\documentclass[a4paper,11pt,oneside,reqno]{amsart}
%
%
\allowdisplaybreaks
%\usepackage[foot]{amsaddr}
%\setlength{\parindent}{0pt}
%\setlength{\parskip}{1pt}
\usepackage[utf8]{inputenc}
\usepackage{tikz}
\usetikzlibrary{arrows.meta,calc}
%\usepackage[pdftex]{graphicx}
\usepackage{amssymb}
\usepackage{esint}
\usepackage{amsmath}
\usepackage{amsfonts}
\usepackage{amsthm}
\usepackage{bm}
\usepackage{mathrsfs}
\usepackage{mathtools}
\usepackage{comment}
\mathtoolsset{showonlyrefs=true}
\usepackage{appendix}
\usepackage{enumerate}
\usepackage[a4paper, total={6.6in, 9.7in}]{geometry}
%\usepackage[left=2.7cm, right=2.7cm, bottom=3.4cm]{geometry}
\usepackage[stretch=30,shrink=30]{microtype}
\linespread{1.1}
\newenvironment{thisnote}{\par\color{blue}}{\par}

% Custom Colours
\usepackage{xcolor}
\definecolor{antiquefuchsia}{rgb}{0.57, 0.36, 0.51}
\definecolor{azure}{rgb}{0.0, 0.5, 1.0}

% Hyperref, add linktocpage inside to change the hyperef to the number instead of the word 
\usepackage[backref=page, colorlinks = true, linkcolor = black, citecolor = antiquefuchsia]{hyperref}

\renewcommand*{\backref}[1]{}
\renewcommand*{\backrefalt}[4]{%
    \ifcase #1 (Not cited.)%
    \or        (Cited on page~#2.)%
    \else      (Cited on pages~#2.)%
    \fi}

% Bold optional argument in statements
\makeatletter
\def\th@plain{%
	\thm@notefont{}% same as heading font
	\itshape % body font
}
\def\th@definition{%
	\thm@notefont{}% same as heading font
	\normalfont % body font
}
\makeatother

\newcommand\res{\mathop{\hbox{\vrule height 7pt width .3pt depth 0pt\vrule height .3pt width 5pt depth 0pt}}\nolimits}

% Number equations according to section
\numberwithin{equation}{section}

\newtheorem{theorem}{Theorem}[section]
\newtheorem{lemma}[theorem]{Lemma}
\newtheorem{proposition}[theorem]{Proposition}
\newtheorem{corollary}[theorem]{Corollary}
\newtheorem{claim}[theorem]{Claim}

\newtheorem*{question}{Question}
\newtheorem{conjecture}[theorem]{Conjecture}
\newtheorem*{open}{Open problem}
\theoremstyle{definition}
\newtheorem{definition}[theorem]{Definition}
\newtheorem{assumption}[theorem]{Assumption}
%\theoremstyle{remark}
\newtheorem{remark}[theorem]{Remark}
\newtheorem{example}[theorem]{Example}

%newcommand
\newcommand{\de}{\partial}
\newcommand{\n}{\mathbb{N}}
\newcommand{\s}{\mathbb{S}}
\renewcommand{\b}{\mathbb{B}}
\newcommand{\z}{\mathbb{Z}}
\renewcommand{\r}{\mathbb{R}}
\let\textcedilla\c
\renewcommand{\c}{\mathbb{C}}
\newcommand{\ind}{\mathrm{ind}}
\newcommand{\nul}{\mathrm{nul}}
\newcommand{\dom}{\mathrm{dom}}
\newcommand{\area}{\mathrm{Area}}
\newcommand{\mres}{\mathbin{\vrule height 1.6ex depth 0pt width
0.13ex\vrule height 0.13ex depth 0pt width 1.3ex}}
\newcommand{\ang}[1]{\langle #1 \rangle}
%\renewcommand{\bar}{\overline}

%DeclareMathOperator
\DeclareMathOperator{\spt}{spt}
\DeclareMathOperator{\HH}{HH}
\DeclareMathOperator{\supp}{supp} 
\DeclareMathOperator{\dvol}{dvol}
\DeclareMathOperator{\dist}{dist}
\DeclareMathOperator{\Per}{Per}
\DeclareMathOperator{\BV}{BV}
\DeclareMathOperator{\curl}{curl}
\DeclareMathOperator{\Maps}{Maps}
\DeclareMathOperator{\vol}{vol}
\DeclareMathOperator{\Area}{Area}
\DeclareMathOperator{\divergence}{div}
\DeclareMathOperator{\Sym}{Sym} 
\DeclareMathOperator{\tr}{tr}
\DeclareMathOperator{\Ric}{Ric}
\DeclareMathOperator{\Hom}{Hom}
\DeclareMathOperator{\Id}{Id}
\DeclareMathOperator{\Hess}{Hess}
\DeclareMathOperator{\loc}{loc}
\DeclareMathOperator{\graph}{graph}
\DeclareMathOperator{\Sing}{Sing}
\DeclareMathOperator{\osc}{osc}
\DeclareMathOperator{\Lip}{Lip}
\DeclareMathOperator{\YM}{YM}
\DeclareMathOperator{\Ind}{Ind}
\DeclareMathOperator{\Jac}{Jac} 
\DeclareMathOperator{\Stab}{Stab}  
\DeclareMathOperator{\id}{id}  
\DeclareMathOperator{\Real}{Re}
\DeclareMathOperator{\Imag}{Im}
\DeclareMathOperator{\End}{End}
\DeclareMathOperator{\E}{E}
\DeclareMathOperator{\hyp}{hyp}
\DeclareMathOperator{\CAT}{CAT}
\DeclareMathOperator{\SL}{SL}
\DeclareMathOperator{\Gr}{Gr}

\newcommand{\flatS}{\mathfrak{F}}

\def\Z{\ensuremath\mathbb{Z}}

\NewDocumentCommand{\sff}{}{\mathrm{I\!I}}

%additional commands
\newcommand{\mz}{\frac{1}{2}}
\newcommand{\inj}{\operatorname{inj}(M)}
\newcommand{\brak}{\mathcal{B}}
\newcommand{\nin}{\not\in}
\newcommand{\uno}{\bm{1}}

\newcommand{\note}[1]{\strut{\color{red}[#1]}}

\newcommand{\toweakstar}{\overset{*}\rightharpoonup}
\DeclareMathOperator{\diverg}{div}


\newcommand{\Reg}{\mathrm{Reg}}
\renewcommand{\H}{\mathrm{H}}
\newcommand{\PSL}{\operatorname{PSL}}
\newcommand{\PSO}{\operatorname{PSO}}
\newcommand{\GL}{\operatorname{GL}}
\newcommand{\U}{\operatorname{U}}
\newcommand{\eps}{\varepsilon}
\newcommand{\pbf}{\mathbf{p}}
\newcommand{\D}{\mathbb{D}}
\newcommand{\cp}{\mathbb{CP}}
\newcommand{\mbf}{\boldsymbol{m}}
\newcommand{\bphi}{\boldsymbol{\varphi}}
\newcommand{\bPhi}{\boldsymbol{\Phi}}
\newcommand{\TV}{\mathrm{TV}}
\newcommand{\hbf}{\mathbf{h}}
\newcommand{\Err}{\mathrm{Err}}
\newcommand{\spn}{\mathrm{span}}
\newcommand{\proofstep}[1]{\textit{#1}}
\setcounter{tocdepth}{1}

\renewcommand\labelitemi{$\vcenter{\hbox{\tiny$\bullet$}}$}

\title[Negatively Curved Minimal Surfaces in the Round Four-Sphere]{Negatively Curved Minimal Surfaces in the Round Four-Sphere}

\author[Riccardo Caniato]{Riccardo Caniato}
\address{Warwick Mathematics Institute, Zeeman Building, University of Warwick, Coventry, CV4 7AL, United Kingdom}
\email{riccardo.caniato@warwick.ac.uk}
%
\date{\today} 
\subjclass[2020]{Primary 53A10; Secondary 53C42, 53C21}

\begin{document}
%
\begin{abstract}
    We construct smooth closed embedded minimal surfaces with everywhere negative Gaussian curvature in the unit round four-sphere. More precisely, for every sufficiently large integer $m$, our construction gives a minimal surface $\Sigma_m$ 
    of genus $m^2+1$ such that $\operatorname{Area}(\Sigma_m)\to 4\pi^2$ as $m\to+\infty$. This result yields a complete solution to Problem 101 in Yau's 1982 list of open problems.
\end{abstract}
%
\maketitle
%
\tableofcontents
%

\section{Introduction}
\label{sec:s4-introduction}
%
Understanding how the curvature of an ambient manifold constrains the intrinsic geometry of its minimal surfaces is a classical and fascinating problem in Riemannian geometry. By the Gauss equation, the Gaussian curvature of a minimal surface equals the ambient sectional curvature along its tangent plane minus one half of the squared norm of its second fundamental form. Nonpositive ambient curvature therefore forces nonpositive Gaussian curvature. On the other hand, positive ambient curvature presents a more subtle question, since the sign of the Gaussian curvature of a minimal surface depends on the balance between the ambient curvature and the second fundamental form. Round spheres, as the model spaces of constant positive sectional curvature, provide a natural setting in which to study this interaction.

Let $\Sigma$ be a smooth Riemannian surface and let $\iota:\Sigma\to\mathbb{S}^n$ be a smooth
minimal immersion, where $g_{\mathbb{S}^n}$ denotes the round
metric of sectional curvature one. Write
$g_\iota:=\iota^*g_{\mathbb{S}^n}$ for its induced metric,
$K_\iota$ for the Gaussian curvature of $g_\iota$, and
$A_\iota$ for its normal-bundle-valued second fundamental
form in $\mathbb{S}^n$. The Gauss equation gives
%
\begin{align*}
    K_\iota=1-\frac{|A_\iota|_{g_\iota}^2}{2},
\end{align*}
%
where the norm uses $g_\iota$ on the tangent bundle and the
restriction of $g_{\mathbb{S}^n}$ on the normal bundle.
Thus minimality gives the upper bound $K_\iota\leq1$,
while everywhere negative Gaussian curvature requires
$|A_\iota|_{g_\iota}^2>2$ at every point.
If $\Sigma$ is closed, connected and orientable, of genus
$\gamma$, the Gauss--Bonnet theorem gives
%
\begin{align*}
    \int_\Sigma K_\iota\,d\operatorname{vol}_{g_\iota}
    =4\pi(1-\gamma).
\end{align*}
%
For $\gamma\geq2$, this forces the curvature to be negative
somewhere, but gives no pointwise upper bound by a negative
number. The existence of minimal surfaces of large genus
therefore does not by itself settle the existence of examples
with everywhere negative curvature.

The classical theory provides several different types of
intrinsic geometry. Bor{\accent23 u}vka constructed minimal two-spheres
of constant positive Gaussian curvature in even-dimensional
round spheres \cite{BoruvkaSphericalSurfaces}, and Calabi
developed a systematic theory of minimal two-spheres in
spheres of arbitrary dimension
\cite{CalabiMinimalImmersionsSpheres}.
The Clifford torus gives a closed flat minimal surface in
$\mathbb{S}^3$. Allowing the curvature to vary, Lawson
\cite{LawsonCompleteMinimalSurfaces} constructed closed
embedded minimal surfaces in $\mathbb{S}^3$ of every
orientable genus. These results show that neither constant
positive curvature, flatness, nor arbitrarily large genus
is incompatible with minimality in a round sphere.
The requirement of everywhere negative curvature is more
restrictive.

In his 1982 problem list, Yau asked whether a round sphere
contains a closed minimal surface whose Gaussian curvature
is everywhere negative
\cite[Problem~101]{YauProblemSection}.
A subsequent theorem of Bryant
\cite[Theorem~2.3]{BryantConstantCurvature}
shows that a surface of constant negative Gaussian curvature
cannot admit a minimal isometric immersion into any
finite-dimensional round sphere, even locally.
Consequently, any solution of Yau's problem must have
nonconstant Gaussian curvature. The question is whether
one can allow the curvature to vary while keeping it
strictly negative at every point. 

An analogous existence problem has been solved in complex
projective geometry. Mohsen
\cite{MohsenNegativelyCurvedCompleteIntersections}
proved that every smooth complex projective manifold of
complex dimension at least three, equipped with a prescribed
Hermitian metric, contains a smooth closed complex curve
whose induced Gaussian curvature is everywhere negative.
When the prescribed metric is K\"ahler, the inclusion of
this curve is minimal, since complex submanifolds of
K\"ahler manifolds are minimal. His construction therefore
provides embedded examples for every K\"ahler metric on
such a projective manifold. The curves are constructed
as complete intersections using high tensor powers of
an ample line bundle and the asymptotic methods of
Donaldson and Auroux; see \cite{Donaldson1996}, \cite{Auroux1997}, \cite{Auroux2001Estimated}.

For round spheres, recent progress has come from the connection
between harmonic maps, unitary representations and random
matrix theory developed by Song
\cite{SongRandomHarmonicMaps}.
In particular, \cite[Theorem~0.3]{SongRandomHarmonicMaps}
constructs closed branched minimal immersions
$\iota_j:\Sigma_j\to\mathbb{S}^{n_j}$, with
$n_j\to+\infty$, such that
%
\begin{align*}
    \lim_{j\to\infty}
    \frac{1}{\operatorname{Area}(\iota_j)}
    \int_{\Sigma_j}
    |K_{\iota_j}+8|\,
    d\operatorname{vol}_{g_{\iota_j}}=0.
\end{align*}
%
The constant $-8$ corresponds to the normalization in which
the ambient spheres have sectional curvature one.
For every fixed $\varepsilon>0$, the proportion of the area
on which $|K_{\iota_j}+8|\geq\varepsilon$ therefore tends
to zero. However, this does not imply negativity at every point, as the convergence of the Gaussian curvatures to $-8$ can be shown to be locally smooth only on compact subsets away from the cusp regions; no analogous uniform statement holds near the images of the cusps. 

Building on Song's method,
Ancona--Labourie--Roig-Sanchis--Toulisse
\cite{AnconaLabourieRoigSanchisToulisse}
obtained closed minimal immersions into all sufficiently
high-dimensional round spheres with everywhere negative
Gaussian curvature. More precisely, their curvatures
converge uniformly to $-8$ as the ambient dimension tends
to infinity. Their construction uses compact
hyperbolic orbifolds with trivial Teichm\"uller space in place of
the thrice-punctured sphere used in Song's construction.
This permits uniform geometric control over the entire
surface. Their argument does not provide an explicit
dimension threshold or establish embeddedness.

In this paper we construct closed embedded examples in
the round four-sphere, the smallest possible ambient
dimension.
Throughout the paper, all round spheres have sectional
curvature one, and the curvature and area of an embedded
surface refer to its induced metric. When a surface is
specified as a submanifold rather than by a parametrization,
we write $K_\Sigma$ for the Gaussian curvature associated
with its inclusion.

\begin{theorem}
    \label{thm:s4-main}
    There exists $m_0\in\mathbb{N}$ such that, for every integer $m\geq m_0$, the unit round sphere $\mathbb{S}^4$ contains a smooth closed connected embedded minimal surface $\Sigma_m$ of genus $m^2+1$ satisfying
    %
    \begin{align*}
        &K_{\Sigma_m}<0 \quad\text{ everywhere on }\Sigma_m \qquad\mbox{ and }\qquad \lim_{m\to\infty}\operatorname{Area}(\Sigma_m)=4\pi^2.
    \end{align*}
\end{theorem}
%
These surfaces are obtained by deforming two nearby toroidal sheets joined by $m^2$ small necks modeled on catenoids. Their areas tend to twice the area of the flat Clifford torus, while their genera tend to infinity. In particular, the Gauss--Bonnet theorem gives
%
\begin{align*}
    \int_{\Sigma_m}K_{\Sigma_m}\,
    d\operatorname{vol}_{\Sigma_m}=-4\pi m^2.
\end{align*}
%
The construction and the estimates that preserve negative
curvature under the deformations are described below.

\begin{corollary}
    \label{cor:s4-higher-spheres}
    Every unit round sphere $\mathbb{S}^n$, $n\geq4$, contains
    a smooth closed connected embedded minimal surface with
    everywhere negative Gaussian curvature.
\end{corollary}
\begin{proof}
    Compose an embedding supplied by Theorem~\ref{thm:s4-main}
    with a totally geodesic inclusion
    $\mathbb{S}^4\hookrightarrow\mathbb{S}^n$.
    This preserves the induced metric, minimality and
    embeddedness.
\end{proof}
%
The ambient dimension in Theorem~\ref{thm:s4-main} is the
smallest possible. Indeed, a closed minimal immersion
$\iota:\Sigma\to\mathbb{S}^3$ cannot have
$K_\iota<0$ everywhere; see
\cite[Proposition~1.5]{LawsonCompleteMinimalSurfaces}.
For an orientable surface, the Hopf quadratic differential
of its second fundamental form is holomorphic, and its
zeros are precisely the umbilic points. Since the immersion
is minimal, an umbilic point satisfies $A_\iota=0$ and
therefore $K_\iota=1$. Everywhere negative curvature would
thus require the Hopf differential to be nowhere zero.
On the other hand, Gauss--Bonnet would force the genus
$\gamma$ to be at least two, and a nonzero holomorphic
quadratic differential on such a surface has
$4\gamma-4$ zeros counted with multiplicity.
This is a contradiction. A nonorientable example would
give an orientable one by passing to its orientable double
cover. Finally, the metric induced by an immersion of a
surface into $\mathbb{S}^2$ has Gaussian curvature one.

Twistor geometry provides a different approach to minimal
surfaces in even-dimensional spheres. Bryant
\cite{BryantMinimalS4} proved that every compact Riemann
surface admits a conformal superminimal immersion into
$\mathbb{S}^4$. In \cite{CaniatoHighDimensionalFamilies},
we constructed families of linearly full branched
superminimal immersions of arbitrarily large dimension,
with both the conformal structure of the domain and the
target sphere $\mathbb{S}^{2n}$, $n\geq2$, fixed.
Here linearly full means that the image is not contained
in any proper totally geodesic subsphere.
In $\mathbb{S}^4$, superminimality means that, at every
regular point, the image of the unit tangent circle under
$v\mapsto A_\iota(v,v)$ is a circle, possibly reduced
to a point, in the normal plane.
This additional condition imposes an obstruction to
negative curvature: a nonconstant branched superminimal
immersion of a closed Riemann surface into
$\mathbb{S}^4$ cannot have nonpositive Gaussian curvature
throughout its regular locus
\cite[Remark~1.4]{CaniatoHighDimensionalFamilies}.
Thus the surfaces of Theorem~\ref{thm:s4-main} lie outside
the superminimal class.
%
\subsection{Outline of the proof}
%
We construct our surfaces in a suitable minimal hypersurface of $\mathcal{N}\subset\mathbb{S}^4$, by joining two nearby copies of the Clifford torus with small catenoidal necks and then correcting their mean curvature. We use the doubling construction of Kapouleas--Yang \cite{KapouleasYangClifford} and the linearized doubling theory of Kapouleas--McGrath \cite{KapouleasMcGrathLD}.

The choice of $\mathcal N$ serves two purposes.
First, the Clifford torus is totally geodesic within
$\mathcal N$, and its scalar Jacobi operator is
$\Delta_{\mathbb T}+2$. We use this operator to
choose the heights of the two sheets so that their
mean curvature within $\mathcal N$ vanishes to
first order. Second, we must control the mean-curvature component
normal to $\mathcal N$: minimality within $\mathcal N$
alone would not imply minimality in $\mathbb S^4$.
For every constant-height torus, this additional
component vanishes. For a sheet defined by a height function over the
central torus, the absolute value of the
mean-curvature component normal to $\mathcal N$
is bounded by a constant times the squared norm
of the gradient of that function. The gradient
and its norm are computed using the flat metric
on $\mathbb T$. Thus, sheets whose height functions vary slowly
have a small additional mean-curvature component.
Together with separate estimates near the necks,
this bound allows us to control the displacement
normal to $\mathcal N$ needed to obtain minimality
in $\mathbb S^4$.

For each sufficiently large integer $m$, we choose
a square lattice $L_m\subset\mathbb T$ with $m^2$
points. We choose $\psi=\tau\Phi_m$ as the principal term of the height function defining the upper sheet, where
%
\begin{align*}
    (\Delta_{\mathbb T}+2)\Phi_m
    =2\pi\sum_{p\in L_m}\delta_p.
\end{align*}
%
Here, $\delta_p$ is the unit Dirac measure at $p$. Since the Dirac measures are supported at the lattice
points, the height function $\psi$ satisfies
$(\Delta_{\mathbb T}+2)\psi=0$ away from $L_m$.
This cancels the linear term in the scalar mean
curvature within $\mathcal N$ of the graph with
height $\psi$. The nonlinear terms will be estimated
and corrected later.
Near a lattice point $p\in L_m$, the Green function $\Phi_m$
has the expansion
%
\begin{align*}
    \Phi_m(q)=\log r+A_m+o(1),
    \qquad r:=d_{\mathbb T}(q,p)\to0.
\end{align*}
%
Thus the corresponding terms in the sheet's height
function $\psi=\tau\Phi_m$ are
$\tau\log r+\tau A_m$.
For a Euclidean catenoid whose smallest circle has
radius $\tau$ and lies at height zero, the height
of its upper end is approximately
$\tau\log r+\tau\log(2/\tau)$ when $r$ is large
compared with $\tau$.
The coefficients of $\log r$ therefore already
agree. To make the additive constants agree as well,
we choose the reference radius $\tau_0$ so that
%
\begin{align*}
    A_m=\log(2/\tau_0),
    \qquad\text{ or equivalently}\qquad
    \tau_0=2e^{-A_m}.
\end{align*}
%
The Green-function estimates give
$A_m=m^2/(2\pi)+\log m+O(1)$, and hence they force
%
\begin{align*}
    \tau_0\asymp
    m^{-1}\exp\left(-\frac{m^2}{2\pi}\right),
\end{align*}
%
with comparison constants independent of $m$. To encode this multiplicative freedom in our construction, we set the common minimum radius of the catenoidal necks to be $\tau_\zeta:=e^\zeta\tau_0$, with $\zeta\in[-1,1]$.
Since, by construction, $\pm\tau_\zeta\Phi_m$ are the principal terms of the two sheet height functions, $\zeta$ varies the neck radius and the
separation of the sheets together. This parameter will be chosen when solving the remaining scalar mean-curvature equation, to cancel the coefficient not determined by the scalar correction alone.

In Section~\ref{sec:s4-scalar-family}, we join the two sheets
to the catenoids at distances comparable to
$\tau_\zeta^\alpha$, for a fixed sufficiently small
$\alpha>0$. These distances are much larger than
$\tau_\zeta$ and much smaller than the lattice spacing,
which is of order $m^{-1}$.
The resulting surface $M_\zeta\subset\mathcal N$
is smooth, closed, connected and embedded, with genus
$m^2+1$.

Two functions are needed to determine suitable normal deformation of $M_{\zeta}$ which will make it minimal.
Writing $\nu_\zeta$ for the outward unit normal of $M_{\zeta}$
within $\mathcal N$, we first move each point within
$\mathcal N$ by a normal displacement $\phi$, and
then move the resulting point in the spherical
normal direction $n$ by a displacement $b$:
%
\begin{align*}
    G_{\zeta,\phi}(x)
    &:=
    \exp_x^{\mathcal N}
    \bigl(\phi(x)\nu_\zeta(x)\bigr),\\
    \widehat X_{\zeta,\phi,b}(x)
    &:=
    \exp_{G_{\zeta,\phi}(x)}^{\mathbb S^4}
    \bigl(b(x)n(G_{\zeta,\phi}(x))\bigr)
    \qquad\forall\,x\in M_\zeta.
\end{align*}
%
The final surface is the image of
$\widehat X_{\zeta,\phi,b}$ for a suitable choice of
$\zeta,\phi,b$.

We impose the symmetries defined in Section~\ref{sec:s4-geometry} on the initial surfaces and all their deformations to make the linear problem solvable with suitable estimates. On the central torus,
lattice invariance excludes the kernel of
$\Delta_{\mathbb T}+2$, making $\Phi_m$ uniquely
determined. On each model neck, the conditions on
$b$ force zero average around every circular
cross-section. The resulting angular Poincar\'e
inequality is essential to obtaining inverse
estimates independent of the neck radius.
The conditions on $\phi$ allow a constant angular
component, which requires a separate linear argument.

To test minimality after deformation, we parallel
transport the two normals of
$G_{\zeta,\phi}(M_\zeta)$ along the spherical
displacement geodesics and project onto the new
normal planes. For small deformations, these
projections remain independent. Minimality is
therefore equivalent to the vanishing of the
mean-curvature vector's inner products with both
normal fields.

For every $\zeta\in[-1,1]$ and every $\phi$ satisfying
the bounds in Section~\ref{sec:s4-transverse},
we first choose $b$ to vanish the component associated with $n$.
The angular inequality, metric comparisons and
elliptic estimates give a uniformly bounded inverse
for the linearized equation.
A contraction argument then solves the nonlinear
equation, with
%
\begin{align*}
    \|b\|_{L^\infty(M_\zeta)}\leq C\tau_\zeta^2.
\end{align*}
%
Section~\ref{sec:s4-transverse} also controls the derivatives of $b$ up to order too, which appear in the curvature estimates.

The displacement $b$ changes the other mean-curvature
component. Section~\ref{sec:s4-coupled} handles this change when solving for $\phi$ and $\zeta$.
The scalar linear construction determines a
correction up to an error proportional to one
prescribed function. Adjusting $\zeta$ cancels
the remaining coefficient.
Since $b$ depends on $(\phi,\zeta)$, these choices
are made by solving a fixed-point problem.
The estimates give a continuous map preserving
a compact convex set of corrections and parameters;
the polynomial losses in $m$ are absorbed by
positive powers of $\tau_\zeta$.
Schauder's fixed-point theorem supplies a choice
for which both normal components vanish.
The resulting surface
$\Sigma_m:=\widehat X_{\zeta,\phi,b}(M_\zeta)$
is therefore minimal in $\mathbb S^4$.

Strictly negative curvature requires a further
comparison. Write $\tau:=\tau_\zeta$ for the final
neck parameter. For a sheet defined by a height
function $w$, the quadratic curvature expression is
%
\begin{align*}
    \mathcal K_2(w)
    :=\det D^2w+w\Delta_{\mathbb T}w+|Dw|^2,
\end{align*}
%
where derivatives, norms and determinants use the
flat metric on $\mathbb T$.
For $\psi=\tau\Phi_m$, the equation
$\Delta_{\mathbb T}\psi=-2\psi$ and the
Green-function estimates give
%
\begin{align*}
    \mathcal K_2(\psi)
    =|D\psi|^2-\frac12|D^2\psi|^2
    \leq-c\tau^2\bigl(m^4+\rho^{-4}\bigr),
    \qquad
    \rho(q):=d_{\mathbb T}(q,L_m),
\end{align*}
%
for $q\in\mathbb T\smallsetminus L_m$, with
$c>0$ independent of $m$.
Thus the Hessian of the chosen height function
provides the negative curvature bound.
At distances comparable to the lattice spacing,
this bound has size $m^4\tau^2\to0$.
Absolute smallness of the corrections is therefore
insufficient: the curvature errors must be smaller
than the negative bound.

Section~\ref{sec:s4-negative-curvature} proves
precisely this comparison on the sheet portions
outside a sufficiently large fixed multiple of
the neck radius. The interpolation, the nonlinear
curvature terms and both deformations together
produce an error smaller than half the negative
bound, once $m$ is sufficiently large.

On the remaining neck portions, we divide lengths
by $\tau$. Both deformations tend to zero in
$C^2$ in the rescaled coordinates, and the final
necks converge to the Euclidean catenoid of minimum
radius one. The catenoid has curvature bounded
above by a strictly negative constant on each
fixed bounded portion. Consequently, the final
necks satisfy $K_{\Sigma_m}\leq-c_R\tau^{-2}$
there. The sheet and neck estimates cover the
entire surface.

Finally, the neck and interpolation regions have
total area $O(m^2\tau^{2\alpha})\to0$.
Each remaining sheet has area tending to the
Clifford torus area $2\pi^2$, and the corrections
preserve these limits. Hence
$\operatorname{Area}(\Sigma_m)\to4\pi^2$.

\subsection{Open problems}
\label{subsec:s4-open-problems}

The results above leave several open questions concerning general ambient metrics, the possible genera of the surfaces, and the existence of negatively curved examples within more specific subclasses of minimal surfaces.

In particular, Song has proposed the following extension of Yau's question
to general closed Riemannian manifolds
\cite[p.~4]{SongQuestions}.

\begin{question}[Song]
    \label{que:s4-general-ambient}
    Let $(M,g)$ be a closed Riemannian manifold with
    $\dim P\geq4$. Does there exist a closed connected
    surface $\Sigma$ and a smooth minimal immersion
    $\iota:\Sigma\to(M,g)$ such that the Gaussian
    curvature of $g_\iota:=\iota^*g$ is strictly
    negative at every point?
\end{question}
%
Corollary~\ref{cor:s4-higher-spheres} answers this question in full generality for round spheres. Mohsen's theorem gives another class of examples, namely complex projective manifolds of complex dimension at least three equipped with K\"ahler metrics. For a general minimal immersion into $(M,g)$, the Gauss
equation reads
%
\begin{align*}
    K_\iota(x)
    =\sec_g\bigl(d\iota_x(T_x\Sigma)\bigr)
    -\frac{|A_\iota(x)|_{g_\iota}^2}{2}
    \qquad\forall\,x\in\Sigma,
\end{align*}
%
where $\sec_g$ denotes ambient sectional curvature and
$A_\iota$ is the second fundamental form in $(M,g)$.
Thus one must construct the immersion while controlling
its second fundamental form relative to the curvature
of its tangent planes in the ambient manifold.
Whether the immersion can always be chosen to be an
embedding is a stronger question.

A more restricted problem is to establish the same
existence result for every smooth metric in some
$C^\infty$-neighborhood of the round metric on
$\mathbb{S}^4$. Strict negativity of Gaussian curvature
persists under sufficiently small $C^2$ perturbations
of a fixed immersion and the ambient metric. However,
one must also construct a minimal immersion for the
perturbed metric, as curvature continuity alone does not
provide this.

Even for the round metric, the present construction does
not determine all genera for which such examples exist.
Theorem~\ref{thm:s4-main} gives the genera $m^2+1$ for
all sufficiently large integers $m$.

\begin{question}
    \label{que:s4-prescribed-genus}
    Which integers $\gamma\geq2$ occur as the genus of
    a smooth closed connected orientable embedded minimal
    surface in $\mathbb{S}^4$ with everywhere negative
    Gaussian curvature? In particular, does every
    sufficiently large integer occur?
\end{question}
%
The lower bound $\gamma\geq2$ follows from Gauss--Bonnet. Determining the smallest possible genus is already a
separate problem.

Finally, the superminimal existence problem discussed
in \cite[Conjecture~1.3]{CaniatoHighDimensionalFamilies}
remains distinct from the existence result proved here.
Its first possible ambient sphere is $\mathbb{S}^6$.
In particular, one can ask for an unbranched embedded
example as follows.

\begin{question}
    \label{que:s4-superminimal-six-sphere}
    Does $\mathbb{S}^6$ contain a smooth closed connected
    embedded linearly full superminimal surface with
    everywhere negative Gaussian curvature?
\end{question}
%
The examples obtained from
Corollary~\ref{cor:s4-higher-spheres} lie in a totally
geodesic $\mathbb{S}^4\subset\mathbb{S}^6$ and therefore
do not answer this question. A construction within the
superminimal class would have to impose the curvature
inequality together with the holomorphicity and
horizontality conditions on the corresponding twistor
lifts.
\renewcommand\labelitemi{$\vcenter{\hbox{\tiny$\bullet$}}$}
\begingroup
\renewcommand{\addcontentsline}[3]{}
\subsection*{On the use of AI}
The author used ChatGPT-6 Astra solely to improve the clarity and style of the prose, to assist in identifying potential errors or gaps in the mathematical arguments, and to create the figures appearing in the manuscript. The research ideas and all the proofs were developed without the use of AI tools. All mathematical arguments were independently checked by the author, who takes full responsibility for the content of this paper.
\par
\endgroup
%
\renewcommand\labelitemi{$\vcenter{\hbox{\tiny$\bullet$}}$}
\begingroup
\renewcommand{\addcontentsline}[3]{}
\subsection*{Acknowledgements}
I am very grateful to Antoine Song, for introducing me to this problem during our first meeting at Caltech more than three years ago, and for many enlightening discussions we had since then on this topic.

I would also like to thank Xingzhe Li, Ben Lowe, and Daniel Stern, for the interesting exchanges we had over time on (equivariant) minimal surfaces in round spheres.
\par
\endgroup
%
\section{Geometry and symmetries near the Clifford torus}
\label{sec:s4-geometry}

We first embed the Clifford torus inside an explicit minimal
hypersurface $\mathcal{N}$ of the round four-sphere and compute the geometric quantities needed in the construction. For a surface in $\mathcal{N}$, we then express the two components of its mean-curvature vector within $\mathbb{S}^4$ and their normal linearization, and provide suitable tubular coordinates for deformations out of the hypersurface. Finally, we specify the finite group of symmetries used in the construction and derive the resulting restrictions on the normal variations and their angular Fourier terms near each neck.

\subsection{A minimal hypersurface containing the Clifford torus}

The hypersurface $\mathcal{N}$ constructed below will contain the initial surfaces used in the perturbation argument. Its central Clifford torus is totally geodesic within it, and the scalar Jacobi operator of that torus within $\mathcal{N}$ has a particularly simple expression.

We regard the unit round four-sphere as
%
\begin{align*}
    \mathbb{S}^4=\{(w_1,w_2,x_5)\in\mathbb{C}^2\oplus\mathbb{R}\mbox{ : }|w_1|^2+|w_2|^2+x_5^2=1\}.
\end{align*}
%
Set $\ell:=\sqrt2\pi$ and let $\mathbb{T}:=(\mathbb{R}/\ell\mathbb{Z})^2$, equipped with the flat metric $g_{\mathbb{T}}:=ds^2+dt^2$. Define $f:\mathbb{T}\to\mathbb{S}^4$ by
%
\begin{align}
    \label{eq:s4-clifford-frame}
    f(s,t):=\frac{1}{\sqrt{2}}
    \bigl(e^{i\sqrt{2}s},e^{i\sqrt{2}t},0\bigr) \qquad\forall\,(s,t)\in\mathbb{T}.
\end{align}
%
Note that $f$ is an isometric embedding of $\mathbb{T}$ into $\mathbb{S}^4$ and $\operatorname{Area}(\mathbb{T})=2\pi^2$. Fix $z_0\in(0,\pi/4)$ and define the smooth embedding $\mathcal F:\mathbb{T}\times(-z_0,z_0)\to\mathbb{S}^4$ by
%
\begin{align}
    \label{eq:s4-suspension}
    \mathcal{F}(s,t,z):=\cos z\,f(s,t)+\sin z\,e_5 \qquad\forall\,(s,t,z)\in\mathbb{T}\times(-z_0,z_0),
\end{align}
%
where $e_5=(0,0,1)\in\mathbb{C}^2\oplus\mathbb{R}$. We let 
%
\begin{align*}
    \mathcal{N}:=\mathcal{F}\bigl(\mathbb{T}\times(-z_0,z_0)\bigr)
\end{align*}
%
and we define the projection $\pi:\mathcal{N}\to\mathbb{T}$ as
%
\begin{align*}
    \pi(\mathcal{F}(s,t,z)):=(s,t) \qquad\forall(s,t,z)\in\mathbb{T}\times(-z_0,z_0).
\end{align*}
%
Let $g_{\mathbb{S}^4}$ denote the unit round metric on $\mathbb{S}^4$,
let $\iota_{\mathcal{N}}:\mathcal{N}\hookrightarrow\mathbb{S}^4$
denote the inclusion, and equip $\mathcal{N}$ with the induced
metric $g_{\mathcal{N}}:=\iota_{\mathcal{N}}^*g_{\mathbb{S}^4}$.
Note that 
%
\begin{align}
    \label{eq:s4-suspension-metric}
    g_{\mathcal{N}}=(\mathcal{F}^{-1})^*\bigl(dz^2+\cos^2z\,(ds^2+dt^2)\bigr).
\end{align}
%
Define $N:\mathbb{T}\to\mathbb{C}^2\oplus\mathbb{R}$ by
%
\begin{align*}
    N(s,t):=\frac{1}{\sqrt{2}}\bigl(e^{i\sqrt{2}s},-e^{i\sqrt{2}t},0\bigr)\qquad\forall\,(s,t)\in\mathbb{T}.
\end{align*}
%
The vector $N(s,t)$ has unit Euclidean norm and is orthogonal to $\mathcal{F}(s,t,z)$ and to each of $\partial_s\mathcal{F}$, $\partial_t\mathcal{F}$ and $\partial_z\mathcal{F}$ at $(s,t,z)$. Consequently, the field $n$ defined along $\mathcal{N}$ by
%
\begin{align*}
    n\bigl(\mathcal{F}(s,t,z)\bigr):=N(s,t) \qquad\forall\,(s,t,z)\in\mathbb{T}\times(-z_0,z_0)
\end{align*}
%
is a smooth global unit normal to $\mathcal{N}$ within $\mathbb{S}^4$.

We call $\mathcal{N}$ the \textit{suspension band} of the Clifford torus and refer to $f(\mathbb{T})=\mathcal{F}(\mathbb{T}\times\{0\})$ as its \textit{central torus}. 

\begin{remark}
 \label{rem:s4-signed-distance}
    Note that the function on $\mathcal{N}$ defined by $\mathcal{F}(s,t,z)\mapsto z$ is the signed distance to the central torus in $(\mathcal{N},g_{\mathcal{N}})$, i.e.
    %
    \begin{align*}
        \operatorname{dist}_{g_{\mathcal{N}}}
        \bigl(\mathcal{F}(s,t,z),f(\mathbb{T})\bigr)=|z|
        \qquad\forall\,(s,t,z)\in
        \mathbb{T}\times(-z_0,z_0).
    \end{align*}
    %
    Indeed, the coordinate segment joining $\mathcal{F}(s,t,z)$ to $\mathcal{F}(s,t,0)$ has length $|z|$, while every curve joining $\mathcal{F}(s,t,z)$ to the central torus has length at least $|z|$, by \eqref{eq:s4-suspension-metric}.
\end{remark}
%
% Figure~\ref{fig:s4-band-coordinates} shows the parameter
% domain and the normal geodesics used in this definition.
%
\begin{figure}[htbp]
    \centering
    \begingroup
    \definecolor{bandteal}{RGB}{0,112,116}
    \definecolor{bandorange}{RGB}{190,102,36}
    \begin{tikzpicture}[
        >=Stealth,
        every node/.style={font=\small}
    ]
        \begin{scope}[
            x={(1cm,0cm)},y={(0.42cm,0.28cm)},z={(0cm,0.9cm)}
        ]
            \node[font=\normalsize] at (1.65,1,2.28)
                {The parameter domain of $\mathcal F$};
            \fill[bandteal!12]
                (0,0,0)--(3.3,0,0)--(3.3,2.3,0)--(0,2.3,0)--cycle;
            \foreach \height in {-1.1,1.1} {
                \draw[black!35,dashed]
                    (0,0,\height)--(3.3,0,\height)--(3.3,2.3,\height)
                    --(0,2.3,\height)--cycle;
            }
            \foreach \x in {0,3.3} {
                \foreach \y in {0,2.3} {
                    \draw[black!40] (\x,\y,-1.1)--(\x,\y,1.1);
                }
            }
            \foreach \x in {0.825,1.65,2.475} {
                \draw[bandteal!35] (\x,0,0)--(\x,2.3,0);
            }
            \foreach \y in {0.575,1.15,1.725} {
                \draw[bandteal!35] (0,\y,0)--(3.3,\y,0);
            }
            \draw[bandteal,line width=0.9pt]
                (0,0,0)--(3.3,0,0)--(3.3,2.3,0)--(0,2.3,0)--cycle;
            \foreach \y in {0,2.3} {
                \draw[bandteal,line width=1pt,->]
                    (1.42,\y,0)--(1.93,\y,0);
            }
            \foreach \x in {0,3.3} {
                \draw[bandorange,line width=1pt,->]
                    (\x,0.85,0)--(\x,1.4,0);
            }
            \draw[bandorange,line width=1.1pt]
                (1.03,0.72,-0.87)--(1.03,0.72,0.87);
            \fill[bandorange] (1.03,0.72,0) circle[radius=1.8pt];
            \node[anchor=south east,text=bandorange] at (1.03,0.72,0.1)
                {$(q,0)$};
            \node[anchor=west] at (3.48,2.3,1.1) {$z=z_0$};
            \node[anchor=west] at (3.48,2.3,-1.1) {$z=-z_0$};
            \node[anchor=north] at (2.45,0,-0.14)
                {$\mathbb T\times\{0\}$};
            \draw[black!50,->] (0,0,-1.1)--(0,0,1.67)
                node[above] {$z$};
        \end{scope}
        \begin{scope}[shift={(9.1,0.55)}]
            \node[font=\normalsize] at (0,2.25)
                {A fixed $q\in\mathbb T$};
            \draw[black!35] (0,0) circle[radius=1.5];
            \draw[black!45,->] (-1.75,0)--(1.92,0);
            \draw[black!45,->] (0,-1.65)--(0,1.73)
                node[right] {$e_5$};
            \coordinate (base) at (1.5,0);
            \coordinate (point) at ({1.5*cos(32)},{1.5*sin(32)});
            \draw[black!40] (0,0)--(point);
            \draw[bandorange,line width=1.5pt]
                ({1.5*cos(-32)},{1.5*sin(-32)})
                arc[start angle=-32,end angle=32,radius=1.5];
            \fill[bandteal] (base) circle[radius=1.8pt];
            \fill[bandorange] (point) circle[radius=1.8pt];
            \node[below right=2pt] at (base) {$f(q)$};
            \node[above right=3pt,text=bandorange] at (point)
                {$\mathcal F(q,z)$};
            \draw[black!65,->]
                (0.48,0) arc[start angle=0,end angle=32,radius=0.48];
            \node at (0.69,0.18) {$z$};
            \node[align=center,text=black!65] at (0,-2.02)
                {$\mathcal F(q,z)=\cos z\,f(q)+\sin z\,e_5$};
        \end{scope}
    \end{tikzpicture}
    \endgroup
    \caption{The hypersurface $\mathcal N$ is the image under
    $\mathcal F$ of the parameter domain on the left.
    Opposite vertical faces are identified in the $s$ and
    $t$ directions; the central square, with its opposite
    edges identified, represents $\mathbb T\times\{0\}$.
    The boundary heights $\pm z_0$ are not included.
    The orange segment keeps $q$ fixed.
    On the right, its image is an arc of the unit circle
    in the plane spanned by $f(q)$ and $e_5$.
    % The box depicts coordinates, not the induced metric
    % or an embedding of $\mathcal N$ into $\mathbb R^3$.
    }
    \label{fig:s4-band-coordinates}
\end{figure}

%
We fix the following conventions for immersions, which will
also be used in the subsequent sections.
Let $\iota:M\to(P,\overline g)$ be a smooth immersion and
equip $M$ with the induced metric $g:=\iota^*\overline g$.
We denote the Levi-Civita connection of $g$ by $\nabla^g$.
For a smooth function $u$ on $M$, the notation $\nabla^g u$
denotes its gradient with respect to $g$. We use the convention
%
\begin{align*}
    \Delta_g u:=\operatorname{div}_g(\nabla^g u)
    \qquad\forall\,u\in C^\infty(M)
\end{align*}
%
for the Laplace--Beltrami operator.
We omit the metric subscript when it is clear from the context.

Denote by $\overline\nabla$ the connection on $\iota^*TP$
induced by the Levi-Civita connection of $(P,\overline g)$.
Let
%
\begin{align*}
    T^\perp_\iota M
    :=\bigl(d\iota(TM)\bigr)^\perp\subset\iota^*TP
\end{align*}
%
be the normal bundle of $\iota$, where ${\cdot}^\perp$ indicates the $\overline g$-orthogonal complement of $d\iota(TM)$ in $\iota^*TP$. For a vector in $\iota^*TP$, the superscript ${}^\perp$
denotes its $\overline g$-orthogonal projection onto
$T^\perp_\iota M$.
The \textit{second fundamental form} of $\iota$ is the tensor $A_\iota\in\Gamma(\operatorname{Sym}^2T^*M\otimes
T^\perp_\iota M)$ defined by
%
\begin{align*}
    A_\iota(X,Y)
    :=\bigl(\overline\nabla_X(d\iota(Y))\bigr)^\perp
    \qquad\forall\,X,Y\in\Gamma(TM).
\end{align*}
%
For a smooth unit normal field $\eta\in\Gamma(T^\perp_\iota M)$, its scalar
component in the direction $\eta$ is the bilinear form
$h_\iota^\eta\in\Gamma(\operatorname{Sym}^2T^*M)$, defined by
%
\begin{align*}
    h_\iota^\eta(X,Y)
    :=\overline g\bigl(A_\iota(X,Y),\eta\bigr)
    \qquad\forall\,X,Y\in\Gamma(TM).
\end{align*}
%
The \textit{shape operator} of $\iota$ in the direction $\eta$ is the $g$-self-adjoint endomorphism
$S_\iota^\eta\in\Gamma(T^*M\otimes TM)$ characterized by
%
\begin{align*}
    g(S_\iota^\eta(X),Y)=h_\iota^\eta(X,Y)
    \qquad\forall\,X,Y\in\Gamma(TM).
\end{align*}
%
In particular, $|S_\iota^\eta|_g=|h_\iota^\eta|_g$.
The \textit{mean-curvature vector} of $\iota$ and its scalar component in the direction $\eta$ are defined using the trace convention:
%
\begin{align*}
    \mathbf{H}_\iota
    &:=\operatorname{tr}_g A_\iota
    =\sum_{a=1}^{\dim M}A_\iota(e_a,e_a),\\
    H_{\iota,\eta}
    &:=\operatorname{tr}_g h_\iota^\eta
    =\sum_{a=1}^{\dim M}h_\iota^\eta(e_a,e_a)
    =\overline g(\mathbf{H}_\iota,\eta),
\end{align*}
%
where $(e_a)_{a=1}^{\dim M}$ is any local
$g$-orthonormal frame of $TM$.

The \textit{normal connection} of $\iota$ is given by
%
\begin{align*}
    \nabla^{\perp,\iota}_X\xi
    :=(\overline\nabla_X\xi)^\perp
    \qquad\forall\,X\in\Gamma(TM),\quad
    \xi\in\Gamma(T^\perp_\iota M).
\end{align*}
%
With these conventions, the Weingarten identity reads
%
\begin{align}
    \label{eq:s4-weingarten-convention}
    \overline\nabla_X\eta
    =-d\iota(S_\iota^\eta(X))
    +\nabla^{\perp,\iota}_X\eta
    \qquad\forall\,X\in\Gamma(TM).
\end{align}
%
Indeed, differentiating $\overline g(\eta,d\iota(Y))=0$
gives
%
\begin{align*}
    \overline g(\overline\nabla_X\eta,d\iota(Y))
    &=-\overline g(\eta,\overline\nabla_X(d\iota(Y)))\\
    &=-h_\iota^\eta(X,Y)
    =-g(S_\iota^\eta(X),Y)
    \qquad\forall\,X,Y\in\Gamma(TM),
\end{align*}
%
which determines the tangential component of
$\overline\nabla_X\eta$; its normal component is
$\nabla^{\perp,\iota}_X\eta$ by definition.
In codimension one, this normal component vanishes because
$\eta$ has unit length and spans the normal bundle.

For the fixed metrics on $\mathbb{S}^4$ and $\mathcal{N}$,
we use the abbreviations
%
\begin{align*}
    \nabla^{\mathbb{S}^4}:=\nabla^{g_{\mathbb{S}^4}},
    \qquad
    \nabla^{\mathcal{N}}:=\nabla^{g_{\mathcal{N}}},
\end{align*}
%
for their respective Levi-Civita connections.

\begin{lemma}
    \label{lem:s4-suspension-geometry}
    Write $S:=S_{\iota_{\mathcal{N}}}^n\in\Gamma(T^*\mathcal{N}\otimes T\mathcal{N})$ for the shape
    operator of the inclusion $\iota_{\mathcal{N}}:\mathcal{N}\hookrightarrow\mathbb{S}^4$ with respect to the unit normal $n$, using the preceding conventions.
    Define $E_1,E_2,E_3\in\Gamma(T\mathcal{N})$ by
    %
    \begin{align*}
        E_1(\mathcal{F}(s,t,z))&:=\frac{1}{\cos z}\partial_s\mathcal F(s,t,z),\\
        E_2(\mathcal{F}(s,t,z))&:=\frac{1}{\cos z}\,
        \partial_t\mathcal F(s,t,z),\\
        E_3(\mathcal{F}(s,t,z))&:=\partial_z\mathcal F(s,t,z),
    \end{align*}
    %
    for every $(s,t,z)\in\mathbb{T}\times(-z_0,z_0)$.
    Then $(E_1,E_2,E_3)$ is a global
    $g_{\mathcal{N}}$-orthonormal frame of $T\mathcal{N}$. For every $(s,t,z)\in\mathbb{T}\times(-z_0,z_0)$, writing $q:=\mathcal{F}(s,t,z)$, one has
    %
    \begin{align}
        \label{eq:s4-S}
        \begin{split}
            S_q(E_1(q))&=-\sec z\,E_1(q),\\
            S_q(E_2(q))&=\sec z\,E_2(q),\\
            S_q(E_3(q))&=0.
        \end{split}
    \end{align}
    %
    In particular, $\mathcal{N}$ is a minimal hypersurface of $\mathbb{S}^4$.

    The central torus $f(\mathbb{T})\subset\mathcal N$ is totally geodesic as a hypersurface of $(\mathcal{N},g_{\mathcal{N}})$. The scalar Jacobi operator
    $L_{\mathbb{T}}:C^\infty(\mathbb{T})\to C^\infty(\mathbb{T})$
    of the isometric embedding
    $f:(\mathbb{T},g_{\mathbb{T}})\to(\mathcal{N},g_{\mathcal{N}})$,
    with respect to the unit normal field $E_3\circ f$,
    is given by
    %
    \begin{align}
        \label{eq:s4-base-Jacobi}
        L_{\mathbb{T}}u=\Delta_{\mathbb{T}}u+2u
        \qquad\forall\,u\in C^\infty(\mathbb{T}),
    \end{align}
    %
    where $\Delta_{\mathbb{T}}:=\Delta_{g_{\mathbb{T}}}$.

    Finally, there exists a constant $C=C(z_0)>0$ such that, for every $(s,t,z)\in\mathbb{T}\times(-z_0,z_0)$ and $q:=\mathcal{F}(s,t,z)$,
    %
    \begin{align}
        \label{eq:s4-S-derivative}
        |S_q|_{g_{\mathcal{N}}}\leq C, 
        \qquad
        |(\nabla^{\mathcal{N}}S)_q|_{g_{\mathcal{N}}}
        \leq C|z|.
    \end{align}
    %
    Here $\nabla^{\mathcal{N}}S$ is the covariant derivative of
    the tensor field $S$, and both norms are induced by
    $g_{\mathcal{N}}$.
\end{lemma}

\begin{proof}
    Fix $(s,t,z)\in\mathbb{T}\times(-z_0,z_0)$ and set $q:=\mathcal{F}(s,t,z)$.
    Differentiating $\mathcal{F}$ and using the
    definitions of the frame fields gives
    %
    \begin{align*}
        E_1(q)&=\partial_sf(s,t),\\
        E_2(q)&=\partial_tf(s,t),\\
        E_3(q)&=-\sin z\,f(s,t)+\cos z\,e_5.
    \end{align*}
    %
    These three vectors are orthonormal in
    $\mathbb{R}^5$, and therefore also with respect to
    the induced metric $g_{\mathcal{N}}$.
    All inner products below are Euclidean inner products
    in $\mathbb{R}^5$.
    Viewing $n$ as a map from $\mathcal{N}$ to
    $\mathbb{R}^5$, the spherical covariant derivative
    satisfies
    %
    \begin{align*}
        \nabla^{\mathbb{S}^4}_v n
        =dn_q(v)+\langle v,n(q)\rangle q
        =dn_q(v)
        \qquad\forall\,v\in T_q\mathcal{N}.
    \end{align*}
    %
    The last equality follows from the orthogonality
    of $n(q)$ to $T_q\mathcal{N}$.
    Consequently,
    %
    \begin{align*}
        \bigl(\nabla^{\mathbb{S}^4}_{E_1}n\bigr)(q)
        &=dn_q(E_1(q))
        =\frac{1}{\cos z}\,
        \partial_s(n\circ\mathcal{F})(s,t,z)\\
        &=\frac{1}{\cos z}\,\partial_sN(s,t)
        =\frac{1}{\cos z}\,\partial_sf(s,t)
        =\sec z\,E_1(q).
    \end{align*}
    %
    Similarly, using $\partial_tN=-\partial_tf$
    and the independence of $n\circ\mathcal{F}$ from $z$,
    we obtain
    %
    \begin{align*}
        \bigl(\nabla^{\mathbb{S}^4}_{E_2}n\bigr)(q)
        &=-\sec z\,E_2(q),\\
        \bigl(\nabla^{\mathbb{S}^4}_{E_3}n\bigr)(q)
        &=0.
    \end{align*}
    %
    By \eqref{eq:s4-weingarten-convention}, this proves \eqref{eq:s4-S}.
    In particular, $\operatorname{tr}S=0$, so
    $\mathcal{N}$ is minimal in $\mathbb{S}^4$.

    We next compute the coefficients of the Levi-Civita connection of
    $g_{\mathcal{N}}$ in the orthonormal frame $(E_1,E_2,E_3)$.
    Viewing tangent vector fields on $\mathcal{N}$
    as $\mathbb{R}^5$-valued maps, orthogonal projection
    of their Euclidean derivatives onto $T_q\mathcal{N}$
    gives
    %
    \begin{align*}
        \bigl(\nabla^{\mathcal{N}}_X Y\bigr)(q)
        &=dY_q(X(q))+\langle X(q),Y(q)\rangle q\\
        &\quad-\langle S_q(X(q)),Y(q)\rangle n(q)
        \qquad\forall\,X,Y\in\Gamma(T\mathcal{N}).
    \end{align*}
    %
    Indeed, the Euclidean normal space to $\mathcal{N}$
    at $q$ is spanned by the orthonormal vectors
    $q$ and $n(q)$.
    %
    For example, since
    $\partial_{ss}f=-f-N$, we obtain
    %
    \begin{align*}
        \bigl(\nabla^{\mathcal{N}}_{E_1}E_1\bigr)(q)
        &=
        \frac{1}{\cos z}\,
        \partial_s(E_1\circ\mathcal{F})(s,t,z)
        +q+\sec z\,n(q)\\
        &=
        \sec z\,\partial_{ss}f(s,t)
        +q+\sec z\,N(s,t)\\
        &=
        -\sec z\bigl(f(s,t)+N(s,t)\bigr)
        +q+\sec z\,N(s,t)\\
        &=q-\sec z\,f(s,t)\\
        &=
        \tan z\bigl(-\sin z\,f(s,t)+\cos z\,e_5\bigr)\\
        &=\tan z\,E_3(q).
    \end{align*}
    %
    Applying the same projection formula to the
    remaining pairs of frame fields gives
    %
    \begin{align*}
        \bigl(\nabla^{\mathcal{N}}_{E_i}E_j\bigr)(q)
        &=\tan z\,\delta_{ij}E_3(q),\\
        \bigl(\nabla^{\mathcal{N}}_{E_i}E_3\bigr)(q)
        &=-\tan z\,E_i(q),\\
        \bigl(\nabla^{\mathcal{N}}_{E_3}E_i\bigr)(q)
        &=0,\\
        \bigl(\nabla^{\mathcal{N}}_{E_3}E_3\bigr)(q)
        &=0,
    \end{align*}
    %
    for every $i,j\in\{1,2\}$. Hence, at $z=0$ the normal components of $\nabla^{\mathcal{N}}_{E_i}E_j$ vanish for $i,j\in\{1,2\}$. Thus, the central torus is totally geodesic in $\mathcal{N}$.

    Since $S(E_3)=0$, the Gauss equation for
    $\mathcal{N}\subset\mathbb{S}^4$ shows that
    each plane spanned by $E_i$ and $E_3$,
    with $i\in\{1,2\}$, has sectional curvature $1$.
    Consequently,
    %
    \begin{align*}
        \operatorname{Ric}_{g_{\mathcal{N}}}(E_3,E_3)=2.
    \end{align*}
    %
    Let $\nu:=E_3\circ f$ be the unit normal field along
    the isometric embedding
    $f:(\mathbb{T},g_{\mathbb{T}})\to(\mathcal{N},g_{\mathcal{N}})$.
    Since the central torus is totally geodesic in $\mathcal{N}$,
    we have $h_f^\nu\equiv0$. The scalar Jacobi formula gives
    %
    \begin{align*}
        L_{\mathbb{T}}u
        &=\Delta_{\mathbb{T}}u+
        \Bigl(
        |h_f^\nu|_{g_{\mathbb{T}}}^2
        +\bigl(\operatorname{Ric}_{g_{\mathcal{N}}}(E_3,E_3)\bigr)
        \circ f
        \Bigr)u\\
        &=\Delta_{\mathbb{T}}u+2u
        \qquad\forall\,u\in C^\infty(\mathbb{T}),
    \end{align*}
    %
    which proves \eqref{eq:s4-base-Jacobi}.

    Finally, set
    %
    \begin{align*}
        \lambda_1(z):=-\sec z,
        \qquad
        \lambda_2(z):=\sec z
        \qquad\forall\,z\in(-z_0,z_0).
    \end{align*}
    %
    Using the connection identities above and the
    formula
    %
    \begin{align*}
        (\nabla^{\mathcal{N}}_X S)(Y)
        =
        \nabla^{\mathcal{N}}_X(S(Y))
        -S(\nabla^{\mathcal{N}}_X Y)
        \qquad\forall\,X,Y\in\Gamma(T\mathcal{N}),
    \end{align*}
    %
    we obtain, for every $i,j\in\{1,2\}$,
    %
    \begin{align*}
        \bigl((\nabla^{\mathcal{N}}_{E_3}S)(E_i)\bigr)(q)
        &=\lambda_i(z)\tan z\,E_i(q),\\
        \bigl((\nabla^{\mathcal{N}}_{E_i}S)(E_j)\bigr)(q)
        &=\lambda_j(z)\tan z\,\delta_{ij}E_3(q),\\
        \bigl((\nabla^{\mathcal{N}}_{E_i}S)(E_3)\bigr)(q)
        &=\lambda_i(z)\tan z\,E_i(q),\\
        \bigl((\nabla^{\mathcal{N}}_{E_3}S)(E_3)\bigr)(q)
        &=0.
    \end{align*}
    %
    Since $(E_1,E_2,E_3)$ is orthonormal, these identities
    and \eqref{eq:s4-S} imply
    %
    \begin{align*}
        |S_q|_{g_{\mathcal{N}}}^2=2\sec^2z,
        \qquad
        |(\nabla^{\mathcal{N}}S)_q|_{g_{\mathcal{N}}}^2
        =6\sec^2z\,\tan^2z.
    \end{align*}
    %
    For every $z\in(-z_0,z_0)$, one has
    $\sec z\leq\sec z_0$ and
    $|\tan z|\leq\sec^2z_0\,|z|$.
    The estimates \eqref{eq:s4-S-derivative} follow.
\end{proof}

\subsection{Normal variations and tubular coordinates}

For a surface $M$ contained in $\mathcal{N}$, minimality in $\mathbb{S}^4$ requires both its scalar mean curvature within $\mathcal{N}$ and the component of its mean-curvature vector within $\mathbb{S}^4$ in the direction $n$ to vanish. We compute the latter component and the normal connection in terms of the shape operator $S$ of $\mathcal{N}\subset\mathbb{S}^4$ and the chosen unit normal of $M$ within $\mathcal{N}$, and specialize
these formulas to height graphs over the central torus. We then represent the normal linearization of the mean curvature map as a $2\times2$ matrix of scalar differential operators, acting on the coefficients of a normal variation in these two directions. Its off-diagonal entries describe how a variation in either direction affects the mean-curvature component in the other direction. Finally, we construct convenient tubular coordinates around $\mathcal{N}$ in $\mathbb{S}^4$.

Let $M\subset\mathcal{N}$ be a smooth embedded two-sided surface in $\mathcal{N}$, and let
%
\begin{align*}
    \jmath:M\hookrightarrow\mathcal{N},
    \qquad
    \iota:M\hookrightarrow\mathbb{S}^4
\end{align*}
%
denote its inclusions in $\mathcal N$ and $\mathbb S^4$, respectively. Equip $M$ with the induced metric
$g_M:=\jmath^*g_{\mathcal{N}}=\iota^*g_{\mathbb{S}^4}$, and fix a smooth global unit normal field $\nu$ to $M$
within $\mathcal{N}$. No compactness or minimality assumption is imposed on $M$ in this subsection. We write $\nabla^M:=\nabla^{g_M}$ and $\Delta_M:=\Delta_{g_M}$, and denote the restriction of $n$ to $M$ again by $n$. Then $(\nu,n)$ is a global orthonormal frame of the normal bundle $T^\perp_\iota M$.

Using the conventions above, the mean-curvature vectors of the two inclusions decompose as
%
\begin{align*}
    \mathbf{H}_\jmath&=H_{\jmath,\nu}\nu,\\
    \mathbf{H}_\iota&=H_{\iota,\nu}\nu+H_{\iota,n}n.
\end{align*}
%
The equality $h_\iota^\nu=h_\jmath^\nu$, verified in the proof of Lemma \ref{lem:s4-normal-identities} below, implies $H_{\iota,\nu}=H_{\jmath,\nu}$. Thus minimality in $\mathcal{N}$ requires $H_{\iota,\nu}=0$, while minimality in $\mathbb{S}^4$ requires both $H_{\iota,\nu}=0$ and $H_{\iota,n}=0$.

Inner products of vectors are taken with respect to $g_{\mathbb{S}^4}$; tensor norms and contractions use the corresponding induced metrics.

For the normal connection of $\iota$, we abbreviate $\nabla^\perp:=\nabla^{\perp,\iota}$. In the frame $(\nu,n)$, the connection $\nabla^{\perp}$ is determined by the one-form $\omega\in\Gamma(T^*M)$ defined by
%
\begin{align}
    \label{eq:s4-omega-definition}
    \omega(X):=\langle\nabla^\perp_X\nu,n\rangle
    =\langle\nabla^{\mathbb{S}^4}_X\nu,n\rangle
    \qquad\forall\,X\in\Gamma(TM).
\end{align}
%
Indeed, differentiating
$\langle\nu,\nu\rangle=\langle n,n\rangle=1$ and
$\langle\nu,n\rangle=0$ gives
%
\begin{align}
    \label{eq:s4-normal-connection-frame}
    \nabla^\perp_X\nu=\omega(X)n,
    \qquad
    \nabla^\perp_Xn=-\omega(X)\nu
    \qquad\forall\,X\in\Gamma(TM).
\end{align}
%
Thus $\omega$ records the coupling of the two normal directions.
The following lemma expresses this one-form and the second
fundamental form in the direction $n$ in terms of $S$. To state
the estimates, write
%
\begin{align*}
    \nu_h(x):=\sum_{a=1}^2
    \langle\nu(x),E_a(x)\rangle E_a(x)
    \qquad\forall\,x\in M
\end{align*}
%
for the component of $\nu$ in the plane spanned by $E_1,E_2$.

\begin{lemma}
    \label{lem:s4-normal-identities}
    For every $X,Y\in\Gamma(TM)$, one has
    %
    \begin{align}
        \label{eq:s4-normal-identities}
        \begin{split}
            h_\iota^n(X,Y)&=\langle S(X),Y\rangle,\\
            H_{\iota,n}&=-\langle S(\nu),\nu\rangle,\\
            \omega(X)&=\langle S(X),\nu\rangle.
        \end{split}
    \end{align}
    %
    Moreover,
    %
    \begin{align}
        \label{eq:s4-omega-derivative}
        \begin{split}
            (\nabla^M_X\omega)(Y)&=\langle(\nabla^{\mathcal{N}}_XS)(Y),\nu\rangle+h_\iota^\nu(X,Y)\langle S(\nu),\nu\rangle-\langle S(Y),S_\iota^\nu (X)\rangle
            \qquad\forall\,X,Y\in\Gamma(TM).
        \end{split}
    \end{align}
    %
    There exists $C=C(z_0)>0$ such that, for every
    $(s,t,z)\in\mathbb{T}\times(-z_0,z_0)$ with
    $x:=\mathcal{F}(s,t,z)\in M$,
    %
    \begin{align}
        \label{eq:s4-normal-bounds}
        \begin{aligned}
            |(h_\iota^n)_x|&\leq C,
            &\qquad|H_{\iota,n}(x)|&\leq C|\nu_h(x)|^2,\\
            |\omega_x|&\leq C|\nu_h(x)|,
            &\qquad|(\nabla^M\omega)_x|
            &\leq C\bigl(|z|+|(S_\iota^\nu)_x|\bigr).
        \end{aligned}
    \end{align}
    %
    The norms of $h_\iota^n$, $\omega$, $\nabla^M\omega$ and $S_\iota^\nu$ are induced by $g_M$, while $|\nu_h|$ is computed using $g_{\mathcal{N}}$.
\end{lemma}

\begin{proof}
    For every $X,Y\in\Gamma(TM)$, differentiating
    $\langle Y,n\rangle=0$ gives
    %
    \begin{align*}
        h_\iota^n(X,Y)=\langle\nabla^{\mathbb{S}^4}_XY,n\rangle=-\langle Y,\nabla^{\mathbb{S}^4}_Xn\rangle
        =\langle S(X),Y\rangle.
    \end{align*}
    %
    Let $(e_1,e_2)$ be a local $g_M$-orthonormal frame of $TM$.
    Since $(e_1,e_2,\nu)$ is an orthonormal frame of
    $T\mathcal{N}|_M$ and $\operatorname{tr}S=0$, taking the
    trace yields
    %
    \begin{align*}
        H_{\iota,n}
        =\sum_{i=1}^2\langle S(e_i),e_i\rangle
        =\operatorname{tr}S-\langle S(\nu),\nu\rangle
        =-\langle S(\nu),\nu\rangle.
    \end{align*}
    %
    Likewise, differentiation of $\langle\nu,n\rangle=0$
    gives
    %
    \begin{align*}
        \omega(X)
        =-\langle\nu,\nabla^{\mathbb{S}^4}_Xn\rangle
        =\langle S(X),\nu\rangle
        \qquad\forall\,X\in\Gamma(TM).
    \end{align*}
    %
    This proves \eqref{eq:s4-normal-identities}.

    To differentiate $\omega$, we also need the geometry of
    $M$ within $\mathcal{N}$. The difference between
    $\nabla^{\mathbb{S}^4}_XY$ and $\nabla^{\mathcal{N}}_XY$
    is a multiple of $n$, which is orthogonal to $\nu$.
    Hence
    %
    \begin{align*}
        h_\iota^\nu(X,Y)
        =\langle\nabla^{\mathcal{N}}_XY,\nu\rangle
        =h_\jmath^\nu(X,Y)
        \qquad\forall\,X,Y\in\Gamma(TM).
    \end{align*}
    %
    Since both inclusions induce the metric $g_M$, it
    follows that $S_\iota^\nu=S_\jmath^\nu$ as well.
    The tangential and normal decompositions in $\mathcal{N}$
    therefore give
    %
    \begin{align*}
        \nabla^{\mathcal{N}}_XY
        &=\nabla^M_XY+h_\jmath^\nu(X,Y)\nu,\\
        \nabla^{\mathcal{N}}_X\nu
        &=-S_\jmath^\nu(X)
        \qquad\forall\,X,Y\in\Gamma(TM).
    \end{align*}
    %
    Substituting these identities into the covariant
    derivative of $\omega$ gives
    %
    \begin{align*}
        (\nabla^M_X\omega)(Y)
        &=X\bigl(\langle S(Y),\nu\rangle\bigr)
        -\langle S(\nabla^M_XY),\nu\rangle\\
        &=\langle(\nabla^{\mathcal{N}}_XS)(Y),\nu\rangle
        +\langle S(\nabla^{\mathcal{N}}_XY
        -\nabla^M_XY),\nu\rangle
        +\langle S(Y),\nabla^{\mathcal{N}}_X\nu\rangle\\
        &=\langle(\nabla^{\mathcal{N}}_XS)(Y),\nu\rangle
        +h_\iota^\nu(X,Y)\langle S(\nu),\nu\rangle
        -\langle S(Y),S_\iota^\nu(X)\rangle
    \end{align*}
    %
    for every $X,Y\in\Gamma(TM)$. This proves
    \eqref{eq:s4-omega-derivative}.

    We now deduce the estimates. Fix
    $(s,t,z)\in\mathbb{T}\times(-z_0,z_0)$ with
    $x:=\mathcal{F}(s,t,z)\in M$.
    The self-adjointness of $S$ and the identity $S(E_3)=0$
    imply
    %
    \begin{align*}
        S_x(\nu(x))&=S_x(\nu_h(x)),\\
        \langle S_x(\nu(x)),\nu(x)\rangle
        &=\langle S_x(\nu_h(x)),\nu_h(x)\rangle.
    \end{align*}
    %
    Therefore \eqref{eq:s4-normal-identities} gives
    %
    \begin{align*}
        |(h_\iota^n)_x|\leq |S_x|,
        \qquad
        |H_{\iota,n}(x)|\leq |S_x|\,|\nu_h(x)|^2,
        \qquad
        |\omega_x|\leq |S_x|\,|\nu_h(x)|.
    \end{align*}
    %
    Taking tensor norms in \eqref{eq:s4-omega-derivative}
    and using $|(h_\iota^\nu)_x|=|(S_\iota^\nu)_x|$ also gives
    %
    \begin{align*}
        |(\nabla^M\omega)_x|
        \leq |(\nabla^{\mathcal{N}}S)_x|
        +2|S_x|\,|(S_\iota^\nu)_x|.
    \end{align*}
    %
    The bounds \eqref{eq:s4-normal-bounds} follow from
    \eqref{eq:s4-S-derivative}.
\end{proof}
%
\begin{remark}
    For a height graph, these geometric identities give explicit
    expressions in terms of its defining function. Let
    $U\subset\mathbb{T}$ be open and let
    $w:U\to(-z_0,z_0)$ be smooth. Define
    $F_w:U\to\mathcal{N}$ and its image $M_w$ by
    %
    \begin{align*}
        F_w(s,t):=\mathcal{F}(s,t,w(s,t))
        \qquad\forall\,(s,t)\in U,
        \qquad M_w:=F_w(U).
    \end{align*}
    %
    On $U$, set $w_s:=\partial_sw$, $w_t:=\partial_tw$,
    $Dw:=(w_s,w_t)$, $c:=\cos w$ and $Q:=c^2+|Dw|^2$,
    where $|Dw|^2=w_s^2+w_t^2$. Apply the preceding notation
    to $M_w$, choosing its unit normal $\nu$ within $\mathcal{N}$
    so that $\langle\nu,E_3\rangle>0$. Then the following
    identities hold on $U$:
    %
    \begin{align}
        \label{eq:s4-height-graph-identities}
        \begin{split}
            \nu\circ F_w
            &=\frac{-w_s(E_1\circ F_w)-w_t(E_2\circ F_w)
            +c(E_3\circ F_w)}{\sqrt Q},\\
            H_{\iota,n}\circ F_w&=\frac{w_s^2-w_t^2}{cQ},\\
            F_w^*\omega&=\frac{w_s\,ds-w_t\,dt}{\sqrt Q}.
        \end{split}
    \end{align}
    %
    To verify these formulas, fix $(s,t)\in U$ and put
    $x:=F_w(s,t)$. The chain rule gives
    %
    \begin{align*}
        d(F_w)_{(s,t)}(\partial_s)
        &=c(s,t)E_1(x)+w_s(s,t)E_3(x),\\
        d(F_w)_{(s,t)}(\partial_t)
        &=c(s,t)E_2(x)+w_t(s,t)E_3(x).
    \end{align*}
    %
    The numerator in the first line of
    \eqref{eq:s4-height-graph-identities} is orthogonal to
    both tangent vectors and has squared norm $Q$. Since
    $c>0$, its $E_3$-component is positive. This proves the
    formula for $\nu$. Applying \eqref{eq:s4-S} at height
    $z=w(s,t)$ then gives
    %
    \begin{align*}
        S_x(\nu(x))
        &=\frac{w_s(s,t)E_1(x)-w_t(s,t)E_2(x)}
        {c(s,t)\sqrt{Q(s,t)}},\\
        H_{\iota,n}(x)
        &=-\langle S_x(\nu(x)),\nu(x)\rangle
        =\frac{w_s(s,t)^2-w_t(s,t)^2}{c(s,t)Q(s,t)}.
    \end{align*}
    %
    Finally, using the definition of pullback and
    \eqref{eq:s4-normal-identities}, we obtain
    %
    \begin{align*}
        (F_w^*\omega)_{(s,t)}(\partial_s)
        &=\left\langle
        S_x\bigl(d(F_w)_{(s,t)}(\partial_s)\bigr),\nu(x)
        \right\rangle=\langle-E_1(x),\nu(x)\rangle
        =\frac{w_s(s,t)}{\sqrt{Q(s,t)}},\\
        (F_w^*\omega)_{(s,t)}(\partial_t)
        &=\left\langle
        S_x\bigl(d(F_w)_{(s,t)}(\partial_t)\bigr),\nu(x)
        \right\rangle=\langle E_2(x),\nu(x)\rangle
        =-\frac{w_t(s,t)}{\sqrt{Q(s,t)}}.
    \end{align*}
    %
    These are the coefficients of the one-form in the last
    line of \eqref{eq:s4-height-graph-identities}.
\end{remark}
%
We next consider normal variations of a general surface
$M\subset\mathcal{N}$ as above. The geometric identities
already established will allow us to express the
linearization of its mean-curvature vector in the frame
$(\nu,n)$.
For $V\in\Gamma(T^\perp_\iota M)$, fix an open set
$U\subset M$ whose closure in $M$ is compact.
There exists $r_U>0$ such that the maps
%
\begin{align*}
    G_r:U&\rightarrow\mathbb{S}^4,\\
    x&\mapsto\exp_{\iota(x)}^{\mathbb{S}^4}(rV(x))
\end{align*}
%
are smooth immersions for every $r\in(-r_U,r_U)$. Indeed, $G_0=\iota|_U$,
and the differential of $G_r$ depends smoothly on
$(x,r)$. Since $\iota$ is an immersion, compactness
of $\overline U$ ensures that this differential
remains injective for all $x\in U$ when $|r|$
is sufficiently small. Moreover,
%
\begin{align*}
    \left.\frac{\partial}{\partial r}G_r(x)\right|_{r=0}
    =V(x)
    \qquad\forall\,x\in U.
\end{align*}
%
Thus these maps define a normal variation of $\iota$
near every point of $M$, although the interval for $r$
may depend on $U$. The local variations agree on
overlapping domains whenever both are defined, so
their linearizations at $r=0$ define an operator
on all of $M$. If $M$ is compact, we may take $U=M$.
% We use this local interpretation when $M$ is not compact. 
Define the \textit{normal linearization of the mean curvature map} at $\iota$ to be the operator
$J:\Gamma(T^\perp_\iota M)\to\Gamma(T^\perp_\iota M)$ by
%
\begin{align*}
    (JV)(x):=\left(
    \left.\frac{D}{dr}\right|_{r=0}\mathbf{H}_{G_r}(x)
    \right)^\perp
    \qquad\forall\,x\in M,
\end{align*}
%
where $D/dr$ denotes covariant differentiation in
$\mathbb{S}^4$ along the geodesic $r\mapsto G_r(x)$ and the projection
is onto $(T^\perp_\iota M)_x$.

The general normal variation formula expresses $J$ in
terms of the Laplacian associated with the normal
connection, given by
%
\begin{align*}
    \Delta^\perp V
    :=\sum_{i=1}^2\left(
    \nabla^\perp_{e_i}\nabla^\perp_{e_i}V
    -\nabla^\perp_{\nabla^M_{e_i}e_i}V\right)
    \qquad\forall\,V\in\Gamma(T^\perp_\iota M),
\end{align*}
%
where $(e_1,e_2)$ is any local $g_M$-orthonormal frame
of $TM$. For a surface in the unit round sphere, this
formula reads
%
\begin{align}
    \label{eq:s4-normal-variation-formula}
    JV=\Delta^\perp V+2V+
    \sum_{i,j=1}^2
    \langle A_\iota(e_i,e_j),V\rangle A_\iota(e_i,e_j)
    \qquad\forall\,V\in\Gamma(T^\perp_\iota M);
\end{align}
%
see \cite[Lemma~3.1]{LagaceLynchCanonicalFoliation}.
This local identity applies without assuming that $M$
is minimal. The term $2V$ is the ambient curvature
contribution, with one copy of $V$ for each orthonormal
tangent direction.

In the next lemma, we compute the scalar operators representing $J$.
In particular, we compare its component in the direction
$\nu$ with the scalar mean-curvature linearization for
$M\subset\mathcal{N}$.

\begin{lemma}
    \label{lem:s4-normal-linearization}
    Define $L_{M/\mathcal{N}},B,B^*:C^\infty(M)\to
    C^\infty(M)$ and $q\in C^\infty(M)$ by
    %
    \begin{align}
        L_{M/\mathcal{N}}u
        &:=\Delta_Mu+
        \bigl(|S_\jmath^\nu|^2+
        \operatorname{Ric}_{g_{\mathcal{N}}}(\nu,\nu)
        \bigr)u,
        \label{eq:s4-scalar-Jacobi}\\
        q&:=2+|S_\iota^n|^2-|\omega|^2,\nonumber\\
        Bb&:=-2\langle\omega,db\rangle
        -(\operatorname{div}_M\omega)b
        +\langle h_\iota^\nu,h_\iota^n\rangle b,\nonumber\\
        B^*u&:=2\langle\omega,du\rangle
        +(\operatorname{div}_M\omega)u
        +\langle h_\iota^\nu,h_\iota^n\rangle u\nonumber
    \end{align}
    %
    for every $u,b\in C^\infty(M)$. Here
    %
    \begin{align*}
        \operatorname{div}_M\omega
        &:=\sum_{i=1}^2(\nabla^M_{e_i}\omega)(e_i),\\
        \langle h_\iota^\nu,h_\iota^n\rangle
        &:=\sum_{i,j=1}^2h_\iota^\nu(e_i,e_j)h_\iota^n(e_i,e_j),
    \end{align*}
    %
    where $(e_1,e_2)$ is any local $g_M$-orthonormal
    frame of $TM$.
    Then, for every $u,b\in C^\infty(M)$,
    %
    \begin{align}
        \label{eq:s4-J-matrix}
        \begin{pmatrix}
            \langle J(u\nu+bn),\nu\rangle\\
            \langle J(u\nu+bn),n\rangle
        \end{pmatrix}
        =
        \begin{pmatrix}
            L_{M/\mathcal{N}}+H_{\iota,n}^2&B\\
            B^*&\Delta_M+q
        \end{pmatrix}
        \begin{pmatrix}u\\b\end{pmatrix},
    \end{align}
    %
    where $H_{\iota,n}^2$ and $q$ act by multiplication.
    If $M$ is closed, then $B^*$ is the formal
    $L^2(M,d\operatorname{vol}_M)$ adjoint of $B$, where $d\operatorname{vol}_M$
    is the area measure induced by $g_M$.
\end{lemma}

\begin{proof}
    We compute the two terms in
    \eqref{eq:s4-normal-variation-formula} that depend
    on the geometry of $M$. Fix $u,b\in C^\infty(M)$
    and put $V:=u\nu+bn$. The connection identities
    \eqref{eq:s4-normal-connection-frame} give
    %
    \begin{align*}
        \nabla^\perp_XV
        =\bigl(Xu-b\omega(X)\bigr)\nu
        +\bigl(Xb+u\omega(X)\bigr)n
        \qquad\forall\,X\in\Gamma(TM).
    \end{align*}
    %
    To compute $\Delta^\perp V$ at $x\in M$, choose
    a local $g_M$-orthonormal frame $(e_1,e_2)$ with
    $(\nabla^M_{e_i}e_j)(x)=0$ for
    $i,j\in\{1,2\}$. Set $\omega_i:=\omega(e_i)$
    for $i\in\{1,2\}$ on the domain of this frame.
    Applying the preceding identity twice gives
    %
    \begin{align*}
        \nabla^\perp_{e_i}\nabla^\perp_{e_i}V
        &={}
        \bigl(e_i(e_iu-b\omega_i)
        -(e_ib+u\omega_i)\omega_i\bigr)\nu\\
        &\quad+
        \bigl(e_i(e_ib+u\omega_i)
        +(e_iu-b\omega_i)\omega_i\bigr)n
        \qquad\forall\,i\in\{1,2\}.
    \end{align*}
    %
    Expanding and summing at $x$ yields
    %
    \begin{align*}
        \Delta^\perp V
        &={}
        \bigl(\Delta_Mu-2\langle\omega,db\rangle
        -b\operatorname{div}_M\omega
        -u|\omega|^2\bigr)\nu\\
        &\quad+
        \bigl(\Delta_Mb+2\langle\omega,du\rangle
        +u\operatorname{div}_M\omega
        -b|\omega|^2\bigr)n.
    \end{align*}
    %
    Since $x$ was arbitrary, this identity holds on $M$.

    For the second-fundamental-form term, the decomposition
    of $A_\iota$ in the orthonormal frame $(\nu,n)$ gives
    %
    \begin{align*}
        A_\iota(e_i,e_j)
        =h_\iota^\nu(e_i,e_j)\nu+h_\iota^n(e_i,e_j)n
        \qquad\forall\,i,j\in\{1,2\}.
    \end{align*}
    %
    Consequently,
    %
    \begin{align*}
        \sum_{i,j=1}^2
        \langle A_\iota(e_i,e_j),V\rangle A_\iota(e_i,e_j)
        &=
        \bigl(|h_\iota^\nu|^2u+\langle h_\iota^\nu,h_\iota^n\rangle b\bigr)\nu+
        \bigl(\langle h_\iota^\nu,h_\iota^n\rangle u+|h_\iota^n|^2b\bigr)n.
    \end{align*}
    %
    Combining these computations with
    \eqref{eq:s4-normal-variation-formula} and using
    $|h_\iota^\eta|=|S_\iota^\eta|$ for $\eta\in\{\nu,n\}$,
    we obtain
    %
    \begin{align*}
        \langle JV,\nu\rangle
        &=\Delta_Mu+
        \bigl(2+|S_\iota^\nu|^2-|\omega|^2\bigr)u+Bb,\\
        \langle JV,n\rangle
        &=B^*u+(\Delta_M+q)b.
    \end{align*}
    %
    It remains to express the first diagonal entry in
    terms of $L_{M/\mathcal{N}}$. The Gauss equation for
    $\mathcal{N}\subset\mathbb{S}^4$, together with
    $\operatorname{tr}S=0$ and $|\nu|=1$, gives
    %
    \begin{align*}
        \operatorname{Ric}_{g_{\mathcal{N}}}(\nu,\nu)
        =2+(\operatorname{tr}S)\langle S(\nu),\nu\rangle
        -|S(\nu)|^2
        =2-|S(\nu)|^2.
    \end{align*}
    %
    By \eqref{eq:s4-normal-identities} and the
    self-adjointness of $S$, the components of $S(\nu)$
    in the orthonormal frame $(e_1,e_2,\nu)$ are
    %
    \begin{align*}
        S(\nu)=\sum_{i=1}^2\omega(e_i)e_i-H_{\iota,n}\nu,
        \qquad
        |S(\nu)|^2=|\omega|^2+H_{\iota,n}^2.
    \end{align*}
    %
    Using $S_\iota^\nu=S_\jmath^\nu$, we conclude that
    %
    \begin{align*}
        2+|S_\iota^\nu|^2-|\omega|^2
        =|S_\jmath^\nu|^2+
        \operatorname{Ric}_{g_{\mathcal{N}}}(\nu,\nu)
        +H_{\iota,n}^2,
    \end{align*}
    %
    which proves \eqref{eq:s4-J-matrix}.

    Finally, if $M$ is closed, integration by parts
    gives, for every $u,b\in C^\infty(M)$,
    %
    \begin{align*}
        \int_M u\,Bb\,d\operatorname{vol}_M
        &=-2\int_Mu\langle\omega,db\rangle\,d\operatorname{vol}_M
        -\int_Mub\operatorname{div}_M\omega\,d\operatorname{vol}_M\\
        &\quad+\int_Mub\langle h_\iota^\nu,h_\iota^n\rangle\,d\operatorname{vol}_M\\
        &=\int_Mb\bigl(
        2\langle\omega,du\rangle
        +u\operatorname{div}_M\omega
        +u\langle h_\iota^\nu,h_\iota^n\rangle\bigr)\,d\operatorname{vol}_M\\
        &=\int_M b\,B^*u\,d\operatorname{vol}_M.
    \end{align*}
    %
    This proves the adjoint assertion.
\end{proof}
%
We finally construct an explicit parametrization of a
neighbourhood of $\mathcal{N}$ in $\mathbb{S}^4$ by following the geodesics normal to $\mathcal{N}$. The following lemma provides a uniform range of the normal parameter on which this parametrization is an embedding.
%
\begin{lemma}
    \label{lem:s4-tubular}
    Choose $b_0\in(0,\pi/2)$ such that
    $\tan b_0<\cos z_0$. Define
    $\Psi:\mathbb{T}\times(-z_0,z_0)\times(-b_0,b_0)
    \to\mathbb{S}^4$ by
    %
    \begin{align}
        \label{eq:s4-tubular-map}
        \Psi(s,t,z,b)
        :=\cos b\,\mathcal{F}(s,t,z)+\sin b\,N(s,t)
    \end{align}
    %
    for every $(s,t,z,b)$ in its domain.
    Then $\Psi$ is a smooth embedding, and its image is
    an open neighbourhood of $\mathcal{N}$ in $\mathbb{S}^4$.
    Moreover,
    %
    \begin{align*}
        \Psi(s,t,z,b)
        =\exp^{\mathbb{S}^4}_{\mathcal{F}(s,t,z)}
        \bigl(b\,n(\mathcal{F}(s,t,z))\bigr)
    \end{align*}
    %
    for every $(s,t,z,b)$ in its domain.
\end{lemma}
\begin{proof}
    The vectors $\mathcal F(s,t,z)$ and $N(s,t)$ are
    orthonormal in $\mathbb R^5$. Hence the formula
    for $\Psi$ defines a smooth map into $\mathbb S^4$.
    Written in the coordinates of $\mathbb C^2\oplus\mathbb R$,
    this map has the form
    %
    \begin{align*}
        \Psi(s,t,z,b)
        &=\bigl(r_1e^{i\sqrt2s},r_2e^{i\sqrt2t},
        \cos b\sin z\bigr),\\
        r_1&:=\frac{\cos b\cos z+\sin b}{\sqrt2},
        \qquad
        r_2:=\frac{\cos b\cos z-\sin b}{\sqrt2},
    \end{align*}
    %
    for every $(s,t,z,b)$ in the domain.
    Since $|z|<z_0$ and $|b|<b_0<\pi/2$, the choice of
    $b_0$ implies
    %
    \begin{align*}
        r_1,r_2
        \geq\frac{\cos b}{\sqrt2}
        \bigl(\cos z-|\tan b|\bigr)
        >\frac{\cos b}{\sqrt2}
        \bigl(\cos z_0-\tan b_0\bigr)>0.
    \end{align*}
    %
    Thus, if $(w_1,w_2,x_5)=\Psi(s,t,z,b)$, then
    $r_1=|w_1|$ and $r_2=|w_2|$. The phase factors
    $w_1/|w_1|$ and $w_2/|w_2|$ recover $s$ and $t$
    uniquely as elements of $\mathbb R/\ell\mathbb Z$:
    the map $s\mapsto e^{i\sqrt2s}$ is a diffeomorphism
    from this quotient circle onto the unit circle in
    $\mathbb C$. The other two parameters are recovered by
    %
    \begin{align*}
        b&=\arcsin\left(\frac{|w_1|-|w_2|}{\sqrt2}\right),\\
        z&=\arcsin\left(\frac{x_5}{\cos b}\right),
    \end{align*}
    %
    using the branch of $\arcsin$ with values in
    $(-\pi/2,\pi/2)$.
    These formulas recover the original parameters because
    $\sin b=(r_1-r_2)/\sqrt2$ and $x_5=\cos b\sin z$,
    while both $b$ and $z$ belong to $(-\pi/2,\pi/2)$.

    All the recovery maps are smooth on a neighbourhood
    of each point of the image: neither complex coordinate
    vanishes, $\cos b>0$, and the arguments of $\arcsin$
    lie strictly between $-1$ and $1$.
    Their composition with $\Psi$ is the identity on its
    domain. Consequently, $d\Psi$ is injective everywhere,
    and the recovered parameters give a continuous smooth
    inverse on the image. This proves that $\Psi$ is an
    embedding. Its domain and target both have dimension
    four, so its image is open by the inverse function
    theorem. The image contains $\mathcal N$, since
    $\Psi(s,t,z,0)=\mathcal F(s,t,z)$.

    Finally, for each $(s,t,z)\in\mathbb T\times(-z_0,z_0)$,
    the curve
    $b\mapsto\cos b\,\mathcal F(s,t,z)+\sin b\,N(s,t)$
    is the unit-speed great circle with initial point
    $\mathcal F(s,t,z)$ and initial velocity
    $N(s,t)=n(\mathcal F(s,t,z))$. This is precisely the
    defining geodesic of the spherical exponential map
    in the last assertion.
\end{proof}
%
\subsection{Symmetries of the construction}
%
We will construct the initial surfaces inside $\mathcal N$ by
joining two sheets on opposite sides of the central torus with
small necks. Each neck will be parametrized using a Euclidean catenoid.
The neck positions will be indexed by a square lattice
$L_m\subset\mathbb T$, defined below: for each $p\in L_m$,
we will insert one neck near the point $f(p)$ of the
central torus $f(\mathbb T)$. We will impose the same finite group of
symmetries on the initial surfaces and on the maps used to
deform them in $\mathbb S^4$.

The symmetry conditions have two uses in the later linear
analysis. Invariance under lattice translations removes the
kernel of $L_{\mathbb T}=\Delta_{\mathbb T}+2$ on the functions
used to construct the two sheets. Near each neck, the other
symmetries impose different conditions on the two coefficients
of a normal variation. The coefficient in the direction $n$
has no constant angular term and satisfies an angular
Poincar\'e inequality with a constant independent of the neck
size. This inequality is used to estimate the correction in
the direction $n$, uniformly as the neck radii tend to zero.
The coefficient in the normal direction
within $\mathcal N$ can have a constant angular term; that
term must be treated in the analysis of deformations within
$\mathcal N$. We establish these facts below and verify that
the normal linearization from the preceding subsection
preserves the selected symmetry conditions.

Fix an integer $m\geq2$, and define
%
\begin{align}
    \label{eq:s4-lattice}
    L_m:=
    \left\{\left(\frac{j\ell}{m},\frac{k\ell}{m}\right):
    j,k\in\{0,\ldots,m-1\}\right\}\subset\mathbb T.
\end{align}
%
The corresponding points of the central torus are $f(L_m)$.
Define the following isometries of $\mathbb S^4$:
%
\begin{align*}
    R_1(w_1,w_2,x_5)&:=(e^{2\pi i/m}w_1,w_2,x_5),\\
    R_2(w_1,w_2,x_5)&:=(w_1,e^{2\pi i/m}w_2,x_5),\\
    C_1(w_1,w_2,x_5)&:=(\overline{w_1},w_2,x_5),\\
    C_2(w_1,w_2,x_5)&:=(w_1,\overline{w_2},x_5),\\
    \sigma(w_1,w_2,x_5)&:=(w_2,w_1,x_5),\\
    Z(w_1,w_2,x_5)&:=(w_1,w_2,-x_5)
    \qquad\forall\,(w_1,w_2,x_5)\in\mathbb S^4.
\end{align*}
%
Set
%
\begin{align*}
    \Gamma_m^0:=\langle R_1,R_2,C_1,C_2,\sigma\rangle,
    \qquad
    \Gamma_m:=\langle\Gamma_m^0,Z\rangle.
\end{align*}
%
These groups are finite. Indeed, every element permutes the
two complex coordinates, possibly conjugates each of them,
and multiplies them by $m$th roots of unity; an element of
$\Gamma_m$ may also change the sign of the fifth coordinate.
There are only finitely many maps of this form.

The definition of $\mathcal F$ gives the following actions
on its parameters:
%
\begin{align*}
    \begin{array}{c|c}
        \text{isometry}&
        \text{parameters of its value at }\mathcal F(s,t,z)\\ \hline
        R_1&(s+\ell/m,t,z)\\
        R_2&(s,t+\ell/m,z)\\
        C_1&(-s,t,z)\\
        C_2&(s,-t,z)\\
        \sigma&(t,s,z)\\
        Z&(s,t,-z)
    \end{array}
    \qquad
    \forall\,(s,t,z)\in\mathbb T\times(-z_0,z_0).
\end{align*}
%
For example,
%
\begin{align*}
    R_1\bigl(\mathcal F(s,t,z)\bigr)
    &=
    \frac{\cos z}{\sqrt2}
    \bigl(e^{i\sqrt2(s+\ell/m)},e^{i\sqrt2t},0\bigr)
    +\sin z\,e_5=\mathcal F(s+\ell/m,t,z),
\end{align*}
%
because $\sqrt2\ell/m=2\pi/m$. The other entries follow
by conjugating or interchanging the displayed coordinates,
or by changing the sign of $\sin z$.
In particular, every $g\in\Gamma_m$ preserves $\mathcal N$
and the central torus. Write $\widehat g:\mathbb T\to\mathbb T$
for the map induced by $g\in\Gamma_m$ on the first two parameters, i.e.
%
\begin{align*}
    \hat g:=f^{-1}\circ (g|_{f(\mathbb{T})})\circ f. 
\end{align*}
%
Then
%
\begin{align*}
    \pi\circ g=\widehat g\circ\pi,
    \qquad
    \widehat g(L_m)=L_m
    \qquad\forall\,g\in\Gamma_m.
\end{align*}
%
The two translations induced by $R_1,R_2$ are denoted by
%
\begin{align*}
    T_s(s,t):=(s+\ell/m,t),\qquad
    T_t(s,t):=(s,t+\ell/m)
    \qquad\forall\,(s,t)\in\mathbb T.
\end{align*}
%
% Figure~\ref{fig:s4-square-lattice} shows the lattice and
% these translations.

\begin{figure}[htbp]
    \centering
    \begingroup
    \definecolor{symteal}{RGB}{0,112,116}
    \definecolor{symorange}{RGB}{190,102,36}
    \begin{tikzpicture}[
        x=0.87cm,y=0.87cm,
        >=Stealth,
        every node/.style={font=\small},
        lattice point/.style={circle,fill=symteal,inner sep=1.8pt},
        repeated point/.style={circle,draw=black!40,fill=white,
                              line width=0.55pt,inner sep=1.65pt}
    ]
        % A fundamental square, with one elementary lattice cell shaded.
        \fill[symteal!8] (0,0) rectangle (1,1);
        \draw[step=1,black!17,line width=0.35pt] (0,0) grid (5,5);
        \draw[black!65,line width=0.8pt] (0,0) rectangle (5,5);
        \foreach \j in {0,...,4} {
            \foreach \k in {0,...,4} {
                \node[lattice point] at (\j,\k) {};
            }
        }
        % These boundary marks are copies, not additional elements of L_m.
        \foreach \j in {0,...,4} {
            \node[repeated point] at (\j,5) {};
        }
        \foreach \k in {0,...,5} {
            \node[repeated point] at (5,\k) {};
        }
        % Matching arrow marks record the two pairs of identified edges.
        \foreach \y in {0,5} {
            \draw[symteal,line width=1pt,->] (2.20,\y) -- (2.75,\y);
        }
        \foreach \x in {0,5} {
            \draw[symorange,line width=1pt,->] (\x,2.20) -- (\x,2.75);
        }
        % The two generators of the translation subgroup.
        \draw[symorange,line width=1.05pt,->,shorten <=4pt,shorten >=4pt]
            (2,1) -- (3,1) node[midway,above=4pt] {$T_s$};
        \draw[symorange,line width=1.05pt,->,shorten <=4pt,shorten >=4pt]
            (1,2) -- (1,3) node[midway,left=4pt] {$T_t$};
        \node[below left=3pt] at (0,0) {$0$};
        \node[below=4pt] at (5,0) {$\ell$};
        \node[left=4pt] at (0,5) {$\ell$};
        \node[right=9pt] at (5,0) {$s$};
        \node[above=9pt] at (0,5) {$t$};
        \draw[black!65,|-|] (0,-0.57) -- (1,-0.57)
            node[midway,below=3pt] {$\ell/m$};
        % An unobtrusive key keeps the lattice itself uncluttered.
        \node[anchor=west,font=\normalsize] at (6.1,4.60)
            {$\mathbb T=(\mathbb R/\ell\mathbb Z)^2$};
        \node[anchor=west,text=black!65] at (6.1,4.05)
            {The case $m=5$.};
        \node[lattice point] at (6.25,3.10) {};
        \node[anchor=west] at (6.60,3.10) {points of $L_m$};
        \node[repeated point] at (6.25,2.48) {};
        \node[anchor=west] at (6.60,2.48) {periodic copies};
        \node[anchor=west,align=left] at (6.1,1.43)
            {$T_s(s,t)=(s+\ell/m,t)$,\\[4pt]
             $T_t(s,t)=(s,t+\ell/m)$.};
        \node[anchor=west,align=left,text=black!65] at (6.1,0.25)
            {Matching edge arrows indicate\\identified sides.};
    \end{tikzpicture}
    \endgroup
    \caption{The square lattice $L_m$ in a fundamental square of
    $\mathbb T$, with its two elementary translations.
    The top and right edges repeat points already represented on
    the bottom and left edges.}
    \label{fig:s4-square-lattice}
\end{figure}

We next show that the scalar Jacobi operator of the central
torus has trivial kernel on the space of smooth functions
invariant under the lattice translations.
For integers $a,b$, define the complex-valued Fourier
function
%
\begin{align*}
    e_{a,b}(s,t):=
    \exp\left(\frac{2\pi i}{\ell}(as+bt)\right)
    \qquad\forall\,(s,t)\in\mathbb T.
\end{align*}
%
Since $\ell^2=2\pi^2$, differentiation gives
%
\begin{align*}
    \Delta_{\mathbb T}e_{a,b}
    &=-2(a^2+b^2)e_{a,b},\\
    L_{\mathbb T}e_{a,b}
    &=2(1-a^2-b^2)e_{a,b}
    \qquad\forall\,(a,b)\in\mathbb Z^2.
\end{align*}
%
Moreover,
%
\begin{align*}
    e_{a,b}\circ T_s=e^{2\pi ia/m}e_{a,b},
    \qquad
    e_{a,b}\circ T_t=e^{2\pi ib/m}e_{a,b}.
\end{align*}
%
Uniqueness of Fourier coefficients therefore implies that
a smooth function invariant under both translations can
have a nonzero coefficient at $(a,b)$ only if $a,b\in m\mathbb Z$.
For $m\geq2$, no such pair satisfies $a^2+b^2=1$. Hence
%
\begin{align}
    \label{eq:s4-translation-kernel}
    \left.
    \begin{gathered}
        v\in C^\infty(\mathbb T),\qquad
        v\circ T_s=v,\qquad v\circ T_t=v,\\
        L_{\mathbb T}v=0
    \end{gathered}
    \right\}
    \quad\rightarrow\quad v=0.
\end{align}
%
The constant functions are permitted by these symmetries,
but they are not in the kernel, since
$L_{\mathbb T}1=2$. The later construction of the sheet
heights uses the operator $L_{\mathbb T}$ on this
translation-invariant space.

We next determine how the group $\Gamma_m$ acts on the two normal
directions of a surface in $\mathcal N$.
Every $g\in\Gamma_m$ either preserves the two complex
coordinate planes individually or interchanges them. Define
%
\begin{align*}
    \chi(g):=
    \begin{cases}
        +1,&\text{if $g$ preserves the two planes individually},\\
        -1,&\text{if $g$ interchanges the two planes}
    \end{cases}
    \qquad\forall\,g\in\Gamma_m.
\end{align*}
%
Interchanging the planes twice preserves them individually,
so $\chi(gh)=\chi(g)\chi(h)$ for every $g,h\in\Gamma_m$.
Thus $\chi:\Gamma_m\to\{1,-1\}$ is a group homomorphism,
also called a \textit{character}. It equals $-1$ on $\sigma$ and
$+1$ on each of the other five generators.

The value $\chi(g)$ determines whether the differential
of $g$ preserves or reverses the chosen unit normal field $n$:
it sends $n(x)$ to $n(g(x))$ when $\chi(g)=+1$, and to
$-n(g(x))$ when $\chi(g)=-1$, for every $x\in\mathcal N$.
For every isometry $g:\mathbb S^4\to\mathbb S^4$ and
$x\in\mathbb S^4$, let $dg_x:T_x\mathbb S^4\rightarrow T_{g(x)}\mathbb S^4$ denote its differential at $x$. We first compute how the differentials of the six generators transform the normal field $n$. Fix
$(s,t,z)\in\mathbb T\times(-z_0,z_0)$ and set
$x:=\mathcal F(s,t,z)$. Since $n(x)=N(s,t)$,
differentiating the coordinate formulas for the generators gives
%
\begin{align*}
    (dR_1)_x(n(x))
    &=N(s+\ell/m,t)=n(R_1(x)),\\
    (dR_2)_x(n(x))
    &=N(s,t+\ell/m)=n(R_2(x)),\\
    (dC_1)_x(n(x))
    &=N(-s,t)=n(C_1(x)),\\
    (dC_2)_x(n(x))
    &=N(s,-t)=n(C_2(x)),\\
    (d\sigma)_x(n(x))
    &=-N(t,s)=-n(\sigma(x))\\
    (dZ)_x(n(x))
    &=N(s,t)=n(Z(x)).\\
\end{align*}
%
Since $\chi(\sigma)=-1$ and $\chi$ equals $+1$ on the
other five generators, these computations prove
%
\begin{align*}
    dg_x(n(x))=\chi(g)n(g(x))
    \qquad
    \forall\,g\in\{R_1,R_2,C_1,C_2,\sigma,Z\},
    \quad x\in\mathcal N.
\end{align*}
Indeed, $\chi(\sigma)=-1$, whereas $\chi$ equals $+1$
on each of the other five generators. We now show that the same identity holds for every element of $\Gamma_m$. Suppose that it holds for two elements
$g,h\in\Gamma_m$. Applying first $h$ and then $g$, the
chain rule gives
%
\begin{align*}
    d(g\circ h)_x(n(x))
    &=dg_{h(x)}\bigl(dh_x(n(x))\bigr)\\
    &=\chi(h)\,dg_{h(x)}(n(h(x)))\\
    &=\chi(h)\chi(g)\,n(g(h(x)))\\
    &=\chi(g\circ h)\,n((g\circ h)(x))
    \qquad\forall\,x\in\mathcal N,
\end{align*}
%
where the last equality uses that $\chi$ is a group
homomorphism. Thus the identity is preserved under composition.
Moreover, $R_i^{-1}=R_i^{m-1}$ for $i\in\{1,2\}$, and
each of $C_1,C_2,\sigma,Z$ is its own inverse. Every element
of $\Gamma_m$ can therefore be written as a finite composition
of these six generators. We conclude that
%
\begin{align}
    \label{eq:s4-character}
    dg_x(n(x))=\chi(g)n(g(x))
    \qquad\forall\,g\in\Gamma_m,\quad x\in\mathcal N.
\end{align}
%

Let $\Omega\subset\mathcal N$ be an open set with smooth
boundary $M:=\partial_{\mathcal N}\Omega$, and assume
$g(\Omega)=\Omega$ for every $g\in\Gamma_m$. Use the
notation of the preceding subsection for the inclusions
$\jmath:M\hookrightarrow\mathcal N$ and
$\iota:M\hookrightarrow\mathbb S^4$, their induced metric
$g_M$, and their second fundamental forms. Choose $\nu$
to be the outward unit normal to $\Omega$ within
$\mathcal N$. Each $g$ preserves $M$, the metric, and
the side occupied by $\Omega$. Therefore,
%
\begin{align*}
    d g_x\bigl(\nu(x)\bigr)=\nu(g(x))
    \qquad\forall\,g\in\Gamma_m,\,\forall\,x\in M.
\end{align*}
%
For $k\in\mathbb Z_{\geq0}$ and $\beta\in(0,1)$, define
%
\begin{align}
    \label{eq:s4-symmetry-spaces}
    \begin{split}
        C^{k,\beta}_{\mathrm{inv}}(M)
        &:=
        \{u\in C^{k,\beta}(M):
        u\circ g=u\ \text{for every }g\in\Gamma_m\},\\
        C^{k,\beta}_{\chi}(M)
        &:=
        \{b\in C^{k,\beta}(M):
        b\circ g=\chi(g)b\ \text{for every }g\in\Gamma_m\}
    \end{split}
\end{align}
%
to be the spaces of the \textit{invariant} and $\chi$-\textit{equivariant} functions on $M$, respectively. In these compositions, $g$ denotes its restriction to $M$.
The same subscripts on $C^\infty(M)$ impose the same
conditions on smooth functions.
If $V=u\nu+bn$ is a normal vector field along $\iota$, then
%
\begin{align*}
    d g_x\bigl(V(x)\bigr)
    &=u(x)\nu(g(x))+\chi(g)b(x)n(g(x)),\\
    V(g(x))
    &=u(g(x))\nu(g(x))+b(g(x))n(g(x))
\end{align*}
%
for every $g\in\Gamma_m$ and $x\in M$. The two vectors
$\nu(g(x)),n(g(x))$ are linearly independent. Thus the
condition
%
\begin{align*}
    d g_x\bigl(V(x)\bigr)=V(g(x))
    \qquad\forall\,g\in\Gamma_m,\,\forall\,x\in M,
\end{align*}
%
which is the definition of \textit{equivariance} for a vector field $V\in\Gamma(T^\perp_\iota M)$, is
equivalent to $u\circ g=u$ and $b\circ g=\chi(g)b$.

These conditions also explain how the symmetries are
preserved by deformations of the surface.
For $u,b\in C^\infty(M)$, set $V:=u\nu+bn$ and define
%
\begin{align*}
    G_r:M&\longrightarrow\mathbb{S}^4,\\
    G_r(x)&:=\exp_{\iota(x)}^{\mathbb{S}^4}
    \bigl(rV(x)\bigr)
    \qquad\forall\,x\in M,\quad r\in\mathbb{R}.
\end{align*}
%
These maps are well defined for every $r\in\mathbb{R}$
because the round sphere is complete.
Their being immersions requires a restriction on $r$.
Indeed, let $U\subset M$ be an open set whose closure
in $M$ is compact. Since $G_0=\iota$ is an immersion,
continuity of the differential of $G_r$ and compactness
of $\overline U$ give $r_U>0$ such that $G_r|_U$
is an immersion whenever $|r|<r_U$.
If $M$ is compact, this assertion applies with $U=M$.

An isometry maps a geodesic to the geodesic whose initial
point and initial velocity are the images of the original
initial point and initial velocity.
Since $\iota$ is the inclusion and $g$ preserves $M$,
one has $\iota(g(x))=g(\iota(x))$.
Consequently, if $V$ is equivariant, then
%
\begin{align*}
    G_r(g(x))
    &=\exp_{\iota(g(x))}^{\mathbb{S}^4}
      \bigl(rV(g(x))\bigr)\\
    &=\exp_{g(\iota(x))}^{\mathbb{S}^4}
      \bigl(r\,d g_{\iota(x)}(V(x))\bigr)\\
    &=g\left(
      \exp_{\iota(x)}^{\mathbb{S}^4}
      \bigl(rV(x)\bigr)\right)\\
    &=g(G_r(x))
\end{align*}
%
for every $g\in\Gamma_m$, $x\in M$ and $r\in\mathbb{R}$.
Since $g(M)=M$, this gives $g(G_r(M))=G_r(M)$.
Whenever $G_r$ is an embedding, its image is therefore
an embedded surface preserved by $\Gamma_m$.
The same calculation applies to deformations defined
using $\exp^{\mathcal N}$, wherever these maps are
defined, because each $g|_{\mathcal N}$ is an isometry
of $(\mathcal N,g_{\mathcal N})$

We now verify that the linearized operator $J$ preserves invariance and $\chi$-equivariance. Recall that, for a covariant two-tensor $h$
on $M$ and for the previously defined one-form $\omega$,
the pullbacks by $g$ are given by
%
\begin{align*}
    (g^*h)_x(X,Y)&:=h_{g(x)}(d g_x(X),d g_x(Y)),\\
    (g^*\omega)_x(X)&:=\omega_{g(x)}(d g_x(X))
    \qquad\forall\,g\in\Gamma_m,\,\forall\,x\in M,\,\forall\,X,Y\in T_xM.
\end{align*}
%
An ambient isometry preserves the Levi-Civita connection
and maps the tangent and normal spaces of $M$ at $x$
onto those at $g(x)$. Applying these facts to the
definition of $A_\iota$ gives
%
\begin{align*}
    (A_\iota)_{g(x)}(d g_x(X),d g_x(Y))
    =d g_x\bigl((A_\iota)_x(X,Y)\bigr)
    \qquad\forall\,g\in\Gamma_m,\quad x\in M,\quad
    X,Y\in T_xM.
\end{align*}
%
Pairing this identity with $\nu(g(x))=d g_x(\nu(x))$
and $n(g(x))=\chi(g)d g_x(n(x))$ yields
%
\begin{align*}
    (g^*h_\iota^\nu)_x(X,Y)
    &=
    \left\langle d g_x\bigl((A_\iota)_x(X,Y)\bigr),
    d g_x(\nu(x))\right\rangle
    =(h_\iota^\nu)_x(X,Y),\\
    (g^*h_\iota^n)_x(X,Y)
    &=
    \chi(g)\left\langle
    d g_x\bigl((A_\iota)_x(X,Y)\bigr),d g_x(n(x))
    \right\rangle
    =\chi(g)(h_\iota^n)_x(X,Y).
\end{align*}
%
Taking traces in an orthonormal frame, whose image under
$d g_x$ is again orthonormal, gives
%
\begin{align*}
    H_{\iota,\nu}\circ g=H_{\iota,\nu},
    \qquad
    H_{\iota,n}\circ g=\chi(g)H_{\iota,n}
    \qquad\forall\,g\in\Gamma_m.
\end{align*}
%
Similarly, the one-form $\omega$ defined in
\eqref{eq:s4-omega-definition} satisfies
%
\begin{align*}
    (g^*\omega)_x(X)
    &=
    \left\langle
    \left(\nabla^{\mathbb S^4}_{d g_xX}\nu\right)(g(x)),
    n(g(x))\right\rangle\\
    &=
    \chi(g)\left\langle
    d g_x\left(\left(\nabla^{\mathbb S^4}_X\nu\right)(x)\right),
    d g_x(n(x))\right\rangle\\
    &=\chi(g)\omega_x(X)
    \qquad\forall\,g\in\Gamma_m,\quad x\in M,\quad X\in T_xM.
\end{align*}
%
Here a covariant derivative in a tangent vector at a point
is evaluated using any local extension of that vector.

Isometries commute with contraction, covariant
differentiation and the scalar Laplacian. The preceding
identities therefore imply
%
\begin{align*}
    \bigl(\langle h_\iota^\nu,h_\iota^n\rangle\bigr)\circ g
    &=\chi(g)\langle h_\iota^\nu,h_\iota^n\rangle,\\
    (\operatorname{div}_M\omega)\circ g
    &=\operatorname{div}_M(g^*\omega)
    =\chi(g)\operatorname{div}_M\omega,\\
    q\circ g&=q,\\
    (\Delta_Mv)\circ g&=\Delta_M(v\circ g)
    \qquad\forall\,g\in\Gamma_m,\quad v\in C^\infty(M).
\end{align*}
%
For the third identity, recall that
$q=2+|S_\iota^n|^2-|\omega|^2$ and
$|S_\iota^n|=|h_\iota^n|$, so both squared norms are
unchanged by the sign $\chi(g)$.
The function
$|S_\jmath^\nu|^2+
\operatorname{Ric}_{g_{\mathcal N}}(\nu,\nu)$ in
$L_{M/\mathcal N}$ is invariant as well, since
$g|_{\mathcal N}$ is an isometry and preserves $\nu$.
Hence both diagonal operators in \eqref{eq:s4-J-matrix}
commute with composition by $g$; here
$H_{\iota,n}^2\circ g=H_{\iota,n}^2$.

For completeness, the transformation of the
off-diagonal operator $B$ follows directly from its
formula in Lemma~\ref{lem:s4-normal-linearization}. For every
$v\in C^\infty(M)$ and $g\in\Gamma_m$,
%
\begin{align*}
    (Bv)\circ g
    &=-2\langle g^*\omega,d(v\circ g)\rangle
    -\bigl((\operatorname{div}_M\omega)\circ g\bigr)(v\circ g)
    +\bigl(\langle h_\iota^\nu,h_\iota^n\rangle\circ g\bigr)
    (v\circ g)\\
    &=\chi(g)\left(
    -2\langle\omega,d(v\circ g)\rangle
    -(\operatorname{div}_M\omega)(v\circ g)
    +\langle h_\iota^\nu,h_\iota^n\rangle(v\circ g)
    \right)\\
    &=\chi(g)B(v\circ g).
\end{align*}
%
The same computation with the signs in the definition
of $B^*$ gives
%
\begin{align*}
    (B^*v)\circ g=\chi(g)B^*(v\circ g)
    \qquad\forall\,v\in C^\infty(M),\,\forall\,g\in\Gamma_m.
\end{align*}
%
In particular, if $b\circ g=\chi(g)b$ and $u\circ g=u$,
then
%
\begin{align*}
    (Bb)\circ g&=\chi(g)B(\chi(g)b)=Bb,\\
    (B^*u)\circ g&=\chi(g)B^*u.
\end{align*}
%
Thus, the off-diagonal operators satisfy
%
\begin{align*}
    B &: C^\infty_\chi(M)\rightarrow
         C^\infty_{\mathrm{inv}}(M),\\
    B^* &: C^\infty_{\mathrm{inv}}(M)\rightarrow
           C^\infty_\chi(M).
\end{align*}
%
The first diagonal entry of $J$ preserves
$C^\infty_{\mathrm{inv}}(M)$, and the second preserves
$C^\infty_\chi(M)$. Consequently, applying $J$ to
$u\nu+bn$ with $u$ invariant and $b$ $\chi$-equivariant
produces a normal field whose $\nu$-coefficient is invariant
and whose $n$-coefficient is $\chi$-equivariant.

We finish by describing the consequences near a neck.
Choose $\tau>0$ and $\xi_*>0$ such that
%
\begin{align}\label{Equation: conditions symmetries}
    \tau\cosh\xi_*<\frac{\ell}{2m},
    \qquad
    \tau\xi_*<z_0.
\end{align}
%
For each $p=(p_s,p_t)\in L_m$, define
%
\begin{align}
    \label{eq:s4-symmetry-neck-map}
    \begin{split}
        F_{\tau,p}:
        (-\xi_*,\xi_*)\times(\mathbb R/2\pi\mathbb Z)
        &\rightarrow\mathcal N,\\
        (\xi,\theta)
        &\mapsto
        \mathcal F\bigl(
        p+\tau\cosh\xi(\cos\theta,\sin\theta),\tau\xi
        \bigr).
    \end{split}
\end{align}
%
The first argument of $\mathcal F$ is computed in the
flat torus $\mathbb T$. The first inequality in \eqref{Equation: conditions symmetries} places it
in an embedded coordinate disk about $p$ containing no
other lattice point, and the second keeps the height
inside $(-z_0,z_0)$. In these coordinates, the map is the usual catenoid parametrization, followed by $\mathcal F$. Its coordinate derivatives before applying $\mathcal F$
are
%
\begin{align*}
    \partial_\xi
    \bigl(p+\tau\cosh\xi(\cos\theta,\sin\theta),\tau\xi\bigr)
    &=
    \bigl(\tau\sinh\xi\cos\theta,
    \tau\sinh\xi\sin\theta,\tau\bigr),\\
    \partial_\theta
    \bigl(p+\tau\cosh\xi(\cos\theta,\sin\theta),\tau\xi\bigr)
    &=
    \bigl(-\tau\cosh\xi\sin\theta,
    \tau\cosh\xi\cos\theta,0\bigr).
\end{align*}
%
These two coordinate derivatives are linearly independent:
the last component of the first is $\tau>0$, whereas
the last component of the second is zero and the second
vector is nonzero. For a point $x\in\mathcal N$, write
%
\begin{align*}
    \mathcal F^{-1}(x)=(q,z),
    \qquad q=\pi(x)\in\mathbb T,\quad z\in(-z_0,z_0).
\end{align*}
%
and recall that $z$ is the signed distance coordinate
from the central torus within $(\mathcal N,g_{\mathcal N})$;
see Remark~\ref{rem:s4-signed-distance}. Therefore, we say that $z$ is the \textit{height} of $x$. For $x=F_{\tau,p}(\xi,\theta)$, the defining formula gives
%
\begin{align*}
    z=\tau\xi,
    \qquad
    q=p+\tau\cosh\xi(\cos\theta,\sin\theta).
\end{align*}
%
In particular, we recover the parameters $(\xi,\theta)$ from $x$ by
%
\begin{align*}
    \xi=\frac z\tau,
    \qquad
    (\cos\theta,\sin\theta)=\frac{q-p}{|q-p|},
\end{align*}
%
where $q-p$ is computed in the flat coordinate disk
about $p$. The denominator is nonzero because
$|q-p|=\tau\cosh\xi\geq\tau>0$.
These formulas give a smooth inverse on the image,
with $\theta$ understood modulo $2\pi$.
Together with the linear independence of the coordinate
derivatives and the fact that $\mathcal F$ is an embedding,
this proves that $F_{\tau,p}$ is an embedding. We call the image of $F_{\tau,p}$ the
\textit{catenoidal model neck} at $p$. Note that a catenoidal model need not be minimal in $(\mathcal N,g_{\mathcal N})$.

For each fixed $\xi\in(-\xi_*,\xi_*)$, varying $\theta$
gives a closed curve at height $z=\tau\xi$.
Its projection under $\pi$ is the circle in the flat
torus centered at $p$ with radius
%
\begin{align*}
    r(\xi):=\tau\cosh\xi.
\end{align*}
%
Since $\cosh\xi\geq1$, with equality only at $\xi=0$,
this projected radius is smallest when $\xi=0$.
We call the corresponding closed curve the
\textit{waist} of the neck. Explicitly,
%
\begin{align*}
    F_{\tau,p}(0,\theta)
    =f\bigl(p+\tau(\cos\theta,\sin\theta)\bigr)
    \qquad\forall\,\theta\in\mathbb R/2\pi\mathbb Z.
\end{align*}
%
Thus the waist lies in the central torus, and its
projection is the circle of radius $\tau$ about $p$.
The portions with $\xi>0$ and $\xi<0$ lie on opposite
sides of the central torus and meet along the waist.

Only group elements fixing $p$ on the base torus
preserve this individual neck. Denote this subgroup by
%
\begin{align*}
    \Gamma_{m,p}
    :=
    \{g\in\Gamma_m:\widehat g(p)=p\}.
\end{align*}
%
For $p=(j\ell/m,k\ell/m)\in L_m$, set
$T_p:=R_1^jR_2^k$. Its induced translation sends $0$
to $p$. Define
%
\begin{align*}
    C_{2,p}:=T_pC_2T_p^{-1},\qquad
    \sigma_p:=T_p\sigma T_p^{-1},\qquad
    Q_p:=\sigma_pC_{2,p}.
\end{align*}
%
These maps and $Z$ belong to $\Gamma_{m,p}$. At $p=0$,
$C_2$ changes $(r\cos\theta,r\sin\theta)$ to
$(r\cos\theta,-r\sin\theta)$, whereas $\sigma$
changes it to $(r\sin\theta,r\cos\theta)$.
Their composition $\sigma C_2$ changes it to
$(-r\sin\theta,r\cos\theta)$.
Conjugation by $T_p$ gives the same actions on the
displacement from any $p\in L_m$.
Finally, $\cosh(-\xi)=\cosh\xi$, so $Z$ changes
$\xi$ to $-\xi$ without changing $\theta$.
We have therefore proved the following table:
%
\begin{align*}
    \begin{array}{c|c|c}
        g&\text{map }\vartheta_g
        \text{ satisfying }g\circ F_{\tau,p}
        =F_{\tau,p}\circ\vartheta_g&\chi(g)\\ \hline
        C_{2,p}&(\xi,\theta)\mapsto(\xi,-\theta)&+1\\
        \sigma_p&(\xi,\theta)\mapsto(\xi,\pi/2-\theta)&-1\\
        Q_p&(\xi,\theta)\mapsto(\xi,\theta+\pi/2)&-1\\
        Z&(\xi,\theta)\mapsto(-\xi,\theta)&+1
    \end{array}
\end{align*}
%
All angular expressions are understood modulo $2\pi$.
The signs are unchanged by conjugation, because
$\chi(T_p)=1$.

Let $b$ be a smooth function on this model neck
satisfying $b\circ g=\chi(g)b$ for $g\in\Gamma_{m,p}$.
For the following calculations, write $b(\xi,\theta)$
for its pullback by $F_{\tau,p}$. The table above gives
%
\begin{align*}
    b(\xi,-\theta)&=b(\xi,\theta),\\
    b(\xi,\pi/2-\theta)&=-b(\xi,\theta),\\
    b(\xi,\theta+\pi/2)&=-b(\xi,\theta),\\
    b(-\xi,\theta)&=b(\xi,\theta)
\end{align*}
%
for every $\xi\in(-\xi_*,\xi_*)$ and
$\theta\in\mathbb R/2\pi\mathbb Z$.
To determine the angular Fourier terms, write
%
\begin{align*}
    b(\xi,\theta)=\sum_{k\in\mathbb Z}
    \widehat b_k(\xi)e^{ik\theta},
    \qquad
    \widehat b_k(\xi):=
    \frac1{2\pi}\int_0^{2\pi}
    b(\xi,\theta)e^{-ik\theta}\,d\theta.
\end{align*}
%
The quarter-turn identity implies
$(e^{ik\pi/2}+1)\widehat b_k(\xi)=0$.
Thus, a nonzero coefficient requires $k\equiv2\pmod4$.
The reflection $\theta\mapsto-\theta$ makes the
coefficients at $k$ and $-k$ equal, removing the sine
terms. The reflection in $\xi$ makes each remaining
coefficient even. Hence
%
\begin{align}
    \label{eq:s4-neck-modes}
    b(\xi,\theta)=
    \sum_{j\geq0}b_j(\xi)\cos((4j+2)\theta),
    \qquad
    b_j(-\xi)=b_j(\xi)
\end{align}
%
for every $(\xi,\theta)$ in the parameter domain.
Here
%
\begin{align*}
    b_j(\xi):=\frac1\pi\int_0^{2\pi}
    b(\xi,\theta)\cos((4j+2)\theta)\,d\theta
    \qquad\forall\,j\in\mathbb Z_{\geq0},\quad
    \xi\in(-\xi_*,\xi_*).
\end{align*}
%
The Fourier series and its derivatives converge on
compact subsets of this domain, since $b$ is smooth.
Orthogonality of the trigonometric functions gives
%
\begin{align*}
    \int_0^{2\pi}|b(\xi,\theta)|^2\,d\theta
    &=\pi\sum_{j\geq0}|b_j(\xi)|^2,\\
    \int_0^{2\pi}|\partial_\theta b(\xi,\theta)|^2\,d\theta
    &=\pi\sum_{j\geq0}(4j+2)^2|b_j(\xi)|^2
\end{align*}
%
for every $\xi\in(-\xi_*,\xi_*)$. Since $(4j+2)^2\geq4$,
%
\begin{align}
    \label{eq:s4-symmetry-angular-poincare}
    \int_0^{2\pi}|\partial_\theta b(\xi,\theta)|^2\,d\theta
    \geq
    4\int_0^{2\pi}|b(\xi,\theta)|^2\,d\theta
    \qquad\forall\,\xi\in(-\xi_*,\xi_*).
\end{align}
%
The constant $4$ depends only on the permitted angular
orders. In particular, it is independent of $\tau$
and $\xi_*$.

The same symmetries determine the conditions on the
parameter rectangle $[0,\xi_*)\times[0,\pi/4]$.
Differentiating $b(\xi,-\theta)=b(\xi,\theta)$ at
$\theta=0$ gives
$-\partial_\theta b(\xi,0)=\partial_\theta b(\xi,0)$.
Evaluating
$b(\xi,\pi/2-\theta)=-b(\xi,\theta)$ at $\theta=\pi/4$
gives $b(\xi,\pi/4)=-b(\xi,\pi/4)$.
Differentiating $b(-\xi,\theta)=b(\xi,\theta)$ at
$\xi=0$ gives the corresponding equality for
$\partial_\xi b$. Therefore
%
\begin{align*}
    \partial_\theta b(\xi,0)=0, \quad b(\xi,\pi/4)=0 \qquad \forall\,\xi\in[0,\xi_*)
\end{align*}
%
and
%
\begin{align*}
    \partial_\xi b(0,\theta)=0 \qquad\forall\,\theta\in[0,\pi/4].
\end{align*}
%
The zero-value condition is a homogeneous Dirichlet
condition. The zero normal-derivative conditions are
homogeneous Neumann conditions. To relate the displayed
coordinate derivatives to normal derivatives, direct
differentiation of $F_{\tau,p}$ gives
%
\begin{align*}
    F_{\tau,p}^*g_{\mathcal N}
    =
    \tau^2\bigl(1+\cos^2(\tau\xi)\sinh^2\xi\bigr)\,d\xi^2
    +
    \tau^2\cos^2(\tau\xi)\cosh^2\xi\,d\theta^2.
\end{align*}
%
This metric is diagonal and positive definite.
Consequently, $\partial_\theta$ is perpendicular to
each angular boundary and $\partial_\xi$ is perpendicular
to the boundary $\xi=0$. These boundary conditions
follow from restrictions of smooth functions with the
stated symmetries. We impose no boundary condition at
the artificial truncation $\xi=\xi_*$.
Figure~\ref{fig:s4-neck-symmetries} illustrates the
three reflection conditions.

For comparison, let $u$ be a smooth function on the
model neck invariant under $\Gamma_{m,p}$, and again use its pullback by
$F_{\tau,p}$. Its quarter-turn identity is
$u(\xi,\theta+\pi/2)=u(\xi,\theta)$.
The Fourier calculation above now permits only
$k\equiv0\pmod4$. The two evenness conditions give
%
\begin{align*}
    u(\xi,\theta)
    =u_0(\xi)+\sum_{j\geq1}u_j(\xi)\cos(4j\theta),
    \qquad
    u_j(-\xi)=u_j(\xi)
    \quad\forall\,j\in\mathbb Z_{\geq0},
\end{align*}
%
where
%
\begin{align*}
    u_0(\xi)&:=\frac1{2\pi}\int_0^{2\pi}u(\xi,\theta)\,d\theta,\\
    u_j(\xi)&:=\frac1\pi\int_0^{2\pi}
    u(\xi,\theta)\cos(4j\theta)\,d\theta
    \qquad\forall\,j\geq1,\quad \xi\in(-\xi_*,\xi_*).
\end{align*}
%
The expansion holds for every $(\xi,\theta)$ in the
parameter domain.
Invariance under $\theta\mapsto\pi/2-\theta$ gives
zero angular derivative at $\theta=\pi/4$, by
differentiating that identity at its fixed point.
The other two reflections are the same as for $b$.
Thus
%
\begin{align*}
    \partial_\theta u(\xi,0)=0, \quad \partial_\theta u(\xi,\pi/4)&=0 \qquad\forall\,\xi\in[0,\xi_*)
\end{align*}
%
and
%
\begin{align*}
    \partial_\xi u(0,\theta)&=0 \qquad\forall\,\theta\in[0,\pi/4].
\end{align*}
%
Both $u$ and $b$ therefore have zero coefficients of
$\cos\theta$ and $\sin\theta$, which are the first
angular harmonics. However, $u_0$ need not vanish.
For example, a nonzero constant function satisfies
all the conditions on $u$ and has zero angular derivative,
so no inequality of the form
\eqref{eq:s4-symmetry-angular-poincare} can hold for all
invariant $u$.
This explains why the later estimates for the correction
in the direction $n$ use angular Poincar\'e inequalities,
whereas the analysis within $\mathcal N$ must also handle
the angular average $u_0$. The model calculation here
concerns a single catenoidal neck; the later analysis
also treats its attachment to the two sheets.

\begin{figure}[htbp]
    \centering
    \begingroup
    \definecolor{symteal}{RGB}{0,112,116}
    \definecolor{symorange}{RGB}{190,102,36}
    \begin{tikzpicture}[
        x=1cm,y=1cm,>=Stealth,
        every node/.style={font=\small},
        axis/.style={black!40,line width=0.55pt},
        symmetry axis/.style={dashed,line width=0.7pt}
    ]
        % Angular picture in the (s,t)-plane centered at a lattice point.
        \begin{scope}[shift={(0,0)}]
            \node[font=\normalsize] at (0,2.75) {Angular action};
            \fill[symteal!11] (0,0) -- (1.7,0)
                arc[start angle=0,end angle=45,radius=1.7] -- cycle;
            \draw[black!45,line width=0.75pt] (0,0) circle[radius=1.7];
            \draw[axis,->] (0,-1.95) -- (0,2.05)
                node[above left=1pt] {$t-p_t$};
            \draw[symmetry axis,symteal] (-2.10,0) -- (2.10,0);
            \draw[symmetry axis,symorange] (-1.45,-1.45) -- (1.56,1.56);
            \node[below=4pt] at (2.02,0) {$s-p_s$};
            \draw[symteal,line width=1.2pt] (0,0) -- (1.7,0);
            \draw[symorange,line width=1.2pt] (0,0) -- (1.2021,1.2021);
            \fill[black!75] (0,0) circle[radius=1.6pt];
            \node[text=symteal,above=5pt] at (1.55,0)
                {$\mathrm N$};
            \node[text=symorange,right=4pt] at (1.22,1.22)
                {$\mathrm D$};
            \node[text=symteal] at (1.02,0.43) {$\mathcal S$};
            \draw[black!60,->] (0.58,0)
                arc[start angle=0,end angle=45,radius=0.58];
            \node[font=\scriptsize] at (0.77,0.18) {$\pi/4$};
            % A positive quarter-turn; its character is -1.
            \draw[black!65,line width=0.85pt,->]
                ({2.14*cos(25)},{2.14*sin(25)})
                arc[start angle=25,end angle=115,radius=2.14];
            \node[fill=white,inner sep=2pt] at (0.85,2.22)
                {$Q_p:\ \chi=-1$};
            \node[anchor=west,text=symteal] at (-2.35,-2.26)
                {$C_{2,p}:\ \theta\mapsto-\theta,\quad\chi=+1$};
            \node[anchor=west,text=symorange] at (-2.35,-2.76)
                {$\sigma_p:\ \theta\mapsto\pi/2-\theta,\quad\chi=-1$};
            \node[anchor=west] at (-2.35,-3.26)
                {$Q_p:\ \theta\mapsto\theta+\pi/2$.};
        \end{scope}
        % A schematic catenoid. Heights are proportional to xi.
        \begin{scope}[shift={(7.3,0)}]
            \node[font=\normalsize] at (0,2.75) {Axial action};
            \fill[symteal!8]
                (0.60,0) -- plot[domain=0:1.65,samples=45] ({0.60*cosh(\x)},\x)
                -- ({-0.60*cosh(1.65)},1.65)
                -- plot[domain=1.65:0,samples=45] ({-0.60*cosh(\x)},\x)
                -- cycle;
            \draw[axis,->] (0,-2.02) -- (0,2.16)
                node[above] {$\xi$};
            \draw[black!65,line width=0.9pt]
                plot[domain=-1.65:1.65,samples=70] ({0.60*cosh(\x)},\x);
            \draw[black!65,line width=0.9pt]
                plot[domain=-1.65:1.65,samples=70] ({-0.60*cosh(\x)},\x);
            \draw[black!40] (0,1.65)
                ellipse[x radius=1.6196,y radius=0.23];
            \draw[black!40] (0,-1.65)
                ellipse[x radius=1.6196,y radius=0.23];
            \draw[symteal,line width=0.85pt] (0,0)
                ellipse[x radius=0.60,y radius=0.11];
            \draw[symteal,symmetry axis] (-1.95,0) -- (1.95,0);
            \node[anchor=east,text=symteal,align=right]
                at (-0.82,0.30) {waist: $\xi=0$\\[-1pt] $\mathrm N$};
            \draw[black!65,line width=0.85pt,<->]
                (1.62,1.05) to[bend left=38] (1.62,-1.05);
            \node[anchor=west,align=left] at (2.10,0)
                {$Z$\\[2pt] $\chi=+1$};
            \node at (0,-2.26) {$Z:\ \xi\mapsto-\xi$};
            \node at (0,-2.76) {$b(-\xi,\theta)=b(\xi,\theta)$};
            \node[text=black!65] at (0,-3.26)
                {$r=\tau\cosh\xi,\qquad z=\tau\xi$.};
        \end{scope}
    \end{tikzpicture}
    \endgroup
    \caption{The neck symmetries and the boundary conditions they impose
    on a $\chi$-equivariant coefficient $b$.
    Left: the sector $\mathcal S=\{0\leq\theta\leq\pi/4\}$
    has Neumann data ($\mathrm N$) on $\theta=0$ and
    Dirichlet data ($\mathrm D$) on $\theta=\pi/4$.
    Right: reflection across the waist makes $b$ even in $\xi$,
    giving Neumann data at $\xi=0$ on the upper half-neck.
    Shading only identifies the chosen sector and half-neck.}
    \label{fig:s4-neck-symmetries}
\end{figure}
%
%
\section{Construction of the equivariant initial surfaces}
\label{sec:s4-scalar-family}

For each sufficiently large integer $m$ and each
$\zeta\in[-1,1]$, we construct a smooth closed embedded
surface $M_\zeta\subset\mathcal N$. It consists of two
sheets on opposite sides of the central torus, joined by
one catenoidal neck near each of the $m^2$ lattice points
in $L_m$. We then establish the estimates needed to deform
this surface within $\mathcal N$. Throughout this section,
scalar mean curvature and normal displacements of a surface
refer to its inclusion into $(\mathcal N,g_{\mathcal N})$.

We first construct and estimate a Green function on a fixed
flat torus. After rescaling, it prescribes the height of the
sheets away from the lattice points. Its logarithmic terms
agree with those of the ends of a catenoid; comparing the
constant terms determines the neck radius $\tau_\zeta$.
We subtract the remaining constant difference and interpolate
between these graph heights and the catenoids. The resulting
surface is denoted by $M_\zeta$, and we prove its embeddedness,
topology, and geometric estimates directly.

The final two subsections define the function spaces and
establish the scalar estimates used later. We identify all
$M_\zeta$ with a fixed reference surface $M_0$ so that the
unknown functions have a fixed domain when $\zeta$ varies.
Using the linear construction of Kapouleas--McGrath, we solve
the scalar linearized equation up to a multiple of one
specified function. We retain that coefficient and estimate
the nonlinear error for every allowed scalar displacement.
Section~\ref{sec:s4-transverse} will add a displacement
normal to $\mathcal N$, and Section~\ref{sec:s4-coupled}
will choose the scalar displacement and $\zeta$ so that
both components of the spherical mean-curvature vector vanish.
No exact minimality is imposed in the present section.

All constants below are independent of $m$ and
$\zeta\in[-1,1]$ after the fixed cutoffs and exponents have
been chosen. Their dependence on a fixed derivative order
or H\"older exponent is allowed. Every assertion for
sufficiently large $m$ uses a threshold independent of
$\zeta$; we may increase that threshold finitely many times.

\subsection{A Green function on the fixed square torus}

We first construct a function with one logarithmic singularity
on a fixed flat torus. Rescaling this function will give a
solution of $L_{\mathbb T}\Phi_m=0$ away from $L_m$, with the
same logarithmic singularity at every lattice point. We need
estimates for its value and derivatives to construct the
initial surfaces. We also prove a lower bound for its Hessian;
this bound will later be used to establish negative Gaussian
curvature between the necks.

Let $\mathbb T_0:=(\mathbb R/2\pi\mathbb Z)^2$, equipped with
the flat metric $dX_1^2+dX_2^2$. Write $dX$ for its area
measure, $D_0$ for its flat covariant derivative, and $\Delta_0$
for its Laplace--Beltrami operator. Thus
$\operatorname{Area}(\mathbb T_0)=4\pi^2$. Let $\delta_0$ be
the Dirac measure of mass one at the origin of $\mathbb T_0$.
For every $a\in[0,1/4]$, define
%
\begin{align}
    \label{eq:s4-cell-Green}
    G_a(X):=\sum_{n\in\mathbb Z^2\smallsetminus\{0\}}
    \frac{e^{in\cdot X}}{a-|n|^2}.
\end{align}
%
Here $X=(X_1,X_2)$ and $n\cdot X:=n_1X_1+n_2X_2$ for
$n=(n_1,n_2)\in\mathbb Z^2$. The series converges in
$L^2(\mathbb T_0)$ because its coefficients are square
summable. We use the Fourier normalization
%
\begin{align*}
    \widehat v(n):=\frac1{4\pi^2}
    \int_{\mathbb T_0}v(X)e^{-in\cdot X}\,dX
    \qquad\forall\,n\in\mathbb Z^2
\end{align*}
%
for integrable functions $v$. In this normalization,
$4\pi^2\delta_0=\sum_{n\in\mathbb Z^2}e^{in\cdot X}$ distributionally. Multiplication of each Fourier coefficient
in \eqref{eq:s4-cell-Green} by $a-|n|^2$ therefore gives
%
\begin{align}
    \label{eq:s4-cell-equation}
    (\Delta_0+a)G_a=4\pi^2\delta_0-1,
    \qquad \int_{\mathbb T_0}G_a\,dX=0.
\end{align}
%
These equations uniquely determine $G_a$: the mean-zero
condition fixes its coefficient at $n=0$, and
$a-|n|^2\neq0$ for every $n\neq0$. The coefficients at $n$
and $-n$ are equal and real, so $G_a$ is real valued. Since
they depend only on $|n|^2$, the function is invariant under
both coordinate reflections and interchange of the coordinates.
By standard elliptic regularity, we have
$G_a\in C^\infty(\mathbb T_0\smallsetminus\{0\})$.

To state estimates near the singularity, let
$\mathcal A_R:=\{X:R/2<|X|<2R\}$, where $R>0$ is small
enough that this annulus lies in the flat coordinate disk
centred at the origin. For a nonnegative integer $k$,
$\beta\in(0,1)$ and $v\in C^{k,\beta}(\mathcal A_R)$, define the rescaled H\"older norm
%
\begin{align*}
    \|v\|_{k,\beta;R}
    :=\sum_{j=0}^kR^j\|D_0^jv\|_{L^\infty(\mathcal A_R)}
    +R^{k+\beta}[D_0^kv]_{\beta;\mathcal A_R},
\end{align*}
%
where $D_0^0v:=v$ and, for a tensor field $T$ on the annulus,
%
\begin{align*}
    [T]_{\beta;\mathcal A_R}
    :=\sup_{\substack{X,Y\in\mathcal A_R\\X\neq Y}}
    \frac{|T(X)-T(Y)|}{|X-Y|^\beta}.
\end{align*}
%
Tensor components are compared in the parallel orthonormal
frame of the flat coordinates, and all pointwise norms use
the flat metric. By the chain rule, $\|v\|_{k,\beta;R}$ is
the corresponding $C^{k,\beta}$ norm of $Y\mapsto v(RY)$
on $\{Y:1/2<|Y|<2\}$. For annuli in $\mathbb T$, we use
the same definition in flat coordinates centred at the
specified point, with $D_0$ replaced by the flat covariant
derivative $D$ of $g_{\mathbb T}$.

\begin{lemma}[Uniform estimates for $G_a$]
    \label{lem:s4-Green}
    Set $r_*:=1/4$. There exist constants
    $a_*\in(0,1/4]$ and $c,C>0$ with the following
    properties. For every $a\in[0,1/4]$, there are a number
    $c_a\in\mathbb R$ and a smooth function $E_a$ on
    $B(0,r_*)\smallsetminus\{0\}$ such that
    %
    \begin{align}
        \label{eq:s4-Green-local}
        G_a(X)=2\pi\log r+c_a+E_a(X),
        \qquad |c_a|\leq C,\qquad |c_a-c_0|\leq Ca,
    \end{align}
    %
    for every $X\in B(0,r_*)\smallsetminus\{0\}$ and
    $r:=|X|$. For every nonnegative integer $k$ and every
    $\beta\in(0,1)$, there exists $C_{k,\beta}>0$ such that
    %
    \begin{align}
        \label{eq:s4-Green-remainder}
        \|E_a\|_{k,\beta;R}
        \leq C_{k,\beta}R^2(1+a|\log R|)
        \qquad\forall\,a\in[0,1/4],\quad 0<R<r_*/2.
    \end{align}
    %
    On the whole punctured torus, for every
    $a\in[0,1/4]$, $X\in\mathbb T_0\smallsetminus\{0\}$
    and $r:=d_{\mathbb T_0}(X,0)$, one has
    %
    \begin{align}
        \label{eq:s4-Green-global}
        |G_a(X)-2\pi\log r|\leq C,
        \qquad |D_0^jG_a(X)|\leq C_j(1+r^{-j})
        \quad\forall\,j\geq1.
    \end{align}
    %
    Finally, for every $a\in[0,a_*]$ and
    $X\in\mathbb T_0\smallsetminus\{0\}$, with the same
    definition of $r$, one has
    %
    \begin{align}
        \label{eq:s4-Green-Hessian}
        |D_0^2G_a(X)|^2&\geq c(1+r^{-4}),\nonumber\\
        2a|D_0G_a(X)|^2-2|D_0^2G_a(X)|^2
        &\leq-c(1+r^{-4}).
    \end{align}
    %
    All constants are independent of $a$ in the stated ranges.
\end{lemma}

\begin{proof}
    We first obtain uniform bounds away from the origin and
    compare $G_a$ with $G_0$. For every
    $n\in\mathbb Z^2\smallsetminus\{0\}$ and $a\in[0,1/4]$,
    %
    \begin{align*}
        \left|\frac1{a-|n|^2}\right|
        &\leq\frac4{3|n|^2},\\
        \widehat{G_a-G_0}(n)
        &=-\frac{a}{|n|^2(|n|^2-a)}.
    \end{align*}
    %
    Parseval's identity and the first inequality give a uniform
    $L^2$ bound for $G_a$. The absolute values in the second
    line have sum at most
    $(4a/3)\sum_{n\neq0}|n|^{-4}\leq Ca$. Thus $G_a-G_0$
    has a uniformly convergent Fourier series and
    %
    \begin{align*}
        \|G_a-G_0\|_{C^0(\mathbb T_0)}\leq Ca.
    \end{align*}
    %
    On a fixed compact subset of
    $\mathbb T_0\smallsetminus\{0\}$, the equation
    $(\Delta_0+a)G_a=-1$ and the $L^2$ bound give uniform
    bounds for every derivative by interior elliptic estimates.

    To identify the logarithmic singularity and its remainder,
    we construct a radial fundamental solution. The following
    power series are defined for every $a\geq0$; only
    $a\in[0,1/4]$ is needed for the uniform estimates in this
    lemma. Set $h_0:=0$ and
    $h_j:=\sum_{i=1}^j i^{-1}$ for every integer $j\geq1$,
    and define
    %
    \begin{align*}
        d_j(a)&:=\frac{(-a/4)^j}{(j!)^2}
            \qquad\forall\,j\geq0,\\
        \mathcal B_a(r)&:=\sum_{j=0}^{\infty}d_j(a)r^{2j}
            =J_0(\sqrt a\,r),\\
        W_a(r)&:=-\sum_{j=1}^{\infty}h_jd_j(a)r^{2j}
            \qquad\forall\,r\geq0,
    \end{align*}
    %
    where
    $J_0(x):=\sum_{j=0}^{\infty}(-1)^j(x^2/4)^j/(j!)^2$
    is the Bessel function of order zero. These series and
    all their derivatives converge uniformly when $a$ and $r$
    range over compact intervals. For the radial operator
    $\mathcal D_a:=\partial_r^2+r^{-1}\partial_r+a$, the identity
    %
    \begin{align*}
        4(j+1)^2d_{j+1}(a)+ad_j(a)=0
        \qquad\forall\,j\geq0
    \end{align*}
    %
    gives $\mathcal D_a\mathcal B_a=0$ on $(0,\infty)$.
    The coefficient of $r^{2j}$ in $\mathcal D_aW_a$ is
    %
    \begin{align*}
        -4(j+1)^2h_{j+1}d_{j+1}(a)-ah_jd_j(a)
        &=a(h_{j+1}-h_j)d_j(a)\\
        &=\frac{a}{j+1}d_j(a)
        =-4(j+1)d_{j+1}(a)
        \qquad\forall\,j\geq0.
    \end{align*}
    %
    These are precisely the coefficients of
    $-2\mathcal B_a'(r)/r$. The product rule also gives
    %
    \begin{align*}
        \mathcal D_a(\mathcal B_a\log r)
        =(\mathcal D_a\mathcal B_a)\log r
        +\frac{2\mathcal B_a'(r)}r
        =\frac{2\mathcal B_a'(r)}r
        \qquad\forall\,r>0.
    \end{align*}
    %
    Consequently the function
    %
    \begin{align}
        \label{eq:s4-radial-fundamental}
        P_a(r):=2\pi\left[
        J_0(\sqrt a\,r)\log r+
        \sum_{j\geq1}\frac{(-1)^{j+1}h_j}{(j!)^2}
        \left(\frac{ar^2}{4}\right)^j
        \right]
        \qquad\forall\,r>0
    \end{align}
    %
    satisfies $\mathcal D_aP_a=0$ away from the origin.

    The series give, as $r\downarrow0$ and uniformly for
    $a\in[0,1/4]$,
    %
    \begin{align*}
        P_a(r)&=2\pi\log r+O\bigl(ar^2(1+|\log r|)\bigr),\\
        P_a'(r)&=\frac{2\pi}{r}
            +O\bigl(ar(1+|\log r|)\bigr).
    \end{align*}
    %
    We also regard $X\mapsto P_a(|X|)$ as a locally
    integrable function on the flat disk $B(0,1)$.
    For $\varphi\in C_c^\infty(B(0,1))$, integration by
    parts on $B(0,1)\smallsetminus B(0,\epsilon)$ gives
    %
    \begin{align*}
        \int_{B(0,1)\smallsetminus B(0,\epsilon)}
        P_a(\Delta_0+a)\varphi\,dX
        =\int_{\partial B(0,\epsilon)}
        \bigl(\varphi P_a'-P_a\partial_r\varphi\bigr)\,d\sigma,
    \end{align*}
    %
    where $d\sigma$ is Euclidean arc length and
    $\partial_r$ points away from the origin. The boundary
    term containing $P_a\partial_r\varphi$ tends to zero,
    since its absolute value is $O(\epsilon|\log\epsilon|)$.
    The other boundary term tends to
    $(2\pi)(2\pi)\varphi(0)=4\pi^2\varphi(0)$.
    It follows that $(\Delta_0+a)P_a=4\pi^2\delta_0$
    distributionally on $B(0,1)$.

    Define $Q_a:=G_a-P_a$ on this disk. Subtracting the
    distributional equations gives $(\Delta_0+a)Q_a=-1$.
    Both terms defining $Q_a$ have uniformly bounded
    $L^2$ norms on $B(0,1)$. Interior elliptic regularity
    therefore shows that $Q_a$ is smooth across the origin
    and, for every nonnegative integer $j$,
    %
    \begin{align*}
        \sup_{a\in[0,1/4]}
        \|D_0^jQ_a\|_{L^\infty(B(0,1/2))}<\infty.
    \end{align*}
    %
    Set $c_a:=Q_a(0)$. Since $P_a-P_0$ tends to zero at
    the origin, evaluating the uniformly convergent series
    for $G_a-G_0$ there gives
    %
    \begin{align*}
        c_a-c_0
        =-a\sum_{n\in\mathbb Z^2\smallsetminus\{0\}}
        \frac1{|n|^2(|n|^2-a)}.
    \end{align*}
    %
    In particular, $|c_a-c_0|\leq Ca$ and $|c_a|\leq C$.
    The coordinate reflection symmetries of $G_a$ and the
    radiality of $P_a$ imply $D_0Q_a(0)=0$. Taylor's
    formula gives
    %
    \begin{align*}
        Q_a(RY)-Q_a(0)
        =R^2\sum_{i,j=1}^2Y_iY_j
        \int_0^1(1-t)(\partial_{ij}Q_a)(tRY)\,dt
    \end{align*}
    %
    whenever $0<R<r_*/2$ and $1/2<|Y|<2$.
    Differentiating with respect to $Y$ and using the
    bounds for the derivatives of $Q_a$ proves
    %
    \begin{align*}
        \|Q_a-Q_a(0)\|_{k,\beta;R}\leq C_{k,\beta}R^2.
    \end{align*}
    %
    Substituting $r=R|Y|$ in the power series for $P_a$
    similarly gives
    %
    \begin{align*}
        \|P_a-2\pi\log r\|_{k,\beta;R}
        \leq C_{k,\beta}aR^2(1+|\log R|).
    \end{align*}
    %
    Thus \eqref{eq:s4-Green-local} and
    \eqref{eq:s4-Green-remainder} follow upon setting
    $E_a:=Q_a-Q_a(0)+P_a-2\pi\log r$.
    The same estimates give \eqref{eq:s4-Green-global}
    near the origin. On the complement of a fixed small
    disk, they follow from the previously established
    derivative bounds and the fact that the distance to
    the origin is bounded above and bounded away from zero.
    No derivatives of the distance function are taken
    in \eqref{eq:s4-Green-global}.

    It remains to prove the Hessian estimates. The local
    expansion and its second-derivative bound give
    %
    \begin{align*}
        D_0^2G_a(X)
        =\frac{2\pi}{r^2}(I-2e_r\otimes e_r)
        +O(1+a|\log r|),
        \qquad e_r:=X/r,
    \end{align*}
    %
    for $0<r=|X|<r_*$. Here $I$ is the identity matrix
    in the flat orthonormal coordinates. The first term
    has norm $2\sqrt2\pi/r^2$. Choose
    $r_1\in(0,r_*)$ sufficiently small that the error
    has at most half this norm for every $0<r<r_1$ and
    $a\in[0,1/4]$. Such a choice is possible because
    $r^2(1+|\log r|)\to0$. Therefore
    $|D_0^2G_a|^2\geq cr^{-4}$ on that disk.

    On $\mathbb T_0\smallsetminus B(0,r_1)$, let
    $M_1$ be a uniform upper bound for $|G_a|$.
    Choose $a_*\in(0,1/4]$ so that $a_*M_1\leq1/2$.
    For a symmetric $2\times2$ matrix $H$, the inequality
    $|H|^2\geq\frac12(\operatorname{tr}H)^2$ follows
    from Cauchy--Schwarz. Applying it to $D_0^2G_a$ gives
    %
    \begin{align*}
        |D_0^2G_a|^2
        \geq\frac12|\Delta_0G_a|^2
        =\frac12|1+aG_a|^2\geq\frac18
        \qquad\forall\,a\in[0,a_*]
    \end{align*}
    %
    on the complement of the disk. Combining the two
    regions yields
    $|D_0^2G_a|^2\geq c_1(1+r^{-4})$ for some $c_1>0$.
    By \eqref{eq:s4-Green-global}, there is also $C_1>0$
    such that $|D_0G_a|^2\leq C_1(1+r^{-4})$.
    Decrease $a_*$, if necessary, so that
    $2a_*C_1\leq c_1$. Then
    %
    \begin{align*}
        2a|D_0G_a|^2-2|D_0^2G_a|^2
        \leq(2aC_1-2c_1)(1+r^{-4})
        \leq-c_1(1+r^{-4}),
    \end{align*}
    %
    as required.
\end{proof}

\subsection{The lattice solution and the choice of neck radii}

We now construct the function on $\mathbb T$ that determines
the two sheets away from the necks. We first rescale $G_a$
so that its singularities occur precisely at $L_m$ and its
equation becomes $L_{\mathbb T}\Phi_m=0$ away from $L_m$.
We then compare its expansion near a lattice point with the
height of a catenoid. This comparison determines the reference
neck radius $\tau_0$ and the error represented by $\zeta$.

For an integer $m\geq2$, set
%
\begin{align}
    \label{eq:s4-physical-Green}
    \begin{split}
        a&:=m^{-2},\qquad \kappa:=\sqrt2m,\\
        U_m(s,t)&:=1+aG_a(\kappa s,\kappa t),\\
        \Phi_m(s,t)&:=\frac{m^2}{2\pi}U_m(s,t)
        =\frac{m^2}{2\pi}+\frac1{2\pi}G_a(\kappa s,\kappa t),
    \end{split}
\end{align}
%
for every $(s,t)\in\mathbb T\smallsetminus L_m$.
The arguments of $G_a$ are taken modulo $2\pi$ in each
coordinate. Since $\kappa\ell=2\pi m$, the map
$(s,t)\mapsto(\kappa s,\kappa t)$ is a well-defined
covering of $\mathbb T_0$ by $\mathbb T$. The inverse
image of the origin is exactly $L_m$. The symmetries of
$G_a$ imply that $U_m$ and $\Phi_m$ are invariant under
the action of $\Gamma_m^0$ on $\mathbb T$.

The normalization in \eqref{eq:s4-physical-Green} and Lemma \ref{lem:s4-Green} give
%
\begin{align}
    \label{eq:s4-physical-equation}
    L_{\mathbb T}\Phi_m=2\pi\sum_{p\in L_m}\delta_p,
    \qquad
    \frac1{2\pi^2}\int_{\mathbb T}\Phi_m\,ds\,dt
    =\frac{m^2}{2\pi},
\end{align}
%
where the first equality holds distributionally and
$\delta_p$ is the Dirac measure of mass one at $p$.
Indeed, for $q=(s,t)\notin L_m$ and
$X:=(\kappa s,\kappa t)\in\mathbb T_0$, the chain rule gives
%
\begin{align*}
    L_{\mathbb T}U_m(q)
    &=a\kappa^2\Delta_0G_a(X)+2+2aG_a(X)\\
    &=2\bigl(-1-aG_a(X)\bigr)+2+2aG_a(X)=0.
\end{align*}
%
Near each $p\in L_m$, use the flat variable $q-p$ and
write $r_p:=|q-p|$. The local expansion of $G_a$ gives
$\Phi_m(q)=\log r_p+O(1)$, with remainder after subtraction
of the constant of order $r_p^2(1+|\log r_p|)$ for fixed $m$.
The first derivative of this remainder has order
$r_p(1+|\log r_p|)$. Integration by parts around $p$,
as in the preceding proof, therefore gives a point mass
$2\pi\delta_p$ and no other distribution supported at $p$.
This proves the first identity in
\eqref{eq:s4-physical-equation}. The average of
$G_a(\kappa s,\kappa t)$ on $\mathbb T$ is zero, because
the covering consists of $m^2$ identical square cells.
This proves the second identity.

For $q\in\mathbb T$, let
$\rho(q):=d_{\mathbb T}(q,L_m)$. Each point lies within
distance $\ell/(\sqrt2m)=\pi/m$ of a lattice point.
Moreover, the covering and the flat metrics give the exact
distance relation
%
\begin{align*}
    d_{\mathbb T_0}(X,0)=\kappa\rho(q)
    \qquad\forall\,q=(s,t)\in\mathbb T,
    \quad X:=(\kappa s,\kappa t).
\end{align*}
%
Hence Lemma~\ref{lem:s4-Green} implies
%
\begin{align}
    \Phi_m(q)&=\frac{m^2}{2\pi}
        +\log(\kappa\rho(q))+O(1),
        \label{eq:s4-physical-values}\\
    |D^j\Phi_m(q)|&\leq C_j\rho(q)^{-j}
        \qquad\forall\,j\geq1,
        \label{eq:s4-physical-derivatives}\\
    |D^2\Phi_m(q)|^2&\geq c\bigl(m^4+\rho(q)^{-4}\bigr)
        \label{eq:s4-physical-Hessian}
\end{align}
%
for every $q\in\mathbb T\smallsetminus L_m$ and all
sufficiently large $m$, chosen so that $m^{-2}\leq a_*$.
Here and below $D$ is the flat covariant derivative of
$g_{\mathbb T}$. To see each scaling factor explicitly,
in the parallel coordinate frames one has
%
\begin{align*}
    D^j\Phi_m(q)&=\frac{\kappa^j}{2\pi}D_0^jG_a(X)
        \qquad\forall\,j\geq1,\\
    |D^j\Phi_m(q)|
        &\leq C_j\bigl(\kappa^j+\rho(q)^{-j}\bigr)
        \leq C_j'\rho(q)^{-j},\\
    |D^2\Phi_m(q)|^2
        &=\frac{m^4}{\pi^2}|D_0^2G_a(X)|^2
        \geq c\bigl(m^4+\rho(q)^{-4}\bigr).
\end{align*}
%
The second inequality uses $m\rho(q)\leq\pi$.
The derivative bounds also give local H\"older estimates.
Precisely, if $B_q:=B(q,\rho(q)/8)$, then
%
\begin{align*}
    \sum_{i=0}^k\rho(q)^i
        \|D^{j+i}\Phi_m\|_{L^\infty(B_q)}
    +\rho(q)^{k+\beta}[D^{j+k}\Phi_m]_{\beta;B_q}
    \leq C_{j,k,\beta}\rho(q)^{-j}
\end{align*}
%
for all integers $j\geq1$, $k\geq0$ and
$\beta\in(0,1)$. Indeed, $\rho$ is bounded below by
$7\rho(q)/8$ on $B_q$. Apply
\eqref{eq:s4-physical-derivatives} there, and use one more
derivative and the mean-value estimate to bound the H\"older
seminorm. All tensor comparisons are in the flat frame.

The combination of derivatives that will enter the curvature
estimate also has a uniform sign. The chain rule gives
$DU_m=(\sqrt2/m)D_0G_a$ and $D^2U_m=2D_0^2G_a$,
where the right-hand sides are evaluated at $X$.
Consequently, with $r:=\kappa\rho(q)$,
%
\begin{align}
    \label{eq:s4-height-quadratic-margin}
    |DU_m(q)|^2-\tfrac12|D^2U_m(q)|^2
    =2a|D_0G_a(X)|^2-2|D_0^2G_a(X)|^2
    \leq-c(1+r^{-4}).
\end{align}
%
This inequality is stronger than the assertion that the
Hessian does not vanish, as it compares the negative Hessian
term directly with the positive gradient term.

We next fix the scales used to join the sheets to the
catenoids. Fix $\beta\in(0,1)$ and choose $\alpha>0$
sufficiently small for Convention~2.10 of
\cite{KapouleasMcGrathLD}, with the additional restrictions
%
\begin{align}
    \label{eq:s4-alpha}
    \alpha<\min\left\{\frac18,\frac15,
    \frac1{9(3+\beta)}\right\}.
\end{align}
%
A further restriction on $\alpha$, depending only on a
fixed cutoff, will be specified when we define the
diffeomorphisms between the initial surfaces. Every choice
of $\alpha$ is made before choosing $m$.
Define
%
\begin{align}
    \label{eq:s4-scales}
    \begin{split}
        \delta&:=(100m)^{-1},\qquad
        A_m:=\frac{m^2}{2\pi}+\log(\sqrt2m)
            +\frac{c_{m^{-2}}}{2\pi},\\
        \tau_0&:=2e^{-A_m},\qquad
        \tau_\zeta:=\tau_0e^\zeta
            \quad\forall\,\zeta\in[-1,1],\qquad
        \varepsilon_\zeta:=\frac{m^2\tau_\zeta}{2\pi}.
    \end{split}
\end{align}
%
The radius $\delta$ is a fixed small fraction of the
lattice spacing. The number $\tau_\zeta$ will be the
waist radius of each catenoidal neck. The quantity
$\varepsilon_\zeta$ is the constant term in the expression
$\tau_\zeta\Phi_m=\varepsilon_\zeta+
(\tau_\zeta/2\pi)G_{m^{-2}}(\kappa s,\kappa t)$.
It gives the principal size of the sheet height away
from the lattice points.

By \eqref{eq:s4-Green-local},
$c_{m^{-2}}=c_0+O(m^{-2})$. Therefore
%
\begin{align*}
    \tau_\zeta
    =\frac{\sqrt2}{m}
    \exp\left(-\frac{m^2}{2\pi}
        -\frac{c_{m^{-2}}}{2\pi}+\zeta\right)
    \qquad\forall\,\zeta\in[-1,1].
\end{align*}
%
In particular, there are constants $c,C>0$ such that
$cm^2\leq|\log\tau_\zeta|\leq Cm^2$ for all sufficiently
large $m$ and all $\zeta\in[-1,1]$. More generally,
%
\begin{align*}
    \sup_{\zeta\in[-1,1]}m^p\tau_\zeta^q\rightarrow0
    \qquad\text{as }m\to\infty
\end{align*}
%
for every fixed pair of real numbers $p\geq0$ and $q>0$.
We will use this fact whenever a fixed power of $m$
multiplies a positive power of $\tau_\zeta$.

For $p\in L_m$, write $r_p(q):=d_{\mathbb T}(q,p)$ in
the flat coordinate disk centred at $p$. Rescaling the local
Green expansion gives
%
\begin{align}
    \label{eq:s4-physical-expansion}
    \Phi_m(q)=\log r_p(q)+A_m+\mathcal E_{m,p}(q),
    \qquad
    \|\mathcal E_{m,p}\|_{k,\beta;R}
    \leq C_{k,\beta}R^2(m^2+|\log R|)
\end{align}
%
for $0<r_p(q)<2\delta$ and $0<R<\delta$, with the norm
taken on the annulus about $p$ of radii $R/2$ and $2R$.
In fact,
$\mathcal E_{m,p}(q)=(2\pi)^{-1}
E_{m^{-2}}(\kappa(q-p))$. Since
$\kappa R<\sqrt2/100<r_*/2$, the scaled norm and
\eqref{eq:s4-Green-remainder} give
%
\begin{align*}
    \|\mathcal E_{m,p}\|_{k,\beta;R}
    \leq C_{k,\beta}(\kappa R)^2
        \bigl(1+m^{-2}|\log(\kappa R)|\bigr)\leq C_{k,\beta}'R^2(m^2+|\log R|).
\end{align*}
%
In the last step we used
$|\log(\kappa R)|\leq|\log R|+\log\kappa$ and
$\log\kappa\leq Cm^2$.

Define the proposed upper sheet height before interpolation by
%
\begin{align*}
    \psi_\zeta:=\tau_\zeta\Phi_m
    \quad\text{on }\mathbb T\smallsetminus L_m,
    \qquad\forall\,\zeta\in[-1,1].
\end{align*}
%
This function satisfies $L_{\mathbb T}\psi_\zeta=0$
away from $L_m$ and has logarithmic coefficient
$\tau_\zeta$ at each lattice point. Thus, it is a linearized
doubling solution in the terminology of
\cite[Definition~3.6]{KapouleasMcGrathLD}.
For $\tau>0$, the upper half of a Euclidean catenoid of
waist radius $\tau$ is the radial graph
%
\begin{align*}
    C_\tau(r):=\tau\operatorname{arcosh}(r/\tau)
    \qquad\forall\,r\geq\tau.
\end{align*}
%
Using $\operatorname{arcosh}x=\log(x+\sqrt{x^2-1})$,
we obtain
%
\begin{align*}
    C_\tau(r)
    &=\tau\log r+\tau\log(2/\tau)
        +O(\tau^3r^{-2})
        \qquad\text{as }r/\tau\to\infty.
\end{align*}
%
The coefficients of $\log r$ in $\tau\Phi_m$ and
$C_\tau$ already agree. Their constant terms agree
precisely when $A_m=\log(2/\tau)$, or
$\tau=2e^{-A_m}=\tau_0$. This explains the reference
radius in \eqref{eq:s4-scales}.
Figure~\ref{fig:s4-catenoid-matching} compares the
catenoidal heights with their logarithmic terms after
dividing both radius and height by $\tau$.

\begin{figure}[htbp]
    \centering
    \begingroup
    \definecolor{matchteal}{RGB}{0,112,116}
    \begin{tikzpicture}[
        x=1cm,y=0.93cm,>=Stealth,
        every node/.style={font=\small}
    ]
        \draw[black!45,->] (0.55,0)--(5.45,0)
            node[right] {$r/\tau$};
        \draw[black!45,->] (0.55,-2.62)--(0.55,2.75)
            node[above] {$z/\tau$};
        \draw[black!75,line width=1.1pt]
            plot[domain=-2.27:2.27,samples=110]
                ({cosh(\x)},\x);
        \draw[matchteal,dashed,line width=0.95pt]
            plot[domain=1:4.89,samples=70] (\x,{ln(2*\x)});
        \draw[matchteal,dashed,line width=0.95pt]
            plot[domain=1:4.89,samples=70] (\x,{-ln(2*\x)});
        \fill[black!75] (1,0) circle[radius=1.8pt];
        \node[below=4pt] at (1,0) {$1$};
        \node[anchor=east,align=right] at (0.33,-0.50)
            {waist\\$r=\tau,\ z=0$};
        \draw[black!45] (0.39,-0.43)--(0.96,-0.04);
        \draw[black!75,line width=1.1pt] (6.2,1.62)--(6.92,1.62);
        \node[anchor=west,align=left] at (7.08,1.62)
            {Catenoidal height\\[3pt]
             $z=\pm\tau\operatorname{arcosh}(r/\tau)$};
        \draw[matchteal,dashed,line width=0.95pt]
            (6.2,0.32)--(6.92,0.32);
        \node[anchor=west,align=left] at (7.08,0.32)
            {Logarithmic term\\[3pt]
             $z=\pm\tau\log(2r/\tau)$};
        \node[anchor=west,align=left] at (6.2,-1.20)
            {$\tau\log(2r/\tau)$\\[3pt]
             $\qquad=\tau\log r+\tau\log(2/\tau)$};
    \end{tikzpicture}
    \endgroup
    \caption{The radial section of a catenoid, shown in
    the dimensionless coordinates $r/\tau$ and $z/\tau$,
    and the logarithmic terms used to match its two ends
    to the sheets. The solid curve is parametrized by
    $(r/\tau,z/\tau)=(\cosh\xi,\xi)$.
    As $r/\tau$ tends to infinity, the upper height
    differs from $\tau\log(2r/\tau)$ by
    $O(\tau^3r^{-2})$.
    The coefficient of $\log r$ fixes the neck radius
    $\tau$, and matching the remaining constant
    $\tau\log(2/\tau)$ determines its relation to the
    Green function.}
    \label{fig:s4-catenoid-matching}
\end{figure}

To relate this comparison to the local matching conditions
in the cited construction, we specify both its constant
and its first-order part. Suppose that $\tau>0$ and that
$u$ is a smooth solution of $L_{\mathbb T}u=0$ on a
punctured flat disk centred at $p$, with
$u-\tau\log r_p=O(1)$ as $q\to p$. There are unique
$c_p\in\mathbb R$ and $\lambda_p\in T_p^*\mathbb T$
such that
%
\begin{align*}
    u(q)=\tau\log(2r_p(q)/\tau)
        +c_p+\lambda_p(q-p)
        +O\bigl(r_p(q)^2|\log r_p(q)|\bigr)
        \qquad\text{as }q\to p.
\end{align*}
%
The variable $q-p$ is understood in the flat coordinate
disk. In particular, $u-\tau\log(2r_p/\tau)$ extends as
a $C^1$ function across $p$, with value $c_p$ and
differential $\lambda_p$. The affine function
$q-p\mapsto c_p+\lambda_p(q-p)$ is called the
\textit{mismatch jet}. It records the difference between
the first-order expansion of $u$ after subtraction of
its logarithmic singularity and that of the catenoidal
end.

For completeness, the stated regularity follows from the
radial construction in \eqref{eq:s4-radial-fundamental}
with parameter $a=2$. The function
$P_2(r)/(2\pi)$ solves $(\Delta+2)(P_2/(2\pi))=0$
    for $r>0$
and differs from $\log r$ by $O(r^2|\log r|)$, with
first derivative $O(r|\log r|)$. Thus
$h:=u-\tau P_2(r_p)/(2\pi)$ is bounded and solves
$(\Delta+2)h=0$ away from $p$.
To justify removability, extend $h$ as a bounded
measurable function on the disk and solve
$\Delta w=-2h$ there with zero boundary values.
For every finite $p_0>2$, local elliptic regularity gives
$w\in W^{2,p_0}_{\mathrm{loc}}\subset C^1$.
Then $h-w$ is bounded and harmonic on the punctured
disk, so its singularity is removable. Hence $h$ solves
$(\Delta+2)h=0$ distributionally on the whole disk and
is smooth by elliptic regularity. Taylor's theorem for
$h$, together with the expansion of $P_2$, proves the
claimed expansion and $C^1$ extension.

Apply this definition to $u=\psi_\zeta$ and
$\tau=\tau_\zeta$. The reflections fixing $p$ force
$\lambda_p=0$, and lattice translations make $c_p$
independent of $p$. By \eqref{eq:s4-physical-expansion},
its common value is
%
\begin{align}
    \label{eq:s4-exact-mismatch}
    c_p
    =\tau_\zeta\bigl(A_m+\log(\tau_\zeta/2)\bigr)
    =\tau_\zeta\zeta
    \qquad\forall\,p\in L_m,\quad\zeta\in[-1,1].
\end{align}
%
The last equality uses
$\log(\tau_\zeta/2)=-A_m+\zeta$.
Thus, the symmetries of the construction allow only one common constant
mismatch and no linear mismatch. The parameter $\zeta$
varies precisely this constant. We retain this freedom
because the later perturbation argument must also
eliminate the one-dimensional scalar error in its
linear equation.


\subsection{Joining the two graphs by catenoidal necks}

We now replace the logarithmic singularity of each graph by
one end of a catenoid. There are two separate steps. First,
we subtract the constant difference
$\tau_\zeta\zeta$ between the expansion of
$\psi_\zeta$ and that of the catenoidal height. We make this
subtraction using a function that solves
$L_{\mathbb T}v=0$ near each lattice point. The resulting
linear error is therefore supported away from the necks.
Second, we interpolate between the corrected graph and the
catenoid on a smaller annulus. The estimates below show that
the resulting surface is smooth, embedded, and has controlled
geometry at the scale of its distance from the lattice.

Choose a smooth nondecreasing function
$\vartheta:\mathbb R\to[0,1]$ such that
%
\begin{align*}
    \begin{cases}
        \vartheta(x)=0 & \forall\,x\le-1\\
        \vartheta(x)=1 & \forall\,x\ge 1\\
        \vartheta(-x)=1-\vartheta(x) & \forall\,x\in\mathbb{R}.
    \end{cases}
\end{align*}
%
For $a_1,a_2\in\mathbb R$ with $a_1\ne a_2$, define
%
\begin{align}
    \label{eq:s4-cutoff}
    \eta_{a_1,a_2}(r)
    :=\vartheta\left(6\frac{r-a_1}{a_2-a_1}-3\right)
    \qquad\forall\,r\in\mathbb R.
\end{align}
%
Then $\eta_{a_1,a_2}+\eta_{a_2,a_1}=1$. If $a_1<a_2$,
the function $\eta_{a_1,a_2}$ is zero for
$r\leq(2a_1+a_2)/3$ and one for
$r\geq(a_1+2a_2)/3$. Thus the change from zero to one
occurs strictly between $a_1$ and $a_2$. For every integer
$j\geq1$,
%
\begin{align*}
    \left|\frac{d^j}{dr^j}\eta_{a_1,a_2}(r)\right|
    \leq C_j|a_2-a_1|^{-j}
    \qquad\forall\,r\in\mathbb R,
\end{align*}
%
where $C_j$ depends only on $j$ and $\vartheta$.
This is the cutoff convention in
\cite[Definition~1.5]{KapouleasMcGrathLD}.

Recall that $\delta=(100m)^{-1}$ and that $J_0$ denotes the Bessel function of order zero, as in the proof of Lemma \ref{lem:s4-Green}. Define smooth functions
$\widehat v_m,k_m:\mathbb T\to\mathbb R$ by
%
\begin{align}
    \label{eq:s4-substitute}
    \widehat v_m(q)
    :=\sum_{p\in L_m}
      \eta_{2\delta,\delta}(r_p(q))
      J_0(\sqrt2r_p(q)),
    \qquad
    k_m:=L_{\mathbb T}\widehat v_m,
\end{align}
%
where each summand is defined on $B(p,2\delta)$ and
extended by zero to its complement. The disks
$B(p,2\delta)$, $p\in L_m$, are pairwise disjoint.
The cutoff vanishes on a neighborhood of the boundary of
each disk, and the power series for $J_0(\sqrt2r_p)$
is a smooth power series in $r_p^2$. These observations
justify smoothness both at $p$ and across the boundary
of the disk. Moreover,
%
\begin{align}
    \label{eq:s4-substitute-properties}
    \widehat v_m(p)=1,\qquad
    d\widehat v_m(p)=0
    \quad\forall\,p\in L_m,
    \qquad
    \operatorname{supp}k_m
    \subset\bigcup_{p\in L_m}
           \{q:\delta\leq r_p(q)\leq2\delta\}.
\end{align}
%
Indeed, near $p$ the cutoff is one and
%
\begin{align*}
    J_0(\sqrt2r_p)
    =1-\frac{r_p^2}{2}+O(r_p^4),
    \qquad
    (\Delta_{\mathbb T}+2)J_0(\sqrt2r_p)=0.
\end{align*}
%
The second identity follows from the radial equation
$J''+r^{-1}J'+2J=0$ and holds also at $p$ by smoothness.
On $B(p,2\delta)$, the product rule gives
%
\begin{align*}
    L_{\mathbb T}
    \bigl(\eta_{2\delta,\delta}(r_p)J_0(\sqrt2r_p)\bigr)
    =J_0(\sqrt2r_p)
        \Delta_{\mathbb T}\eta_{2\delta,\delta}(r_p)+2\left\langle
        D\eta_{2\delta,\delta}(r_p),
        DJ_0(\sqrt2r_p)\right\rangle.
\end{align*}
%
Both terms vanish wherever the cutoff is constant, proving
the support assertion. In particular, subtracting
$c\widehat v_m$ subtracts the same constant $c$ from the
regular part at every lattice point, changes none of its
first derivatives there, and adds the linear error $-ck_m$.
The value map $c\widehat v_m\mapsto c$ is an isomorphism
from $\operatorname{span}\{\widehat v_m\}$ to $\mathbb R$.
In the terminology of the scalar linearized doubling
construction, $k_m$ is the substitute function associated
with this one-dimensional space.

We apply this correction with $c=\tau_\zeta\zeta$ and set
%
\begin{align}
    \label{eq:s4-corrected-exterior}
    \widetilde\psi_\zeta
    :=\tau_\zeta\Phi_m
       -\tau_\zeta\zeta\widehat v_m
    \qquad\text{on }\mathbb T\smallsetminus L_m.
\end{align}
%
Since $L_{\mathbb T}\Phi_m=0$ away from $L_m$,
%
\begin{align*}
    L_{\mathbb T}\widetilde\psi_\zeta
    =-\tau_\zeta\zeta k_m
    \qquad\text{on }\mathbb T\smallsetminus L_m.
\end{align*}
%
Thus the correction has removed the constant mismatch,
while its linear error is confined to the annuli
$\delta\leq r_p\leq2\delta$.

For the rest of this subsection, fix $\zeta\in[-1,1]$ and
abbreviate $\tau:=\tau_\zeta$ and $r_c:=\tau^\alpha$.
The radius $r_c$ specifies where we interpolate between the
catenoidal height and $\widetilde\psi_\zeta$. Since
$0<\alpha<1$, the estimates for $\tau_\zeta$ give
%
\begin{align*}
    \frac{\tau}{r_c}=\tau^{1-\alpha}\rightarrow0,
    \qquad
    \frac{r_c}{\delta}=100m\tau^\alpha\rightarrow0
    \qquad\text{as }m\to\infty,
\end{align*}
%
uniformly in $\zeta$. Hence, for all sufficiently large $m$,
we have $9\tau<2r_c<3r_c<\delta$. The interpolation
annuli are therefore much farther from each neck centre
than its waist radius $\tau$, and much closer to that centre
than the annuli supporting $k_m$.

On the torus with disks removed,
%
\begin{align*}
    \Omega_\zeta
    :=\mathbb T\smallsetminus
      \bigcup_{p\in L_m}B(p,9\tau),
\end{align*}
%
define the height function
%
\begin{align}
    \label{eq:s4-glued-height}
    h_\zeta(q):=
    \begin{cases}
        \eta_{3r_c,2r_c}(r_p(q))C_\tau(r_p(q))
        +\eta_{2r_c,3r_c}(r_p(q))\widetilde\psi_\zeta(q),
        &9\tau\leq r_p(q)\leq3r_c,
          \quad p\in L_m,\\
        \widetilde\psi_\zeta(q),
        &q\notin\displaystyle\bigcup_{p\in L_m}B(p,3r_c).
    \end{cases}
\end{align}
%
The disks in the first line are disjoint, so the choice of
$p$ is unambiguous. The two formulas agree on a neighborhood
of $r_p=3r_c$, where the first cutoff is zero and the
second is one. In particular, $h_\zeta$ is smooth up to
the boundary of $\Omega_\zeta$.

For each $p\in L_m$, define
%
\begin{align*}
    \mathcal D_\zeta
    :=\{(\xi,\theta)\in\mathbb R\times
         (\mathbb R/2\pi\mathbb Z):\tau\cosh\xi<2r_c\}
\end{align*}
%
and the catenoidal parametrization
%
\begin{align}
    \label{eq:s4-catenoid-parametrization}
    \mathcal C_{\zeta,p}(\xi,\theta)
    :=\mathcal F\bigl(
         p+\tau\cosh\xi(\cos\theta,\sin\theta),
         \tau\xi\bigr)
    \qquad\forall\,(\xi,\theta)\in\mathcal D_\zeta.
\end{align}
%
Addition to $p$ is taken in its flat coordinate disk.
The first component has distance $\tau\cosh\xi$ from
$p$, and the second component is the signed-distance
coordinate $z=\tau\xi$ in $\mathcal N$.
For large $m$, these distances are less than $2r_c<\delta$
and $|\tau\xi|\leq\tau\operatorname{arcosh}(2r_c/\tau)
\to0$, respectively. Thus the whole parametrization lies
in the domain of $\mathcal F$. Its circle $\xi=0$ lies
at $z=0$ and projects onto the circle of radius $\tau$
about $p$, which is the waist of the neck.

The height estimates proved below give
$|h_\zeta|\leq C\varepsilon_\zeta<z_0$ for all
sufficiently large $m$. Thus, the following graphs lie into $\mathcal N$. Define the initial surface of parameter $\zeta$ by
%
\begin{align}
    \label{eq:s4-initial-double}
    \begin{split}
        M_\zeta:=
        \{\mathcal F(q,h_\zeta(q)):q\in\Omega_\zeta\}
        \cup
        \{\mathcal F(q,-h_\zeta(q)):q\in\Omega_\zeta\}\cup
        \bigcup_{p\in L_m}
        \mathcal C_{\zeta,p}(\mathcal D_\zeta).
    \end{split}
\end{align}
%
On the overlap $9\tau\leq r_p<2r_c$, the height
$h_\zeta$ equals $C_\tau(r_p)$. If
$r_p=\tau\cosh\xi$, then
$C_\tau(r_p)=\tau|\xi|$. Therefore the upper graph agrees
exactly with the part of the catenoid where $\xi>0$,
and the lower graph agrees exactly with the part where
$\xi<0$. The graph pieces and catenoidal pieces
therefore have exactly the same images on these overlaps.

% Figure~\ref{fig:s4-initial-double} shows a local part of this surface in the coordinates of $\mathcal F$.

% Schematic in the parameter domain of \mathcal F, not a Euclidean model of its metric.
\begin{figure}[htbp]
    \centering
    \begingroup
    \definecolor{doubleteal}{RGB}{0,112,116}
    \definecolor{doubleorange}{RGB}{190,102,36}
    \begin{tikzpicture}[
        x=1cm,y=1cm,>=Stealth,
        every node/.style={font=\small},
        sheet edge/.style={black!65,line width=0.8pt},
        neck edge/.style={black!75,line width=0.85pt},
        reference/.style={black!40,dashed,line width=0.55pt}
    ]
        % Each plate is a schematic local piece of a toroidal sheet.
        % The neck radius and the vertical separation are exaggerated.
        \path[fill=doubleteal!8,even odd rule]
            (0,-1.15)--(7.3,-1.15)--(8.25,-0.40)--(0.95,-0.40)--cycle
            (1.6,-0.78) ellipse[x radius=0.66,y radius=0.19]
            (3.98,-0.78) ellipse[x radius=0.66,y radius=0.19]
            (6.36,-0.78) ellipse[x radius=0.66,y radius=0.19];
        \draw[sheet edge]
            (0,-1.15)--(7.3,-1.15)--(8.25,-0.40)--(0.95,-0.40)--cycle;
        \foreach \xc in {1.6,3.98,6.36} {
            \draw[black!35,dashed] (\xc,-0.78)
                ellipse[x radius=0.66,y radius=0.19];
        }
        % The central level is only a reference, not part of M_zeta.
        \draw[reference]
            (0,0)--(7.3,0)--(8.25,0.75)--(0.95,0.75)--cycle;
        \draw[reference] (0.48,0.37)--(7.78,0.37);
        % The upper sheet is drawn before the necks so that the latter remain visible.
        \path[fill=doubleteal!12,even odd rule]
            (0,1.15)--(7.3,1.15)--(8.25,1.90)--(0.95,1.90)--cycle
            (1.6,1.52) ellipse[x radius=0.66,y radius=0.19]
            (3.98,1.52) ellipse[x radius=0.66,y radius=0.19]
            (6.36,1.52) ellipse[x radius=0.66,y radius=0.19];
        \draw[sheet edge]
            (0,1.15)--(7.3,1.15)--(8.25,1.90)--(0.95,1.90)--cycle;
        % Three catenoidal necks in one row of the square lattice.
        \foreach \xc in {1.6,3.98,6.36} {
            \begin{scope}[shift={(\xc,0.37)}]
                \path[fill=doubleteal!7]
                    plot[domain=-2.23:2.23,samples=65]
                        ({0.14*cosh(\x)},{1.15*\x/2.23})
                    -- plot[domain=2.23:-2.23,samples=65]
                        ({-0.14*cosh(\x)},{1.15*\x/2.23}) -- cycle;
                \draw[neck edge]
                    plot[domain=-2.23:2.23,samples=65]
                        ({0.14*cosh(\x)},{1.15*\x/2.23});
                \draw[neck edge]
                    plot[domain=-2.23:2.23,samples=65]
                        ({-0.14*cosh(\x)},{1.15*\x/2.23});
                \draw[black!45] (0,1.15)
                    ellipse[x radius=0.66,y radius=0.19];
                \draw[black!35,dashed] (0,-1.15)
                    ellipse[x radius=0.66,y radius=0.19];
                \draw[doubleteal,line width=1pt] (0,0)
                    ellipse[x radius=0.14,y radius=0.04];
                \fill[doubleorange] (0,0) circle[radius=1.3pt];
            \end{scope}
        }
        \draw[black!50,->] (-0.30,-1.34)--(-0.30,2.13)
            node[above] {$z$};
        \node[anchor=west] at (8.65,1.53) {upper sheet: $z=h_\zeta(q)$};
        \draw[black!45] (8.56,1.53)--(7.91,1.65);
        \node[anchor=west] at (8.65,0.36) {central level: $z=0$};
        \draw[black!45,dashed] (8.56,0.36)--(7.94,0.51);
        \node[anchor=west] at (8.65,-0.81) {lower sheet: $z=-h_\zeta(q)$};
        \draw[black!45] (8.56,-0.81)--(7.91,-0.62);
        \node[anchor=north,align=center] at (1.6,-1.78)
            {waist\\$r_p=\tau_\zeta,\quad z=0$};
        \draw[black!45] (1.6,-1.68)--(1.02,-1.45)--(1.02,0.20)--(1.46,0.35);
        \node[anchor=north,align=center] at (5.65,-1.78)
            {centres $(p,0)$\\with $p\in L_m$};
        \draw[doubleorange,->] (5.50,-1.67)--(5.27,-1.41)--(5.27,0.15)--(6.27,0.35);
        \node[anchor=west,text=black!65] at (0,2.38)
            {A local part of $\mathcal F^{-1}(M_\zeta)$ in the $(s,t,z)$ coordinates};
    \end{tikzpicture}
    \endgroup
    \caption{The initial surface $M_\zeta$ is formed from two graphs
    over the central torus, joined near each $p\in L_m$ by a
    catenoidal neck. Three neighbouring necks in one lattice row
    are shown. The teal circles have radius $\tau_\zeta$ in the
    flat base coordinates and lie at $z=0$; their centres are
    marked in orange and do not belong to the surface. The dashed
    lines indicate the reference level $z=0$, which meets the
    surface precisely along these circles. The drawing represents
    the parameter coordinates
    of $\mathcal F$, not an isometric embedding of $\mathcal N$
    into $\mathbb R^3$. It suppresses the interpolation details
    and exaggerates both the sheet separation and the neck radii.
    The displayed sheet boundaries are artificial: the complete
    surface has two punctured toroidal sheets joined by all $m^2$ necks
    and has no boundary.}
    \label{fig:s4-initial-double}
\end{figure}


We write $\jmath_\zeta:M_\zeta\hookrightarrow\mathcal N$
and $\iota_\zeta:M_\zeta\hookrightarrow\mathbb S^4$ for
the inclusions, once the smoothness and embeddedness asserted
below have been proved. We retain the notation of
Section~\ref{sec:s4-geometry} for their second fundamental
forms and shape operators. In particular,
$S_{\jmath_\zeta}^{\nu_\zeta}$ is the shape operator
within $\mathcal N$, with respect to the unit normal
$\nu_\zeta$ specified in the lemma. Both inclusions induce
the same metric
%
\begin{align*}
    g_\zeta:=\jmath_\zeta^*g_{\mathcal N}
             =\iota_\zeta^*g_{\mathbb S^4}.
\end{align*}

\begin{lemma}
    \label{lem:s4-initial-geometry}
    For all sufficiently large $m$ and every
    $\zeta\in[-1,1]$, $M_\zeta$ is a smooth closed
    connected embedded $\Gamma_m$-invariant surface in
    $\mathcal N$, of genus $m^2+1$. It is the boundary
    of a $\Gamma_m$-invariant open subset of
    $\mathcal N$ with compact closure. Let $\nu_\zeta$
    be its outward unit normal within $\mathcal N$.
    For every $x\in M_\zeta$,
    %
    \begin{align}
        \label{eq:s4-initial-geometry-bounds}
        |z(x)|\leq C\varepsilon_\zeta,\qquad
        |(\nu_\zeta)_h(x)|\leq
          C\frac{\tau_\zeta}{\rho_\zeta(x)},\qquad
        |S_{\jmath_\zeta}^{\nu_\zeta}(x)|_{g_\zeta}
          \leq C\frac{\tau_\zeta}{\rho_\zeta(x)^2},
    \end{align}
    %
    where $(\nu_\zeta)_h$ denotes the orthogonal projection
    onto $\operatorname{span}\{E_1,E_2\}$ and
    %
    \begin{align}
        \label{eq:s4-rho-on-double}
        \rho_\zeta(x)
        :=d_{\mathbb T}(\pi(x),L_m),
        \qquad
        \tau_\zeta\leq\rho_\zeta(x)\leq\pi/m.
    \end{align}
    %
    The norm of $(\nu_\zeta)_h$ is taken with respect
    to $g_{\mathcal N}$. At every $x\in M_\zeta$,
    coordinates on a neighborhood of radius a fixed
    positive fraction of $\rho_\zeta(x)$ give uniformly
    elliptic coefficients for the metric
    $\rho_\zeta(x)^{-2}g_\zeta$, with bounds for every
    fixed finite number of their coordinate derivatives.
    All constants are independent of $m$, $\zeta$, and $x$.
    Moreover,
    %
    \begin{align}
        \label{eq:s4-scaled-injectivity}
        \operatorname{inj}_{g_\zeta}(x)
        \geq\rho_\zeta(x)/5
        \qquad\forall\,x\in M_\zeta.
    \end{align}
\end{lemma}

\begin{proof}
    We first check that the two graphs are separated.
    Write $\rho(q):=d_{\mathbb T}(q,L_m)$ for
    $q\in\mathbb T$. From
    \eqref{eq:s4-physical-values} and \eqref{eq:s4-scales},
    %
    \begin{align}
        \label{eq:s4-truncated-positivity}
        \frac{1-\alpha}{4\pi}m^2
        \leq\Phi_m(q)\leq Cm^2
        \qquad\text{whenever }\rho(q)\geq r_c,
    \end{align}
    %
    for all sufficiently large $m$. To obtain the lower
    bound, use $\log\rho\geq\alpha\log\tau$ and
    $\log\tau=\log2-A_m+\zeta$ to write
    %
    \begin{align*}
        \Phi_m(q)
        &\geq\frac{m^2}{2\pi}+\log\kappa
                     +\alpha\log\tau-C\\
        &=\frac{1-\alpha}{2\pi}m^2
                     +(1-\alpha)\log\kappa+O(1).
    \end{align*}
    %
    The $O(1)$ term is uniform in $\zeta$. The upper
    bound follows from $\kappa\rho\leq\sqrt2\pi$.
    Since $|\widehat v_m|\leq C$ and $|\zeta|\leq1$,
    it follows that $\widetilde\psi_\zeta>0$ on
    $\rho\geq r_c$. Also $C_\tau(r)>0$ for $r>\tau$.
    Every value of $h_\zeta$ is either one of these
    positive heights or a convex combination of them.
    Consequently $h_\zeta>0$ on $\Omega_\zeta$.

    We next estimate the difference between the two heights
    on the interpolation annuli. If $r_p\leq\delta$,
    then $\widehat v_m=J_0(\sqrt2r_p)$, and the local
    expansion of $\Phi_m$ gives
    %
    \begin{align*}
        \widetilde\psi_\zeta
        -\tau\log(2r_p/\tau)
        =\tau\mathcal E_{m,p}
         +\tau\zeta\bigl(1-J_0(\sqrt2r_p)\bigr).
    \end{align*}
    %
    Here we used
    $A_m+\log(\tau/2)=\zeta$ to cancel the constant
    term. On the other hand, for $r>\tau$,
    %
    \begin{align*}
        C_\tau(r)
        &=\tau\log(2r/\tau)
          +\tau\log\left(
                \frac{1+\sqrt{1-\tau^2/r^2}}2\right)\\
        &=\tau\log(2r/\tau)
          -\frac{\tau^3}{4r^2}+O(\tau^5r^{-4})
          \qquad\text{as }\tau/r\to0.
    \end{align*}
    %
    The last expansion follows from the convergent power
    series in $\tau^2/r^2$ and can be differentiated any
    fixed number of times on rescaled annuli. Since
    $m^2+|\log r_c|\leq C|\log\tau|$,
    \eqref{eq:s4-physical-expansion} and the power series
    for $J_0$ imply
    %
    \begin{align}
        \label{eq:s4-matching-estimate}
        \|\widetilde\psi_\zeta-C_\tau\|_{k,\beta;r_c}
        \leq C_{k,\beta}
        \bigl(\tau r_c^2|\log\tau|
                   +\tau^3r_c^{-2}\bigr).
    \end{align}
    %
    In this formula the norm is taken on
    $\{q:2r_c\leq r_p(q)\leq3r_c\}$ using the same
    scaled formula as before, with scale $r_c$.
    Multiplication by either cutoff in
    \eqref{eq:s4-glued-height} preserves this estimate:
    after the substitution $q=p+r_cY$, the cutoffs are
    fixed smooth functions of $|Y|$. In particular, for
    each positive derivative order, the interpolation
    error divided by the scale $\tau r_c^{-j}$ is bounded by
    %
    \begin{align*}
        C_j\left(r_c^2|\log\tau|+\tau^2r_c^{-2}\right)
        =C_j\left(\tau^{2\alpha}|\log\tau|
                        +\tau^{2-2\alpha}\right)
        \rightarrow0.
    \end{align*}

    The graph heights satisfy $|h_\zeta|\leq C\tau m^2$.
    On the part where $h_\zeta=\widetilde\psi_\zeta$ this
    follows from \eqref{eq:s4-truncated-positivity}; on
    the catenoidal part it follows from
    $C_\tau(r)\leq\tau\log(2r/\tau)$ and
    $r\leq3r_c$. The interpolated heights satisfy the same
    estimate. Since $\tau m^2=2\pi\varepsilon_\zeta\to0$,
    all these graphs lie within the prescribed band
    $|z|<z_0$ for large $m$.

    To verify regularity of the catenoidal pieces, put
    $r:=\tau\cosh\xi$, $z:=\tau\xi$, and
    $c:=\cos z$ at a point of $\mathcal D_\zeta$.
    Along that piece, let
    %
    \begin{align*}
        E_r:=\cos\theta\,E_1+\sin\theta\,E_2,
        \qquad
        E_\theta:=-\sin\theta\,E_1+\cos\theta\,E_2.
    \end{align*}
    %
    Differentiating \eqref{eq:s4-catenoid-parametrization}
    gives
    %
    \begin{align*}
        \partial_\xi\mathcal C_{\zeta,p}
        &=\tau c\sinh\xi\,E_r+\tau E_3,\\
        \partial_\theta\mathcal C_{\zeta,p}
        &=\tau c\cosh\xi\,E_\theta.
    \end{align*}
    %
    These vectors are orthogonal and nonzero. In particular,
    the first equals $\tau E_3$ at $\xi=0$.
    Their squared lengths give the exact induced metric
    %
    \begin{align*}
        \mathcal C_{\zeta,p}^*g_{\mathcal N}
        =\tau^2(1+c^2\sinh^2\xi)\,d\xi^2
          +\tau^2c^2\cosh^2\xi\,d\theta^2.
    \end{align*}
    %
    The height recovers $\xi=z/\tau$ and the direction
    of the displacement from $p$ recovers
    $\theta\in\mathbb R/2\pi\mathbb Z$. Hence each
    catenoidal parametrization is an embedding.
    The two graphs are disjoint by positivity of
    $h_\zeta$. The neck disks are pairwise disjoint,
    and each neck meets the graphs only where the formulas
    agree exactly. Since $\mathcal F$ is an embedding,
    the union defining $M_\zeta$ is therefore a smooth
    embedded surface.

    For completeness, extend $h_\zeta$ to the open set
    $\{q\in\mathbb T:\rho(q)>\tau\}$ by setting
    $h_\zeta(q):=C_\tau(r_p(q))$ when
    $\tau<r_p(q)<9\tau$. This agrees with the existing
    definition near $r_p=9\tau$. Define
    %
    \begin{align*}
        \mathcal O_\zeta
        :=\{\mathcal F(q,z):
              \rho(q)>\tau,\ |z|<h_\zeta(q)\}
        \subset\mathcal N.
    \end{align*}
    %
    Its closure is compact and lies in $|z|<z_0$ for
    large $m$. Near each waist the condition defining
    this open set is
    $r_p>\tau\cosh(z/\tau)$. Its boundary is smooth
    there because the derivative with respect to $r_p$
    of $r_p-\tau\cosh(z/\tau)$ equals one.
    Consequently $\partial\mathcal O_\zeta=M_\zeta$,
    which defines the outward unit normal $\nu_\zeta$.
    All functions, disks, and parametrized necks used above
    are preserved by the prescribed symmetries; the
    reflection $z\mapsto-z$ exchanges the graphs.
    Thus $\mathcal O_\zeta$ and $M_\zeta$ are
    $\Gamma_m$-invariant.

    The surface is compact and has no boundary. Indeed,
    the graph pieces have compact domains, and each
    remaining central portion of a neck is the image of
    a compact cylinder. All boundaries of these pieces
    lie in the exact overlaps described above.
    Topologically, $M_\zeta$ consists of two tori with
    $m^2$ disks removed from each, joined by $m^2$
    annuli. The first annulus connects the two components;
    the remaining annuli preserve connectedness. The
    boundary orientation of $\mathcal O_\zeta$ makes
    $M_\zeta$ orientable. Removing the disks gives Euler
    characteristic $-2m^2$, and attaching the annuli
    does not change it. Therefore
    %
    \begin{align*}
        2-2\operatorname{genus}(M_\zeta)=-2m^2,
    \end{align*}
    %
    proving the stated genus. The inequality
    $\tau\leq\rho_\zeta\leq\pi/m$ follows from the
    catenoidal radius $r=\tau\cosh\xi$ and the covering
    radius $\pi/m$ of the square lattice.

    We now prove the differential estimates. The height
    estimate has already been established. For each fixed
    integer $j\geq1$, the extended graph height satisfies
    %
    \begin{align*}
        |D^jh_\zeta(q)|
        \leq C_j\tau\rho(q)^{-j}
        \qquad\text{whenever }\rho(q)\geq2\tau.
    \end{align*}
    %
    On the exterior part, use
    \eqref{eq:s4-physical-derivatives}. On the support of
    derivatives of the cutoff defining $\widehat v_m$,
    the quantities $\rho$ and $\delta$ are comparable,
    so the cutoff derivative bounds give the same estimate.
    Where this cutoff is constant, the power series for
    $J_0$ applies. On the interpolation annuli, use
    \eqref{eq:s4-matching-estimate} and the derivatives
    of $C_\tau$. Finally, on $2\tau\leq r_p\leq2r_c$,
    differentiation of $C_\tau(r_p)$ gives the estimate
    directly. The same arguments give the corresponding
    bounds in local scaled H\"older norms.

    For a graph with height $w=h_\zeta$ or $w=-h_\zeta$,
    put $c:=\cos w$ and $Q:=c^2+|Dw|^2$, and denote
    its upward unit normal by $\nu^+$. Its horizontal
    part has norm $|Dw|/\sqrt Q$. Thus
    $|\nu_h|\leq C\tau/\rho$ on $\rho\geq2\tau$;
    reversing the normal on the lower graph does not
    change this estimate. We also give the computation
    for its second fundamental form. In the coordinates
    $(s,t,z)$ on the band, the nonzero connection
    coefficients are
    %
    \begin{align*}
        \Gamma^z_{ij}=\sin z\cos z\,\delta_{ij},
        \qquad
        \Gamma^i_{jz}=\Gamma^i_{zj}
            =-\tan z\,\delta^i_j
        \qquad(i,j\in\{1,2\}).
    \end{align*}
    %
    Here and in the following computation, coordinate
    vectors on the parameter domain denote their images
    under $d\mathcal F$. With graph tangent vectors
    $T_i:=\partial_i+w_i\partial_z$, these identities give
    %
    \begin{align*}
        \nabla^{\mathcal N}_{T_i}T_j
        =(w_{ij}+\sin w\cos w\,\delta_{ij})\partial_z
          -\tan w\,(w_j\partial_i+w_i\partial_j).
    \end{align*}
    %
    The normal formula for a height graph gives
    $\langle\partial_z,\nu^+\rangle=c/\sqrt Q$ and
    $\langle\partial_i,\nu^+\rangle=-cw_i/\sqrt Q$.
    Taking the inner product of the preceding identity
    with $\nu^+$ therefore gives
    %
    \begin{align*}
        \bigl\langle\nabla^{\mathcal N}_{T_i}T_j,\nu^+\bigr\rangle
        =\frac{c}{\sqrt Q}
          \bigl(w_{ij}+\sin w\cos w\,\delta_{ij}
                       +2\tan w\,w_iw_j\bigr)
        \qquad(i,j\in\{1,2\}).
    \end{align*}
    %
    The graph metric has coefficients
    $c^2\delta_{ij}+w_iw_j$. On $\rho\geq2\tau$ it is
    uniformly comparable to the flat metric, because
    $|w|\to0$ and $|Dw|\leq C$. The preceding formula
    consequently implies
    %
    \begin{align*}
        |S_{\jmath_\zeta}^{\nu_\zeta}|_{g_\zeta}
        &\leq C\bigl(|D^2w|
                    +|w|(1+|Dw|^2)\bigr)\\
        &\leq C\bigl(\tau\rho^{-2}+\tau m^2\bigr)
         \leq C\tau\rho^{-2},
    \end{align*}
    %
    where the last inequality uses $\rho\leq\pi/m$.

    It remains to obtain these bounds where
    $\tau\leq\rho\leq2\tau$. There the surface is exactly
    one of the parametrized catenoids and $|\xi|$ is
    bounded. Its outward unit normal is
    %
    \begin{align*}
        \nu_\zeta
        =\frac{-E_r+c\sinh\xi\,E_3}
               {\sqrt{1+c^2\sinh^2\xi}}.
    \end{align*}
    %
    Since $c\geq\cos z_0$, its horizontal part satisfies
    %
    \begin{align*}
        |(\nu_\zeta)_h|
        =\frac1{\sqrt{1+c^2\sinh^2\xi}}
        \leq\frac{\sec z_0}{\cosh\xi}
        =\sec z_0\frac{\tau}{\rho}.
    \end{align*}
    %
    In coordinates obtained by multiplying the three
    parameters $(q-p,z)$ by $\tau^{-1}$, the ambient
    metric, after multiplication by $\tau^{-2}$, is
    %
    \begin{align*}
        dz'^2+\cos^2(\tau z')\,(dq_1'^2+dq_2'^2).
    \end{align*}
    %
    On bounded coordinate sets it converges with all
    derivatives to the Euclidean metric. The parametrized
    surface in these coordinates is the fixed catenoid
    $(q',z')=(\cosh\xi(\cos\theta,\sin\theta),\xi)$.
    Its second fundamental form in the rescaled metrics
    is therefore uniformly bounded on $|\xi|\leq
    \operatorname{arcosh}2$. Returning to the original
    scale gives $|S_{\jmath_\zeta}^{\nu_\zeta}|\leq C/\tau
    \leq4C\tau/\rho^2$. This completes
    \eqref{eq:s4-initial-geometry-bounds}.

    These computations also give the asserted local
    coordinate estimates. On any fixed region
    $\rho/\tau\leq B$, with $B>1$, use the catenoidal
    coordinates and rescale by $\tau^{-1}$; rescaling
    instead by $\rho^{-1}$ changes lengths by a factor
    between $B^{-1}$ and one. On regions where
    $\rho/\tau$ is larger, use the height graph over a
    flat disk of radius a fixed fraction of $\rho$.
    After dividing lengths by $\rho$, the bounds for
    $D^jh_\zeta$ above give uniform ellipticity and bounds
    for any fixed number of derivatives of the metric
    coefficients. The function $\rho$ itself is never
    differentiated in this argument: the rescaling uses
    its value at the centre of each coordinate disk.

    Finally, we prove the stated lower bound for the
    injectivity radius. We first record a bound on the
    standard Euclidean catenoid. Its horizontal projection
    to $\mathbb R^2\smallsetminus\{0\}$ is length
    nonincreasing and induces an isomorphism on fundamental
    groups. A based loop of length $L$ projects into the
    Euclidean disk of radius $L/2$ about its base point:
    each point of the loop can be reached from the base
    point by one of its two arcs of length at most $L/2$.
    If the base point has horizontal radius $r$ and
    $L<2r$, that disk misses the origin, so the loop is
    contractible. The complete catenoid has nonpositive
    Gaussian curvature. Its simply connected covering
    surface therefore has no cut or conjugate points,
    and the lower bound $2r$ for lengths of
    noncontractible based loops implies that its
    injectivity radius at the base point is at least $r$.

    Fix a large constant $B>4$. On the initial surface
    with $\rho\leq4B\tau$, the parametrization is exactly
    catenoidal for all sufficiently large $m$.
    After rescaling by $\tau^{-1}$, its metric converges
    smoothly to the Euclidean catenoid metric on this
    fixed compact portion. If $\rho(x)\leq B\tau$,
    any path from $x$ to $\rho=4B\tau$ has length at
    least $\cos z_0(4B\tau-\rho(x))$. Indeed, for
    every tangent vector to the band the length of its
    projection under $\pi$ is at most $\sec z_0$ times
    its original length, and distance to the lattice is
    a $1$-Lipschitz function on $\mathbb T$.
    In particular, paths of length at most $\rho(x)/2$
    stay in this catenoidal portion. Smooth dependence of
    the geodesic equation and its differential on the
    metric, together with the catenoid injectivity bound,
    then gives
    $\operatorname{inj}_{g_\zeta}(x)\geq\rho(x)/2$
    for large $m$, uniformly on $\rho(x)\leq B\tau$.
    For clarity, injectivity here follows as well from
    this smooth convergence: a sequence of distinct
    vectors in the closed balls of half the limiting
    injectivity radius with equal exponential images
    would have convergent subsequences. Distinct limits
    contradict injectivity of the limiting exponential
    map; equal limits contradict its nonsingular
    differential and uniform local invertibility.

    If $\rho(x)\geq B\tau$, consider the flat disk
    $B_{\mathbb T}(\pi(x),\rho(x)/2)$. At every point
    of this disk the distance to $L_m$ is at least
    $\rho(x)/2>\tau$, so the sheet containing $x$ is
    a single graph over the whole disk. In coordinates
    rescaled by $\rho(x)^{-1}$, its metric differs from
    the Euclidean metric in $C^2$ norm by
    $O(B^{-2})+o_m(1)$. This follows from the bounds for
    the first three derivatives of $h_\zeta$ and the
    metric formula $\cos^2h_\zeta\,g_{\mathbb T}
    +dh_\zeta\otimes dh_\zeta$.
    A path leaving this graph neighborhood has length
    at least $\cos z_0\rho(x)/2>\rho(x)/3$.
    Therefore the exponential map on tangent vectors
    of length at most $\rho(x)/4$ stays within this
    neighborhood. After the same rescaling, it is
    $C^1$ close to the Euclidean exponential map.
    Choosing $B$ large and then $m$ large makes it
    injective and nonsingular on the tangent ball of
    radius $\rho(x)/5$. This proves
    \eqref{eq:s4-scaled-injectivity} in the remaining
    region and completes the proof.
\end{proof}

\subsection{Weighted norms and a fixed domain for the corrections}

The norm of the second fundamental form near a waist is of
order $\tau_\zeta^{-1}$. The estimates in
Lemma~\ref{lem:s4-initial-geometry} are uniform only after
distances near $x\in M_\zeta$ are divided by $\rho_\zeta(x)$.
We use that same scale to measure derivatives of the functions
that will deform $M_\zeta$. We also specify how to regard
functions on the different surfaces $M_\zeta$ as functions on
the single surface $M_0$. These two definitions are needed for
the linear estimates below and for the fixed-point argument
in Section~\ref{sec:s4-coupled}.

Define the positive function $f_\zeta:M_\zeta\to(0,\infty)$ by
%
\begin{align}
    \label{eq:s4-weight}
    f_\zeta(x)=\max\{\rho_\zeta(x)^{3/2},
    \tau_\zeta^{(1-\alpha)/2}\rho_\zeta(x)^{1/2}\}.
\end{align}
%
The weight $f_\zeta$ specifies the allowed size of a normal
displacement, apart from its overall norm. The two terms agree when
$\rho_\zeta=\tau_\zeta^{(1-\alpha)/2}$. Thus the weight is
$\tau_\zeta^{(1-\alpha)/2}\rho_\zeta^{1/2}$ nearer the
waist and $\rho_\zeta^{3/2}$ farther from it.
Here $\rho_\zeta$ is the function in
\eqref{eq:s4-rho-on-double}. For every continuous positive
function $w:M_\zeta\to(0,\infty)$,
every nonnegative integer $k$, and every function
$u\in C^{k,\beta}(M_\zeta)$, define
%
\begin{align}
    \label{eq:s4-weighted-norm}
    \|u\|_{C^{k,\beta}_\zeta(w)}
    =\sup_{x\in M_\zeta}w(x)^{-1}
    \|u\|_{C^{k,\beta}
    (B_{g_\zeta}(x,\rho_\zeta(x)/100),\,
    \rho_\zeta(x)^{-2}g_\zeta)}.
\end{align}
%
For each centre $x$, the factor $\rho_\zeta(x)^{-2}$ is a
constant on the whole ball in this formula. In particular, it
does not mean that the metric is multiplied by the variable
function $\rho_\zeta^{-2}$. In the rescaled metric this ball
has radius $1/100$.
The $C^{k,\beta}$ norm is the sum of the supremum norms of
the covariant derivatives of orders $0,\ldots,k$ and the
$\beta$-H\"older seminorm of the derivative of order $k$.
For that seminorm we compare tensors by parallel transport
along the unique short minimizing geodesic and divide by the
$\beta$th power of the distance in the rescaled metric.
Such geodesics are well defined for the pairs of points in
these balls. Indeed, $\rho_\zeta$ is Lipschitz with constant
at most $\sec z_0$, because $|d\pi|\leq\sec z_0$, and
\eqref{eq:s4-scaled-injectivity} consequently gives an
injectivity radius larger than the diameter of the ball at
every point of the ball.
The functions $\rho_\zeta$ and $f_\zeta$ need only be continuous;
no differentiation of the distance function at a cell boundary
is involved in this definition.
In particular, a bound $\|u\|_{C^{k,\beta}_\zeta(w)}\leq K$
implies
%
\begin{align*}
    |\nabla^j u(x)|_{g_\zeta}
    \leq K w(x)\rho_\zeta(x)^{-j}
    \qquad\forall\,x\in M_\zeta,\quad 0\leq j\leq k,
\end{align*}
%
as well as the corresponding H\"older estimates on the indicated
balls. Here $\nabla$ is the Levi-Civita connection of $g_\zeta$.
In particular, two derivatives of a displacement of size
$f_\zeta$ have size $f_\zeta/\rho_\zeta^2$.
The potential term in the scalar Jacobi operator is bounded
by $C\rho_\zeta^{-2}$, by
\eqref{eq:s4-initial-geometry-bounds} and the bounded curvature
of $\mathcal N$. Thus its second-order and zeroth-order terms
both send this displacement scale to $f_\zeta/\rho_\zeta^2$.
We therefore use the following domain and target spaces:
%
\begin{align}
    \label{eq:s4-XY}
    X_\zeta=C^{2,\beta}_{\mathrm{inv}}(M_\zeta),\qquad
    Y_\zeta=C^{0,\beta}_{\mathrm{inv}}(M_\zeta),\qquad
    \|u\|_{X_\zeta}=\|u\|_{C^{2,\beta}_\zeta(f_\zeta)},\quad
    \|E\|_{Y_\zeta}=\|E\|_{C^{0,\beta}_\zeta(f_\zeta/\rho_\zeta^2)}.
\end{align}
%
Here $C^{k,\beta}_{\mathrm{inv}}(M_\zeta)$ consists of the
$C^{k,\beta}$ functions invariant under every element of
$\Gamma_m$. For fixed $m$ and $\zeta$, the weights are
positive on the compact surface $M_\zeta$, so these norms
make $X_\zeta$ and $Y_\zeta$ Banach spaces.

We explain the relation with the norms used in the cited
linear construction. For a function $u$ on a parametrized neck,
write $u(\xi,\theta):=u(\mathcal C_{\zeta,p}(\xi,\theta))$.
Its first angular harmonic is
%
\begin{align*}
    (\Pi_1u)(\xi,\theta)
    :=\frac{\cos\theta}{\pi}
       \int_0^{2\pi}u(\xi,\vartheta)\cos\vartheta\,d\vartheta
      +\frac{\sin\theta}{\pi}
       \int_0^{2\pi}u(\xi,\vartheta)\sin\vartheta\,d\vartheta.
\end{align*}
%
The half-turn about the lattice point belongs to the symmetry
group and preserves the axial coordinate. Thus an invariant
function satisfies $u(\xi,\theta+\pi)=u(\xi,\theta)$.
Changing variables by $\vartheta\mapsto\vartheta+\pi$ in
each integral shows that $\Pi_1u=0$.
Thus the additional first-harmonic summand in Definition~4.2
of \cite{KapouleasMcGrathLD} vanishes identically, and the
norms above are uniformly equivalent to the restrictions of its
two-weight norms with $\gamma=3/2$ and $\gamma'=1/2$.
The possible change in the fixed radius used in the local norm
changes only a constant independent of $m$ and $\zeta$.

To use the prescribed error $k_m$ in an equation on $M_\zeta$,
define
%
\begin{align}
    \label{eq:s4-J-lift}
    \mathcal{J}_\zeta=k_m\circ\pi\big|_{M_\zeta}.
\end{align}
%
This is a function on $M_\zeta$. It is smooth, vanishes on
the necks, and on each of the two sheets equals $k_m$ at
the projected point. We use the outward unit normal on both
sheets. The isometry $z\mapsto-z$ exchanges the sheets and
preserves their outward normals as vector fields under its
differential; hence their scalar mean curvatures have the same
transformation rule. This explains why the same function,
rather than its negative, is used on the lower sheet.
In the notation of \cite[Notation~3.25]{KapouleasMcGrathLD},
it is exactly $J_{M_\zeta}(k_m,k_m)$.

The surfaces themselves depend on $\zeta$, so a function on
$M_\zeta$ does not have a fixed domain when $\zeta$ varies.
We remove this difficulty by constructing explicit
diffeomorphisms $F_\zeta:M_0\to M_\zeta$.
The subscript $0$ here refers to $\zeta=0$; the integer $m$
remains fixed. Near a waist these maps rescale the catenoidal
axial coordinate, while away from the necks they preserve the
base point on $\mathbb T$.
Set, for every $\zeta\in[-1,1]$,
%
\begin{align*}
    r_0=\tfrac12\tau_0^\alpha,\qquad
    s_0=\operatorname{arcosh}(r_0/\tau_0),\qquad
    s_\zeta=\operatorname{arcosh}(\tau_\zeta^{\alpha-1}/2),
    \qquad a_\zeta=s_\zeta/s_0.
\end{align*}
%
For sufficiently large $m$, both arguments of
$\operatorname{arcosh}$ exceed one, so $s_0,s_\zeta>0$.
The circles $|\xi|=s_0$ on the reference neck have radius
$r_0$, and the circles $|\xi|=s_\zeta$ on the target neck
have radius $\tau_\zeta^\alpha/2$.
Choose a fixed smooth cutoff $\chi_0:[0,\infty)\to[0,1]$
equal to one on $[0,1]$ and zero on $[2,\infty)$, and for
$r\geq\tau_0$ set
%
\begin{align}
    \label{eq:s4-radius-identification}
    \begin{split}
        A_\zeta(r)&=\tau_\zeta
        \cosh\bigl(a_\zeta\operatorname{arcosh}(r/\tau_0)\bigr),\\
        R_\zeta(r)&=r+\chi_0(r/r_0)(A_\zeta(r)-r).
    \end{split}
\end{align}
%
The function $A_\zeta$ is the change in radius induced by
the linear change $\xi\mapsto a_\zeta\xi$ in the axial
coordinate. The function $R_\zeta$ makes this change equal
to the identity once $r\geq2r_0$.
Near each neck define $F_\zeta:M_0\to M_\zeta$ by preserving
the polar angle and the choice of upper or lower sheet,
replacing the base radius $r$ by $R_\zeta(r)$, and taking the
corresponding point on $M_\zeta$. Outside the radius-$2r_0$
disks it preserves the base point and sheet sign.
On the region $\rho_0\leq r_0$, the definition is equivalently
the smooth formula
%
\begin{align}
    \label{eq:s4-core-identification}
    F_\zeta(\mathcal{C}_{0,p}(\xi,\theta))
    =\mathcal{C}_{\zeta,p}(a_\zeta\xi,\theta),
    \qquad |\xi|\leq s_0.
\end{align}
%

\begin{lemma}[Comparison of the surfaces and their function spaces]
    \label{lem:s4-identifications}
    After taking $\alpha$ sufficiently small in terms of
    $\chi_0$, the maps $F_\zeta$ are symmetry-preserving smooth
    diffeomorphisms for all sufficiently large $m$. They depend
    smoothly on $\zeta$ and $F_0$ is the identity. There is a
    constant $C_F$, independent of $m,\zeta$, such that
    %
    \begin{align}
        \label{eq:s4-norm-identification}
        \begin{split}
            C_F^{-1}\|u\|_{X_\zeta}
            &\leq\|u\circ F_\zeta\|_{X_0}
            \leq C_F\|u\|_{X_\zeta},\\
            C_F^{-1}\|E\|_{Y_\zeta}
            &\leq\|E\circ F_\zeta\|_{Y_0}
            \leq C_F\|E\|_{Y_\zeta}.
        \end{split}
    \end{align}
    %
    The inequalities in the first line hold for every
    $u\in X_\zeta$, and those in the second line hold
    for every $E\in Y_\zeta$.
\end{lemma}

\begin{proof}
    Set $L_0:=|\log\tau_0|$. The expansion
    $\operatorname{arcosh}x=\log(2x)+O(x^{-2})$ gives
    %
    \begin{align*}
        s_0&=(1-\alpha)L_0+O(\tau_0^{2(1-\alpha)}),\\
        s_\zeta-s_0
        &=-(1-\alpha)\zeta+O(\tau_0^{2(1-\alpha)}),\\
        a_\zeta-1&=O(L_0^{-1}),
    \end{align*}
    %
    uniformly for $\zeta\in[-1,1]$. If $r\in[r_0,2r_0]$
    and $\xi:=\operatorname{arcosh}(r/\tau_0)$, then
    $\xi=s_0+\log(r/r_0)+o_m(1)$. Therefore
    %
    \begin{align*}
        \log\frac{A_\zeta(r)}r
        &=\zeta+\log\cosh(a_\zeta\xi)-\log\cosh\xi\\
        &=\zeta+(a_\zeta-1)\xi+o_m(1)
          =\alpha\zeta+O(L_0^{-1}).
    \end{align*}
    %
    Here $o_m(1)$ denotes a quantity tending uniformly to zero
    as $m\to\infty$. Differentiation gives
    %
    \begin{align*}
        A_\zeta'(r)
        &=\frac{\tau_\zeta}{\tau_0}
             a_\zeta\frac{\sinh(a_\zeta\xi)}{\sinh\xi}\\
        &=a_\zeta\frac{A_\zeta(r)}r
             \frac{\tanh(a_\zeta\xi)}{\tanh\xi}.
    \end{align*}
    %
    Consequently, on $[r_0,2r_0]$,
    %
    \begin{align*}
        A_\zeta(r)/r-1=O(\alpha)+o_m(1),
        \qquad A_\zeta'(r)-1=O(\alpha)+o_m(1),
    \end{align*}
    %
    with constants uniform for $\alpha$ in a fixed small
    interval. The derivative of the interpolating radius is
    %
    \begin{align*}
        R_\zeta'(r)
        =1+\chi_0(r/r_0)(A_\zeta'(r)-1)
           +\frac{\chi_0'(r/r_0)}{r_0}(A_\zeta(r)-r).
    \end{align*}
    %
    Choose $\alpha$ so small, in terms of the fixed cutoff,
    that the sum of the $O(\alpha)$ terms has absolute value
    less than $1/4$. Then choose $m$ large enough that the
    remaining terms have absolute value less than $1/4$.
    This gives $R_\zeta'>1/2$ on the transition annulus.
    Outside it, $R_\zeta$ equals either the increasing function
    $A_\zeta$ or the identity. Since
    $R_\zeta(\tau_0)=\tau_\zeta$, it maps the allowed radii
    bijectively onto the corresponding radii of $M_\zeta$.

    Both the transition annulus and its image lie in the
    portions defined by \eqref{eq:s4-catenoid-parametrization}.
    Indeed, their outer radii satisfy
    %
    \begin{align*}
        \max\{2r_0,R_\zeta(2r_0)\}
        \leq\tau_0^\alpha(1+O(\alpha)+o_m(1))
        <2\tau_\zeta^\alpha.
    \end{align*}
    %
    Smoothness at the waist follows from
    \eqref{eq:s4-core-identification}, which is the map
    $(\xi,\theta)\mapsto(a_\zeta\xi,\theta)$ between regular
    catenoidal coordinates. Elsewhere the graph descriptions
    and $R_\zeta'>0$ give smoothness and a smooth inverse.
    All formulas depend smoothly on $\zeta$, and radial
    changes of coordinates commute with the rotations,
    reflections, and factor interchange used in $\Gamma_m$.
    The sheet signs are preserved and height reflection
    exchanges them on both sides. Thus $F_\zeta$ commutes
    with every element of $\Gamma_m$. When $\zeta=0$,
    $a_0=1$ and $R_0(r)=r$, so $F_0$ is the identity.

    For the norm estimates, the ratios
    $\rho_\zeta(F_\zeta(x))/\rho_0(x)$ are bounded above
    and below uniformly. On $\rho_0\leq r_0$, this follows
    from
    %
    \begin{align*}
        \frac{\rho_\zeta(F_\zeta(x))}{\rho_0(x)}
        =e^\zeta\frac{\cosh(a_\zeta\xi)}{\cosh\xi},
        \qquad |\xi|\leq s_0,
    \end{align*}
    %
    because $|(a_\zeta-1)\xi|$ is uniformly bounded.
    The transition-annulus estimates give the same conclusion
    there, and outside those annuli the radii agree exactly.
    Since $\tau_\zeta/\tau_0\in[e^{-1},e]$, the ratios of
    the weights $f_\zeta\circ F_\zeta$ and $f_0$ are also
    bounded above and below uniformly.

    Differentiating the formulas for $A_\zeta$ and $R_\zeta$
    in rescaled graph coordinates gives uniform derivative
    bounds of every fixed order on the transition annuli.
    In the catenoidal coordinates at the waist this follows
    from the linear map in $\xi$ and the uniform local
    geometry of Lemma~\ref{lem:s4-initial-geometry}.
    The same statements hold for the inverses. Hence the
    pullback metrics, after rescaling by the corresponding
    radii, are uniformly comparable with uniformly bounded
    derivatives. The chain rule then compares the derivative
    terms in the local H\"older norms. Distances are uniformly
    comparable as well, giving the comparison of the H\"older
    seminorms. An image of one of the balls in
    \eqref{eq:s4-weighted-norm} is covered by a uniformly
    bounded number of such balls in the target. Taking the
    supremum over the centers proves
    \eqref{eq:s4-norm-identification}.
\end{proof}

\subsection{The scalar linear equation and the remaining corrections}

We now construct a first correction of the scalar mean
curvature within $\mathcal N$, and estimate every further
correction allowed in Sections~\ref{sec:s4-transverse}
and~\ref{sec:s4-coupled}. We first solve the linearized
equation with an additional term $c\mathcal J_\zeta$, where
$c\in\mathbb R$ is part of the output. We apply that solution
operator to the initial mean-curvature error to obtain a
specified function $u_\zeta$. Finally, we allow an additional
function $v$ on $M_0$ and estimate the total displacement
on $M_\zeta$.

Write $H_\zeta:=H_{\jmath_\zeta,\nu_\zeta}$ for the scalar
mean curvature of the inclusion into $\mathcal N$, measured
with its outward unit normal. Its linearization under a normal
variation within $\mathcal N$ is the operator
$\mathcal L_\zeta:C^{2,\beta}(M_\zeta)\to C^{0,\beta}(M_\zeta)$
given by
%
\begin{align}
    \label{eq:s4-L-zeta}
    \mathcal L_\zeta u
    :=\Delta_{g_\zeta}u+
       \bigl(|S_{\jmath_\zeta}^{\nu_\zeta}|_{g_\zeta}^2
       +\operatorname{Ric}_{g_{\mathcal N}}
                         (\nu_\zeta,\nu_\zeta)\bigr)u
    \qquad\forall\,u\in C^{2,\beta}(M_\zeta).
\end{align}
%
No minimality assumption on $M_\zeta$ is needed for this
linearization formula. We use the construction of
\cite{KapouleasMcGrathLD} to solve
$\mathcal L_\zeta u=E$ up to a multiple of
$\mathcal J_\zeta$. The extra term allows the linear construction
to adjust the common value of its auxiliary solution on the
central torus at the lattice points while maintaining the
weighted estimates. The assertion below is
a right inverse for the augmented map
$(u,c)\mapsto\mathcal L_\zeta u-c\mathcal J_\zeta$;
it does not assert invertibility of $\mathcal L_\zeta$ alone.
The total coefficient of $\mathcal J_\zeta$ in the remaining
equation will be set equal to zero by choosing $\zeta$ in
Section~\ref{sec:s4-coupled}.

The estimates are needed not just for one correction, but for
all scalar functions considered in that argument. Using the
maps $F_\zeta$ to express those functions on $M_0$, define
%
\begin{align}
    \label{eq:s4-scalar-ball}
    \mathcal B_m:=\{(v,\zeta)\in X_0\times[-1,1]:
                         \|v\|_{X_0}\leq\tau_0^{1+\alpha}\}.
\end{align}
%
The integer $m$ is fixed when this set is used in a fixed-point
argument. The function $v$ always has domain $M_0$;
$v\circ F_\zeta^{-1}$ has domain $M_\zeta$.
The total displacement, denoted by $\phi$ in the proposition,
is the sum of this function and the preliminary correction
$u_\zeta$. Thus $v$ and $\phi$ are not independent unknowns.

\begin{proposition}[Uniform scalar estimates]
    \label{prop:s4-scalar-input}
    For all sufficiently large $m$, the following hold uniformly
    for $|\zeta|\leq1$.

    \begin{enumerate}
        \item There is a linear map
            $\mathscr{R}_\zeta:Y_\zeta\to X_\zeta\times\mathbb{R}$,
            $\mathscr{R}_\zeta E=(u_E,c_E)$, such that
            %
            \begin{align}
                \label{eq:s4-scalar-inverse}
                \mathcal{L}_\zeta u_E=E+c_E\mathcal{J}_\zeta,
                \qquad
                \|u_E\|_{X_\zeta}+|c_E|\|k_m\|_{C^{0,\beta}(\mathbb{T})}
                \leq C\delta^{-4-2\beta}\|E\|_{Y_\zeta}.
            \end{align}
            %
            In particular,
            %
            \begin{align}
                \label{eq:s4-coefficient-bound}
                |c_E|\leq C\delta^{-2-\beta}\|E\|_{Y_\zeta}.
            \end{align}
            %
            The solution operators can be chosen with the following continuity
            property. Fix $m$ and $0<\beta'<\beta$. Suppose
            $\zeta_j\to\zeta$, $E_j\in Y_{\zeta_j}$ have uniformly
            bounded $Y_{\zeta_j}$ norms, $E\in Y_\zeta$, and
            $E_j\circ F_{\zeta_j}\to E\circ F_\zeta$ in
            $C^{0,\beta'}(M_0)$. If
            $(u_j,c_j)=\mathscr{R}_{\zeta_j}E_j$ and
            $(u,c)=\mathscr{R}_\zeta E$, then
            $u_j\circ F_{\zeta_j}\to u\circ F_\zeta$ in
            $C^{2,\beta'}(M_0)$ and $c_j\to c$.

        \item Define the scalar error $D_\zeta$ and the preliminary
            correction $u_\zeta$ by
            %
            \begin{align}
                \label{eq:s4-initial-correction}
                D_\zeta=H_\zeta+\tau_\zeta\zeta\mathcal{J}_\zeta,
                \qquad
                (u_\zeta,c_{H,\zeta})=-\mathscr{R}_\zeta D_\zeta.
            \end{align}
            %
            Then
            %
            \begin{align}
                \|D_\zeta\|_{Y_\zeta}&\leq C\tau_\zeta^{1+\alpha/3},
                \label{eq:s4-H-error}\\
                \mathcal{L}_\zeta u_\zeta+H_\zeta
                &=(-\tau_\zeta\zeta+c_{H,\zeta})\mathcal{J}_\zeta.
                \label{eq:s4-linear-cancellation}
            \end{align}
            %

        \item For every $(v,\zeta)\in\mathcal{B}_m$, define
            the total scalar displacement $\phi:M_\zeta\to\mathbb R$ by
            %
            \begin{align}
                \label{eq:s4-full-scalar-correction}
                \phi=u_\zeta+v\circ F_\zeta^{-1}.
            \end{align}
            %
            Then
            %
            \begin{align}
                \label{eq:s4-full-ball-bound}
                \|\phi\|_{X_\zeta}
                +|c_{H,\zeta}|\|k_m\|_{C^{0,\beta}(\mathbb{T})}
                \leq\tau_\zeta^{1+\alpha/4}.
            \end{align}
            %
            The map
            %
            \begin{align*}
                G_{\zeta,\phi}:M_\zeta&\rightarrow\mathcal N,\\
                G_{\zeta,\phi}(x)
                &:=\exp_{\jmath_\zeta(x)}^{\mathcal N}
                    (\phi(x)\nu_\zeta(x))
                \qquad\forall\,x\in M_\zeta
            \end{align*}
            %
            is well defined and is an embedding of class
            $C^{2,\beta}$. The same conclusion holds with
            $\phi$ replaced by $t\phi$ for every $t\in[0,1]$.
            Write
            $M_{\zeta,\phi}:=G_{\zeta,\phi}(M_\zeta)$.
            Choose the unit normal $\nu_{\zeta,\phi}$ along
            $G_{\zeta,\phi}$ by continuity from $\nu_\zeta$
            along the family $G_{\zeta,t\phi}$, $0\leq t\leq1$,
            and set
            $H_{\zeta,\phi}:=H_{G_{\zeta,\phi},\nu_{\zeta,\phi}}$.
            This is a function on the fixed domain $M_\zeta$,
            computed using the induced metric
            $G_{\zeta,\phi}^*g_{\mathcal N}$.
            The nonlinear remainder
            %
            \begin{align}
                \label{eq:s4-Q-definition}
                Q_\zeta(\phi)=H_{\zeta,\phi}-H_\zeta-\mathcal{L}_\zeta\phi
            \end{align}
            %
            satisfies
            %
            \begin{align}
                \label{eq:s4-scalar-quadratic}
                \|Q_\zeta(\phi)\|_{Y_\zeta}
                \leq C\tau_\zeta^{-\alpha/2}\|\phi\|_{X_\zeta}^2
                \leq C\tau_\zeta^2.
            \end{align}
            %
            In particular, on every member of this iteration set,
            %
            \begin{align}
                \label{eq:s4-scalar-two-jet}
                |\nabla^j\phi(x)|_{g_\zeta}\leq
                C\tau_\zeta^{1+\alpha/4}f_\zeta(x)\rho_\zeta(x)^{-j},
                \qquad x\in M_\zeta,\quad j\in\{0,1,2\},
            \end{align}
            %
            with the corresponding local scaled H\"older estimates.
    \end{enumerate}
\end{proposition}

\begin{proof}
    The analytic statements we use from
    \cite{KapouleasMcGrathLD} are the linear construction in
    Proposition~4.18 and Corollary~4.21, the initial
    mean-curvature estimate in Lemma~4.6, and the normal-graph
    estimate in Lemma~5.1. We first verify their geometric
    hypotheses and explain the restriction to invariant
    functions. We then derive the estimates in the present
    notation. The background three-manifold is the fixed open
    band $\mathcal{N}$, and the base is the closed two-sided minimal
    torus. Completeness of the background is not required in its
    Convention~2.1. Its Fermi map agrees with $\mathcal{F}$ because
    the curves $z\mapsto\mathcal{F}(q,z)$ are unit-speed geodesics
    normal to the central torus.

    We first check the invertibility required on the central torus.
    On the full torus, the eigenvalues of $L_{\mathbb{T}}$ on
    $e^{i\sqrt2(js+kt)}$ are $2-2(j^2+k^2)$.
    Indeed, translating $s$ by $\ell/m$ multiplies this
    Fourier mode by $e^{2\pi i j/m}$, while translating $t$
    by $\ell/m$ multiplies it by $e^{2\pi i k/m}$.
    Invariance under both translations therefore forces
    $j,k\in m\mathbb Z$.
    For $m\geq2$ none of these eigenvalues is zero. Thus the
    $\Gamma_m^0$-invariant kernel is trivial.
    More explicitly, the full kernel is spanned by
    $\cos(\sqrt2s)$, $\sin(\sqrt2s)$,
    $\cos(\sqrt2t)$, and $\sin(\sqrt2t)$.
    Every invariant source is orthogonal to these functions.
    The inverse of the fixed operator on their orthogonal
    complement has a fixed $C^{0,\beta}\to C^{2,\beta}$
    bound, and uniqueness makes it commute with all the
    isometries under consideration. This supplies the
    invariant version of Assumption~4.1 of the cited paper;
    the unrestricted kernel is not being assumed trivial.

    We next verify that the values and differentials prescribed
    at the lattice points can be corrected with the required bounds.
    Before imposing symmetry, at each $p$ take the cutoff in
    \eqref{eq:s4-substitute} times the three functions
    %
    \begin{align*}
        J_0(\sqrt2r),\qquad
        \sqrt2J_1(\sqrt2r)\cos\theta,\qquad
        \sqrt2J_1(\sqrt2r)\sin\theta,
    \end{align*}
    %
    where
    $J_1(x)=\sum_{j\geq0}(-1)^j(x/2)^{2j+1}/(j!(j+1)!)$.
    Write $x_1=s-s(p)$, $x_2=t-t(p)$ in the flat disk,
    $r=(x_1^2+x_2^2)^{1/2}$, and
    $(x_1,x_2)=r(\cos\theta,\sin\theta)$.
    The Bessel equations give $(\Delta+2)u=0$ for all
    three functions. Their expansions are
    %
    \begin{align*}
        J_0(\sqrt2r)&=1-r^2/2+O(r^4),\\
        \sqrt2J_1(\sqrt2r)\cos\theta
           &=x_1+O(r^3),\\
        \sqrt2J_1(\sqrt2r)\sin\theta
           &=x_2+O(r^3).
    \end{align*}
    %
    Hence their values and differentials at $p$ form a basis
    for $\mathbb R\oplus T_p^*\mathbb T$.
    Their supports and the supports of their Jacobi images obey
    Assumption~3.11 of the cited paper.
    If the prescribed value and differential have norm at
    most one, the corresponding combination of these three
    cutoff functions has $C^{2,\beta}$ norm at most
    $C\delta^{-2-\beta}$. To see this, every derivative of
    a cutoff contributes one factor $\delta^{-1}$;
    taking a H\"older quotient contributes at most
    $\delta^{-\beta}$. The uncut Bessel functions have
    uniformly bounded derivatives on the fixed small disks.
    The difference from each prescribed affine jet is $O(r^2)$,
    with all fixed scaled derivatives, by the displayed Bessel
    series. Disjoint supports give the same estimates for arbitrary
    collections of jets, with the maximum norm on their
    coefficients; there is no factor equal to the number of necks.
    Invariance under the reflections removes the linear jets, and
    lattice transitivity identifies all constant jets. The invariant
    space is therefore exactly $\mathbb{R}\widehat v_m$, with
    Jacobi image $\mathbb{R} k_m$.

    We check the separation and smallness assumptions with
    the scales already chosen.
    Convention~3.8 holds: the $9\delta$ disks are disjoint because
    $18\delta<\ell/m$; the Fermi injectivity radius has a fixed
    positive lower bound; and the Dirichlet operator $\Delta+2$
    on every disk of radius $\delta'\leq3\delta$ is invertible.
    For the last assertion its first $-\Delta$ eigenvalue is
    $j_{0,1}^2/(\delta')^2>2$ for large $m$, where $j_{0,1}$
    is the first positive zero of $J_0$. The same convergent
    radial series \eqref{eq:s4-radial-fundamental} with $a=2$,
    divided by $2\pi$, supplies a local Green function
    $P_2/(2\pi)$ with logarithmic coefficient one on these small
    flat disks, together with the required regularity.

    All logarithmic coefficients are equal to
    $\tau=\tau_\zeta$. For large $m$,
    %
    \begin{align*}
        9\tau^\alpha<\tau^{\alpha/100}<\delta,
        \qquad \tau\leq\tau^{1-\alpha/100}.
    \end{align*}
    %
    The boundary norm in Convention~3.15(iv) satisfies
    %
    \begin{align*}
        \delta^{-2}\|\psi_\zeta\|_{C^{2,\beta}(\partial B(p,\delta))}
        \leq C\tau m^{4+\beta}
        =o(\tau^{1-\alpha/9}).
    \end{align*}
    %
    On $\rho\geq\tau^\alpha$, the local derivative estimates
    and a two-point estimate for the H\"older seminorm give
    %
    \begin{align*}
        \|\psi_\zeta\|_{C^{3,\beta}}
        \leq C\tau^{1-\alpha(3+\beta)}
        =o(\tau^{8/9}).
    \end{align*}
    %
    For the two-point estimate, use a scaled local estimate when
    the points are within a fixed fraction of $\tau^\alpha$,
    and the supremum bound otherwise. If needed near the boundary,
    these local estimates are taken on the larger region
    $\rho\geq\tfrac12\tau^\alpha$.
    Finally \eqref{eq:s4-truncated-positivity} gives
    $\psi_\zeta\geq c\tau m^2\geq\tau^{1+\alpha/5}$ there.
    This verifies all parts of Convention~3.15 uniformly in
    $\zeta$. The inequalities in Convention~4.5 are precisely
    $1/4>2\alpha$ and $(1-\alpha)/2>2\alpha$, which follow
    from \eqref{eq:s4-alpha}. The same convention fixes a
    large constant separating the neck region
    $\rho\leq B\tau$ from its graphical ends. Choose that
    constant once and for all before choosing $m$.

    We now check the definition of the initial surface and
    record the parameter conditions needed in the later
    existence argument.
    The lattice, local disks and jet spaces do not move with
    $\zeta$; the base diffeomorphisms are identities. The
    normalized mismatch map is $c\mapsto c/\tau_\zeta$, so
    \eqref{eq:s4-exact-mismatch} makes the normalized mismatch
    exactly $\zeta$. The neck elevation and tilt parameters are
    zero. The bounded ratio $\tau_\zeta/\tau_0=e^\zeta$ and the
    preceding estimates give the invariant specialization of
    Assumption~5.2. Definitions \eqref{eq:s4-glued-height} and
    \eqref{eq:s4-initial-double} are Definition~3.17 of the cited
    paper with these choices. Since the supports of the Jacobi
    images $k_m$ lie
    in $\delta\leq r_p\leq2\delta$, they are disjoint from all
    regions $\rho\leq3\tau^\alpha$ for large $m$; its lifted
    substitute function is therefore exactly
    \eqref{eq:s4-J-lift}.

    With these verifications, we apply the linear construction
    and explain why its error term has only one coefficient here.
    Proposition~4.18 and Corollary~4.21 of the cited paper
    construct a linear solution operator whose equation has
    one substitute-function output on each of the two sheets.
    We use their proofs on the invariant spaces, as in the
    symmetry restriction in the proof of their Theorem~6.17.
    In particular, the torus equation in their equation~(4.13)
    can be solved explicitly within the invariant spaces.
    Given a $\Gamma_m^0$-invariant source
    $E'\in C^{0,\beta}(\mathbb T)$, let $u^{\mathrm{base}}$
    be the unique invariant solution of
    %
    \begin{align*}
        L_{\mathbb T}u^{\mathrm{base}}=E'.
    \end{align*}
    %
    Lattice transitivity makes its value at every $p\in L_m$
    equal to a common number $c$, and the coordinate reflections
    fixing $p$ imply $du^{\mathrm{base}}(p)=0$.
    Using the cutoff function already constructed, set
    %
    \begin{align*}
        \widehat v:=c\widehat v_m,\qquad
        u':=u^{\mathrm{base}}-\widehat v,\qquad
        w':=-L_{\mathbb T}\widehat v=-c k_m.
    \end{align*}
    %
    Then all these functions are invariant and satisfy
    %
    \begin{align*}
        L_{\mathbb T}u'=E'+w',\qquad
        u'(p)=0,\qquad du'(p)=0
        \qquad\forall\,p\in L_m.
    \end{align*}
    %
    The invariant torus inverse and the cutoff estimates above
    provide exactly the bounds for this step; no inverse on
    the unrestricted torus space is needed.
    The symmetry restriction is likewise legitimate at every
    other step: the base inverse just described,
    the radial cutoffs, the local neck solutions, the jet maps,
    and the sums over the lattice commute with $\Gamma_m^0$.
    If a local solution operator is not initially invariant,
    average its conjugates over the finite group fixing the
    corresponding neck. For a local group $G$ and operator
    $R$, this average is
    $|G|^{-1}\sum_{g\in G}g^*R(g^{-1})^*$.
    Since the group acts isometrically and preserves the
    radial weights, the average has no larger norm and
    satisfies the same local equation. Translating the chosen
    operators to every lattice point then gives an invariant
    global construction. The fixed-torus inverse used in that
    proof is replaced by the invariant inverse justified above;
    all other local estimates are unchanged.
    One then imposes $Z$-invariance. Height reflection exchanges
    the two exterior substitute outputs with the same sign, so
    both are $c_Ek_m$.
    The two-weight norms are uniformly equivalent to those
    in \eqref{eq:s4-XY}, because the first-harmonic projection
    vanishes. Corollary~4.21 therefore gives
    \eqref{eq:s4-scalar-inverse}, with loss
    $C\delta^{-4-2\beta}$.

    For the coefficient estimate, the fixed cutoff satisfies
    %
    \begin{align}
        \label{eq:s4-k-size}
        c\delta^{-2-\beta}\leq\|k_m\|_{C^{0,\beta}(\mathbb{T})}
        \leq C\delta^{-2-\beta}.
    \end{align}
    %
    The upper bound follows by scaling. For the lower bound,
    write $r=\delta x$ on a single annulus and let
    $\eta(x)=\eta_{2,1}(x)$. Then
    %
    \begin{align*}
        k_m(p+\delta x e_\theta)
        =\delta^{-2}\bigl(\eta''(x)+x^{-1}\eta'(x)\bigr)+O(1).
    \end{align*}
    %
    The leading function is not identically zero, since a radial
    harmonic function that is constant near each boundary circle
    cannot change from one to zero. Compare a point where this
    function is nonzero with a point in the inner region where
    $k_m=0$. Their distance is comparable to $\delta$, giving
    the required H\"older quotient. Combining
    \eqref{eq:s4-k-size} with \eqref{eq:s4-scalar-inverse}
    proves \eqref{eq:s4-coefficient-bound}.

    We next apply the initial mean-curvature estimate and
    compute the equation satisfied by $u_\zeta$.
    The two substitute terms in Lemma~4.6 of the cited paper
    are both $-\tau_\zeta\zeta k_m$, because the same height
    correction $-\tau_\zeta\zeta\widehat v_m$ is used on
    the two sheets. That lemma therefore gives
    %
    \begin{align*}
        \|H_\zeta+\tau_\zeta\zeta\mathcal J_\zeta\|_{Y_\zeta}
        \leq C\tau_\zeta^{1+\alpha/3},
    \end{align*}
    %
    which is \eqref{eq:s4-H-error}. By the linearity of
    $\mathscr R_\zeta$, the definition
    $(u_\zeta,c_{H,\zeta})=-\mathscr R_\zeta D_\zeta$
    gives
    %
    \begin{align*}
        \mathcal L_\zeta u_\zeta
        &=-D_\zeta+c_{H,\zeta}\mathcal J_\zeta\\
        &=-H_\zeta+
          (-\tau_\zeta\zeta+c_{H,\zeta})\mathcal J_\zeta.
    \end{align*}
    %
    This proves \eqref{eq:s4-linear-cancellation} with the
    stated signs.

    The solution operators may be chosen continuously in
    the precise sense stated in the proposition.
    The right inverse is the one constructed by the local
    solution operators and convergent correction series in
    Proposition~4.18 and Corollary~4.21. 
    
    We retain the normalizations in equations~(4.13)
    and~(4.16) of \cite{KapouleasMcGrathLD}.
    Each auxiliary torus solution has value and
    differential zero at every point of $L_m$, as in
    the construction of $u'$ above.
    On each catenoid, the solutions for the constant
    and first angular harmonics have both their value
    and their $\xi$-derivative equal to zero at
    $\xi=0$, the waist.
    The remaining angular modes are solved using the
    inverse of $\Delta_{\mathbb{S}^2}+2$ on functions
    with angular Fourier orders at least two, after
    the conformal change of metric given by the
    catenoid Gauss map.
    The restriction to these angular orders excludes
    every nonzero element of the kernel of
    $\Delta_{\mathbb{S}^2}+2$.
    In the invariant spaces used here, the first
    angular harmonics vanish identically.
    
    After conjugation by
    $F_\zeta$, each finite stage of that construction is jointly
    continuous in the parameter and the source with any loss from
    exponent $\beta$ to $\beta'$: the cutoffs and coefficients
    vary smoothly, the local elliptic inverses vary continuously,
    and the moving smooth pullbacks are continuous on bounded
    stronger-H\"older sets. The tails of the correction series
    are uniformly small in the $\beta$ norm, hence also in the
    $\beta'$ norm. Applying these observations first to a finite
    stage and then letting the number of stages tend to infinity
    proves the precise sequential continuity asserted above.

    It remains to estimate the total displacement $\phi$
    and its nonlinear mean-curvature error on all of
    $\mathcal B_m$.
    The preceding bounds give
    %
    \begin{align*}
        \|u_\zeta\|_{X_\zeta}
        +|c_{H,\zeta}|\|k_m\|_{C^{0,\beta}}
        \leq C\delta^{-4-2\beta}\tau_\zeta^{1+\alpha/3}.
    \end{align*}
    %
    Its ratio to $\tau_\zeta^{1+\alpha/4}$ tends to zero
    uniformly because
    $\delta^{-4-2\beta}\tau_\zeta^{\alpha/12}\to0$.
    Also \eqref{eq:s4-norm-identification} gives
    %
    \begin{align*}
        \|v\circ F_\zeta^{-1}\|_{X_\zeta}
        \leq C_F\tau_0^{1+\alpha}
        \leq C\tau_\zeta^{1+\alpha},
    \end{align*}
    %
    whose ratio to $\tau_\zeta^{1+\alpha/4}$ also tends to
    zero. In fact we have proved the stronger uniform bound
    %
    \begin{align*}
        \|\phi\|_{X_\zeta}
        +|c_{H,\zeta}|\|k_m\|_{C^{0,\beta}(\mathbb T)}
        =o\bigl(\tau_\zeta^{1+\alpha/4}\bigr).
    \end{align*}
    %
    The little-$o$ estimate is uniform for
    $(v,\zeta)\in\mathcal B_m$ as $m\to\infty$.
    This proves \eqref{eq:s4-full-ball-bound} and also absorbs
    the fixed constants in the comparison with the norms of
    \cite{KapouleasMcGrathLD}. The smallness hypothesis of
    its Lemma~5.1 is therefore satisfied. That lemma gives the embeddedness
    of $G_{\zeta,\phi}$ and the first inequality in
    \eqref{eq:s4-scalar-quadratic}. We explain the scaling
    of this estimate, since it will be used again for the
    second normal direction. On a coordinate neighborhood
    of radius comparable to $\rho=\rho_\zeta(x)$, divide
    lengths by $\rho$. The rescaled normal displacement is
    $\phi/\rho$, and its $C^{2,\beta}$ norm is at most
    $C\|\phi\|_{X_\zeta}f_\zeta(x)/\rho$.
    The local geometry established above gives uniform bounds
    for the coefficients of the mean-curvature map in these
    coordinates. Taylor's formula at zero therefore bounds
    its remainder, after subtracting its value and derivative,
    by the square of this displacement norm. Returning to
    the original metric contributes the factor $\rho^{-1}$
    to mean curvature, so
    %
    \begin{align*}
        \|Q_\zeta(\phi)\|_{C^{0,\beta}
          (B_{g_\zeta}(x,\rho/100),\rho^{-2}g_\zeta)}
        \leq C\rho^{-3}f_\zeta(x)^2\|\phi\|_{X_\zeta}^2.
    \end{align*}
    %
    Dividing by the target weight $f_\zeta(x)/\rho^2$
    leaves the factor
    %
    \begin{align*}
        \frac{f_\zeta(x)}\rho
        =\max\{\rho^{1/2},
                   \tau_\zeta^{(1-\alpha)/2}\rho^{-1/2}\}
        \leq C\tau_\zeta^{-\alpha/2},
    \end{align*}
    %
    since $\tau_\zeta\leq\rho\leq\pi/m$.
    This is the factor in the cited nonlinear estimate.
    In addition, the rescaled displacement is small uniformly,
    because
    $\|\phi\|_{X_\zeta}f_\zeta/\rho
       \leq C\tau_\zeta^{1-\alpha/4}\to0$.
    Finally,
    %
    \begin{align*}
        \tau_\zeta^{-\alpha/2}\|\phi\|_{X_\zeta}^2
        \leq\tau_\zeta^{-\alpha/2}
                    \tau_\zeta^{2+\alpha/2}
        =\tau_\zeta^2,
    \end{align*}
    %
    proving the second inequality in
    \eqref{eq:s4-scalar-quadratic}.
    Evaluating the derivative terms in the definition of
    $X_\zeta$ at the center of each ball gives
    \eqref{eq:s4-scalar-two-jet}, and the seminorm term gives
    its local H\"older counterpart.
\end{proof}

\begin{remark}
    \label{rem:s4-scalar-scope}
    All estimates in Proposition~\ref{prop:s4-scalar-input} hold
    on the entire scalar iteration set; none assumes that an
    iterate is minimal in $\mathcal{N}$ or in $\mathbb{S}^4$. In particular,
    the available bounds control the scalar correction and
    its first two covariant derivatives, together with the
    local $\beta$-H\"older seminorm of its second derivative.
    Subsequent estimates on arbitrary iterates will use this
    regularity without differentiating the correction a third
    time.
\end{remark}

\section{Correction normal to the suspension band}
\label{sec:s4-transverse}

Section~\ref{sec:s4-scalar-family} provides a family of embeddings
$G_{\zeta,\phi}:M_\zeta\to\mathcal N$, where
$\phi=u_\zeta+v\circ F_\zeta^{-1}$ and
$(v,\zeta)\in\mathcal B_m$. Each image has two independent
normal directions in $\mathbb S^4$: its unit normal $\nu$ within
$\mathcal N$, and the unit normal $n$ to $\mathcal N$.
Its mean-curvature vector has a scalar component in each
direction. To obtain a minimal surface in $\mathbb S^4$, both
components must vanish.

In this section we deform each of these surfaces along the
spherical geodesics with initial velocity $n$. The displacement
is a scalar function $b$ on $M_\zeta$. We choose $b$ so that
one scalar component of the mean-curvature vector of the
displaced surface vanishes. Since the tangent plane changes
under the deformation, the normal direction in which we test
this component must also be specified; we do this by projecting
the parallel transport of $n$ onto the new normal plane.
We then estimate the change in the other scalar component.
Section~\ref{sec:s4-coupled} will use that estimate to choose
$v$ and $\zeta$ so that the remaining component vanishes as well.
No minimality assumption is imposed on $G_{\zeta,\phi}$ here.

The argument has three parts. First, we obtain geometric
estimates valid for every $(v,\zeta)\in\mathcal B_m$.
Next, we use the symmetries of the coefficient $b$ to invert
the scalar operator that describes the first-order change
of the component in the direction $n$. Finally, we estimate
the nonlinear terms, solve for $b$, and bound its effect on
the remaining component. All estimates must be uniform on
$\mathcal B_m$, because the choice of $(v,\zeta)$ is made
only in the next section.

Fix $(v,\zeta)\in\mathcal B_m$, with $\mathcal B_m$ as in
\eqref{eq:s4-scalar-ball}, and abbreviate
%
\begin{align}
    \label{eq:s4-transverse-abbreviations}
    \begin{split}
        \tau:=\tau_\zeta,\qquad
        \varepsilon:=\varepsilon_\zeta,\qquad
        \rho:=\rho_\zeta,\qquad f:=f_\zeta,\qquad\phi:=u_\zeta+v\circ F_\zeta^{-1}.
    \end{split}
\end{align}
%
Recall that $\jmath_\zeta:M_\zeta\hookrightarrow\mathcal N$
is the inclusion, and set
%
\begin{align}
    \label{eq:s4-scalar-graph-map}
    \begin{split}
        G=G_{\zeta,\phi}:M_\zeta&\rightarrow\mathcal N,\\
        G(x)&:=\exp_{\jmath_\zeta(x)}^{\mathcal N}
            \bigl(\phi(x)\nu_\zeta(x)\bigr)
            \qquad\forall\,x\in M_\zeta,\\
        g_\phi&:=G^*g_{\mathcal N}.
    \end{split}
\end{align}
%
Proposition~\ref{prop:s4-scalar-input} shows that $G$ is a
$C^{2,\beta}$ embedding with image $M_{\zeta,\phi}$.
To distinguish the two ambient manifolds in the notation
for its second fundamental forms, let
%
\begin{align*}
    \jmath:=G:M_\zeta\rightarrow\mathcal N,
    \qquad
    \iota:=\iota_{\mathcal N}\circ G:
        M_\zeta\rightarrow\mathbb S^4,
\end{align*}
%
where $\iota_{\mathcal N}$ is the inclusion of $\mathcal N$
into $\mathbb S^4$. Both maps induce $g_\phi$ on $M_\zeta$.
We write $\nu:=\nu_{\zeta,\phi}$ for the unit normal along
$\jmath$ chosen in Proposition~\ref{prop:s4-scalar-input},
by continuity from $\nu_\zeta$ along $G_{\zeta,t\phi}$,
$0\leq t\leq1$. We also write $n$ for the field along
$\iota$ given by
%
\begin{align*}
    n(x):=N\bigl(\pi(G(x))\bigr)
    \qquad\forall\,x\in M_\zeta.
\end{align*}
%
Thus $(\nu,n)$ is an orthonormal frame of
$T^\perp_\iota M_\zeta$. With the notation of
Section~\ref{sec:s4-geometry},
%
\begin{align*}
    h_\iota^\nu=h_\jmath^\nu,\qquad
    S_\iota^\nu=S_\jmath^\nu,\qquad
    H_{\iota,\nu}=H_{\jmath,\nu}=H_{\zeta,\phi},\qquad\mathbf H_\iota
      =H_{\zeta,\phi}\nu+H_{\iota,n}n.
\end{align*}
%
The first three identities hold because the difference
between the spherical and the $\mathcal N$ covariant
derivatives is in the direction $n$, perpendicular to
$\nu$. All the displayed tensors are defined on the
domain $M_\zeta$. In particular, they are the geometric
quantities of the image $M_{\zeta,\phi}$ pulled back by $G$;
they are not the corresponding quantities of the initial
inclusion $\jmath_\zeta$.

The normal-connection one-form $\omega$ is the one in
\eqref{eq:s4-omega-definition}, now computed along $\iota$
in the frame $(\nu,n)$. The field $\nu$ is $C^{1,\beta}$.
The field $n$ is $C^{2,\beta}$, since it is obtained by
composing $G$ with the fixed smooth normal field on
$\mathcal N$. We continue to use the smooth metric
$g_\zeta$ and the function $\rho=\rho_\zeta$ on the
initial surface to define all weighted norms. In particular,
$\rho$ is not replaced by a distance function on the
deformed surface.

Unless a dependence is explicitly stated, every constant
below is independent of $m$ and of $(v,\zeta)\in\mathcal B_m$.
It may depend on the fixed band, cutoffs, $\alpha$, and
$\beta$. Continuity in $(v,\zeta)$ always refers to fixed
$m$. We write $o_m(1)$ for a quantity tending to zero as
$m\to\infty$, uniformly over all parameters and points
under consideration.

\subsection{Uniform geometry after the deformation within $\mathcal N$}

We first specify the norms for the new displacement $b$ and
the scalar component it must cancel. Unlike the displacement
$\phi$, whose allowed size varies according to $f_\zeta$,
the new displacement will have size $O(\tau_\zeta^2)$
everywhere. Its first and second derivatives are still
measured at the local length scale $\rho_\zeta$.
Using \eqref{eq:s4-weighted-norm}, define
%
\begin{align}
    \label{eq:s4-transverse-spaces}
    \begin{aligned}
        \mathcal X_\zeta&:=C^{2,\beta}_\chi(M_\zeta),
        &\qquad\|b\|_{\mathcal X_\zeta}
          &:=\|b\|_{C^{2,\beta}_\zeta(1)},\\
        \mathcal Z_\zeta&:=C^{0,\beta}(M_\zeta),
        &\qquad\|E\|_{\mathcal Z_\zeta}
          &:=\|E\|_{C^{0,\beta}_\zeta(\rho^{-2})},\\
        \mathcal Y_\zeta&:=C^{0,\beta}_\chi(M_\zeta),
        &\qquad\|E\|_{\mathcal Y_\zeta}
          &:=\|E\|_{\mathcal Z_\zeta}.
    \end{aligned}
\end{align}
%
Here the subscript $\chi$ imposes the transformation rule
defined in \eqref{eq:s4-symmetry-spaces}, with the character
$\chi$ of \eqref{eq:s4-character}. The space
$\mathcal Z_\zeta$ has no symmetry restriction; we will
also use its norm for the change in the invariant scalar
component. It differs from the scalar source space
$Y_\zeta$ in \eqref{eq:s4-XY}, whose weight is
$f_\zeta/\rho_\zeta^2$. All three spaces in
\eqref{eq:s4-transverse-spaces} are Banach spaces for
fixed $m$ and $\zeta$.

For example, if $b\in\mathcal X_\zeta$ and
$K_b:=\|b\|_{\mathcal X_\zeta}$, then
%
\begin{align*}
    |b(x)|&\leq K_b,\\
    |\nabla^{g_\zeta}b(x)|_{g_\zeta}
      &\leq K_b\rho(x)^{-1},\\
    |(\nabla^{g_\zeta})^2b(x)|_{g_\zeta}
      &\leq K_b\rho(x)^{-2}
      \qquad\forall\,x\in M_\zeta.
\end{align*}
%
The same norm includes the corresponding H\"older
estimates on the balls in \eqref{eq:s4-weighted-norm}.
Similarly,
%
\begin{align*}
    |E(x)|\leq\|E\|_{\mathcal Z_\zeta}\rho(x)^{-2}
    \qquad\forall\,x\in M_\zeta,
    \quad E\in\mathcal Z_\zeta.
\end{align*}
%
Thus a second derivative of a function of size one has
the size allowed by the source weight $\rho^{-2}$.
The estimate $|H_{\iota,n}|\leq C\tau^2\rho^{-2}$ below
will consequently give a source of norm $O(\tau^2)$.

We also need a precise convention for the local H\"older
estimates on geometric tensors. Given $x\in M_\zeta$,
put $r:=\rho(x)$ and
$U_x:=B_{g_\zeta}(x,r/100)$. When we assert an estimate
$|T|\leq Cw(\rho)$ with local scaled $C^{0,\beta}$
control, we mean that
%
\begin{align*}
    \sup_{U_x}|T|+r^\beta[T]_{\beta;U_x}
        \leq Cw(r)
        \qquad\forall\,x\in M_\zeta.
\end{align*}
%
In the seminorm, points are compared using their
$g_\zeta$ distance, tangent tensors are transported by
the $g_\zeta$ connection, and ambient vector components
are compared in $\mathbb R^5$. Tangent tensor norms use
$g_\phi$, unless another metric is displayed. These are
norms in the original metric, not in the rescaled metric
$r^{-2}g_\zeta$; the factor $r^\beta$ scales only the
H\"older quotient. The same estimates hold on balls with
any other fixed sufficiently small multiple of $r$ as
radius, after adjusting the constant.

\begin{lemma}[Geometry after the deformation within $\mathcal N$]
    \label{lem:s4-full-family-geometry}
    There exists $C>0$, independent of $m$ and of
    $(v,\zeta)\in\mathcal B_m$, such that the following
    estimates hold on $M_\zeta$, with local scaled
    $C^{0,\beta}$ control:
    %
    \begin{align}
        \label{eq:s4-full-family-geometry}
        \begin{aligned}
            |z\circ G|&\leq C\varepsilon,
            &\quad|\nu_h|&\leq C\tau/\rho,
            &\quad|S_\iota^\nu|&\leq C\tau/\rho^2,\\
            |h_\iota^n|&\leq C,
            &\quad|H_{\iota,n}|&\leq C\tau^2/\rho^2,
            &\quad|\omega|&\leq C\tau/\rho,\\
            |\nabla^{g_\phi}\omega|&\leq C\tau/\rho^2,
            &\quad|q|&\leq C,
            &\quad q&:=2+|S_\iota^n|^2-|\omega|^2.
        \end{aligned}
    \end{align}
    %
    Here $z:\mathcal N\to(-z_0,z_0)$ is the height
    coordinate, and
    %
    \begin{align*}
        \nu_h(x):=\sum_{a=1}^2
        \bigl\langle\nu(x),E_a(G(x))\bigr\rangle
             E_a(G(x))
        \qquad\forall\,x\in M_\zeta.
    \end{align*}
    %
    Norms and contractions of tangent tensors use
    $g_\phi$; ambient vector norms and inner products use
    $g_{\mathbb S^4}$. Moreover,
    %
    \begin{align}
        \label{eq:s4-deformed-scale}
        d_{\mathbb T}\bigl(\pi(G(x)),L_m\bigr)
          =\bigl(1+o_m(1)\bigr)\rho(x)
          \qquad\forall\,x\in M_\zeta,
    \end{align}
    %
    uniformly for $(v,\zeta)\in\mathcal B_m$.
\end{lemma}

\begin{proof}
    We first compare the size of $\phi$ with the local
    length scale $\rho$ and with the neck radius $\tau$.
    Put $B_\tau:=\tau^{1+\alpha/4}$. By
    \eqref{eq:s4-full-ball-bound}, the local scaled
    $C^{2,\beta}$ norm of $\phi$ on $U_x$ is at most
    $B_\tau f(x)$. The definition of $f$ gives
    %
    \begin{align*}
        \frac{B_\tau f(x)}{\rho(x)}
        =\max\left\{
          \tau^{1+\alpha/4}\rho(x)^{1/2},\,
          \tau^{3/2-\alpha/4}\rho(x)^{-1/2}
        \right\}
        \qquad\forall\,x\in M_\zeta.
    \end{align*}
    %
    Since $\tau\leq\rho(x)\leq\pi/m$, it follows that
    %
    \begin{align}
        \label{eq:s4-dimensionless-phi}
        \sup_{x\in M_\zeta}\frac{B_\tau f(x)}{\rho(x)}
        \leq C\left(
            \tau^{1-\alpha/4}
            +\tau^{1+\alpha/4}m^{-1/2}\right)
        =o_m(1).
    \end{align}
    %
    A second estimate will ensure that the change in the
    normal and second fundamental form is smaller than
    the bounds already proved for the initial surface.
    Both terms defining $f$ are increasing functions of
    $\rho$. At $\rho=\pi/m$, the ratio of the second
    term to the first is
    %
    \begin{align*}
        \frac{\tau^{(1-\alpha)/2}(\pi/m)^{1/2}}
             {(\pi/m)^{3/2}}
        =\frac{m}{\pi}\tau^{(1-\alpha)/2}=o_m(1).
    \end{align*}
    %
    Therefore
    %
    \begin{align}
        \label{eq:s4-sharp-phi}
        \sup_{x\in M_\zeta}\frac{B_\tau f(x)}{\tau}
        \leq C\tau^{\alpha/4}m^{-3/2}=o_m(1).
    \end{align}
    %
    The limits are uniform in $\zeta$, since
    $e^{-1}\tau_0\leq\tau_\zeta\leq e\tau_0$ and
    \eqref{eq:s4-scales} implies that every positive
    power of $\tau_0$ tends to zero faster than every
    fixed inverse power of $m$.

    We next explain how these two estimates control the
    geometric quantities. Fix $x\in M_\zeta$ and put
    $r:=\rho(x)$. Use one of the charts supplied by
    Lemma~\ref{lem:s4-initial-geometry}, with its
    independent variables rescaled so that they range
    over a fixed disk $D\subset\mathbb R^2$. Denote
    this chart by $\kappa_x:D\to M_\zeta$, with
    $\kappa_x(0)=x$, and write
    $I:=\jmath_\zeta\circ\kappa_x$. The coefficients of
    $r^{-2}\kappa_x^*g_\zeta$ are uniformly elliptic with
    uniformly bounded derivatives. The explicit graph
    and catenoid parametrizations in the proof of that
    lemma also bound every fixed number of derivatives
    of $(I-I(0))/r$ and $\nu_\zeta\circ\kappa_x$.
    By using a smaller fixed disk if necessary, the
    estimates include the ball $U_x$.

    Define
    %
    \begin{align*}
        \widehat\phi(y)&:=r^{-1}\phi(\kappa_x(y)),\\
        \mathcal E_r(y,t)&:=r^{-1}\left[
          \exp_{I(y)}^{\mathcal N}
            \bigl(rt\nu_\zeta(\kappa_x(y))\bigr)-I(0)
          \right]
    \end{align*}
    %
    for $y\in D$ and $|t|$ less than a fixed positive
    constant. The subtraction is in $\mathbb R^5$.
    The initial height tends uniformly to zero and
    $r\leq\pi/m$, so these exponential maps are defined
    inside the fixed band for large $m$.
    On a smaller fixed disk and a fixed interval of $t$,
    all derivatives of $\mathcal E_r$ needed below are
    uniformly bounded. This follows by differentiating
    the exponential map of the fixed smooth metric,
    using the stated bounds for $I$ and $\nu_\zeta$;
    each differentiation with respect to $t$ contributes
    a factor $r$ before the overall division by $r$.
    The rescaled deformed immersion is
    $y\mapsto\mathcal E_r(y,\widehat\phi(y))$, and
    %
    \begin{align*}
        \|\widehat\phi\|_{C^{2,\beta}(D)}
        \leq C\frac{B_\tau f(x)}r=o_m(1),
    \end{align*}
    %
    after restricting the disk once more if necessary.
    A fixed finite covering by the balls used in
    \eqref{eq:s4-weighted-norm} gives this estimate;
    $\rho$ and $f$ are comparable to their central
    values on each such chart.

    For $i,j\in\{1,2\}$, the first two coordinate
    derivatives are
    %
    \begin{align*}
        \partial_i[\mathcal E_r(y,\widehat\phi(y))]
        &=(\mathcal E_r)_i
             +(\mathcal E_r)_t\partial_i\widehat\phi,\\
        \partial_{ij}[\mathcal E_r(y,\widehat\phi(y))]
        &=(\mathcal E_r)_{ij}
             +(\mathcal E_r)_{it}\partial_j\widehat\phi
             +(\mathcal E_r)_{jt}\partial_i\widehat\phi+(\mathcal E_r)_{tt}
                  \partial_i\widehat\phi\partial_j\widehat\phi
             +(\mathcal E_r)_t\partial_{ij}\widehat\phi,
    \end{align*}
    %
    where the derivatives of $\mathcal E_r$ on the
    right are evaluated at $(y,\widehat\phi(y))$.
    The induced metric depends on the first derivatives
    in these formulas. The unit normal is obtained by
    solving its orthogonality equations and normalizing,
    and the second fundamental form is obtained by
    pairing the second covariant derivatives of the
    immersion with this normal. These are smooth
    functions of the displayed derivatives, the ambient
    metric, and its first derivatives, as long as the
    induced metric is uniformly positive definite.
    Estimate~\eqref{eq:s4-dimensionless-phi} preserves
    that positive definiteness.

    Applying the mean value theorem to these smooth
    functions, and using
    $[ab]_\beta\leq\|a\|_\infty[b]_\beta+
    \|b\|_\infty[a]_\beta$, bounds the change in the
    rescaled normal and second fundamental form by
    $C\|\widehat\phi\|_{C^{2,\beta}}$ in
    $C^{0,\beta}$. The same argument bounds the change
    in the rescaled metric coefficients in
    $C^{1,\beta}$. In particular, $g_\phi$ and
    $g_\zeta$ are uniformly equivalent.
    Returning to the original length scale, the changes
    in the normal and shape operator are bounded by
    %
    \begin{align*}
        C\frac{B_\tau f(x)}r,
        \qquad
        C\frac{B_\tau f(x)}{r^2},
    \end{align*}
    %
    respectively, with their local scaled
    $C^{0,\beta}$ bounds. Normals may be compared in
    $\mathbb R^5$; shape operators are compared as
    endomorphisms of the same tangent bundle $TM_\zeta$.

    In view of \eqref{eq:s4-sharp-phi}, these changes
    are $o_m(1)\tau/r$ and $o_m(1)\tau/r^2$.
    The horizontal projection also changes because
    it is evaluated at $G(x)$ instead of
    $\jmath_\zeta(x)$. The fields $E_1,E_2$ are
    smooth with bounded derivatives on the fixed band,
    so this contributes at most $C|\phi|$, with
    the corresponding local scaled bound.
    Since $r<1$ for large $m$, this is bounded by
    $CB_\tau f(x)/r$. The initial estimates
    \eqref{eq:s4-initial-geometry-bounds} therefore give
    %
    \begin{align*}
        |\nu_h|\leq C\tau/\rho,
        \qquad |S_\iota^\nu|\leq C\tau/\rho^2,
    \end{align*}
    %
    with the asserted H\"older control. Here
    $S_\iota^\nu=S_\jmath^\nu$, as explained above.

    The coordinate $z$ has $|\nabla^{\mathcal N}z|=1$.
    The geodesic segment defining $G(x)$ has length
    $|\phi(x)|$, and hence
    %
    \begin{align*}
        |z(G(x))|
        \leq |z(\jmath_\zeta(x))|+|\phi(x)|\leq C\varepsilon+B_\tau f(x)
         \leq C\varepsilon
         \qquad\forall\,x\in M_\zeta.
    \end{align*}
    %
    The last step uses \eqref{eq:s4-sharp-phi} and
    $\varepsilon=m^2\tau/(2\pi)$. On the initial
    charts, the explicit parametrizations in
    Lemma~\ref{lem:s4-initial-geometry} give local
    scaled H\"older bounds of size $C\varepsilon$ for
    the initial height. Composing the smooth height
    function with the exponential map above changes
    these bounds by at most $CB_\tau f(x)$.
    This proves the same estimate for $z\circ G$.

    The other quantities follow from the identities of
    Section~\ref{sec:s4-geometry}. Written on the domain
    $M_\zeta$, they are
    %
    \begin{align*}
        h_\iota^n(Y_1,Y_2)
          &=\bigl\langle S(dG(Y_1)),dG(Y_2)\bigr\rangle,\\
        H_{\iota,n}
          &=-\bigl\langle S(\nu_h),\nu_h\bigr\rangle,\\
        \omega(Y_1)
          &=\bigl\langle S(dG(Y_1)),\nu_h\bigr\rangle
    \end{align*}
    %
    for $Y_1,Y_2\in\Gamma(TM_\zeta)$, with $S$
    evaluated at $G(x)$ in each pointwise identity.
    Since $dG$ is an isometry from $(TM_\zeta,g_\phi)$
    onto the tangent plane of the image and $|S|\leq C$,
    we obtain
    %
    \begin{align*}
        |h_\iota^n|&\leq C,\\
        |H_{\iota,n}|&\leq C|\nu_h|^2
            \leq C\tau^2/\rho^2,\\
        |\omega|&\leq C|\nu_h|
            \leq C\tau/\rho.
    \end{align*}
    %
    The local scaled H\"older bounds follow from the
    same formulas and the product inequality. In
    particular, the estimate for $H_{\iota,n}$ uses
    both factors of $\nu_h$; replacing it by a bound
    linear in $|\nu_h|$ would lose the factor of $\tau$
    needed for the displacement estimate.

    To estimate the covariant derivative of $\omega$,
    use \eqref{eq:s4-omega-derivative} in the explicit
    pullback form
    %
    \begin{align*}
        (\nabla^{g_\phi}_{Y_1}\omega)(Y_2)
        =\bigl\langle
            (\nabla^{\mathcal N}_{dG(Y_1)}S)(dG(Y_2)),
            \nu\bigr\rangle
        +h_\iota^\nu(Y_1,Y_2)
            \langle S(\nu),\nu\rangle
          -\bigl\langle S(dG(Y_2)),
                dG(S_\iota^\nu(Y_1))\bigr\rangle
    \end{align*}
    %
    for $Y_1,Y_2\in\Gamma(TM_\zeta)$. It follows that
    %
    \begin{align*}
        |\nabla^{g_\phi}\omega|
        \leq |(\nabla^{\mathcal N}S)\circ G|
               +2|S\circ G|\,|S_\iota^\nu|
        \leq C\bigl(|z\circ G|+|S_\iota^\nu|\bigr)
        \leq C\bigl(m^2\tau+\tau/\rho^2\bigr)
         \leq C\tau/\rho^2.
    \end{align*}
    %
    The last inequality uses $\rho\leq\pi/m$.
    The explicit connection and shape-operator formulas
    in Lemma~\ref{lem:s4-suspension-geometry} show that
    $\nabla^{\mathcal N}S$ is a smooth tensor with a
    factor $\tan z$. Its composition with $G$ therefore
    has local scaled $C^{0,\beta}$ norm at most
    $C\varepsilon$. Applying the product inequality to
    the preceding identity proves the claimed
    H\"older bound for $\nabla^{g_\phi}\omega$.
    This argument differentiates neither
    $S_\iota^\nu$ nor the second derivatives of $\phi$.
    The bound for $q$ now follows from
    $|S_\iota^n|=|h_\iota^n|\leq C$ and
    $|\omega|\leq C$, since $\rho\geq\tau$.

    Finally, $|d\pi|\leq\sec z_0$ on the fixed band.
    Projecting the geodesic segment used to define
    $G(x)$ gives
    %
    \begin{align*}
        d_{\mathbb T}\bigl(\pi(G(x)),
                         \pi(\jmath_\zeta(x))\bigr)
        \leq\sec z_0\,|\phi(x)|.
    \end{align*}
    %
    Distance to the set $L_m$ is a $1$-Lipschitz
    function on $\mathbb T$. Recalling
    $\rho(x)=d_{\mathbb T}(\pi(\jmath_\zeta(x)),L_m)$,
    we consequently obtain
    %
    \begin{align*}
        \left|d_{\mathbb T}\bigl(\pi(G(x)),L_m\bigr)
                        -\rho(x)\right|
        \leq\sec z_0\,|\phi(x)|\leq CB_\tau f(x)
         =o_m(1)\rho(x)
         \qquad\forall\,x\in M_\zeta.
    \end{align*}
    %
    This proves \eqref{eq:s4-deformed-scale}. No
    differentiability of the distance to $L_m$ at a
    boundary between adjacent lattice cells is needed.
\end{proof}

\subsection{An auxiliary metric and a uniform inverse for its Laplacian}

The linear operator for the correction in the direction $n$
contains a Laplace operator and a bounded zeroth-order term.
We first study a Laplace operator for which the symmetry
conditions give estimates independent of the neck radius.
For this purpose we define a smooth metric $g_*$ on
$M_\zeta$. This metric agrees with the Euclidean catenoid
metric near every waist and with the flat metric pulled
back from $\mathbb T$ away from the necks. We will prove a
uniform inverse estimate for $\Delta_{g_*}$ on functions
with character $\chi$. The next subsection compares this
operator with the operator on the deformed surface.
The metric $g_*$ is used only for these estimates; it is
not asserted to be induced by an immersion into
$\mathbb S^4$.

Set
%
\begin{align*}
    a_m:=\frac{\ell}{2m},\qquad
    r_c:=\tau^\alpha,\qquad
    \theta_*:=\frac\pi4.
\end{align*}
%
Thus $a_m$ is half the side length of a square lattice
cell. Since $\tau/r_c\to0$ and $r_c/a_m\to0$ as
$m\to\infty$, all the changes to the flat metric described
below occur strictly inside each cell.
Choose a fixed smooth function
$\eta_*:[0,\infty)\to[0,1]$ that equals one on
$[0,1/4]$ and zero on $[1/2,\infty)$. Define
%
\begin{align}
    \label{eq:s4-model-radius}
    R(\xi):=\frac\tau2e^\xi,\qquad
    r_\tau(\xi):=R(\xi)+
    \eta_*(R(\xi)/r_c)\frac{\tau^2}{4R(\xi)}
    \qquad\forall\,\xi\in[0,\infty).
\end{align}
%
When $R(\xi)\leq r_c/4$, the cutoff equals one, and
%
\begin{align*}
    r_\tau(\xi)
    =\frac\tau2(e^\xi+e^{-\xi})
    =\tau\cosh\xi.
\end{align*}
%
When $R(\xi)\geq r_c/2$, it equals zero, so
$r_\tau(\xi)=R(\xi)$. In the remaining interval,
differentiation gives
%
\begin{align*}
    r_\tau'(\xi)
    &=R(\xi)+\frac{\tau^2}{4R(\xi)}
      \left[
        \frac{R(\xi)}{r_c}\eta_*'(R(\xi)/r_c)
        -\eta_*(R(\xi)/r_c)
      \right]\\
    &=R(\xi)\bigl(1+O((\tau/r_c)^2)\bigr).
\end{align*}
%
The constants in this estimate depend only on the fixed
cutoff. Consequently, for sufficiently large $m$,
$r_\tau(0)=\tau$, $r_\tau'(0)=0$, and
$r_\tau'(\xi)>0$ for every $\xi>0$.
The variable $\xi$ therefore determines the distance
$r_\tau(\xi)$ from the lattice point on the upper half
of a neck and on its adjoining sheet.

To specify the parametrization on this entire region,
extend the upper height $h_\zeta$ in
\eqref{eq:s4-glued-height} to $r_p>\tau$ by setting it
equal to $C_\tau(r_p)$ for $\tau<r_p\leq9\tau$.
The extension is smooth because the original height
already equals $C_\tau$ on the overlap. For
$p\in L_m$, define
%
\begin{align}
    \label{eq:s4-model-identification}
    \begin{split}
        \Theta_{\zeta,p}(\xi,\theta)
        :=\mathcal F\bigl(p+r_\tau(\xi)
                       (\cos\theta,\sin\theta),
                    h_\zeta(p+r_\tau(\xi)
                       (\cos\theta,\sin\theta))\bigr),
    \end{split}
\end{align}
%
on the domain
%
\begin{align}
    \label{eq:s4-model-wedge}
    \mathcal W_{m,\tau}
    :=\left\{(\xi,\theta):
       0<\theta<\theta_*,\quad
       0<\xi<\log\frac{2a_m}{\tau\cos\theta}
      \right\}.
\end{align}
%
Here the displacement from $p$ is taken in the flat
coordinate disk containing its square cell. The image
is the part of the upper sheet over the triangular
sector $0<\theta<\pi/4$ of that cell. Its outer boundary
is the cell wall
$r_\tau(\xi)\cos\theta=a_m$. This wall lies outside
the cutoff region, where $r_\tau=R$, and hence its
equation is
%
\begin{align*}
    \xi=d_m-\log\cos\theta,
    \qquad d_m:=\log\frac{2a_m}{\tau}.
\end{align*}
%
This explains the upper bound for $\xi$ in
\eqref{eq:s4-model-wedge}.
Near $\xi=0$, the height in
\eqref{eq:s4-model-identification} is
$C_\tau(\tau\cosh\xi)=\tau\xi$. Thus this
parametrization extends smoothly to the waist,
even though the height as a function of the base
radius has an unbounded derivative there.

We define $g_*$ on this parametrized sector by
%
\begin{align}
    \label{eq:s4-model-metric}
    \Theta_{\zeta,p}^*g_*
    :=r_\tau(\xi)^2(d\xi^2+d\theta^2),
\end{align}
%
and extend the definition using the symmetries of
Section~\ref{sec:s4-geometry}. We check that this
procedure defines a smooth metric on all of $M_\zeta$.
Near the waist, extending $\xi$ to negative values
gives the signed catenoidal parametrization, and the
metric is
%
\begin{align*}
    \tau^2\cosh^2\xi\,(d\xi^2+d\theta^2).
\end{align*}
%
This expression is smooth across $\xi=0$ and is
unchanged by $\xi\mapsto-\xi$. Across the sides
$\theta=0$ and $\theta=\pi/4$, smoothness follows
because the coefficient depends only on $\xi$.
Finally, outside the radial cutoff we have
$r_\tau=R$ and $dr_\tau=r_\tau\,d\xi$, so
%
\begin{align*}
    r_\tau^2(d\xi^2+d\theta^2)
    =dr_\tau^2+r_\tau^2d\theta^2
    =\Theta_{\zeta,p}^*(\pi^*g_{\mathbb T}).
\end{align*}
%
In particular, the metrics defined in adjacent cells
agree on a neighborhood of their common wall, where
both are $\pi^*g_{\mathbb T}$. These observations also
apply at intersections of the sides. The resulting
metric is smooth and $\Gamma_m$-invariant.

The scale function $\rho$ remains the one already
defined on $M_\zeta$. Since its value is the distance
of the base point from $L_m$, the present coordinates
give
%
\begin{align}
    \label{eq:s4-model-scale}
    \rho\circ\Theta_{\zeta,p}=r_\tau(\xi),
    \qquad R(\xi)\leq r_\tau(\xi)\leq2R(\xi),
    \qquad \tau\leq\rho\leq\pi/m.
\end{align}
%
The middle inequalities follow from
$0\leq\eta_*\leq1$ and $R\geq\tau/2$.
The last upper bound is the distance from the center
of a square cell to one of its corners.

A subscript $*$ on the spaces and norms in
\eqref{eq:s4-transverse-spaces} means that $g_\zeta$
is replaced by $g_*$. More generally,
$C^{k,\beta}_{*,\rho}(w)$ denotes the norm
\eqref{eq:s4-weighted-norm} with metric $g_*$,
scale $\rho$, and weight $w$. The scale is frozen at
the center of each ball, exactly as in that definition.
In particular,
%
\begin{align*}
    \|b\|_{\mathcal X_{\zeta,*}}
    &=\sup_{x\in M_\zeta}
      \|b\|_{C^{2,\beta}
       (B_{g_*}(x,\rho(x)/100),\,\rho(x)^{-2}g_*)},\\
    \|E\|_{\mathcal Y_{\zeta,*}}
    &=\sup_{x\in M_\zeta}\rho(x)^2
      \|E\|_{C^{0,\beta}
       (B_{g_*}(x,\rho(x)/100),\,\rho(x)^{-2}g_*)}.
\end{align*}
%
Both spaces consist of functions with character $\chi$.

We record the uniform local estimates for $g_*$ needed
below. For every fixed integer $j\geq1$, the $j$th derivative
of $\log r_\tau$ with respect to $\xi$ is uniformly bounded. On the
cutoff interval this follows by differentiating
%
\begin{align*}
    \frac{r_\tau(\xi)}{R(\xi)}-1
    =\eta_*(R(\xi)/r_c)\frac{\tau^2}{4R(\xi)^2},
    \qquad
    \left|\partial_\xi^j
      \left(\frac{r_\tau}{R}-1\right)\right|
      \leq C_j(\tau/r_c)^2.
\end{align*}
%
On the other intervals it follows from
$r_\tau=\tau\cosh\xi$ or $r_\tau=(\tau/2)e^\xi$.
These formulas also give
$|d\rho|_{g_*}\leq2$ wherever $\rho$ is smooth,
for sufficiently large $m$. At cell walls, $\rho$
is the minimum of the adjacent distance functions,
so it is globally $2$-Lipschitz. In particular, its
values are uniformly comparable on balls whose radii
are a sufficiently small fixed multiple of $\rho(x)$.

For completeness, there are uniform coordinate disks
on these balls, including the balls that meet a waist
or a cell wall. To see this near a waist, fix $B>4$.
For sufficiently large $m$, the region $\rho\leq4B\tau$
has exactly the Euclidean catenoid metric. At a point
with $\rho(x)\leq B\tau$, any path reaching
$\rho=4B\tau$ has length at least
$(4B\tau-\rho(x))/2>\rho(x)$, by the Lipschitz
bound just proved. The catenoid injectivity estimate
used in the proof of \eqref{eq:s4-scaled-injectivity}
therefore applies on a ball of radius $\rho(x)/2$.

For $\rho(x)\geq B\tau$, use the base coordinates
on the sheet over
$B_{\mathbb T}(\pi(x),\rho(x)/2)$.
The distance from every point of this disk to $L_m$
is at least $\rho(x)/2$, so the disk avoids all the
waists. In a cell, writing $r=r_\tau(\xi)$ gives
%
\begin{align*}
    g_*=
       \left(\frac{r_\tau(\xi)}{r_\tau'(\xi)}\right)^2
       dr^2+r^2d\theta^2.
\end{align*}
%
In the catenoidal region, the radial coefficient is
$(1-\tau^2/r^2)^{-1}$. In the cutoff region it is
$1+O((\tau/r_c)^2)$, with the same bounds for its
fixed-order derivatives after multiplication by the
corresponding powers of $r_c$. Elsewhere it equals
one. Consequently, after distances in the base disk
are divided by $\rho(x)$, the metric differs from
the Euclidean metric in $C^2$ norm by
$O(B^{-2})+o_m(1)$. The same conclusion holds across
cell walls, where the metric is already flat.
Choosing $B$ large and then $m$ large gives
%
\begin{align*}
    \operatorname{inj}_{g_*}(x)\geq\rho(x)/10
    \qquad\forall\,x\in M_\zeta.
\end{align*}
%
Indeed, a path leaving the indicated base disk has
length at least $\rho(x)/3$, by the metric comparison,
and the exponential map on vectors of length at most
$\rho(x)/10$ is $C^1$ close, after rescaling, to the
Euclidean exponential map. The argument by local
invertibility and injectivity is the same as in the
proof of \eqref{eq:s4-scaled-injectivity}.
The derivative bounds above give uniform bounds of
every fixed order for the metric coefficients in
these scaled coordinate disks. Thus the constants
in the local elliptic estimates used below do not
depend on $m$ or $\zeta$.

The boundary of the sector has a different role from
these local coordinate disks. It is used only to
prove a global estimate for functions with character
$\chi$. Figure~\ref{fig:s4-reflected-domain} shows
its domain in the $(\theta,\xi)$ plane. At fixed
$\xi$, the angular interval ends at $\theta_*$.
It begins at $0$ for $0<\xi\leq d_m$, and at
$\arccos(e^{d_m-\xi})$ for
$d_m<\xi<d_m+\log\sqrt2$. Thus its length is
always at most $\pi/4$. This length bound, including
where the interval is shortened by the cell wall,
is what gives the estimate below.

\begin{figure}[htbp]
    \centering
    \begingroup
    \definecolor{modelteal}{RGB}{0,112,116}
    \definecolor{modelorange}{RGB}{190,102,36}
    \begin{tikzpicture}[
        x=5.1cm,y=1.22cm,>=Stealth,
        every node/.style={font=\small}
    ]
        \fill[modelteal!8]
            (0,0)--({pi/4},0)--({pi/4},{2.6-ln(cos(45))})
            --plot[domain=pi/4:0,samples=50]
                (\x,{2.6-ln(cos(\x*180/pi))})--cycle;
        \draw[modelteal,line width=1pt] (0,0)--({pi/4},0);
        \draw[modelteal,line width=1pt] (0,0)--(0,2.6);
        \draw[modelorange,line width=1.2pt]
            ({pi/4},0)--({pi/4},{2.6-ln(cos(45))});
        \draw[modelteal,line width=1pt]
            plot[domain=0:pi/4,samples=50]
                (\x,{2.6-ln(cos(\x*180/pi))});
        \draw[black!40,dashed] (0,2.6)--({pi/4},2.6);
        \node[left=4pt] at (0,2.6) {$d_m$};
        \draw[black!65,->] (0,-0.10)--(0,3.38)
            node[above] {$\xi$};
        \draw[black!65,->] (-0.025,0)--(0.94,0)
            node[right] {$\theta$};
        \node[below left=3pt] at (0,0) {$0$};
        \node[below=4pt] at ({pi/4},0) {$\pi/4$};
        \node[left=6pt,text=modelteal] at (0,1.35)
            {$\partial_\theta b=0$};
        \node[right=7pt,text=modelorange] at ({pi/4},1.35)
            {$b=0$};
        \node[below=15pt,text=modelteal] at ({pi/8},0)
            {$\partial_\xi b=0$};
        \node[align=left,anchor=west] at (1.02,2.98)
            {Cell wall: $\xi=d_m-\log\cos\theta$\\[3pt]
             $\partial_{\eta_{\mathrm{out}}}b=0$};
        \draw[black!45] (0.61,2.83)--(0.96,2.83);
        \draw[black!55,dashed]
            (0,1.91)--({pi/4},1.91);
        \node[above=2pt,fill=modelteal!8,inner sep=1pt]
            at ({pi/8},1.91) {$0<\theta<\pi/4$};
        \draw[modelorange,line width=1.25pt]
            ({acos(exp(-0.16))*pi/180},2.76)--({pi/4},2.76);
        \fill[modelorange]
            ({acos(exp(-0.16))*pi/180},2.76)
            circle[radius=1.4pt];
        \node[anchor=west,align=left]
            at (1.02,1.94)
            {$\theta=\arccos(e^{d_m-\xi})$\\for $\xi>d_m$};
        \draw[black!45]
            ({acos(exp(-0.16))*pi/180},2.76)
            --(0.91,2.14)--(0.97,2.14);
    \end{tikzpicture}
    \endgroup
    \caption{The domain $\mathcal W_{m,\tau}$ used for the
    model metric, with $d_m=\log(2a_m/\tau)$.
    The boundary conditions shown are those for the
    coefficient $b$.
    Below $\xi=d_m$, a horizontal interval has endpoints
    $0$ and $\pi/4$.
    Above it, the interval begins on the cell wall and
    is shorter, but still has zero value at its endpoint
    $\theta=\pi/4$.
    Its length is therefore at most $\pi/4$, which is
    the fact used in the one-dimensional Poincar\'e
    estimate. The outward derivative on the curved
    boundary is taken in the model metric.}
    \label{fig:s4-reflected-domain}
\end{figure}
%
For the weak formulation used below, let $H^1(M_\zeta)$
denote the Sobolev space of functions with square-integrable
value and weak differential, computed with respect to $g_*$.
Define
%
\begin{align*}
    H^1_\chi(M_\zeta)
    :=\bigl\{
        b\in H^1(M_\zeta):
        b\circ g=\chi(g)b
        \text{ almost everywhere on }M_\zeta
        \text{ for every }g\in\Gamma_m
    \bigr\}.
\end{align*}
%
Each $g\in\Gamma_m$ preserves $g_*$, so the map
$b\mapsto b\circ g$ is an isometry of $H^1(M_\zeta)$.
The defining conditions therefore make
$H^1_\chi(M_\zeta)$ a closed linear subspace of
$H^1(M_\zeta)$.
%
\begin{lemma}[Uniform inverse for the auxiliary Laplacian]
    \label{lem:s4-model-inverse}
    The operator
    %
    \begin{align*}
        \Delta_{g_*}:\mathcal X_{\zeta,*}
        \rightarrow\mathcal Y_{\zeta,*}
    \end{align*}
    %
    is an isomorphism. There is a constant $C>0$,
    independent of $m$ and $\zeta$, such that
    %
    \begin{align}
        \label{eq:s4-model-inverse-bound}
        \|b\|_{\mathcal X_{\zeta,*}}
        \leq C\|\Delta_{g_*}b\|_{\mathcal Y_{\zeta,*}}
        \qquad\forall\,b\in\mathcal X_{\zeta,*}.
    \end{align}
\end{lemma}

\begin{proof}
    We first describe the restrictions of a smooth function
    $b$ with character $\chi$ to one sector. Interchange of
    the two coordinates fixes its diagonal side and has
    character $-1$. Hence $b=0$ on $\theta=\theta_*$.
    Reflection across the axis $\theta=0$ and reflection
    across the waist both have character $+1$. The same
    is true of reflection across the cell wall: on the
    base it is the composition of a lattice translation
    with a coordinate reflection. These three reflections
    preserve $g_*$, and the corresponding evenness of
    $b$ implies zero outward normal derivative on the
    axis, waist, and cell wall. Thus the boundary
    conditions in Figure~\ref{fig:s4-reflected-domain}
    follow from the globally imposed symmetry. They are
    not additional conditions on $b$.

    At a fixed value of $\xi$, denote the angular interval
    by $(a,a+L)$, where $a+L=\theta_*$ and
    $0<L\leq\theta_*$. If $v\in H^1(a,a+L)$ has
    trace zero at $a+L$, its even extension
    $\widetilde v(t):=v(a+|t|)$ belongs to
    $H^1_0(-L,L)$. The Dirichlet Poincar\'e inequality
    on this interval gives
    %
    \begin{align*}
        2\int_a^{a+L}|v'|^2\,d\theta
        =\int_{-L}^L|\widetilde v'|^2\,dt\geq\frac{\pi^2}{4L^2}
          \int_{-L}^L|\widetilde v|^2\,dt
        =\frac{\pi^2}{2L^2}
          \int_a^{a+L}|v|^2\,d\theta.
    \end{align*}
    %
    Dividing by two and using $L\leq\pi/4$ yields
    $\int|v'|^2\geq4\int|v|^2$.
    No condition at the left endpoint was used. This
    is why the same inequality applies above $\xi=d_m$,
    where that endpoint lies on the curved cell wall.

    Apply this inequality to the angular slices of $b$.
    In the coordinates \eqref{eq:s4-model-metric},
    the area measure and the energy density satisfy
    %
    \begin{align*}
        d\mu_*&=\rho^2\,d\xi\,d\theta,\\
        |\nabla b|_{g_*}^2\,d\mu_*
        &=\bigl(|\partial_\xi b|^2
                   +|\partial_\theta b|^2\bigr)
                     \,d\xi\,d\theta.
    \end{align*}
    %
    Integration over $\xi$ therefore gives
    %
    \begin{align}
        \label{eq:s4-angular-coercivity}
        4\int_{\mathcal W_{m,\tau}}\rho^{-2}b^2\,d\mu_*
        \leq\int_{\mathcal W_{m,\tau}}
                  |\nabla b|_{g_*}^2\,d\mu_*.
    \end{align}
    %
    Here, and in the computations on the sector below,
    functions, the metric, and the area measure are
    pulled back by $\Theta_{\zeta,p}$.
    The argument proves \eqref{eq:s4-angular-coercivity}
    for every $H^1$ function with zero trace on the
    diagonal side. Summing over all reflected sectors
    gives the same inequality on $M_\zeta$ for functions
    in $H^1_\chi(M_\zeta)$. In particular,
    %
    \begin{align*}
        \frac{4m^2}{\pi^2}
        \int_{M_\zeta}b^2\,d\mu_*
        \leq\int_{M_\zeta}|\nabla b|_{g_*}^2\,d\mu_*.
    \end{align*}
    %
    This uses $\rho\leq\pi/m$ and excludes a nonzero
    constant function from the space under consideration.

    We next solve the equation before proving the
    uniform H\"older estimate. For
    $E\in\mathcal Y_{\zeta,*}$, the preceding
    inequality makes the Dirichlet bilinear form
    coercive on the closed subspace
    $H^1_\chi(M_\zeta)\subset H^1(M_\zeta)$.
    The Lax--Milgram theorem therefore gives a unique
    $b\in H^1_\chi(M_\zeta)$ satisfying
    %
    \begin{align*}
        \int_{M_\zeta}
             \langle\nabla b,\nabla w\rangle_{g_*}\,d\mu_*
        =-\int_{M_\zeta}Ew\,d\mu_*
        \qquad\forall\,w\in H^1_\chi(M_\zeta).
    \end{align*}
    %
    This is also the weak equation against arbitrary
    $H^1$ test functions. Indeed, for such a function
    $w$, set
    %
    \begin{align*}
        P_\chi w
        :=\frac1{|\Gamma_m|}
          \sum_{g\in\Gamma_m}\chi(g)\,w\circ g.
    \end{align*}
    %
    Then $P_\chi w\in H^1_\chi(M_\zeta)$.
    Since $b,E$ have character $\chi$ and each $g$
    preserves $g_*$, change of variables gives
    %
    \begin{align*}
        \int_{M_\zeta}E(P_\chi w)\,d\mu_*
        &=\int_{M_\zeta}Ew\,d\mu_*,\\
        \int_{M_\zeta}
          \langle\nabla b,\nabla(P_\chi w)\rangle_{g_*}
                 \,d\mu_*
        &=\int_{M_\zeta}
          \langle\nabla b,\nabla w\rangle_{g_*}\,d\mu_*.
    \end{align*}
    %
    Elliptic regularity on the smooth closed surface
    gives $b\in C^{2,\beta}_\chi(M_\zeta)$.
    At this step $m$ is fixed. The uniform estimate
    is established next.

    To bound the supremum of $b$, set
    %
    \begin{align*}
        A_0:=\sup_{x\in M_\zeta}\rho(x)^2|E(x)|,
        \qquad
        \Phi(\theta):=\theta_*^2-\theta^2
        \quad\forall\,\theta\in[0,\theta_*].
    \end{align*}
    %
    The conformal formula for the Laplacian gives
    %
    \begin{align*}
        -\Delta_{g_*}\Phi
        =-\rho^{-2}(\partial_\xi^2\Phi
                         +\partial_\theta^2\Phi)
        =2\rho^{-2}.
    \end{align*}
    %
    The function $\Phi$ vanishes on the diagonal.
    Its outward derivative is zero on the axis and
    waist, because $\partial_\theta\Phi(0)=0$ and
    $\partial_\xi\Phi=0$.
    On the outer wall, a defining function whose
    positive side is outside the sector is
    $R(\xi)\cos\theta-a_m$. Its differential is
    $R\cos\theta\,d\xi-R\sin\theta\,d\theta$.
    Thus the outward unit normal to the outer wall, computed within the sector with respect to $g_*$, is
    %
    \begin{align*}
        \eta_{\mathrm{out}}
        =\rho^{-1}\bigl(\cos\theta\,\partial_\xi
                         -\sin\theta\,\partial_\theta\bigr),
        \qquad
        \partial_{\eta_{\mathrm{out}}}\Phi
        =\frac{2\theta\sin\theta}{\rho}\geq0.
    \end{align*}
    %
    Consider either of the two functions
    $w:=A_0\Phi\pm b$. Since $\Delta_{g_*}b=E$,
    %
    \begin{align*}
        -\Delta_{g_*}w
        =2A_0\rho^{-2}\mp E
        \geq A_0\rho^{-2}\geq0.
    \end{align*}
    %
    Moreover, $w=0$ on the diagonal and
    $\partial_{\eta_{\mathrm{out}}}w\geq0$ on
    every other side. Let $w^-:=\max\{-w,0\}$.
    Its trace vanishes on the diagonal. Integration
    by parts therefore gives
    %
    \begin{align*}
        -\int_{\mathcal W_{m,\tau}}|\nabla w^-|_{g_*}^2\,d\mu_*
        &=\int_{\mathcal W_{m,\tau}}\langle\nabla w,\nabla w^-\rangle_{g_*}\,d\mu_*\\
        &=\int_{\mathcal W_{m,\tau}}(-\Delta_{g_*}w)w^-\,d\mu_*+\int_{\partial\mathcal W_{m,\tau}}(\partial_{\eta_{\mathrm{out}}}w)w^-\,d\sigma_*\geq0.
    \end{align*}
    %
    Here $d\sigma_*$ is the boundary length measure.
    The boundary term on the diagonal is zero, and
    the terms on the other sides are nonnegative.
    The identity is valid for the $H^1$ test function
    $w^-$ on this piecewise smooth domain; the corner
    points have boundary measure zero. It follows that
    $\nabla w^-=0$. Applying
    \eqref{eq:s4-angular-coercivity} to $w^-$ gives
    $w^-=0$. Performing this argument for both signs
    yields $|b|\leq A_0\Phi$ on the sector. Absolute
    values are preserved under all reflections, so
    %
    \begin{align}
        \label{eq:s4-model-barrier}
        \|b\|_{L^\infty(M_\zeta)}
        \leq\theta_*^2 A_0
        \leq\theta_*^2\|E\|_{\mathcal Y_{\zeta,*}}.
    \end{align}

    Finally, fix $x\in M_\zeta$ and put
    $\rho_x:=\rho(x)$ and $\widehat g_x:=\rho_x^{-2}g_*$.
    This is a constant rescaling of the metric, so
    %
    \begin{align*}
        \Delta_{\widehat g_x}b=\rho_x^2 E.
    \end{align*}
    %
    In the uniform coordinate disks constructed above,
    the coefficients of $\widehat g_x$ are uniformly
    elliptic and have uniform $C^{1,\beta}$ bounds.
    Interior Schauder estimates on two concentric
    disks, with radii bounded below independently
    of $x,m,\zeta$ in the rescaled metric, give
    %
    \begin{align*}
        \|b\|_{C^{2,\beta}
          (B_{g_*}(x,\rho_x/100),\widehat g_x)}\leq C\left(
            \|b\|_{L^\infty(M_\zeta)}
            +\rho_x^2\|E\|_{C^{0,\beta}
              (B_{g_*}(x,\rho_x/20),\widehat g_x)}
                       \right).
    \end{align*}
    %
    The larger ball is covered by a uniformly bounded
    number of balls used in the definition of
    $\mathcal Y_{\zeta,*}$. On these balls $\rho$
    is comparable to $\rho_x$. The source term on
    the right is therefore bounded by
    $C\|E\|_{\mathcal Y_{\zeta,*}}$. This includes
    its H\"older seminorm: sufficiently close pairs
    lie in one of the smaller coordinate balls, while
    pairs separated by a fixed positive rescaled
    distance are controlled by the supremum bound.
    Taking the supremum over $x$ and using
    \eqref{eq:s4-model-barrier} proves
    \eqref{eq:s4-model-inverse-bound}.
    All these estimates are interior estimates on
    the full smooth surface. The sides of the sector
    introduce no boundary regularity problem.
\end{proof}

\begin{remark}
    \label{rem:s4-forcing-symmetry}
    The source in Lemma~\ref{lem:s4-model-inverse} is
    a globally defined $C^{0,\beta}$ function with
    character $\chi$. A H\"older function specified
    only on the sector need not have this property:
    for example, its trace on the diagonal must vanish
    in order that its odd reflection be continuous.
    The sources used below are defined on the full
    surface and have the required character by
    equivariance of the geometric operators.
\end{remark}

\subsection{An inverse for the linearized operator in the direction $n$}

The preceding lemma gives an inverse for $\Delta_{g_*}$.
We now transfer that inverse to
$\Delta_{g_\phi}+q$, which is the linearization of the
mean-curvature component in the direction $n$.
There are two differences to estimate: the change from the
auxiliary metric $g_*$ to the induced metric $g_\phi$, and
the addition of the function $q$. The metric difference
will tend to zero in local rescaled coordinates. The
function $q$ is uniformly bounded, so its contribution to
the weighted operator norm is bounded by a constant times
$\sup_{M_\zeta}\rho^2\leq\pi^2/m^2$.

To state the metric comparison precisely, for a
$C^{1,\beta}$ covariant two-tensor $k$ on $M_\zeta$ define
%
\begin{align}
    \label{eq:s4-metric-comparison-norm}
    \mathfrak{M}_*(k)
    :=\sup_{x\in M_\zeta}
    \|\rho(x)^{-2}k\|_{C^{1,\beta}
    (B_{g_*}(x,\rho(x)/100),\,\rho(x)^{-2}g_*)}.
\end{align}
%
At each centre $x$, the number $\rho(x)$ is held fixed
throughout the ball. Both the tensor and the metric are
rescaled: if $r=\rho(x)$, then
%
\begin{align*}
    |r^{-2}k|_{r^{-2}g_*}=|k|_{g_*},
    \qquad
    |\nabla^{r^{-2}g_*}(r^{-2}k)|_{r^{-2}g_*}
    =r|\nabla^{g_*}k|_{g_*}.
\end{align*}
%
Thus $\mathfrak M_*(k)$ controls the relative size of the
metric coefficients, their first derivatives at scale $r$,
and the corresponding H\"older seminorms.

\begin{lemma}[Comparison of the induced and auxiliary metrics]
    \label{lem:s4-metric-comparison}
    For $\zeta\in[-1,1]$, set
    %
    \begin{align}
        \label{eq:s4-comparison-smallness}
        \eta_m(\zeta)
        :=\varepsilon^2+\tau^{2-2\alpha}
        +\tau^{1-\alpha/4}
        +\tau^{1+\alpha/4}m^{-1/2}.
    \end{align}
    %
    Then $\eta_m(\zeta)$ tends to zero uniformly in
    $\zeta$ as $m\to\infty$. Uniformly for
    $(v,\zeta)\in\mathcal B_m$, one has
    %
    \begin{align}
        \label{eq:s4-metric-comparison}
        \begin{split}
            \mathfrak{M}_*(g_\zeta-g_*)
            &\leq C(\varepsilon^2+\tau^{2-2\alpha}),\\
            \mathfrak{M}_*(g_\phi-g_*)
            &\leq C\eta_m(\zeta).
        \end{split}
    \end{align}
    %
    The norms based on $g_*$ and $g_\zeta$ defining
    $\mathcal X_\zeta$ and $\mathcal Z_\zeta$ are uniformly
    equivalent. Moreover, for every $b\in\mathcal X_\zeta$,
    %
    \begin{align}
        \label{eq:s4-Laplacian-comparison}
        \| (\Delta_{g_\phi}-\Delta_{g_*})b\|_{\mathcal{Z}_\zeta}
        \leq C\eta_m(\zeta)\|b\|_{\mathcal{X}_\zeta}.
    \end{align}
\end{lemma}

\begin{proof}
    We first compare $g_\zeta$ with $g_*$ in the three
    regions used to define the auxiliary metric.
    On the part of each neck where
    $r_\tau(\xi)=\tau\cosh\xi$, the initial immersion is
    the catenoidal parametrization with $z=\tau\xi$.
    Its induced metric is
    %
    \begin{align*}
        \Theta_{\zeta,p}^*g_\zeta
        =\tau^2\bigl(\cos^2z\sinh^2\xi+1\bigr)d\xi^2
        +\tau^2\cos^2z\cosh^2\xi\,d\theta^2,
    \end{align*}
    %
    where the pullback is by \eqref{eq:s4-model-identification}.
    Since $\Theta_{\zeta,p}^*g_*=\tau^2\cosh^2\xi
    (d\xi^2+d\theta^2)$, the
    relative coefficient matrix of $g_\zeta-g_*$ in these
    coordinates is
    %
    \begin{align*}
        (\cos^2z-1)
        \begin{pmatrix}
            \tanh^2\xi&0\\
            0&1
        \end{pmatrix}.
    \end{align*}
    %
    For $j\in\{0,1,2\}$, its $j$th derivative in $\xi$
    has norm at most
    %
    \begin{align*}
        C\bigl(z^2+\tau|z|+\tau^2\bigr)
        \leq C\varepsilon^2.
    \end{align*}
    %
    Here $|z|\leq C\varepsilon$ by
    Lemma~\ref{lem:s4-initial-geometry},
    $\partial_\xi z=\tau$, and $\tau\leq C\varepsilon$.
    The coordinates rescaled at a centre of radius $\rho$
    have uniformly bounded coefficients, as proved above
    for $g_*$. The displayed derivative bounds therefore
    give a local rescaled $C^{1,\beta}$ bound
    $C\varepsilon^2$. The same calculation holds in the
    signed coordinate $\xi$ through the waist.

    Next consider the radial cutoff region
    $r_c/4\leq R\leq r_c/2$ in
    \eqref{eq:s4-model-radius}. Put
    %
    \begin{align*}
        a(\xi):=\eta_*(R(\xi)/r_c),
        \qquad e(\xi):=\frac{\tau^2}{4R(\xi)^2}.
    \end{align*}
    %
    The notation in this paragraph is local to this
    calculation. Since $R'=R$ and $e'=-2e$, one has
    %
    \begin{align*}
        r_\tau=R(1+ae),\qquad
        r_\tau'=R\bigl(1+e(a'-a)\bigr),\qquad
        r_\tau^2-\tau^2
        =R^2\bigl((1+ae)^2-4e\bigr).
    \end{align*}
    %
    Every fixed-order derivative of $a$ is uniformly
    bounded on this region, whereas
    $|\partial_\xi^j e|\leq C_j(\tau/r_c)^2$.
    Consequently,
    %
    \begin{align*}
        \frac{r_\tau'^2}{r_\tau^2-\tau^2}
        &=\frac{\bigl(1+e(a'-a)\bigr)^2}
        {(1+ae)^2-4e}
        =1+O\bigl((\tau/r_c)^2\bigr),
    \end{align*}
    %
    and the same bound holds after any fixed positive
    number of $\xi$ derivatives of the difference from
    $1$. The denominator is bounded away from zero for
    all sufficiently large $m$.

    The initial surface in this region is still the exact
    catenoidal parametrization. On its upper half,
    $z=\tau\operatorname{arcosh}(r_\tau/\tau)$.
    The radial coefficient of its metric, divided by
    $r_\tau^2$, can be written as
    %
    \begin{align*}
        \frac{r_\tau'^2}{r_\tau^2}
        \left(\cos^2z+
        \frac{\tau^2}{r_\tau^2-\tau^2}\right)
        =\frac{r_\tau'^2}{r_\tau^2-\tau^2}
        +(\cos^2z-1)\frac{r_\tau'^2}{r_\tau^2}.
    \end{align*}
    %
    Its angular coefficient, divided by $r_\tau^2$, is
    $\cos^2z$. In addition to $|z|\leq C\varepsilon$,
    differentiation of the height gives
    %
    \begin{align*}
        z'=\frac{\tau r_\tau'}
        {\sqrt{r_\tau^2-\tau^2}},
        \qquad |\partial_\xi^j z|\leq C_j\tau
        \quad\text{for every fixed }j\geq1.
    \end{align*}
    %
    These formulas bound the relative metric error and
    its first two coordinate derivatives by
    $C(\varepsilon^2+\tau^2/r_c^2)$. They give the
    corresponding local $C^{1,\beta}$ estimate, including
    the derivatives of the radial cutoff. The lower half
    has the same metric because its height is $-z$.

    On the rest of each sheet, where $R\geq r_c/2$,
    the auxiliary metric is the pullback of the flat
    metric on $\mathbb T$. In the base coordinates the
    initial metric is
    %
    \begin{align*}
        g_\zeta=\cos^2h_\zeta\,g_{\mathbb T}
        +dh_\zeta\otimes dh_\zeta.
    \end{align*}
    %
    The graph estimates in the proof of
    Lemma~\ref{lem:s4-initial-geometry}, including the
    matching estimate \eqref{eq:s4-matching-estimate},
    give
    %
    \begin{align*}
        |h_\zeta|\leq C\varepsilon,
        \qquad
        \rho^j|D^jh_\zeta|\leq C_j\tau
        \quad\text{for every fixed }j\geq1,
    \end{align*}
    %
    together with their local scaled H\"older bounds.
    For example, the first scaled derivative of
    $\cos^2h_\zeta-1$ is bounded by
    $C|h_\zeta|\rho|Dh_\zeta|\leq C\varepsilon\tau$,
    while that of $dh_\zeta\otimes dh_\zeta$ is bounded by
    $C\rho|Dh_\zeta||D^2h_\zeta|
    \leq C\tau^2/\rho^2$.
    The product estimate gives the same orders in the
    scaled H\"older norms. Since $\rho\geq r_c/2$ here,
    the total relative $C^{1,\beta}$ error is at most
    %
    \begin{align*}
        C\left(\varepsilon^2+\frac{\tau^2}{\rho^2}\right)
        \leq C\left(\varepsilon^2+
        \frac{\tau^2}{r_c^2}\right).
    \end{align*}
    %
    The three regions cover $M_\zeta$, and the coefficient
    estimates agree on their overlaps. Because
    $r_c=\tau^\alpha$, this proves the first inequality in
    \eqref{eq:s4-metric-comparison}.

    We next include the deformation $\phi$ within
    $\mathcal N$. Let $B_\tau:=\tau^{1+\alpha/4}$ and set
    %
    \begin{align*}
        d_\phi:=\sup_{x\in M_\zeta}
        \frac{B_\tau f(x)}{\rho(x)}.
    \end{align*}
    %
    In an initial chart of scale $r=\rho(x)$, write
    $\widehat I$ and $\widehat G$ for the initial and
    deformed immersions into $\mathbb R^5$, after subtracting
    the initial chart centre and dividing by $r$.
    The chain-rule calculation in the proof of
    Lemma~\ref{lem:s4-full-family-geometry} gives
    %
    \begin{align*}
        \|\widehat G-\widehat I\|_{C^{2,\beta}}
        \leq Cd_\phi,
        \qquad
        \|\widehat I\|_{C^{2,\beta}}\leq C.
    \end{align*}
    %
    Put $W:=\widehat G-\widehat I$. The coefficients of
    the two rescaled induced metrics satisfy the exact identity
    %
    \begin{align*}
        (\widehat g_\phi-\widehat g_\zeta)_{ij}
        =\langle\partial_i\widehat I,\partial_jW\rangle
        +\langle\partial_iW,\partial_j\widehat I\rangle
        +\langle\partial_iW,\partial_jW\rangle
        \qquad i,j\in\{1,2\}.
    \end{align*}
    %
    The inner products are Euclidean. Differentiating
    this identity once uses at most two derivatives of
    $W$, hence at most two derivatives of $\phi$.
    The $C^{1,\beta}$ product estimate therefore yields
    %
    \begin{align*}
        \mathfrak M_*(g_\phi-g_\zeta)
        \leq Cd_\phi(1+d_\phi).
    \end{align*}
    %
    By \eqref{eq:s4-dimensionless-phi},
    %
    \begin{align*}
        d_\phi
        \leq C\left(\tau^{1-\alpha/4}
        +\tau^{1+\alpha/4}m^{-1/2}\right)=o_m(1).
    \end{align*}
    %
    Adding the initial metric error proves the second
    inequality in \eqref{eq:s4-metric-comparison}.

    The first metric comparison gives uniform ellipticity
    and uniform $C^{1,\beta}$ bounds for the coefficients
    of $g_\zeta$ in the scaled $g_*$ charts. In particular,
    the balls used in the two definitions of the weighted
    norms have uniformly comparable radii. A ball in one
    definition is covered by a uniformly bounded number
    of slightly smaller balls for the other definition.
    The scale $\rho$ is comparable on these balls.
    Finally, for a scalar function $b$, the first covariant
    derivative is $db$, and in local coordinates its
    second covariant derivative is
    %
    \begin{align*}
        (\nabla^g)^2_{ij}b
        =\partial_{ij}b-\Gamma(g)^k_{ij}\partial_kb.
    \end{align*}
    %
    The $C^{0,\beta}$ bounds for the Christoffel symbols
    thus give uniform equivalence of the two
    $C^{2,\beta}$ function norms. The same argument for
    order zero, with the common weight $\rho^{-2}$,
    gives equivalence of the source norms.

    To compare the Laplacians, fix a scaled chart with
    centre $x$ and $r=\rho(x)$. In its coordinates the
    difference is
    %
    \begin{align*}
        r^2(\Delta_{g_\phi}-\Delta_{g_*})b
        =(\widehat g_\phi^{ij}-\widehat g_*^{ij})
        \partial_{ij}b-
        \bigl(\widehat g_\phi^{ij}
        \Gamma(\widehat g_\phi)^k_{ij}
        -\widehat g_*^{ij}\Gamma(\widehat g_*)^k_{ij}\bigr)
        \partial_kb.
    \end{align*}
    %
    Here $b$ denotes its coordinate expression, and the
    repeated indices range from $1$ to $2$.
    Inverting a uniformly positive definite matrix and
    forming its Christoffel symbols are smooth operations
    on the metric coefficients and their first derivatives.
    Therefore the two coefficients on the right have
    $C^{0,\beta}$ norms at most $C\eta_m(\zeta)$.
    The product estimate, followed by the supremum over
    chart centres, proves
    \eqref{eq:s4-Laplacian-comparison}.
\end{proof}

\begin{proposition}[Uniform inverse in the direction $n$]
    \label{prop:s4-transverse-inverse}
    Let $q$ be the function defined in
    \eqref{eq:s4-full-family-geometry}, and set
    %
    \begin{align}
        \label{eq:s4-P-definition}
        P_{\zeta,\phi}:=\Delta_{g_\phi}+q.
    \end{align}
    %
    For all sufficiently large $m$ and every
    $(v,\zeta)\in\mathcal B_m$, this operator is an
    isomorphism from $\mathcal X_\zeta$ to
    $\mathcal Y_\zeta$. Its inverse
    $\mathscr G_{\zeta,\phi}$ satisfies, for every
    $E\in\mathcal Y_\zeta$,
    %
    \begin{align}
        \label{eq:s4-actual-inverse}
        \|\mathscr G_{\zeta,\phi}E\|_{\mathcal X_\zeta}
        \leq C\|E\|_{\mathcal Y_\zeta},
    \end{align}
    %
    with $C$ independent of $(v,\zeta)$ and $m$.
    Furthermore,
    %
    \begin{align}
        \label{eq:s4-forcing-and-mixed}
        \begin{split}
            \|H_{\iota,n}\|_{\mathcal Y_\zeta}
            &\leq C\tau^2,\\
            \|Bb\|_{\mathcal Z_\zeta}
            &\leq C\tau\|b\|_{\mathcal X_\zeta}
            \qquad\forall\,b\in\mathcal X_\zeta,
        \end{split}
    \end{align}
    %
    where $B$ is the off-diagonal operator in
    \eqref{eq:s4-scalar-Jacobi}, evaluated along $\iota$.
\end{proposition}

\begin{proof}
    Equivariance of $G$ implies that $g_\phi$ is
    $\Gamma_m$-invariant. The expression for $q$ in
    \eqref{eq:s4-full-family-geometry} is invariant as
    well. Hence $P_{\zeta,\phi}$ sends functions with
    character $\chi$ to functions with the same character.

    Write $\mathcal D:=P_{\zeta,\phi}-\Delta_{g_*}$.
    The preceding lemma bounds the difference of the
    Laplacians. For the potential term, fix $x\in M_\zeta$
    and put $r=\rho(x)$. The local scaled
    $C^{0,\beta}$ bound for $q$ from
    Lemma~\ref{lem:s4-full-family-geometry} and the
    H\"older product estimate give
    %
    \begin{align*}
        r^2\|qb\|_{0,\beta;r}
        \leq Cr^2\|b\|_{\mathcal X_\zeta}
        \leq Cm^{-2}\|b\|_{\mathcal X_\zeta}.
    \end{align*}
    %
    Here $\|\cdot\|_{0,\beta;r}$ denotes the local
    H\"older norm after distances are divided by $r$;
    tensor sizes are measured in the unrescaled metric,
    as specified at the beginning of this section.
    Taking the supremum over $x$ proves
    %
    \begin{align*}
        \|\mathcal D b\|_{\mathcal Y_\zeta}
        \leq C\bigl(\eta_m(\zeta)+m^{-2}\bigr)
        \|b\|_{\mathcal X_\zeta}
        \qquad\forall\,b\in\mathcal X_\zeta.
    \end{align*}

    Let $K_*:=(\Delta_{g_*})^{-1}$ be the inverse in
    Lemma~\ref{lem:s4-model-inverse}, now regarded as an
    operator $\mathcal Y_\zeta\to\mathcal X_\zeta$
    by the equivalence of norms just proved.
    Its norm is uniformly bounded. Since
    $\eta_m(\zeta)+m^{-2}\to0$ uniformly, we may choose
    the lower threshold for $m$ so that
    %
    \begin{align*}
        \|K_*\mathcal D\|_{\mathcal X_\zeta\to\mathcal X_\zeta}
        \leq\frac12
        \qquad\forall\,(v,\zeta)\in\mathcal B_m.
    \end{align*}
    %
    The identity
    $P_{\zeta,\phi}=\Delta_{g_*}(I+K_*\mathcal D)$ then gives
    %
    \begin{align*}
        \mathscr G_{\zeta,\phi}
        =(I+K_*\mathcal D)^{-1}K_*
        =\sum_{j=0}^{\infty}(-K_*\mathcal D)^jK_*.
    \end{align*}
    %
    The series converges in operator norm. It is an
    inverse on both sides, and its norm is at most
    $2\|K_*\|$. This proves both the isomorphism assertion
    and \eqref{eq:s4-actual-inverse}.

    All operators in this argument act on the same
    spaces of functions on the closed surface $M_\zeta$,
    with the same prescribed character. Their restrictions
    to a reflection sector therefore have the same
    boundary conditions. More explicitly, at a fixed
    curve of a reflection, the $+1$ and $-1$ tangent
    eigenspaces are orthogonal for every invariant metric.
    An even function has zero derivative in the $-1$
    direction, which is the metric-normal direction to
    that curve; an odd function has zero value on it.
    This applies to both $g_*$ and $g_\phi$, including
    the reflection through the waist. At intersections
    of reflection curves the argument is made on the
    full surface in its ordinary local coordinates.

    It remains to estimate the source and the
    off-diagonal term. The transformation rule for $n$
    and equivariance of the mean-curvature vector give
    $H_{\iota,n}\in\mathcal Y_\zeta$.
    At a centre $x$ with $r=\rho(x)$, the geometric
    estimates of Lemma~\ref{lem:s4-full-family-geometry}
    give
    %
    \begin{align*}
        \|H_{\iota,n}\|_{0,\beta;r}
        \leq C\tau^2/r^2.
    \end{align*}
    %
    Multiplication by $r^2$ and the supremum over
    $x\in M_\zeta$ yield the first inequality in
    \eqref{eq:s4-forcing-and-mixed}.

    For the second inequality, the exact formula for
    the off-diagonal operator is
    %
    \begin{align*}
        Bb=-2\langle\omega,db\rangle_{g_\phi}
        -(\operatorname{div}_{g_\phi}\omega)b
        +\langle h_\iota^\nu,h_\iota^n\rangle_{g_\phi}b.
    \end{align*}
    %
    Set $K_b:=\|b\|_{\mathcal X_\zeta}$ and keep the
    same centre and scale. The definition of this norm
    gives local scaled bounds $K_b$ for $b$ and
    $CK_b/r$ for $db$. The geometric lemma gives
    bounds $C\tau/r$ for $\omega$,
    $C\tau/r^2$ for $\nabla^{g_\phi}\omega$ and
    $h_\iota^\nu$, and $C$ for $h_\iota^n$.
    Each of these statements includes its local scaled
    $C^{0,\beta}$ bound. Thus
    %
    \begin{align*}
        \|2\langle\omega,db\rangle_{g_\phi}
        \|_{0,\beta;r}
        &\leq C\frac\tau r\frac{K_b}r,\\
        \|(\operatorname{div}_{g_\phi}\omega)b
        \|_{0,\beta;r}
        &\leq C\frac\tau{r^2}K_b,\\
        \|\langle h_\iota^\nu,h_\iota^n\rangle_{g_\phi}b
        \|_{0,\beta;r}
        &\leq C\frac\tau{r^2}K_b.
    \end{align*}
    %
    Contraction by $g_\phi^{-1}$ preserves these bounds
    because its coefficients have uniform local scaled
    $C^{0,\beta}$ bounds. Multiplication by $r^2$
    and the supremum over chart centres prove the
    second inequality in \eqref{eq:s4-forcing-and-mixed}.
    The output $Bb$ is invariant, as proved in the
    symmetry discussion in Section~\ref{sec:s4-geometry};
    this is why its norm is taken in $\mathcal Z_\zeta$,
    which imposes no character condition.
\end{proof}

The first source estimate allows $H_{\iota,n}$ to have
size independent of $\tau$ at a waist, where $\rho=\tau$.
Indeed, the local estimate there is only
$|H_{\iota,n}|\leq C\tau^2/\rho^2=C$.
Nevertheless, the solution of the linear correction equation
$P_{\zeta,\phi}b=-H_{\iota,n}$ satisfies
%
\begin{align*}
    \|b\|_{\mathcal X_\zeta}\leq C\tau^2,
    \qquad
    |b(x)|+\rho(x)|db(x)|_{g_\zeta}
    \leq C\tau^2
    \quad\forall\,x\in M_\zeta.
\end{align*}
%
At a waist, the displacement divided by the neck radius
and the size of its first derivative are therefore both
$O(\tau)$. The next subsection estimates the nonlinear
error so that this linear correction can be replaced by
one for which the chosen mean-curvature component vanishes
exactly.

\subsection{The displaced surface and its two mean-curvature components}

A function $b$ will specify how far to move each point of
$M_{\zeta,\phi}$ along the spherical geodesic normal to
$\mathcal N$. We first write this new immersion explicitly
and compute its tangent vectors and induced metric. We then
define two scalar functions whose simultaneous vanishing
means that the new immersion is minimal in $\mathbb S^4$.
This requires choosing normal vectors along the new immersion:
parallel transport of $\nu$ and $n$ alone does not in general
give vectors perpendicular to its new tangent plane.

We will project the transported vectors onto that normal
plane and normalize them. This choice has a useful property:
at $b=0$, the derivatives of the two scalar functions are
exactly $B$ and $P_{\zeta,\phi}$, including when
$M_{\zeta,\phi}$ is not minimal. The final estimate in this
subsection bounds the error after subtracting these linear
terms. It will allow us to solve the equation with derivative
$P_{\zeta,\phi}$ by contraction.

Keep $(v,\zeta)\in\mathcal B_m$ fixed, and write
$X:=\iota:M_\zeta\to\mathbb S^4\subset\mathbb R^5$
and $g:=g_\phi$. Thus $X$ and $G$ have the same values
as maps into $\mathbb R^5$, but the specified target of
$X$ is $\mathbb S^4$. For $b\in\mathcal X_\zeta$, define
$X_b:M_\zeta\to\mathbb S^4$ and a vector field $V_b$
along it by
%
\begin{align}
    \label{eq:s4-transverse-graph}
    \begin{split}
        X_b(x)&:=\exp_{X(x)}^{\mathbb S^4}(b(x)n(x))
        =\cos(b(x))X(x)+\sin(b(x))n(x),\\
        V_b(x)&:=-\sin(b(x))X(x)+\cos(b(x))n(x)
        \qquad\forall\,x\in M_\zeta.
    \end{split}
\end{align}
%
These expressions define a map for every real-valued $b$.
Below we restrict its norm so that the map is an immersion
and the normal projections we use are nonzero. For each
$x\in M_\zeta$, set
%
\begin{align*}
    \gamma_x(t):=\cos t\,X(x)+\sin t\,n(x)
    \qquad\forall\,t\in\mathbb R.
\end{align*}
%
Then $\gamma_x$ is a unit-speed spherical geodesic,
$X_b(x)=\gamma_x(b(x))$, and
$V_b(x)=\dot\gamma_x(b(x))$. Indeed,
$\ddot\gamma_x=-\gamma_x$, and the spherical connection gives
%
\begin{align*}
    \nabla^{\mathbb S^4}_{\dot\gamma_x}\dot\gamma_x
    =\ddot\gamma_x+|\dot\gamma_x|^2\gamma_x=0.
\end{align*}
%
Thus $V_b(x)$ is the parallel transport of $n(x)$ from
$t=0$ to $t=b(x)$. The vector $\nu(x)$, kept constant in
$\mathbb R^5$, is perpendicular to $\gamma_x(t)$ and
$\dot\gamma_x(t)$ for every $t$, since it is perpendicular
to both $X(x)$ and $n(x)$. Consequently
%
\begin{align*}
    \nabla^{\mathbb S^4}_{\dot\gamma_x}\nu(x)
    =0+\langle\dot\gamma_x,\nu(x)\rangle\gamma_x=0.
\end{align*}
%
We therefore use the same vector $\nu(x)$ for its parallel
transport to $X_b(x)$. All inner products in the explicit
$\mathbb R^5$ formulas below are Euclidean; on vectors
tangent to $\mathbb S^4$ they equal the spherical metric.

Fix a coordinate chart on $M_\zeta$, write $\partial_i$
for its coordinate vector fields, and use indices
$i,j\in\{1,2\}$, summing over repeated indices. Set
$X_i:=dX(\partial_i)$, $b_i:=\partial_i b$ and
$\omega_i:=\omega(\partial_i)$. The Weingarten equation is
%
\begin{align*}
    \partial_i n
    =\nabla^{\mathbb S^4}_{\partial_i}n
    =-dX\bigl(S_\iota^n(\partial_i)\bigr)-\omega_i\nu.
\end{align*}
%
The first equality holds because $\langle X_i,n\rangle=0$,
and the sign in the last term is the one fixed by
$\nabla^\perp_{\partial_i}\nu=\omega_i n$.
Define the tangent-bundle endomorphism
%
\begin{align*}
    \mathsf P_b:=\cos b\,\operatorname{Id}_{TM_\zeta}
                      -\sin b\,S_\iota^n.
\end{align*}
%
Differentiating \eqref{eq:s4-transverse-graph} now gives
%
\begin{align*}
    \partial_iX_b
    &=\cos b\,X_i+\sin b\,\partial_i n
                         +b_i(-\sin b\,X+\cos b\,n)\\
    &=\cos b\,dX(\partial_i)
       -\sin b\,dX\bigl(S_\iota^n(\partial_i)\bigr)
       -\sin b\,\omega_i\nu+b_iV_b.
\end{align*}
%
Equivalently,
%
\begin{align}
    \label{eq:s4-graph-tangent}
    \partial_iX_b
    =dX\bigl(\mathsf P_b(\partial_i)\bigr)
    -\sin b\,\omega_i\nu+b_iV_b.
\end{align}
%
The three terms in \eqref{eq:s4-graph-tangent} lie in
mutually orthogonal subspaces of $\mathbb R^5$: the first
is tangent to $X$, the second is proportional to $\nu$,
and the third belongs to $\operatorname{span}\{X,n\}$.
Thus the pullback metric $g_b:=X_b^*g_{\mathbb S^4}$ has
coefficients
%
\begin{align}
    \label{eq:s4-graph-metric-exact}
    (g_b)_{ij}
    =g\bigl(\mathsf P_b(\partial_i),\mathsf P_b(\partial_j)\bigr)
    +\sin^2b\,\omega_i\omega_j+b_ib_j.
\end{align}
%
There is a constant $c_0>0$, independent of $(v,\zeta)$
and $m$, such that $X_b$ is an immersion whenever
$\|b\|_{\mathcal X_\zeta}\leq c_0\tau$. To see this, the
norm definition, $\rho\geq\tau$, and metric equivalence give
%
\begin{align*}
    |b|\leq c_0\tau,\qquad
    |db|_g\leq Cc_0\tau/\rho\leq Cc_0.
\end{align*}
%
Since $|S_\iota^n|_g\leq C$, decreasing $c_0$ ensures
$|\mathsf P_b(Y)|_g\geq |Y|_g/2$ for every tangent
vector $Y$. Formula \eqref{eq:s4-graph-metric-exact} then
gives $g_b(Y,Y)\geq g(Y,Y)/4$. It also gives a uniform
upper bound $g_b\leq Cg$. The small constant $c_0$ will be
decreased below to satisfy the other stated requirements.

We next project the transported vectors $V_b$ and $\nu$
onto the normal bundle of $X_b$ within $\mathbb S^4$.
If $(g_b)^{ij}$ denotes the inverse matrix to $(g_b)_{ij}$,
the tangential projection of a vector $Z$ along $X_b$ is
$ (g_b)^{ij}\langle Z,\partial_iX_b\rangle\partial_jX_b$.
In the present case, \eqref{eq:s4-graph-tangent} gives
%
\begin{align*}
    \langle V_b,\partial_iX_b\rangle=b_i,
    \qquad
    \langle\nu,\partial_iX_b\rangle=-\sin b\,\omega_i.
\end{align*}
%
The required normal projections are therefore
%
\begin{align}
    \label{eq:s4-testing-projections}
    \begin{split}
        U_b&=V_b-(g_b)^{ij}b_i\partial_jX_b,\\
        W_b&=\nu+\sin b\,(g_b)^{ij}\omega_i\partial_jX_b.
    \end{split}
\end{align}
%
Both projections remain tangent to $\mathbb S^4$ and are
perpendicular to $dX_b(TM_\zeta)$. They are defined
independently of the coordinate chart because orthogonal
projection is intrinsic. To calculate their lengths, set
$T_b:=(g_b)^{ij}b_i\partial_jX_b$ and
$T_\omega:=(g_b)^{ij}\omega_i\partial_jX_b$.
Then
$|T_b|^2=\langle V_b,T_b\rangle=|db|_{g_b}^2$,
$|T_\omega|^2=|\omega|_{g_b}^2$, and
$\langle\nu,T_\omega\rangle=-\sin b\,|\omega|_{g_b}^2$.
Also
$\langle V_b,T_\omega\rangle
=\langle T_b,T_\omega\rangle=\langle db,\omega\rangle_{g_b}$
and
$\langle T_b,\nu\rangle=-\sin b\langle db,\omega\rangle_{g_b}$.
For example,
%
\begin{align*}
    |U_b|^2
    &=|V_b-T_b|^2
      =1-2\langle V_b,T_b\rangle+|T_b|^2
      =1-|db|_{g_b}^2,\\
    \langle U_b,W_b\rangle
    &=\sin b\,\langle V_b,T_\omega\rangle
      -\langle T_b,\nu\rangle
      -\sin b\,\langle T_b,T_\omega\rangle=\sin b\,\langle db,\omega\rangle_{g_b}.
\end{align*}
%
Together with the same expansion for $W_b$, this proves
%
\begin{align}
    \label{eq:s4-testing-lengths}
    \begin{split}
        |U_b|^2&=1-|db|_{g_b}^2,\\
        |W_b|^2&=1-\sin^2b\,|\omega|_{g_b}^2,\\
        \langle U_b,W_b\rangle
        &=\sin b\,\langle db,\omega\rangle_{g_b}.
    \end{split}
\end{align}
%
On the indicated ball, $|db|_{g_b}\leq Cc_0$ and
$|\sin b|\,|\omega|_{g_b}\leq Cc_0\tau$.
It follows that
%
\begin{align*}
    |U_b|^2|W_b|^2-\langle U_b,W_b\rangle^2
    \geq 1-Cc_0^2\geq\frac12
\end{align*}
%
after decreasing $c_0$. Hence $U_b(x)$ and $W_b(x)$
are linearly independent and span the normal plane of $X_b$
at every $x\in M_\zeta$. Define the unit normal fields
%
\begin{align}
    \label{eq:s4-testing-normals}
    n_b(x):=\frac{U_b(x)}{|U_b(x)|},\qquad
    \nu_b(x):=\frac{W_b(x)}{|W_b(x)|}
    \qquad\forall\,x\in M_\zeta.
\end{align}
%
They form a basis of the normal plane, but need not be
orthogonal. At $b=0$ they reduce to $n$ and $\nu$.

Let $\mathbf H(X_b)$ denote the mean-curvature vector of
$X_b$ in $\mathbb S^4$, as a vector field along the map
with domain $M_\zeta$. We retain the trace convention from
Section~\ref{sec:s4-geometry}. For
$b\in\mathcal X_\zeta$ with
$\|b\|_{\mathcal X_\zeta}\leq c_0\tau$, define functions
on $M_\zeta$ by
%
\begin{align}
    \label{eq:s4-component-maps}
    \mathscr{P}_{\zeta,\phi}(b)
    :=\langle\mathbf{H}(X_b),n_b\rangle,
    \qquad
    \mathscr{T}_{\zeta,\phi}(b)
    :=\langle\mathbf{H}(X_b),\nu_b\rangle.
\end{align}
%
These are scalar products with the two unit normals; they
are not being used as coefficients of a decomposition in
the generally nonorthogonal basis $(n_b,\nu_b)$. Since
that basis spans the normal plane,
%
\begin{align*}
    \mathbf H(X_b)=0
    \quad\Leftrightarrow\quad
    \mathscr P_{\zeta,\phi}(b)=0
    \ \text{and}\ \mathscr T_{\zeta,\phi}(b)=0.
\end{align*}
%
At zero the two functions are $H_{\iota,n}$ and
$H_{\zeta,\phi}$, respectively. In particular the latter
agrees with the scalar mean curvature defined in
Proposition~\ref{prop:s4-scalar-input}.
Figure~\ref{fig:s4-normal-displacement} illustrates the
displacement defining $X_b$.

% A schematic of normal geodesics to the fixed hypersurface, not a metric model.
\begin{figure}[htbp]
    \centering
    \begingroup
    \definecolor{normalteal}{RGB}{0,112,116}
    \definecolor{normalorange}{RGB}{190,102,36}
    \begin{tikzpicture}[
        x=1cm,y=1cm,>=Stealth,
        every node/.style={font=\small},
        reference/.style={black!40,line width=0.6pt},
        base edge/.style={normalteal,line width=0.9pt},
        new edge/.style={normalorange,line width=0.9pt}
    ]
        % One tangential coordinate is suppressed.  The plane is schematic.
        \path[fill=black!3]
            (0,0)--(7.0,0)--(8.25,1.55)--(1.25,1.55)--cycle;
        \draw[reference]
            (0,0)--(7.0,0)--(8.25,1.55)--(1.25,1.55)--cycle;
        \node[anchor=north west,text=black!65] at (0.12,-0.06)
            {fixed hypersurface $\mathcal N$};
        % A local portion of the surface already deformed within N.
        \path[fill=normalteal!12]
            (0.85,0.30)--(6.50,0.30)--(7.10,1.02)--(1.45,1.02)--cycle;
        \draw[base edge]
            (0.85,0.30)--(6.50,0.30)--(7.10,1.02)--(1.45,1.02)--cycle;
        % The corresponding portion after a nonconstant normal displacement.
        \path[fill=normalorange!12]
            (1.03,2.05)
            .. controls (2.70,2.48) and (4.20,2.04) .. (6.75,2.35)
            --(7.34,3.07)
            .. controls (4.82,2.76) and (3.28,3.20) .. (1.62,2.77)
            --cycle;
        \draw[new edge]
            (1.03,2.05)
            .. controls (2.70,2.48) and (4.20,2.04) .. (6.75,2.35)
            --(7.34,3.07)
            .. controls (4.82,2.76) and (3.28,3.20) .. (1.62,2.77)
            --cycle;
        % Two further geodesics indicate correspondence between points.
        \draw[black!35,dashed]
            (1.45,1.02) .. controls (1.45,1.65) and (1.48,2.40) .. (1.62,2.77);
        \draw[black!35,dashed]
            (6.50,0.30) .. controls (6.50,1.04) and (6.55,1.98) .. (6.75,2.35);
        % A marked point and its normal geodesic, for b(x)>0.
        \coordinate (basepoint) at (3.65,0.63);
        \coordinate (endpoint) at (3.93,2.61);
        \draw[black!55,->,line width=0.65pt]
            (basepoint)--(3.65,1.29);
        \node[anchor=east,text=black!70] at (3.52,1.13) {$n(x)$};
        \draw[normalorange,->,line width=1.05pt]
            (basepoint) .. controls (3.65,1.35) and (3.69,2.09) .. (endpoint);
        \fill[normalteal] (basepoint) circle[radius=1.6pt];
        \fill[normalorange] (endpoint) circle[radius=1.6pt];
        \node[anchor=north] at (3.65,-0.38) {$X(x)=G(x)$};
        \draw[black!40] (3.65,-0.28)--(3.65,0.51);
        \node[anchor=south] at (3.93,3.27) {$X_b(x)$};
        \draw[black!40] (3.93,3.17)--(3.93,2.72);
        \node[anchor=west] at (4.11,1.86)
            {geodesic of length $|b(x)|$};
        % Labels are outside the drawing to keep the geometric comparison clear.
        \node[anchor=west,align=left,text=normalteal] at (8.57,0.62)
            {$M_{\zeta,\phi}=G(M_\zeta)$\\surface in $\mathcal N$};
        \draw[normalteal!70] (8.43,0.66)--(7.02,0.89);
        \node[anchor=west,align=left,text=normalorange] at (8.57,2.64)
            {$X_b(M_\zeta)$\\displaced surface};
        \draw[normalorange!70] (8.43,2.68)--(7.31,2.88);
    \end{tikzpicture}
    \endgroup
    \caption{The displacement normal to the fixed hypersurface
    $\mathcal N$. The teal patch belongs to the surface
    $M_{\zeta,\phi}$ obtained by the preceding deformation within
    $\mathcal N$. Each point $X(x)=G(x)$ moves along the spherical
    geodesic with initial unit velocity $n(x)$ to
    $X_b(x)=\cos(b(x))X(x)+\sin(b(x))n(x)$.
    A point with $b(x)>0$ is shown; the function $b$ may have
    either sign. The gray plane represents $\mathcal N$
    schematically, and the displacement is exaggerated.
    The drawing does not represent the spherical metric or
    all three dimensions of $\mathcal N$.}
    \label{fig:s4-normal-displacement}
\end{figure}

\begin{lemma}[Linearization without a minimality assumption]
    \label{lem:s4-component-linearization}
    For every $h\in\mathcal{X}_\zeta$,
    %
    \begin{align}
        \label{eq:s4-component-linearization}
        D\mathscr{P}_{\zeta,\phi}(0)h=P_{\zeta,\phi}h,
        \qquad D\mathscr{T}_{\zeta,\phi}(0)h=Bh.
    \end{align}
    %
    These formulas require no minimality of $M_{\zeta,\phi}$.
\end{lemma}

\begin{proof}
    Fix $h\in\mathcal X_\zeta$ and consider $b=th$ for
    $t$ in an interval about zero satisfying
    $|t|\|h\|_{\mathcal X_\zeta}<c_0\tau$. At $t=0$
    the velocity of $X_{th}$ is $hn$. The ordinary derivative
    of $V_{th}$ is $-hX$, so its spherical covariant
    derivative is $-hX+\langle hn,n\rangle X=0$.

    In the tangential term of $U_{th}$ there is a factor
    $t\partial_i h$. At zero, its derivative is consequently
    $g^{ij}(\partial_i h)X_j=dX(\nabla^g h)$; the derivatives
    of the inverse metric and of $\partial_jX_{th}$ are
    multiplied by $t=0$. Likewise, differentiating the
    tangential term of $W_{th}$ differentiates only
    $\sin(th)$ at zero. Thus
    %
    \begin{align*}
        \left.\nabla^{\mathbb{S}^4}_t U_{th}\right|_{t=0}
        =-dX(\nabla^g h),
        \qquad
        \left.\nabla^{\mathbb{S}^4}_t W_{th}\right|_{t=0}
        =h\,dX(\omega^{\sharp_g}).
    \end{align*}
    %
    Here $\omega^{\sharp_g}$ is defined by
    $g(\omega^{\sharp_g},Y)=\omega(Y)$ for every tangent
    vector $Y$. Both displayed derivatives are tangent to
    $X$. Hence, differentiating the squared lengths gives
    %
    \begin{align*}
        \left.\frac{d}{dt}\right|_{t=0}|U_{th}|^2
        &=2\langle-dX(\nabla^g h),n\rangle=0,\\
        \left.\frac{d}{dt}\right|_{t=0}|W_{th}|^2
        &=2\langle h\,dX(\omega^{\sharp_g}),\nu\rangle=0.
    \end{align*}
    %
    Since $U_0=n$ and $W_0=\nu$ are unit vectors, normalization
    does not change these first derivatives. Differentiating
    each scalar component now yields
    %
    \begin{align*}
        D\mathscr P_{\zeta,\phi}(0)h
        &=\left\langle
        \left.\nabla_t^{\mathbb S^4}\right|_{t=0}
        \mathbf H(X_{th}),n\right\rangle
        -\langle\mathbf H(X),dX(\nabla^g h)\rangle,\\
        D\mathscr T_{\zeta,\phi}(0)h
        &=\left\langle
        \left.\nabla_t^{\mathbb S^4}\right|_{t=0}
        \mathbf H(X_{th}),\nu\right\rangle
        +\langle\mathbf H(X),h\,dX(\omega^{\sharp_g})\rangle.
    \end{align*}
    %
    The last term in each line is zero because the
    mean-curvature vector is normal to $X$. The first terms
    are the $n$ and $\nu$ components of $J(hn)$.
    Lemma~\ref{lem:s4-normal-linearization} gives
    $J(hn)=(Bh)\nu+(P_{\zeta,\phi}h)n$, proving both formulas.
    No step has used $\mathbf H(X)=0$.
\end{proof}

We now estimate the terms left after subtracting the value
and derivative at $b=0$. Define
%
\begin{align}
    \label{eq:s4-transverse-remainders}
    \begin{split}
        Q_n(b)&=\mathscr{P}_{\zeta,\phi}(b)-H_{\iota,n}-P_{\zeta,\phi}b,\\
        Q_\nu(b)&=\mathscr{T}_{\zeta,\phi}(b)-H_{\zeta,\phi}-Bb.
    \end{split}
\end{align}
%
The fixed pair $(\zeta,\phi)$ is suppressed in this notation.
Both remainders vanish at zero and have zero derivative there.
The next estimate quantifies their quadratic dependence on
$b$, uniformly over the scalar deformations allowed in
$\mathcal B_m$.

\begin{lemma}[Quadratic estimates for the two components]
    \label{lem:s4-transverse-quadratic}
    There are constants $c_0,C>0$ such that, for every
    $b,\widetilde b\in\mathcal X_\zeta$ with
    $\|b\|_{\mathcal{X}_\zeta},
    \|\widetilde b\|_{\mathcal{X}_\zeta}\leq c_0\tau$,
    %
    \begin{align}
        \label{eq:s4-transverse-quadratic}
        \begin{split}
            \|Q_\eta(b)-Q_\eta(\widetilde b)\|_{\mathcal{Z}_\zeta}
            \leq{}& C\tau^{-1}
            \bigl(\|b\|_{\mathcal{X}_\zeta}
            +\|\widetilde b\|_{\mathcal{X}_\zeta}\bigr)
            \|b-\widetilde b\|_{\mathcal{X}_\zeta},
            \qquad \eta\in\{n,\nu\}.
        \end{split}
    \end{align}
    %
    The constants are uniform over $\mathcal{B}_m$.
    Only the $C^{2,\beta}$ bound for the scalar correction
    $\phi$ is used. In particular, setting $\widetilde b=0$
    gives
    %
    \begin{align*}
        \|Q_\eta(b)\|_{\mathcal Z_\zeta}
        \leq C\tau^{-1}\|b\|_{\mathcal X_\zeta}^2
        \qquad\forall\,\eta\in\{n,\nu\}.
    \end{align*}
\end{lemma}

\begin{proof}
    Fix $x\in M_\zeta$ and set $r:=\rho(x)$. Choose one
    of the uniformly controlled charts for the initial
    surface from Lemma~\ref{lem:s4-initial-geometry}. After
    rescaling its independent variables, denote its map by
    $\kappa_x:D\to M_\zeta$, where $D\subset\mathbb R^2$
    is a fixed disk and $\kappa_x(0)=x$. We use a smaller
    concentric disk for the local norm; all estimates are
    obtained on a slightly larger disk of the same fixed
    size. Set
    %
    \begin{align*}
        \overline X:=\frac{X\circ\kappa_x-X(x)}r,\qquad
        \overline X_b:=\frac{X_b\circ\kappa_x-X(x)}r,\qquad
        \overline g_b:=r^{-2}\kappa_x^*g_b.
    \end{align*}
    %
    The rescaled immersion $\overline X$ has bounded
    $C^{2,\beta}$ norm and uniformly nondegenerate first
    derivative. This follows from
    the initial geometry and \eqref{eq:s4-dimensionless-phi}.
    In the computations in this chart, write $X,n,b$ also
    for their compositions with $\kappa_x$, and let
    $\partial_i$ denote differentiation with respect to its
    coordinates. The field $n$ is obtained by composing $X$
    with a fixed smooth field on $\mathcal N$. Extend that
    field smoothly to a fixed neighborhood in $\mathbb R^5$
    of the smaller band containing all the images of $X$,
    and denote the extension by $\widetilde n$. The chain
    rule gives
    %
    \begin{align*}
        \partial_i n&=d\widetilde n_{X}(\partial_iX),\\
        \partial_{ij}n
        &=d^2\widetilde n_{X}(\partial_iX,\partial_jX)
        +d\widetilde n_{X}(\partial_{ij}X).
    \end{align*}
    %
    Thus the value and derivatives of $n$ through order two
    are controlled by those of $X$. The derivatives of
    $\widetilde n$ are bounded because the extension and
    neighborhood are fixed. These formulas compute the derivatives of the original
    normal field along the parametrized surface and do not depend on the choice
    of extension. In particular, no second derivative of
    the other normal $\nu$ is needed.

    Differentiating \eqref{eq:s4-transverse-graph} twice gives
    %
    \begin{align*}
        \partial_iX_b
        &=\cos b\,\partial_iX+\sin b\,\partial_i n
        +b_iV_b,\\
        \partial_{ij}X_b
        &=\cos b\,\partial_{ij}X+\sin b\,\partial_{ij}n
        +b_{ij}V_b-b_ib_jX_b\\
        &\quad +b_i(-\sin b\,\partial_jX
        +\cos b\,\partial_j n)
        +b_j(-\sin b\,\partial_iX
        +\cos b\,\partial_i n).
    \end{align*}
    %
    These expressions involve no derivative of $b$ or $X$
    of order greater than two. The coefficients of the metric
    involve only first derivatives of $X_b$. The vector $\nu$
    can be computed from $X,n,\partial_1\overline X,
    \partial_2\overline X$ and the chosen orientation:
    it is the oriented unit vector orthogonal to these four
    vectors. Their Gram matrix is uniformly nondegenerate,
    because $X,n$ are orthonormal and perpendicular to the
    rescaled tangent vectors, whose Gram matrix is the
    uniformly elliptic rescaled induced metric. Thus this
    algebraic formula for $\nu$ has uniformly bounded
    derivatives with respect to its entries. The
    one-form $\omega$ occurring in $W_b$ can also be computed
    from first derivatives, using
    $\omega_i=\langle S(X_i),\nu\rangle$.
    Thus the testing normals in
    \eqref{eq:s4-testing-normals} are algebraic functions of
    these first derivatives, the fixed ambient tensors, and
    the values of $b,db$.

    For either $\eta_b=n_b$ or $\eta_b=\nu_b$,
    %
    \begin{align}
        \label{eq:s4-tested-curvature-formula}
        \langle\mathbf{H}(X_b),\eta_b\rangle
        =(g_b)^{ij}\langle\partial_{ij}X_b,\eta_b\rangle.
    \end{align}
    %
    Indeed, the Euclidean formula for the spherical
    mean-curvature vector is
    %
    \begin{align*}
        \mathbf H(X_b)
        =(g_b)^{ij}
        \bigl(\partial_{ij}X_b-\Gamma(g_b)^k_{ij}
                       \partial_kX_b\bigr)+2X_b.
    \end{align*}
    %
    The terms containing the Christoffel symbols are tangent
    to $X_b$, and $\eta_b$ is perpendicular to $X_b$ itself.
    Pairing with $\eta_b$ therefore gives
    \eqref{eq:s4-tested-curvature-formula}. In the rescaled
    chart it becomes
    %
    \begin{align*}
        r\langle\mathbf{H}(X_b),\eta_b\rangle
        =\overline g_b^{ij}
        \langle\partial_{ij}\overline X_b,\eta_b\rangle.
    \end{align*}
    %
    Thus the rescaled scalar functions are obtained by smooth
    algebraic expressions in the values and derivatives of
    $X$ and $b$ through order two, and the first derivatives
    determining $\nu$. The determinant of $\overline g_b$
    and the lengths of $U_b,W_b$ are uniformly bounded away
    from zero,
    by \eqref{eq:s4-graph-metric-exact} and
    \eqref{eq:s4-testing-lengths}.

    More explicitly, write $b=r\widehat b$ on this chart and
    multiply mean curvature by $r$. Substituting
    $X_b=\cos(r\widehat b)X+\sin(r\widehat b)n$ into
    \eqref{eq:s4-tested-curvature-formula} yields a smooth map
    of $\widehat b$ from $C^{2,\beta}$ to $C^{0,\beta}$ whose
    second Fr\'echet derivative is uniformly bounded on a fixed
    small ball. To see that dilation introduces no singular
    coefficient, write the rescaled graph difference as
    %
    \begin{align*}
        \frac{X_b-X}{r}
        =\frac{\cos(r\widehat b)-1}{r}X
        +\frac{\sin(r\widehat b)}{r}n.
    \end{align*}
    %
    The two scalar coefficients extend smoothly to $r=0$,
    with limits $0$ and $\widehat b$. One can see this without
    dividing a Taylor remainder by $r$: for a scalar $a$,
    %
    \begin{align*}
        \frac{\sin(ra)}r
        &=a\int_0^1\cos(tra)\,dt,\\
        \frac{\cos(ra)-1}r
        &=-ra^2\int_0^1(1-t)\cos(tra)\,dt.
    \end{align*}
    %
    Both formulas and every fixed derivative in $a$ are
    bounded uniformly for bounded $r,a$. All other
    coefficients are the rescaled derivatives of $X,n$
    already estimated, or inverses and square roots of
    uniformly positive quantities.

    To state the resulting Taylor estimate precisely, let
    $\mathcal M_{r,\eta}(\widehat b)$ denote $r$ times the
    scalar product with $\eta_b$, where $\eta\in\{n,\nu\}$.
    The formulas above express it as a function
    $\mathcal F_{r,\eta}(y,\mathbf z_{\widehat b}(y))$,
    where $y\in D$ and
    %
    \begin{align*}
        \mathbf z_{\widehat b}
        :=(\widehat b,\partial_1\widehat b,
          \partial_2\widehat b,\partial_{11}\widehat b,
          \partial_{12}\widehat b,\partial_{22}\widehat b).
    \end{align*}
    %
    The function $\mathcal F_{r,\eta}$ is smooth in these
    six variables. Its derivatives through order three in
    those variables have uniformly bounded local
    $C^{0,\beta}$ norms on the fixed small ball. This follows
    from the explicit formulas: the operations are products,
    trigonometric functions, inverses of nondegenerate metric
    matrices, and square roots of quantities bounded away
    from zero. The H\"older product and composition estimates
    consequently give
    %
    \begin{align*}
        \|D^2\mathcal M_{r,\eta}(\widehat b)
        [h,k]\|_{C^{0,\beta}}
        \leq C\|h\|_{C^{2,\beta}}\|k\|_{C^{2,\beta}}.
    \end{align*}
    %
    For example, set
    $(\mathcal D_1,\ldots,\mathcal D_6)
    :=(\operatorname{Id},\partial_1,\partial_2,
         \partial_{11},\partial_{12},\partial_{22})$.
    The second Fr\'echet derivative is exactly
    %
    \begin{align*}
        D^2\mathcal M_{r,\eta}(\widehat b)[h,k]
        =\sum_{a,c=1}^6
        (\partial_{z_a z_c}\mathcal F_{r,\eta})
            (y,\mathbf z_{\widehat b}(y))
        \mathcal D_a h\,\mathcal D_c k.
    \end{align*}
    %
    Estimating each of the finitely many products proves the
    preceding bound. The derivatives in this formula are
    derivatives with respect to the six variables, not
    additional spatial derivatives of $X$ or $\phi$.

    If $\mathcal Q_{r,\eta}$ is the remainder obtained by
    subtracting the value and derivative of $\mathcal M_{r,\eta}$
    at zero, then
    %
    \begin{align*}
        \mathcal Q_{r,\eta}(a)-\mathcal Q_{r,\eta}(c)
        =\int_0^1\bigl(
        D\mathcal M_{r,\eta}(c+t(a-c))
        -D\mathcal M_{r,\eta}(0)\bigr)(a-c)\,dt.
    \end{align*}
    %
    Applying the second-derivative bound along the line
    segment from zero to $c+t(a-c)$ bounds the norm of
    this expression by
    $C(\|a\|_{C^{2,\beta}}+\|c\|_{C^{2,\beta}})
    \|a-c\|_{C^{2,\beta}}$.
    Take $a=b/r$ and $c=\widetilde b/r$.
    Since $\mathcal Q_{r,\eta}(b/r)=rQ_\eta(b)$, multiplication
    by $r$ gives on the original chart
    %
    \begin{align*}
        \begin{split}
            r^2\|Q_\eta(b)-Q_\eta(\widetilde b)\|_{0,\beta;r}
            \leq{}&\frac Cr\bigl(\|b\|_{2,\beta;r}
            +\|\widetilde b\|_{2,\beta;r}\bigr)
            \|b-\widetilde b\|_{2,\beta;r}.
        \end{split}
    \end{align*}
    %
    Here $\|\cdot\|_{k,\beta;r}$ denotes the ordinary
    $C^{k,\beta}$ norm after rescaling the independent variable
    by $r^{-1}$, with no rescaling of function values. Passing
    to the frozen-scale norms and using $r\geq c\tau$ proves
    \eqref{eq:s4-transverse-quadratic}.

    The normal $\nu$ is used only algebraically in
    \eqref{eq:s4-testing-projections}; it is not differentiated
    twice. No bound on $\nabla S_\iota^\nu$, $\nabla^2\nu$, or
    $\nabla^3\phi$ has entered the argument. The normalization
    of the testing fields is also included: because
    $\mathbf{H}(X_b)$ is normal, its two scalar products equal
    %
    \begin{align*}
        \frac{\langle\mathbf{H}(X_b),V_b\rangle}
        {\sqrt{1-|db|_{g_b}^2}},
        \qquad
        \frac{\langle\mathbf{H}(X_b),\nu\rangle}
        {\sqrt{1-\sin^2b\,|\omega|_{g_b}^2}}.
    \end{align*}
    %
    Their normalization corrections start at quadratic order
    and are covered by the same Taylor estimate.
\end{proof}

We also need the nonlinear maps to preserve the symmetry
conditions of the linear problem. Fix $a\in\Gamma_m$ and
let $b\circ a=\chi(a)b$. Equivariance of $X$ and the
transformation rule for $n$ give
$X_b\circ a=a\circ X_b$: the cosine is even, and the
signs from $\sin(b\circ a)$ and $n\circ a$ cancel.
The same computation in the formula for $V_b$ gives
%
\begin{align*}
    V_b(a(x))
    &=\chi(a)\,d a_{X_b(x)}(V_b(x)),\\
    \nu(a(x))&=d a_{X_b(x)}(\nu(x))
    \qquad\forall\,x\in M_\zeta.
\end{align*}
%
In the second line $\nu$ denotes the transported field
along $X_b$. Since $a$ is an isometry and $X_b$ is
equivariant, its differential maps the tangent and normal
planes of $X_b$ at $x$ onto the corresponding planes at
$a(x)$. Orthogonal projection and normalization therefore
give the same two transformation rules for $n_b$ and
$\nu_b$, respectively. The mean-curvature vector itself
transforms by $d a$. Taking scalar products proves
%
\begin{align}
    \label{eq:s4-component-parities}
    \mathscr{P}_{\zeta,\phi}(b)\circ a
    =\chi(a)\mathscr{P}_{\zeta,\phi}(b),
    \qquad
    \mathscr{T}_{\zeta,\phi}(b)\circ a
    =\mathscr{T}_{\zeta,\phi}(b).
\end{align}
%
for every $a\in\Gamma_m$. The constant and linear terms
subtracted in \eqref{eq:s4-transverse-remainders} have the
same symmetries by Section~\ref{sec:s4-geometry}. Thus
$Q_n(b)\in\mathcal Y_\zeta$, whereas $Q_\nu(b)$ is an
invariant element of $\mathcal Z_\zeta$.

\subsection{The normal correction and the remaining scalar equation}

For each fixed $(v,\zeta)\in\mathcal B_m$, we now solve
$\mathscr P_{\zeta,\phi}(b)=0$. Its source at $b=0$ has
weighted size $O(\tau^2)$, its linearization has a uniformly
bounded inverse, and its quadratic remainder has Lipschitz
constant $O(\tau)$ on a ball of radius $O(\tau^2)$.
A contraction therefore selects a unique small function
$b(v,\zeta)$. We then retain the change in the other
component of mean curvature as a function of the original
choice $(v,\zeta)$.

\begin{proposition}[Correction in the direction $n$]
    \label{prop:s4-transverse-elimination}
    There exist constants $K>0$ and $m_0$ such that for every
    $m\geq m_0$ and every $(v,\zeta)\in\mathcal{B}_m$ there is
    a unique solution $b=b(v,\zeta)$ in the ball
    %
    \begin{align}
        \label{eq:s4-transverse-ball}
        \{b\in\mathcal{X}_\zeta:\|b\|_{\mathcal{X}_\zeta}
        \leq K\tau_\zeta^2\}
    \end{align}
    %
    of
    %
    \begin{align}
        \label{eq:s4-transverse-equation}
        \mathscr{P}_{\zeta,\phi}(b)=0.
    \end{align}
    %
    The change in the other scalar mean-curvature component caused
    by this displacement is the following invariant function
    on $M_\zeta$:
    %
    \begin{align}
        \label{eq:s4-feedback-definition}
        \mathcal{R}_\zeta(\phi)
        =\mathscr{T}_{\zeta,\phi}(b(v,\zeta))-H_{\zeta,\phi}.
    \end{align}
    %
    It obeys
    %
    \begin{align}
        \label{eq:s4-feedback-estimates}
        \|\mathcal{R}_\zeta(\phi)\|_{\mathcal{Z}_\zeta}
        \leq C\tau_\zeta^3,
        \qquad
        \|\mathcal{R}_\zeta(\phi)\|_{Y_\zeta}
        \leq C\tau_\zeta^{2+\alpha/2}.
    \end{align}
    %
    In particular $\mathcal{R}_\zeta(\phi)$ has zero first
    angular harmonic on every neck.

    For fixed $m$, these outputs are continuous with any loss
    of H\"older exponent. Precisely, if $0<\beta'<\beta$,
    $(v_j,\zeta_j),(v,\zeta)\in\mathcal{B}_m$,
    $\zeta_j\to\zeta$ and $v_j\to v$ in
    $C^{2,\beta'}(M_0)$, then
    %
    \begin{align}
        \label{eq:s4-transverse-continuity}
        \begin{aligned}
            b(v_j,\zeta_j)\circ F_{\zeta_j}
            &\rightarrow b(v,\zeta)\circ F_\zeta
            &&\qquad\text{in }C^{2,\beta'}(M_0),\\
            \mathcal{R}_{\zeta_j}(\phi_j)\circ F_{\zeta_j}
            &\rightarrow\mathcal{R}_\zeta(\phi)\circ F_\zeta
            &&\qquad\text{in }C^{0,\beta'}(M_0),
        \end{aligned}
    \end{align}
    %
    where $\phi_j=u_{\zeta_j}+v_j\circ F_{\zeta_j}^{-1}$.
\end{proposition}

\begin{proof}
    Let $C_L$ be the inverse bound in
    \eqref{eq:s4-actual-inverse}, and let $C_H$ be the forcing
    constant in \eqref{eq:s4-forcing-and-mixed}. Fix
    $K>2C_LC_H$. Consider
    %
    \begin{align}
        \label{eq:s4-transverse-contraction}
        \mathcal{A}_{\zeta,\phi}(b)
        =-\mathscr{G}_{\zeta,\phi}\bigl(H_{\iota,n}+Q_n(b)\bigr).
    \end{align}
    %
    Choose $m_0$ large enough that $K\tau\leq c_0$ for
    every allowed $\zeta$ and $m\geq m_0$. The ball
    \eqref{eq:s4-transverse-ball} is then contained in
    the domain of Lemma~\ref{lem:s4-transverse-quadratic}.
    Since $Q_n(0)=0$, that lemma gives
    %
    \begin{align*}
        \|\mathcal{A}_{\zeta,\phi}(b)\|_{\mathcal{X}_\zeta}
        \leq C_LC_H\tau^2+CC_LK^2\tau^3.
    \end{align*}
    %
    The first term is less than $K\tau^2/2$ by the choice
    of $K$. Increasing $m_0$ makes the second at most
    $K\tau^2/2$. Thus $\mathcal A_{\zeta,\phi}$ maps the
    closed ball into itself. For any two points of the ball,
    the same estimates give
    %
    \begin{align*}
        \|\mathcal A_{\zeta,\phi}(b)
                  -\mathcal A_{\zeta,\phi}(\widetilde b)
                       \|_{\mathcal X_\zeta}
        \leq2CC_LK\tau
                  \|b-\widetilde b\|_{\mathcal X_\zeta}.
    \end{align*}
    %
    This factor is at most $1/2$ for large $m$.
    Equation~\eqref{eq:s4-component-parities} ensures that
    $H_{\iota,n}+Q_n(b)\in\mathcal Y_\zeta$, so the
    inverse in the definition of $\mathcal A_{\zeta,\phi}$
    is applied to its proper domain. The closed ball is a
    complete metric space, and the contraction theorem gives
    a unique fixed point in it. Applying $P_{\zeta,\phi}$
    to the fixed-point identity gives
    %
    \begin{align*}
        P_{\zeta,\phi}b=-H_{\iota,n}-Q_n(b),
    \end{align*}
    %
    which is precisely \eqref{eq:s4-transverse-equation}.
    Conversely, any solution in this ball satisfies the
    fixed-point identity, so the uniqueness is uniqueness
    for that equation as well.

    By \eqref{eq:s4-transverse-remainders}, the remaining
    component is
    %
    \begin{align*}
        \mathcal{R}_\zeta(\phi)=Bb+Q_\nu(b).
    \end{align*}
    %
    Using \eqref{eq:s4-forcing-and-mixed} and
    \eqref{eq:s4-transverse-quadratic}, we obtain
    %
    \begin{align*}
        \|\mathcal{R}_\zeta(\phi)\|_{\mathcal{Z}_\zeta}
        \leq C\tau\|b\|_{\mathcal{X}_\zeta}
        +C\tau^{-1}\|b\|_{\mathcal{X}_\zeta}^2
        \leq C\tau^3.
    \end{align*}
    %
    This proves the first bound in
    \eqref{eq:s4-feedback-estimates}. To convert it to the
    scalar source norm, use the definitions with the same
    frozen-scale balls:
    %
    \begin{align*}
        \|\mathcal{R}_\zeta(\phi)\|_{Y_\zeta}
        \leq C\tau^3\sup_{M_\zeta}f(\rho)^{-1}.
    \end{align*}
    %
    The weight $f$ is increasing and $\rho\geq\tau$; since
    $0<\tau<1$,
    %
    \begin{align*}
        f(\tau)=\max\{\tau^{3/2},\tau^{1-\alpha/2}\}
        =\tau^{1-\alpha/2}.
    \end{align*}
    %
    This gives the second bound. The function $\mathcal R_\zeta(\phi)$ is invariant by
    \eqref{eq:s4-component-parities}. In particular, a
    quarter-turn around each neck preserves the axial
    coordinate and the function. If $a_1(\xi),a_2(\xi)$
    were its coefficients of $\cos\theta,\sin\theta$,
    this invariance would require the vector $(a_1,a_2)$
    to be fixed by a rotation through $\pi/2$, and hence
    to be zero. Thus its first angular harmonic vanishes.
    It therefore belongs to the precise scalar source space
    used in Proposition~\ref{prop:s4-scalar-input}.

    It remains to prove continuity. Fix $m$ and $\beta'<\beta$.
    After pullback by $F_{\zeta_j}$, the scalar corrections
    $\phi_j$ converge in $C^{2,\beta'}$ to the correction at
    $(v,\zeta)$, by the continuity statement in
    Proposition~\ref{prop:s4-scalar-input} and the smooth
    parameter dependence of the initial surfaces. They are
    uniformly bounded in $C^{2,\beta}$, since $m$ is fixed
    and the norms in Section~\ref{sec:s4-scalar-family} are uniformly equivalent over
    the compact parameter interval. The solutions for the displacement in the direction $n$
    have the same uniform stronger bound by
    \eqref{eq:s4-transverse-ball}.

    Every subsequence of $b(v_j,\zeta_j)\circ F_{\zeta_j}$
    therefore has a further subsequence converging in
    $C^{2,\beta'}(M_0)$. The explicit formulas
    \eqref{eq:s4-testing-projections} and
    \eqref{eq:s4-tested-curvature-formula} allow passage to the
    limit in the equation $\mathscr P_{\zeta,\phi}(b)=0$. The limit retains the
    bound $\|b\|_{\mathcal{X}_\zeta}\leq K\tau_\zeta^2$:
    the stronger H\"older seminorm is lower semicontinuous,
    and the smooth parameter identifications and metrics converge.
    Uniqueness in the contraction ball identifies this limit
    with $b(v,\zeta)\circ F_\zeta$. Thus the whole sequence
    converges. Applying the same formulas, which use
    derivatives only through order two, to the other scalar
    component proves the second convergence in
    \eqref{eq:s4-transverse-continuity}.
\end{proof}

\begin{proposition}[The remaining equation]
    \label{prop:s4-remaining-scalar-equation}
    For every $(v,\zeta)\in\mathcal B_m$ and all
    sufficiently large $m$, the map
    %
    \begin{align*}
        \widehat X_{v,\zeta}:M_\zeta&\rightarrow\mathbb S^4,\\
        \widehat X_{v,\zeta}(x)
        &:=\cos(b(v,\zeta)(x))G_{\zeta,\phi}(x)
          +\sin(b(v,\zeta)(x))n(x)
          \qquad\forall\,x\in M_\zeta
    \end{align*}
    %
    is a $C^{2,\beta}$ embedding. Its image is a closed
    connected surface, and it is minimal in $\mathbb S^4$
    if and only if
    %
    \begin{align}
        \label{eq:s4-remaining-scalar-equation}
        H_{\zeta,\phi}+\mathcal{R}_\zeta(\phi)=0
        \quad\text{on }M_\zeta.
    \end{align}
    %
    Its displacement in the direction $n$ satisfies
    %
    \begin{align}
        \label{eq:s4-transverse-two-jet-final}
        |(\nabla^{g_\zeta})^j b(v,\zeta)(x)|_{g_\zeta}
        \leq C\tau_\zeta^2\rho_\zeta(x)^{-j}
        \qquad\forall\,x\in M_\zeta,\quad j\in\{0,1,2\},
    \end{align}
    %
    with the corresponding local scaled H\"older bounds.
\end{proposition}

\begin{proof}
    Write $b:=b(v,\zeta)$ and $G:=G_{\zeta,\phi}$.
    By Proposition~\ref{prop:s4-scalar-input}, the map $G$
    is an embedding. Consider
    %
    \begin{align*}
        x\mapsto
        \bigl(\mathcal F^{-1}(G(x)),b(x)\bigr)
        \quad\text{from }M_\zeta\text{ into }
        \mathbb T\times(-z_0,z_0)\times\mathbb R.
    \end{align*}
    %
    This map is an embedding: projection onto the first
    factor $\mathbb T\times(-z_0,z_0)$ gives the embedding
    $\mathcal F^{-1}\circ G$. Its injectivity and injective
    differential follow immediately from this projection,
    and its domain is compact.

    The bounds $|z\circ G|\leq C\varepsilon$ and
    $\|b\|_\infty\leq K\tau^2$ tend uniformly to zero.
    For large $m$, the displayed graph therefore lies in
    the domain of the fixed tubular embedding $\Psi$ from
    Lemma~\ref{lem:s4-tubular}. The definition of $\Psi$
    and $n(x)=N(\pi(G(x)))$ give the exact identity
    %
    \begin{align*}
        \widehat X_{v,\zeta}(x)
        =\Psi\bigl(\mathcal F^{-1}(G(x)),b(x)\bigr)
        \qquad\forall\,x\in M_\zeta.
    \end{align*}
    %
    Thus $\widehat X_{v,\zeta}$ is a composition of two
    embeddings and is itself an embedding. Its image is
    closed and connected because $M_\zeta$ is closed
    and connected.

    The component tested against $n_b$ vanishes by construction.
    Both $n_b$ and $\nu_b$ belong to the normal plane and are
    linearly independent. A vector in that plane that is
    orthogonal to both must be zero, whether or not this
    pair is orthogonal. Thus $\mathbf H(\widehat X_{v,\zeta})=0$
    is equivalent to the vanishing of its component tested
    against $\nu_b$. By definition that
    component is the left side of
    \eqref{eq:s4-remaining-scalar-equation}. Finally,
    \eqref{eq:s4-transverse-two-jet-final} is the pointwise and
    local H\"older content of the bound
    \eqref{eq:s4-transverse-ball}.
\end{proof}
%
\section{Construction of the final minimal surfaces}
\label{sec:s4-coupled}

For each pair $(v,\zeta)\in\mathcal B_m$, the preceding section
constructs an embedded surface in $\mathbb S^4$ in two steps.
First, the initial surface $M_\zeta$ is displaced within
$\mathcal N$ by the function
$\phi=u_\zeta+v\circ F_\zeta^{-1}$. Then its image is displaced
along the spherical geodesics initially normal to $\mathcal N$
by the function $b(v,\zeta)$. The second displacement is chosen
so that the mean-curvature vector of the resulting surface has
zero inner product with the normal field $n_b$ defined in
Section~\ref{sec:s4-transverse}. Its inner product with the other
normal field $\nu_b$ is
$H_{\zeta,\phi}+\mathcal R_\zeta(\phi)$, expressed as a function
on $M_\zeta$. By Proposition~\ref{prop:s4-remaining-scalar-equation},
the resulting surface is minimal if and only if this function
vanishes.

We now choose $v$ and $\zeta$ to obtain that vanishing. The function
$b$ is determined by these choices, so it is not an additional
unknown in this section. We also prove that the final embedding
is smooth, identify the topology of its image, and record the
estimates for both displacements needed in the curvature argument
of Section~\ref{sec:s4-negative-curvature}.

The choice of $\phi$ and $\zeta$ is made by a fixed-point argument.
We work on the set $\mathcal B_m$ defined in
\eqref{eq:s4-scalar-ball}, whose function coordinate belongs to the
fixed reference surface $M_0$. Thus the function space does not change
when the parameter $\zeta$ changes. The constants $\alpha,\beta$ and
the cutoff functions are fixed as in the preceding sections.
All constants denoted by $C$ in this section are independent of $m$
and of $(v,\zeta)\in\mathcal B_m$. Statements concerning compactness
and continuity are made with $m$ fixed; the equivalence constants
between weighted and unweighted H\"older norms in those arguments
may depend on $m$.

The definitions in \eqref{eq:s4-scales} and the bound for
$c_{m^{-2}}$ in Lemma~\ref{lem:s4-Green} give
%
\begin{align}
    \label{eq:s4-coupled-scale-comparison}
    \begin{split}
        \delta&=(100m)^{-1},\qquad
        e^{-1}\tau_0\leq\tau_\zeta\leq e\tau_0
        \quad\forall\,\zeta\in[-1,1],\\
        \log\tau_0&=-\frac{m^2}{2\pi}-\log m+O(1)
        \qquad\text{as }m\to\infty.
    \end{split}
\end{align}
%
In particular, for every fixed $p\geq0$ and $q>0$,
%
\begin{align*}
    \delta^{-p}\tau_0^q\rightarrow0
    \qquad\mbox{as }m\rightarrow\infty.
\end{align*}
%
Indeed, the logarithm of this expression is
$p\log(100m)+q\log\tau_0$, which tends to $-\infty$.
Consequently, a fixed power of $\delta^{-1}$ multiplied by a
positive power of $\tau_0$ still tends to zero. We will use this
fact when applying the scalar inverse estimates, whose constants
contain powers of $\delta^{-1}$.

\subsection{The remaining equation and the fixed-point map}

For $(v,\zeta)\in\mathcal B_m$, recall the displacement
$\phi(v,\zeta)$ and define the combined nonlinear remainder
$E_\zeta(v)$ by
%
\begin{align}
    \label{eq:s4-coupled-source}
    \phi(v,\zeta):=u_\zeta+v\circ F_\zeta^{-1},\qquad
    E_\zeta(v):=Q_\zeta(\phi(v,\zeta))
    +\mathcal R_\zeta(\phi(v,\zeta)).
\end{align}
%
Here $u_\zeta$ is the preliminary correction defined in
\eqref{eq:s4-initial-correction}. The function $Q_\zeta(\phi)$ is
the nonlinear remainder in the scalar mean curvature of the graph
within $\mathcal N$, defined in \eqref{eq:s4-Q-definition}.
The function $\mathcal R_\zeta(\phi)$, defined in
\eqref{eq:s4-feedback-definition}, is the change in the remaining
mean-curvature component caused by the subsequent displacement
$b(v,\zeta)$ normal to $\mathcal N$. Thus $E_\zeta(v)$ contains
exactly the terms not included in
$H_\zeta+\mathcal L_\zeta\phi$:
%
\begin{align*}
    H_{\zeta,\phi}+\mathcal R_\zeta(\phi)
    =H_\zeta+\mathcal L_\zeta\phi+E_\zeta(v)
    \qquad\text{on }M_\zeta,
    \quad \phi:=\phi(v,\zeta).
\end{align*}
%
The two summands defining $E_\zeta(v)$ are
$\Gamma_m$-invariant and belong to $Y_\zeta$.
In particular, their projections onto the angular functions
$\cos\theta$ and $\sin\theta$ on every neck vanish, as required
for the scalar estimates of
Proposition~\ref{prop:s4-scalar-input}.

The geometric meaning of $\phi$ is shown in
Figure~\ref{fig:s4-internal-displacement}. For fixed $\zeta$,
the map $G_{\zeta,\phi}$ moves each point of $M_\zeta$ along a
geodesic of $\mathcal N$ whose initial velocity is normal to
$M_\zeta$ within $\mathcal N$. Changing $\zeta$ also changes
the initial surface and the neck radius $\tau_\zeta$.
For every such choice, the second displacement $b(v,\zeta)$
is the one constructed in the preceding section.

\begin{figure}[htbp]
    \centering
    \begingroup
    \definecolor{internalteal}{RGB}{0,112,116}
    \definecolor{internalorange}{RGB}{190,102,36}
    \begin{tikzpicture}[
        x=1cm,y=1cm,>=Stealth,
        every node/.style={font=\small},
        reference/.style={black!30,line width=0.6pt},
        initial edge/.style={internalteal,line width=0.9pt},
        deformed edge/.style={internalorange,line width=0.9pt}
    ]
        % The box represents a coordinate neighborhood in the three-manifold N.
        \path[fill=black!2]
            (0,0)--(7.7,0)--(8.9,0.9)--(8.9,4.2)--(1.2,4.2)--(0,3.3)--cycle;
        \draw[reference,dashed]
            (0,0)--(1.2,0.9)--(8.9,0.9)
            (1.2,0.9)--(1.2,4.2);
        \draw[reference]
            (0,0)--(7.7,0)--(8.9,0.9)--(8.9,4.2)--(1.2,4.2)--(0,3.3)--cycle
            (0,3.3)--(7.7,3.3)--(8.9,4.2)
            (7.7,0)--(7.7,3.3);
        \node[anchor=south west,text=black!65] at (1.35,4.24)
            {coordinate neighborhood in $\mathcal N$};
        % Initial local surface patch, with artificial boundary.
        \path[fill=internalteal!12]
            (0.72,0.63)--(6.53,0.63)--(7.30,1.22)--(1.49,1.22)--cycle;
        \draw[initial edge]
            (0.72,0.63)--(6.53,0.63)--(7.30,1.22)--(1.49,1.22)--cycle;
        % The corresponding patch after displacement inside N.
        \path[fill=internalorange!12]
            (0.84,2.00)
            .. controls (2.61,2.37) and (4.13,1.92) .. (6.73,2.14)
            --(7.50,2.73)
            .. controls (4.90,2.51) and (3.38,2.96) .. (1.61,2.59)
            --cycle;
        \draw[deformed edge]
            (0.84,2.00)
            .. controls (2.61,2.37) and (4.13,1.92) .. (6.73,2.14)
            --(7.50,2.73)
            .. controls (4.90,2.51) and (3.38,2.96) .. (1.61,2.59)
            --cycle;
        % Two additional correspondences, also entirely within N.
        \draw[black!35,dashed]
            (1.49,1.22) .. controls (1.49,1.69) and (1.53,2.13) .. (1.61,2.59);
        \draw[black!35,dashed]
            (6.53,0.63) .. controls (6.53,1.08) and (6.61,1.68) .. (6.73,2.14);
        % One point with phi(x)>0 and its normal geodesic in N.
        \coordinate (startpoint) at (3.51,0.91);
        \coordinate (endpoint) at (3.81,2.48);
        \draw[black!65,->,line width=0.65pt]
            (startpoint)--(3.51,1.58);
        \node[anchor=east,text=black!70] at (3.35,1.48)
            {$\nu_\zeta(x)$};
        \draw[internalorange,->,line width=1.05pt]
            (startpoint) .. controls (3.51,1.42) and (3.58,1.99) .. (endpoint);
        \fill[internalteal] (startpoint) circle[radius=1.6pt];
        \fill[internalorange] (endpoint) circle[radius=1.6pt];
        \node[anchor=north] at (3.51,0.38) {$\jmath_\zeta(x)$};
        \draw[black!40] (3.51,0.43)--(3.51,0.80);
        \node[anchor=south] at (3.81,2.83) {$G_{\zeta,\phi}(x)$};
        \draw[black!40] (3.81,2.79)--(3.81,2.58);
        \node[anchor=west,align=left] at (4.22,1.66)
            {geodesic in $\mathcal N$\\length $|\phi(x)|$};
        \node[anchor=west,align=left,text=internalteal] at (9.18,1.08)
            {initial surface\\$M_\zeta$};
        \draw[internalteal!70] (9.04,1.12)--(7.28,1.12);
        \node[anchor=west,align=left,text=internalorange] at (9.18,2.72)
            {deformed surface\\$M_{\zeta,\phi}$};
        \draw[internalorange!70] (9.04,2.76)--(7.47,2.66);
    \end{tikzpicture}
    \endgroup
    \caption{For a fixed $\zeta$, the deformation within $\mathcal N$
    sends $\jmath_\zeta(x)$ to
    $G_{\zeta,\phi}(x)=\exp^{\mathcal N}_{\jmath_\zeta(x)}
    (\phi(x)\nu_\zeta(x))$. Both surface patches and the connecting
    geodesics lie in $\mathcal N$. The geodesic starts in a direction
    normal to $M_\zeta$ within $\mathcal N$; this direction need not
    be the coordinate direction $E_3$. A point with $\phi(x)>0$ is
    shown. The box is a schematic coordinate neighborhood, and
    the drawing does not represent the metric of $\mathcal N$.
    The patch boundaries are artificial and the displacement is
    exaggerated. In this section both $\phi$ and $\zeta$ are chosen;
    for each choice, the correction $b(v,\zeta)$ from
    Section~\ref{sec:s4-transverse} is recomputed and applied
    afterwards. That second deformation, which generally leaves
    $\mathcal N$, is not shown here.}
    \label{fig:s4-internal-displacement}
\end{figure}

Apply the linear map $\mathscr R_\zeta$ from
Proposition~\ref{prop:s4-scalar-input} with a minus sign, and write
%
\begin{align}
    \label{eq:s4-coupled-inverse-output}
    (w_{v,\zeta},d_{v,\zeta})
    :=-\mathscr R_\zeta E_\zeta(v)
    \in X_\zeta\times\mathbb R.
\end{align}
%
By \eqref{eq:s4-scalar-inverse}, these quantities satisfy
%
\begin{align}
    \label{eq:s4-coupled-linear-equation}
    \mathcal L_\zeta w_{v,\zeta}
    =-E_\zeta(v)+d_{v,\zeta}\mathcal J_\zeta
    \qquad\mbox{on }M_\zeta.
\end{align}
%
Thus $w_{v,\zeta}$ cancels $E_\zeta(v)$ in the linearized equation
up to a multiple of the prescribed function $\mathcal J_\zeta$.
The preliminary correction already satisfies
%
\begin{align*}
    H_\zeta+\mathcal L_\zeta u_\zeta
    =(-\tau_\zeta\zeta+c_{H,\zeta})\mathcal J_\zeta.
\end{align*}
%
Substituting $\phi=u_\zeta+v\circ F_\zeta^{-1}$ and
$E_\zeta(v)=-\mathcal L_\zeta w_{v,\zeta}
 +d_{v,\zeta}\mathcal J_\zeta$ into the remaining component
therefore gives the exact identity
%
\begin{align}
    \label{eq:s4-coupled-residual-identity}
    \begin{split}
        H_{\zeta,\phi}+\mathcal R_\zeta(\phi)
        &=\mathcal L_\zeta
            \bigl(v\circ F_\zeta^{-1}-w_{v,\zeta}\bigr)\\
        &\quad+
            (-\tau_\zeta\zeta+c_{H,\zeta}+d_{v,\zeta})
            \mathcal J_\zeta
            \qquad\text{on }M_\zeta.
    \end{split}
\end{align}
%
It is thus enough to choose $(v,\zeta)$ so that the function
in parentheses in the first line vanishes and the real
coefficient in the second line is zero. The nonlinear terms
used to construct $w_{v,\zeta}$ are evaluated at the input
$\phi(v,\zeta)$. We must therefore require that the returned
correction agrees with the correction in that same input.
This is the reason for using a fixed point.

Define the map, initially with values in $X_0\times\mathbb R$, by
%
\begin{align}
    \label{eq:s4-coupled-map}
    \mathcal T_m(v,\zeta):=
    \left(
        w_{v,\zeta}\circ F_\zeta,
        \frac{c_{H,\zeta}+d_{v,\zeta}}{\tau_\zeta}
    \right)
    \qquad\forall\,(v,\zeta)\in\mathcal B_m.
\end{align}
%
The first coordinate is expressed on $M_0$ using the diffeomorphism
$F_\zeta:M_0\to M_\zeta$. At a fixed point, it gives
$v\circ F_\zeta^{-1}=w_{v,\zeta}$. The second coordinate gives
$\tau_\zeta\zeta=c_{H,\zeta}+d_{v,\zeta}$.
These are precisely the two equalities that make the right-hand
side of \eqref{eq:s4-coupled-residual-identity} vanish.
We first verify that $\mathcal T_m$ takes values in
$\mathcal B_m$, so that it can be used in a fixed-point argument.

\begin{lemma}
    \label{lem:s4-coupled-self-map}
    For every sufficiently large integer $m$, the map
    $\mathcal T_m$ takes $\mathcal B_m$ into itself.
    Uniformly for $(v,\zeta)\in\mathcal B_m$, one has
    %
    \begin{align}
        \|E_\zeta(v)\|_{Y_\zeta}
        &\leq C\tau_\zeta^2,
        \label{eq:s4-coupled-source-bound}\\
        \|w_{v,\zeta}\|_{X_\zeta}
        &\leq C\delta^{-4-2\beta}\tau_\zeta^2,
        \label{eq:s4-coupled-w-bound}\\
        |d_{v,\zeta}|
        &\leq C\delta^{-2-\beta}\tau_\zeta^2,
        \label{eq:s4-coupled-d-bound}\\
        |c_{H,\zeta}|
        &\leq C\delta^{-2-\beta}\tau_\zeta^{1+\alpha/3}.
        \label{eq:s4-coupled-cH-bound}
    \end{align}
    %
    Writing $\mathcal T_m=(\mathcal T_m^1,\mathcal T_m^2)$,
    it follows that
    %
    \begin{align}
        \frac{\|\mathcal T_m^1(v,\zeta)\|_{X_0}}
        {\tau_0^{1+\alpha}}
        &\leq C\delta^{-4-2\beta}\tau_0^{1-\alpha}
        \rightarrow0,
        \label{eq:s4-coupled-function-smallness}\\
        |\mathcal T_m^2(v,\zeta)|
        &\leq C\delta^{-2-\beta}
        (\tau_0^{\alpha/3}+\tau_0)
        \rightarrow0
        \label{eq:s4-coupled-parameter-smallness}
    \end{align}
    %
    as $m\to\infty$, uniformly on $\mathcal B_m$.
\end{lemma}

\begin{proof}
    By \eqref{eq:s4-scalar-quadratic} and
    \eqref{eq:s4-feedback-estimates},
    %
    \begin{align*}
        \|E_\zeta(v)\|_{Y_\zeta}
        &\leq \|Q_\zeta(\phi(v,\zeta))\|_{Y_\zeta}
        +\|\mathcal R_\zeta(\phi(v,\zeta))\|_{Y_\zeta}\\
        &\leq C\tau_\zeta^2+C\tau_\zeta^{2+\alpha/2}
        \leq C\tau_\zeta^2,
    \end{align*}
    %
    where the last inequality uses $0<\tau_\zeta<1$.
    Apply the function estimate in \eqref{eq:s4-scalar-inverse}
    and the coefficient estimate \eqref{eq:s4-coefficient-bound}
    to \eqref{eq:s4-coupled-inverse-output}. They give
    %
    \begin{align*}
        \|w_{v,\zeta}\|_{X_\zeta}
        &\leq C\delta^{-4-2\beta}\|E_\zeta(v)\|_{Y_\zeta}
        \leq C\delta^{-4-2\beta}\tau_\zeta^2,\\
        |d_{v,\zeta}|
        &\leq C\delta^{-2-\beta}\|E_\zeta(v)\|_{Y_\zeta}
        \leq C\delta^{-2-\beta}\tau_\zeta^2.
    \end{align*}
    %
    Similarly, from
    $(u_\zeta,c_{H,\zeta})=-\mathscr R_\zeta D_\zeta$ and
    \eqref{eq:s4-H-error},
    %
    \begin{align*}
        |c_{H,\zeta}|
        \leq C\delta^{-2-\beta}\|D_\zeta\|_{Y_\zeta}
        \leq C\delta^{-2-\beta}\tau_\zeta^{1+\alpha/3}.
    \end{align*}
    %
    This proves the four asserted bounds.

    The comparison of norms in \eqref{eq:s4-norm-identification}
    and the scale comparison \eqref{eq:s4-coupled-scale-comparison}
    imply
    %
    \begin{align*}
        \|\mathcal T_m^1(v,\zeta)\|_{X_0}
        &=\|w_{v,\zeta}\circ F_\zeta\|_{X_0}\\
        &\leq C_F\|w_{v,\zeta}\|_{X_\zeta}
        \leq C\delta^{-4-2\beta}\tau_0^2.
    \end{align*}
    %
    Dividing by the radius $\tau_0^{1+\alpha}$ of the function
    ball gives \eqref{eq:s4-coupled-function-smallness}.
    For the second coordinate, dividing the estimates for
    $c_{H,\zeta}$ and $d_{v,\zeta}$ by $\tau_\zeta$ gives
    %
    \begin{align*}
        |\mathcal T_m^2(v,\zeta)|
        &\leq\frac{|c_{H,\zeta}|+|d_{v,\zeta}|}{\tau_\zeta}\\
        &\leq C\delta^{-2-\beta}
        (\tau_\zeta^{\alpha/3}+\tau_\zeta)
        \leq C\delta^{-2-\beta}
        (\tau_0^{\alpha/3}+\tau_0).
    \end{align*}
    %
    Because $0<\alpha<1$, all powers of $\tau_0$ in
    \eqref{eq:s4-coupled-function-smallness} and
    \eqref{eq:s4-coupled-parameter-smallness} are positive.
    Both right-hand sides therefore tend to zero uniformly.
    For sufficiently large $m$, they are at most $1/2$.
    In particular,
    %
    \begin{align*}
        \|\mathcal T_m^1(v,\zeta)\|_{X_0}
        \leq\tau_0^{1+\alpha},
        \qquad
        |\mathcal T_m^2(v,\zeta)|\leq1,
    \end{align*}
    %
    which are precisely the defining inequalities for
    membership in $\mathcal B_m$.

    The same calculation isolates the contribution of the
    displacement normal to $\mathcal N$. Its contribution to
    the function coordinate is at most
    $C\delta^{-4-2\beta}\tau_\zeta^{2+\alpha/2}$.
    Relative to the radius $\tau_0^{1+\alpha}$, this is bounded by
    %
    \begin{align*}
        C\delta^{-4-2\beta}\tau_0^{1-\alpha/2}
        \rightarrow0.
    \end{align*}
    %
    Thus the estimate in the preceding section is strong enough
    to be included in the same scalar fixed-point argument.
\end{proof}

\subsection{Compactness and continuity in a weaker H\"older topology}

The preceding lemma gives the bounds needed for the fixed-point
argument. We must also specify a topology in which its domain is
compact and the map is continuous. The bound defining
$\mathcal B_m$ controls two derivatives with H\"older exponent
$\beta$. We retain that bound and use convergence with a smaller
exponent. All arguments in this subsection are for a fixed $m$;
none requires a compactness estimate uniform as $m\to\infty$.

Fix $\beta'\in(0,\beta)$ and let
%
\begin{align*}
    \mathcal E_m
    &:=C^{2,\beta'}_{\mathrm{inv}}(M_0)\times\mathbb R,\\
    \|(v,\zeta)\|_{\mathcal E_m}
    &:=\|v\|_{C^{2,\beta'}(M_0,g_0)}+|\zeta|.
\end{align*}
%
Here $g_0$ is the induced metric of the initial surface $M_0$,
and the subscript $\mathrm{inv}$ has the same meaning as in
\eqref{eq:s4-XY}: the functions are invariant under $\Gamma_m$.
Invariance is preserved under uniform convergence, so this is
a closed subspace of the ordinary H\"older space. In particular,
$\mathcal E_m$ is a Banach space. We regard $\mathcal B_m$ as
a subset of $\mathcal E_m$, with its original defining bound
$\|v\|_{X_0}\leq\tau_0^{1+\alpha}$.

\begin{lemma}
    \label{lem:s4-coupled-compactness}
    For every fixed sufficiently large $m$, the set
    $\mathcal B_m$ is a nonempty compact convex subset of
    $\mathcal E_m$. The map
    $\mathcal T_m:\mathcal B_m\to\mathcal B_m$ is continuous
    in the topology induced by $\mathcal E_m$.
\end{lemma}

\begin{proof}
    The pair $(0,0)$ belongs to $\mathcal B_m$. Its function
    coordinate ranges over a norm ball in the linear space
    $X_0$, and its real coordinate ranges over $[-1,1]$.
    Both sets are convex, hence so is $\mathcal B_m$.

    We first prove compactness. With $m$ fixed, the surface
    $M_0$ is compact and the continuous weights $\rho_0,f_0$
    are strictly positive. Their minima are therefore
    positive and their maxima finite. Consequently, the
    radii of the local balls in \eqref{eq:s4-weighted-norm}
    are bounded below by a positive number, and all the
    constant metric-rescaling factors and function weights
    are bounded above and below. A finite covering by smaller
    such balls shows that $\|\cdot\|_{X_0}$ is equivalent
    to the ordinary $C^{2,\beta}(M_0,g_0)$ norm. The constants
    in this equivalence may depend on $m$.

    Thus the function coordinates of $\mathcal B_m$ are
    bounded in $C^{2,\beta}(M_0,g_0)$. The compact inclusion
    into $C^{2,\beta'}(M_0,g_0)$, together with compactness
    of $[-1,1]$, shows that every sequence in $\mathcal B_m$
    has a subsequence converging in $\mathcal E_m$. To prove
    compactness of $\mathcal B_m$ itself, we must check that
    the limit still satisfies the original weighted bound.

    Suppose $(v_j,\zeta_j)\in\mathcal B_m$ and
    $(v_j,\zeta_j)\to(v,\zeta)$ in $\mathcal E_m$.
    Fix a centre $x\in M_0$ and the corresponding ball in
    \eqref{eq:s4-weighted-norm}. On this fixed ball, the
    covariant derivatives of $v_j$ of orders zero, one and
    two converge uniformly to those of $v$. For each pair
    of distinct points, the $\beta$-H\"older quotient of the
    second derivative converges as well: the metric, the
    distance and the parallel transport used to compare the
    tensors are all fixed. Taking the supremum over these
    pairs gives lower semicontinuity of the H\"older
    seminorm. The three supremum-norm terms converge.
    Multiplying by $f_0(x)^{-1}$ and then taking the supremum
    over $x$ therefore gives
    %
    \begin{align*}
        \|v\|_{X_0}
        \leq\liminf_{j\to\infty}\|v_j\|_{X_0}
        \leq\tau_0^{1+\alpha}.
    \end{align*}
    %
    In particular, the finite local H\"older bounds show that
    $v\in C^{2,\beta}(M_0)$. Invariance passes to the limit
    pointwise, and $\zeta\in[-1,1]$. Hence the limit belongs
    to $\mathcal B_m$. This proves compactness.

    We next prove continuity of $\mathcal T_m$. Let
    $(v_j,\zeta_j)\to(v,\zeta)$ in $\mathcal E_m$, with all
    these pairs in $\mathcal B_m$. Functions on different
    initial surfaces will be compared on $M_0$ by composition
    with the corresponding $F_\zeta$. We follow the order of
    the operations defining \eqref{eq:s4-coupled-map}: first
    the preliminary correction, then the nonlinear source,
    and finally the scalar solution operator.

    The initial embeddings, their induced metrics and the
    diffeomorphisms $F_\zeta$ depend smoothly on $\zeta$.
    Therefore the functions
    $D_\zeta=H_\zeta+\tau_\zeta\zeta\mathcal J_\zeta$
    satisfy
    %
    \begin{align*}
        D_{\zeta_j}\circ F_{\zeta_j}
        \rightarrow D_\zeta\circ F_\zeta
        \qquad\text{in }C^{0,\beta'}(M_0).
    \end{align*}
    %
    Their stronger source norms are bounded by
    \eqref{eq:s4-H-error}. We may therefore apply the
    sequential continuity assertion in
    Proposition~\ref{prop:s4-scalar-input} to
    $(u_\zeta,c_{H,\zeta})=-\mathscr R_\zeta D_\zeta$.
    It gives
    %
    \begin{align}
        \label{eq:s4-initial-correction-continuity}
        \begin{split}
            u_{\zeta_j}\circ F_{\zeta_j}
            &\rightarrow u_\zeta\circ F_\zeta
            \qquad\text{in }C^{2,\beta'}(M_0),\\
            c_{H,\zeta_j}&\rightarrow c_{H,\zeta}.
        \end{split}
    \end{align}
    %
    Setting $\phi_j:=\phi(v_j,\zeta_j)$ and
    $\phi:=\phi(v,\zeta)$, we obtain
    %
    \begin{align*}
        \phi_j\circ F_{\zeta_j}
        &=u_{\zeta_j}\circ F_{\zeta_j}+v_j\\
        &\rightarrow u_\zeta\circ F_\zeta+v
        =\phi\circ F_\zeta
        \qquad\text{in }C^{2,\beta'}(M_0).
    \end{align*}

    The scalar mean curvature of a graph is, in local
    coordinates, a smooth function of the graph height and
    its derivatives of orders one and two, on the set where
    the induced metric is positive definite. The embeddings
    supplied by Proposition~\ref{prop:s4-scalar-input} have
    this property. With $m$ fixed, their metrics are uniformly
    comparable to the initial metrics, and the latter form
    a smooth family on the compact parameter interval.
    The H\"older product and composition estimates therefore
    apply to the preceding convergence. Using also the
    smooth dependence of $H_\zeta$ and $\mathcal L_\zeta$,
    the definition \eqref{eq:s4-Q-definition} gives
    %
    \begin{align*}
        Q_{\zeta_j}(\phi_j)\circ F_{\zeta_j}
        \rightarrow Q_\zeta(\phi)\circ F_\zeta
        \qquad\text{in }C^{0,\beta'}(M_0).
    \end{align*}
    %
    Proposition~\ref{prop:s4-transverse-elimination} gives
    convergence in the same space for
    $\mathcal R_{\zeta_j}(\phi_j)\circ F_{\zeta_j}$.
    Adding these two convergences proves
    %
    \begin{align}
        \label{eq:s4-coupled-source-continuity}
        E_{\zeta_j}(v_j)\circ F_{\zeta_j}
        \rightarrow E_\zeta(v)\circ F_\zeta
        \qquad\text{in }C^{0,\beta'}(M_0).
    \end{align}
    %
    The bound \eqref{eq:s4-coupled-source-bound} also gives
    a uniform bound for their $Y_{\zeta_j}$ norms, with $m$
    fixed. These are precisely the two hypotheses needed
    to apply the sequential continuity of $\mathscr R_\zeta$
    again, this time to
    $(w_{v,\zeta},d_{v,\zeta})=-\mathscr R_\zeta E_\zeta(v)$.
    We obtain
    %
    \begin{align*}
        w_{v_j,\zeta_j}\circ F_{\zeta_j}
        &\rightarrow w_{v,\zeta}\circ F_\zeta
        \qquad\text{in }C^{2,\beta'}(M_0),\\
        d_{v_j,\zeta_j}&\rightarrow d_{v,\zeta}.
    \end{align*}
    %
    Finally, $\tau_\zeta=\tau_0e^\zeta$ is positive and
    continuous. Together with
    \eqref{eq:s4-initial-correction-continuity}, the last two
    limits prove convergence of both coordinates of
    $\mathcal T_m(v_j,\zeta_j)$ in $\mathcal E_m$.
\end{proof}

\subsection{Existence, regularity, and bounds for the deformations}

The two preceding lemmas provide a fixed point. We now verify
that its defining identities make the entire mean-curvature vector
vanish. We also record stronger estimates for this fixed point
than those required to define $\mathcal B_m$, since the curvature
comparison in the next section will use those stronger estimates.

\begin{proposition}
    \label{prop:s4-coupled-existence}
    There exists $m_0\in\mathbb N$ such that, for every integer
    $m\geq m_0$, there is a pair
    $(v_m,\zeta_m)\in\mathcal B_m$ with the following properties.
    Define
    %
    \begin{align*}
        \tau_m&:=\tau_{\zeta_m},\qquad
        \phi_m:=u_{\zeta_m}+v_m\circ F_{\zeta_m}^{-1},
        \qquad b_m:=b(v_m,\zeta_m).
    \end{align*}
    %
    The displacement within $\mathcal N$ is given by
    %
    \begin{align*}
        G_m:M_{\zeta_m}&\rightarrow\mathcal N,\\
        G_m(x)&:=\exp_{\jmath_{\zeta_m}(x)}^{\mathcal N}
        \bigl(\phi_m(x)\nu_{\zeta_m}(x)\bigr)
        \qquad\forall\,x\in M_{\zeta_m}.
    \end{align*}
    %
    Here $\jmath_{\zeta_m}:M_{\zeta_m}\hookrightarrow\mathcal N$
    is the inclusion. Define the unit normal to $\mathcal N$
    along $G_m$ by
    %
    \begin{align*}
        n_m(x):=N\bigl(\pi(G_m(x))\bigr)
        \qquad\forall\,x\in M_{\zeta_m}.
    \end{align*}
    %
    The map $\widehat X_m:M_{\zeta_m}\rightarrow\mathbb S^4$ describing the subsequent displacement normal to $\mathcal N$, given by 
    %
    \begin{align}
        \label{eq:s4-final-embedding}
        \begin{split}
            \widehat X_m(x):=\exp_{G_m(x)}^{\mathbb S^4}
            \bigl(b_m(x)n_m(x)\bigr)=\cos\bigl(b_m(x)\bigr)G_m(x)
            +\sin\bigl(b_m(x)\bigr)n_m(x)
            \quad\forall\,x\in M_{\zeta_m},
        \end{split}
    \end{align}
    %
    is a smooth embedding. Its image
    $\Sigma_m:=\widehat X_m(M_{\zeta_m})$ is a closed, connected,
    $\Gamma_m$-invariant minimal surface in $\mathbb S^4$, with
    %
    \begin{align}
        \label{eq:s4-final-genus}
        \operatorname{genus}(\Sigma_m)=m^2+1.
    \end{align}
    %
    The following estimates hold with constants independent of $m$:
    %
    \begin{align}
        |\zeta_m|&\leq C\delta^{-2-\beta}\tau_m^{\alpha/3},
        \label{eq:s4-final-balance}\\
        \|v_m\|_{X_0}&\leq C\delta^{-4-2\beta}\tau_m^2,
        \label{eq:s4-final-v}\\
        \|\phi_m\|_{X_{\zeta_m}}
        &\leq C\delta^{-4-2\beta}\tau_m^{1+\alpha/3},
        \label{eq:s4-final-phi}\\
        \|b_m\|_{\mathcal X_{\zeta_m}}&\leq C\tau_m^2.
        \label{eq:s4-final-b}
    \end{align}
    %
    In particular, $\zeta_m\to0$ and
    %
    \begin{align}
        \label{eq:s4-final-tau-asymptotic}
        \frac{\tau_m}{\tau_0}
        =1+O\bigl(\delta^{-2-\beta}\tau_0^{\alpha/3}\bigr)
        \rightarrow1
        \qquad\mbox{as }m\rightarrow\infty.
    \end{align}
    %
    Write $\rho_m:=\rho_{\zeta_m}$ and $f_m:=f_{\zeta_m}$,
    and take covariant derivatives in the following estimates
    with respect to $g_{\zeta_m}$. For every $x\in M_{\zeta_m}$
    and $j\in\{0,1,2\}$,
    %
    \begin{align}
        \label{eq:s4-final-two-jets}
        \begin{split}
            |\nabla^j\phi_m(x)|_{g_{\zeta_m}}
            &\leq C\delta^{-4-2\beta}\tau_m^{1+\alpha/3}
            f_m(x)\rho_m(x)^{-j},\\
            |\nabla^j b_m(x)|_{g_{\zeta_m}}
            &\leq C\tau_m^2\rho_m(x)^{-j}.
        \end{split}
    \end{align}
    %
    The corresponding local scaled H\"older estimates hold
    on the balls and with the rescaled metrics specified in
    \eqref{eq:s4-weighted-norm}.
    All estimates in Sections~\ref{sec:s4-scalar-family}
    and~\ref{sec:s4-transverse} that hold for every
    $(v,\zeta)\in\mathcal B_m$ apply to this pair.
\end{proposition}

\begin{proof}
    By Lemmas~\ref{lem:s4-coupled-self-map} and
    \ref{lem:s4-coupled-compactness}, the Schauder fixed-point
    theorem applies to $\mathcal T_m$ on the nonempty compact
    convex subset $\mathcal B_m$ of the Banach space
    $\mathcal E_m$. Choose a fixed point $(v_m,\zeta_m)$.

    We first verify minimality. In this calculation, write
    $(v,\zeta):=(v_m,\zeta_m)$, $w:=w_{v,\zeta}$,
    $d:=d_{v,\zeta}$, and $\phi:=\phi_m$.
    Equality of the two coordinates in
    $\mathcal T_m(v,\zeta)=(v,\zeta)$ means
    %
    \begin{align}
        \label{eq:s4-fixed-point-identities}
        v\circ F_\zeta^{-1}=w,
        \qquad
        \tau_\zeta\zeta=c_{H,\zeta}+d.
    \end{align}
    %
    In particular, $\phi=u_\zeta+w$. Using in order the
    definition of $Q_\zeta$, the definition of $E_\zeta(v)$,
    \eqref{eq:s4-linear-cancellation}, and
    \eqref{eq:s4-coupled-linear-equation}, we obtain
    %
    \begin{align*}
        H_{\zeta,\phi}+\mathcal R_\zeta(\phi)
        &=H_\zeta+\mathcal L_\zeta\phi
        +Q_\zeta(\phi)+\mathcal R_\zeta(\phi)\\
        &=H_\zeta+\mathcal L_\zeta u_\zeta
        +\mathcal L_\zeta w+E_\zeta(v)\\
        &=(-\tau_\zeta\zeta+c_{H,\zeta})\mathcal J_\zeta
        +\bigl(-E_\zeta(v)+d\mathcal J_\zeta\bigr)
        +E_\zeta(v)\\
        &=(-\tau_\zeta\zeta+c_{H,\zeta}+d)
        \mathcal J_\zeta\\
        &=0.
    \end{align*}
    %
    This is the mean-curvature component tested against
    $\nu_{b_m}$ on the final surface. The component tested
    against $n_{b_m}$ is zero by
    \eqref{eq:s4-transverse-equation}. Those two vectors span
    the normal plane at every point. Since the mean-curvature
    vector is normal to the surface and has zero inner product
    with both, it vanishes identically.

    Proposition~\ref{prop:s4-remaining-scalar-equation} already
    gives a $C^{2,\beta}$ embedding in
    \eqref{eq:s4-final-embedding}. The exact relation with the
    tubular map of Lemma~\ref{lem:s4-tubular} is
    %
    \begin{align*}
        \widehat X_m(x)
        =\Psi\bigl(\mathcal F^{-1}(G_m(x)),b_m(x)\bigr)
        \qquad\forall\,x\in M_{\zeta_m}.
    \end{align*}
    %
    The graph $x\mapsto(\mathcal F^{-1}(G_m(x)),b_m(x))$
    is embedded because its projection onto the first factor
    is the embedding $\mathcal F^{-1}\circ G_m$.
    Section~\ref{sec:s4-transverse} ensures that this graph
    lies in the domain on which $\Psi$ is an embedding.
    Thus its image is closed and connected,
    and is diffeomorphic as a $C^{2,\beta}$ surface to
    $M_{\zeta_m}$.

    We also check the symmetries at the level of these maps.
    Fix $g\in\Gamma_m$ and $x\in M_{\zeta_m}$.
    The isometry $g$ preserves $\mathcal N$ and hence commutes
    with its exponential map. Invariance of $\phi_m$ and
    equivariance of $\nu_{\zeta_m}$ therefore give
    %
    \begin{align*}
        G_m(g(x))=g(G_m(x)).
    \end{align*}
    %
    The transformation rules for $b_m$ and $n$ give
    %
    \begin{align*}
        b_m(g(x))&=\chi(g)b_m(x),\\
        n_m(g(x))&=\chi(g)\,d g_{G_m(x)}(n_m(x)).
    \end{align*}
    %
    Since $\chi(g)^2=1$, their product consequently satisfies
    %
    \begin{align*}
        b_m(g(x))n_m(g(x))
        =d g_{G_m(x)}\bigl(b_m(x)n_m(x)\bigr).
    \end{align*}
    %
    Applying equivariance of the spherical exponential map now
    gives
    %
    \begin{align*}
        \widehat X_m(g(x))=\exp_{g(G_m(x))}^{\mathbb S^4}
          \bigl(d g_{G_m(x)}(b_m(x)n_m(x))\bigr)=g\left(\exp_{G_m(x)}^{\mathbb S^4}
          \bigl(b_m(x)n_m(x)\bigr)\right)=g(\widehat X_m(x)).
    \end{align*}
    %
    This proves $\Gamma_m$-invariance of $\Sigma_m$.
    Its genus is the genus of $M_{\zeta_m}$, namely $m^2+1$
    by Lemma~\ref{lem:s4-initial-geometry}.

    We next prove the estimates. The second identity in
    \eqref{eq:s4-fixed-point-identities} gives
    %
    \begin{align*}
        |\zeta_m|
        \leq\tau_m^{-1}
        \bigl(|c_{H,\zeta_m}|+|d_{v_m,\zeta_m}|\bigr)\leq C\delta^{-2-\beta}
        (\tau_m^{\alpha/3}+\tau_m)
        \leq C\delta^{-2-\beta}\tau_m^{\alpha/3}.
    \end{align*}
    %
    The last inequality uses $0<\tau_m<1$ and $\alpha/3<1$.
    This proves \eqref{eq:s4-final-balance}, and its right-hand
    side tends to zero by \eqref{eq:s4-coupled-scale-comparison}.
    Since $\tau_m/\tau_0=e^{\zeta_m}$ and $|\zeta_m|\leq1$,
    %
    \begin{align*}
        \left|\frac{\tau_m}{\tau_0}-1\right|
        =|e^{\zeta_m}-1|
        \leq e|\zeta_m|
        \leq C\delta^{-2-\beta}\tau_0^{\alpha/3},
    \end{align*}
    %
    which proves \eqref{eq:s4-final-tau-asymptotic}.

    The first identity in \eqref{eq:s4-fixed-point-identities}
    and the norm comparison give
    %
    \begin{align*}
        \|v_m\|_{X_0}
        =\|w_{v_m,\zeta_m}\circ F_{\zeta_m}\|_{X_0}
        \leq C\delta^{-4-2\beta}\tau_m^2.
    \end{align*}
    %
    For the preliminary correction, the scalar inverse estimate
    and \eqref{eq:s4-H-error} give
    %
    \begin{align*}
        \|u_{\zeta_m}\|_{X_{\zeta_m}}
        \leq C\delta^{-4-2\beta}\|D_{\zeta_m}\|_{Y_{\zeta_m}}
        \leq C\delta^{-4-2\beta}\tau_m^{1+\alpha/3}.
    \end{align*}
    %
    Adding the estimate for $w_{v_m,\zeta_m}$ proves
    \eqref{eq:s4-final-phi}, because
    $\phi_m=u_{\zeta_m}+w_{v_m,\zeta_m}$ and
    $2>1+\alpha/3$.
    The estimate \eqref{eq:s4-final-b} is the bound in
    Proposition~\ref{prop:s4-transverse-elimination}.

    To obtain \eqref{eq:s4-final-two-jets}, evaluate the local
    norms defining $X_{\zeta_m}$ and
    $\mathcal X_{\zeta_m}$ at a center $x$.
    Multiplication of the metric by $\rho_m(x)^{-2}$ multiplies
    the norm of the $j$-th covariant derivative of a function
    by $\rho_m(x)^j$. The function weight is $f_m(x)$ in
    $X_{\zeta_m}$ and is $1$ in $\mathcal X_{\zeta_m}$.
    Thus the norm bounds just proved give exactly the two
    pointwise inequalities in \eqref{eq:s4-final-two-jets}.
    The H\"older terms in the same local norms give the stated
    local H\"older estimates. Since the fixed point lies in
    $\mathcal B_m$, it also retains all bounds proved on that
    set, including
    $\|\phi_m\|_{X_{\zeta_m}}\leq\tau_m^{1+\alpha/4}$.

    It remains to prove smoothness. We first prove that the image
    $\Sigma_m$ is smooth, and then that its given parametrization
    is smooth. Throughout this part of the argument $m$ is fixed;
    no uniform estimate for derivatives of order greater than two
    is claimed.

    Near each point of the $C^{2,\beta}$ embedded surface
    $\Sigma_m$, choose smooth ambient coordinates
    $(y^1,y^2,y^3,y^4)$ such that the surface is represented by
    %
    \begin{align*}
        Z(y^1,y^2)
        =\bigl(y^1,y^2,V^1(y^1,y^2),V^2(y^1,y^2)\bigr),
    \end{align*}
    %
    with $V=(V^1,V^2)\in C^{2,\beta}$ on an open subset of
    $\mathbb R^2$. Write $\overline g_{BC}$ and
    $\overline\Gamma^A_{BC}$ for the metric components and
    Christoffel symbols of $g_{\mathbb S^4}$ in these coordinates,
    with ambient indices in $\{1,2,3,4\}$.
    For $i,j\in\{1,2\}$, set
    %
    \begin{align*}
        T_i^B:=\partial_i Z^B,
        \qquad
        g_{ij}:=\overline g_{BC}(Z)T_i^BT_j^C,
        \qquad
        (a^{ij}):=(g_{ij})^{-1}.
    \end{align*}
    %
    Repeated indices in these local formulas are summed over
    their indicated ranges. Let $\Gamma^k_{ij}(g)$ be the
    Christoffel symbols of the induced metric.
    Vanishing of the mean-curvature vector means
    %
    \begin{align*}
        a^{ij}\bigl(
        \partial_{ij}Z^A
        +\overline\Gamma^A_{BC}(Z)T_i^BT_j^C
        -\Gamma^k_{ij}(g)\partial_k Z^A
        \bigr)=0
        \qquad\forall\,A\in\{1,2,3,4\}.
    \end{align*}
    %
    Since $Z^k=y^k$ for $k\in\{1,2\}$, the first two
    equations give
    %
    \begin{align*}
        a^{ij}\Gamma^k_{ij}(g)
        =a^{ij}\overline\Gamma^k_{BC}(Z)T_i^BT_j^C
        \qquad\forall\,k\in\{1,2\}.
    \end{align*}
    %
    Substituting these identities into the remaining two
    equations yields, for $\mu\in\{1,2\}$,
    %
    \begin{align*}
        a^{ij}\partial_{ij}V^\mu
        =a^{ij}\bigl(
        \partial_kV^\mu\,\overline\Gamma^k_{BC}(Z)
        -\overline\Gamma^{\mu+2}_{BC}(Z)
        \bigr)T_i^BT_j^C.
    \end{align*}
    %
    Both $a^{ij}$ and the right-hand side are smooth functions
    of $(y,V,DV)$. The matrix $(a^{ij})$ is positive definite
    and uniformly so on a smaller relatively compact coordinate
    neighborhood. Since $V\in C^{2,\beta}$, its coefficients
    and right-hand sides belong to $C^{1,\beta}$.
    Interior Schauder regularity applied to each component
    gives $V\in C^{3,\beta}$ on smaller neighborhoods.
    Repeating this argument raises the regularity by one
    derivative each time. Thus $\Sigma_m$ is smooth.

    Smoothness of the image alone does not yet assert smoothness
    of the original parametrization. To recover it, use the
    smooth inverse of the fixed tubular map
    \eqref{eq:s4-tubular-map}. Define its two relevant components
    by
    %
    \begin{align*}
        P:\operatorname{im}\Psi&\rightarrow\mathcal N,
        &\hspace{-1cm}P(\Psi(s,t,z,b))&:=\mathcal F(s,t,z),\\
        t_\perp:\operatorname{im}\Psi&\rightarrow(-b_0,b_0),
        &\hspace{-1cm}t_\perp(\Psi(s,t,z,b))&:=b,
    \end{align*}
    %
    for every $(s,t,z,b)$ in the domain of $\Psi$.
    By construction,
    %
    \begin{align*}
        P\circ\widehat X_m=G_m,
        \qquad
        t_\perp\circ\widehat X_m=b_m.
    \end{align*}
    %
    The smooth map $P|_{\Sigma_m}$ is injective because $G_m$
    is injective. It is also an immersion: differentiating
    $P\circ\widehat X_m=G_m$ shows that its differential has
    rank two, since $dG_m$ has rank two and $d\widehat X_m$
    maps onto the tangent space of $\Sigma_m$.
    An injective immersion of a compact smooth manifold into
    a Hausdorff manifold is an embedding. Consequently,
    %
    \begin{align*}
        M_{\zeta_m,\phi_m}=P(\Sigma_m)=G_m(M_{\zeta_m})
    \end{align*}
    %
    is a smooth embedded surface in $\mathcal N$.

    To recover the original parametrization of
    $M_{\zeta_m,\phi_m}$, we use
    normal exponential coordinates based on the smooth initial
    surface $M_{\zeta_m}$. Define
    %
    \begin{align*}
        \mathfrak E_m(x,t)
        :=\exp_{\jmath_{\zeta_m}(x)}^{\mathcal N}
          \bigl(t\nu_{\zeta_m}(x)\bigr)
    \end{align*}
    %
    for $x\in M_{\zeta_m}$ and for every $t\in\mathbb R$
    for which the indicated geodesic is defined. This is a smooth
    map on its open domain. We will prove that its differential
    is invertible at $(x,\phi_m(x))$ for every $x$.
    Invertibility at $(x,0)$ alone would not justify this claim,
    because the surfaces and their normal displacements depend
    on $m$.

    Fix $x\in M_{\zeta_m}$ and put $r:=\rho_m(x)$.
    For large $m$, the bound $|z(x)|\leq C\varepsilon_{\zeta_m}$
    places $x$ in the smaller band $|z|<z_0/2$.
    Since $|\nabla^{\mathcal N}z|=1$ and $r\leq\pi/m$,
    the unit-speed normal geodesic
    %
    \begin{align*}
        c_x(t):=\mathfrak E_m(x,t)
    \end{align*}
    %
    is defined for $|t|\leq c_*r$, with any fixed
    $c_*\in(0,1)$ once $m$ is sufficiently large.
    Indeed, it cannot reach the boundary of the band in such
    a time, and the metric is smooth on a neighborhood of
    the closure of every smaller band. The curvature tensor
    of $g_{\mathcal N}$ is uniformly bounded there by the Gauss equation and
    \eqref{eq:s4-S}.

    In the following calculation, vector norms and inner
    products in $T\mathcal N$ are those of $g_{\mathcal N}$.

    Let $\mathsf P_t:T_{\jmath_{\zeta_m}(x)}\mathcal N
    \to T_{c_x(t)}\mathcal N$ denote parallel transport along
    $c_x$. For $\xi\in T_xM_{\zeta_m}$, differentiating the
    family of normal geodesics in its initial point gives the
    Jacobi field
    %
    \begin{align*}
        J_\xi(t):=d\bigl(\mathfrak E_m(\cdot,t)\bigr)_x(\xi).
    \end{align*}
    %
    Set $\xi_0:=d\jmath_{\zeta_m,x}(\xi)$ and
    $T_x:=S_{\jmath_{\zeta_m}}^{\nu_{\zeta_m}}(x)$.
    The initial conditions and the Jacobi equation are
    %
    \begin{align*}
        J_\xi(0)=\xi_0,\qquad
        \nabla_tJ_\xi(0)=-d\jmath_{\zeta_m,x}(T_x(\xi)),\qquad
        \nabla_t^2J_\xi=-R^{\mathcal N}(J_\xi,\dot c_x)\dot c_x.
    \end{align*}
    %
    Here $R^{\mathcal N}$ is the Riemann curvature tensor
    of $g_{\mathcal N}$, and $\nabla_t$ denotes covariant
    differentiation along $c_x$. The formula for the initial
    derivative is the Weingarten equation for
    $M_{\zeta_m}\subset\mathcal N$.
    By \eqref{eq:s4-initial-geometry-bounds} and
    $\tau_m\leq r$, one has $|T_x|\leq C/r$.

    Parallel transport reduces this equation to an ordinary
    differential equation in the fixed vector space
    $T_{\jmath_{\zeta_m}(x)}\mathcal N$. More precisely, set
    %
    \begin{align*}
        K_\xi(t)&:=\mathsf P_t^{-1}J_\xi(t),\\
        \mathcal K_t(w)
        &:=\mathsf P_t^{-1}
          \bigl(R^{\mathcal N}(\mathsf P_t w,\dot c_x(t))
          \dot c_x(t)\bigr)
          \qquad\forall\,w\in T_{\jmath_{\zeta_m}(x)}\mathcal N.
    \end{align*}
    %
    The uniform curvature bound gives $|\mathcal K_t|\leq C$.
    Integrating the Jacobi equation twice gives
    %
    \begin{align*}
        K_\xi(t)
        =\xi_0-t\,d\jmath_{\zeta_m,x}(T_x(\xi))
          -\int_0^t(t-s)\mathcal K_s(K_\xi(s))\,ds.
    \end{align*}
    %
    Since $\jmath_{\zeta_m}$ is isometric, it follows that
    %
    \begin{align*}
        \sup_{|t|\leq c_*r}|K_\xi(t)|
        &\leq (1+Cc_*)|\xi|_{g_{\zeta_m}}
          +Cc_*^2r^2
             \sup_{|t|\leq c_*r}|K_\xi(t)|.
    \end{align*}
    %
    Choose $c_*$ small enough to absorb the last term; this
    choice is independent of $m$ and $x$, since $r\leq1$
    for large $m$. Substituting the resulting bound back into
    the integral identity yields
    %
    \begin{align*}
        |K_\xi(t)-\xi_0|
        \leq C\left(\frac{|t|}{r}+t^2\right)
             |\xi|_{g_{\zeta_m}}
        \qquad\forall\,|t|\leq c_*r.
    \end{align*}
    %
    Decreasing $c_*$ once more makes the right-hand side at
    most $|\xi|_{g_{\zeta_m}}/2$. Thus
    %
    \begin{align*}
        |J_\xi(t)|_{g_{\mathcal N}}
        \geq\tfrac12|\xi|_{g_{\zeta_m}}
        \qquad\forall\,|t|\leq c_*r.
    \end{align*}
    %
    The function $t\mapsto\langle J_\xi(t),\dot c_x(t)\rangle$
    has zero second derivative by the Jacobi equation and
    the symmetries of $R^{\mathcal N}$. Its value and first
    derivative at zero vanish by the initial conditions.
    Hence $J_\xi(t)$ is orthogonal to the unit vector
    $\dot c_x(t)=\partial_t\mathfrak E_m(x,t)$.
    The lower bound just obtained therefore shows that
    %
    \begin{align*}
        d\mathfrak E_{m,(x,t)}(\xi,a)
        =J_\xi(t)+a\dot c_x(t)
        \qquad\forall\,(\xi,a)\in T_xM_{\zeta_m}\times\mathbb R
    \end{align*}
    %
    is an isomorphism for $|t|\leq c_*r$.
    Estimate~\eqref{eq:s4-dimensionless-phi} implies
    $\sup_x|\phi_m(x)|/\rho_m(x)\to0$, so this applies
    at $t=\phi_m(x)$ for every $x$ once $m$ is large.
    The smooth inverse function theorem consequently makes
    $\mathfrak E_m$ a smooth coordinate map on a neighborhood
    of each $(x,\phi_m(x))$.

    After shrinking these coordinate neighborhoods, the inverse
    image of the already smooth surface $M_{\zeta_m,\phi_m}$ under
    $\mathfrak E_m$ is the graph of the $C^{2,\beta}$ function
    $\phi_m$. This graph is therefore a smooth submanifold
    of $M_{\zeta_m}\times\mathbb R$.
    Projection of this graph to the $x$ coordinate has
    invertible differential: in the known parametrization
    $x\mapsto(x,\phi_m(x))$, that projection is the identity.
    The smooth inverse function theorem on this graph therefore
    gives a smooth inverse to this projection, whose last
    coordinate is $\phi_m$. Thus $\phi_m$ is smooth, and so
    is $G_m$. Finally,
    %
    \begin{align*}
        \widehat X_m=(P|_{\Sigma_m})^{-1}\circ G_m,
        \qquad
        b_m=t_\perp\circ\widehat X_m
    \end{align*}
    %
    are smooth. This proves the asserted regularity of the
    parametrization and completes the proof.
\end{proof}
%
\section{Strictly negative Gaussian curvature and the area limit}
\label{sec:s4-negative-curvature}

Let $\Sigma_m=\widehat X_m(M_{\zeta_m})$ be the smooth embedded
minimal surface obtained in
Proposition~\ref{prop:s4-coupled-existence}. We prove that its
Gaussian curvature is strictly negative at every point when $m$
is sufficiently large, and that its area converges to $4\pi^2$.
We use different coordinates in two parts of the surface.
Away from the circles of smallest neck radius, we first express
$G_m(M_{\zeta_m})\subset\mathcal N$ as local height graphs
over $\mathbb T$. We compare their heights with the Green
function used in Section~\ref{sec:s4-scalar-family}, compute
their Gaussian curvatures, and estimate the change caused by
the displacement $b_m$ normal to $\mathcal N$. Near each
waist, the height graph coordinates are
not valid. There we rescale lengths by the neck radius and
compare the final surface with a fixed portion of the
Euclidean catenoid. After proving negativity in both regions,
we estimate the areas of the initial surfaces and compare
their metrics with those of the final surfaces.

The first comparison must be quantitative. On the portions of
the sheets furthest from the lattice, the negative curvature
that we obtain tends to zero as $m\to\infty$. Consequently,
an error in curvature that merely tends to zero would not
suffice. We must bound that error by a small multiple of the
positive function $\mathcal W$ defined below, and prove a
negative upper bound of the form $-c\mathcal W$, with $c>0$
independent of $m$.

For the rest of the section, abbreviate
%
\begin{align*}
    \zeta:=\zeta_m,\qquad \tau:=\tau_m,\qquad
    \varepsilon:=\frac{m^2\tau}{2\pi},\qquad
    r_c:=\tau^\alpha,\qquad \psi:=\tau\Phi_m.
\end{align*}
%
Let $h:=h_\zeta$ be the upper initial height from
\eqref{eq:s4-glued-height}, extended to the domain
$\{q\in\mathbb T:d_{\mathbb T}(q,L_m)>\tau\}$ by
$h(q):=C_\tau(r_p(q))$ wherever
$\tau<r_p(q)\leq9\tau$, for each $p\in L_m$.
Here $r_p(q)=d_{\mathbb T}(q,p)$ and
$C_\tau(r)=\tau\operatorname{arcosh}(r/\tau)$ for $r>\tau$.
This extension agrees with the existing height on its overlap
with the exact catenoidal region. It describes the upper
half of the initial surface arbitrarily close to a waist,
but not the waist itself.
For every $q\in\mathbb T\smallsetminus L_m$, put
%
\begin{align}
    \label{eq:s4-curvature-radius-weight}
    \rho(q):=d_{\mathbb T}(q,L_m),\qquad
    \mathcal W(q):=\varepsilon^2+\tau^2\rho(q)^{-4}.
\end{align}
%
The two sides of each local graph comparison must be evaluated
at the same point of $\mathbb T$. The deformation $G_m$ can
change that point. Accordingly, for every $x\in M_\zeta$,
we use the previously defined initial radius $\rho_\zeta(x)$
and denote the radius after deformation by $\widehat\rho(x)$:
%
\begin{align*}
    \rho_\zeta(x)=\rho(\pi(x)),\qquad
    \widehat\rho(x):=\rho(\pi(G_m(x))).
\end{align*}
%
By \eqref{eq:s4-deformed-scale},
$\widehat\rho(x)=(1+o_m(1))\rho_\zeta(x)$ uniformly in $x$.
Here $o_m(1)$ denotes a quantity tending to zero as
$m\to\infty$, uniformly over the points under consideration.
Although the distance function $\rho$ need not be smooth, we
use it only as a pointwise scale and never differentiate it.
Each square lattice cell has side length $\ell/m$, so its
points have distance at most $\ell/(\sqrt2m)=\pi/m$ from
the lattice. Thus
%
\begin{align*}
    \frac{\varepsilon^2}{\tau^2\rho(q)^{-4}}
    =\frac{m^4\rho(q)^4}{4\pi^2}
    \leq\frac{\pi^2}{4}
    \qquad\forall\,q\in\mathbb T\smallsetminus L_m,
\end{align*}
%
and hence
%
\begin{align}
    \label{eq:s4-curvature-weight-equivalence}
    \tau^2\rho^{-4}\leq\mathcal W
    \leq\left(1+\frac{\pi^2}{4}\right)\tau^2\rho^{-4}
    \qquad\text{on }\mathbb T\smallsetminus L_m.
\end{align}
%
We will first choose a large constant $R$, independently of
$m$, and then choose $m$ sufficiently large in terms of $R$.
There is a fixed lower threshold $R_0>10$ such that all constants
in the graph estimates below are independent of $R\geq R_0$.
Constants in the comparison on a fixed rescaled neck are
allowed to depend on $R$.
The two regions will be specified on the initial surface by
$\rho_\zeta\geq2R\tau$ and $\rho_\zeta\leq2R\tau$.
In the first region, the comparison between $\widehat\rho$
and $\rho_\zeta$ gives $\widehat\rho\geq R\tau$ for large
$m$. In the second, the rescaled neck radius is at most $2R$.
These two conditions cover the initial surface, and hence
their images cover $\Sigma_m$.

\subsection{Height coordinates and comparison with the Green function}

The estimates in Proposition~\ref{prop:s4-coupled-existence}
are stated on the initial surface $M_\zeta$. The curvature
calculation below instead uses the height of
$G_m(M_\zeta)$ over the flat torus. We must therefore
estimate the change of the torus coordinate caused by the
normal displacement $\phi_m$, invert that change, and express
both $\phi_m$ and $b_m$ in the resulting coordinates.
We then compare the new height, together with its first two
derivatives, with $\psi=\tau\Phi_m$. These are precisely the
derivatives that enter the second fundamental form.

For every $r>0$, define
%
\begin{align*}
    f_\tau(r):=\max\{r^{3/2},
    \tau^{(1-\alpha)/2}r^{1/2}\},\qquad
    \mathfrak a_m:=\delta^{-4-2\beta}\tau^{1+\alpha/3},
    \qquad \delta:=(100m)^{-1}.
\end{align*}
%
Thus $f_\zeta(x)=f_\tau(\rho_\zeta(x))$ for every
$x\in M_\zeta$. By \eqref{eq:s4-scales}, there are
constants $c,C>0$ such that
$cm^2\leq|\log\tau|\leq Cm^2$ for all sufficiently large
$m$. In particular, $m^p\tau^q\to0$ for every fixed
$p\geq0$ and $q>0$.
Both terms in the maximum defining $f_\tau$ are increasing
functions of $r$. Moreover,
%
\begin{align*}
    \frac{\tau^{(1-\alpha)/2}(\pi/m)^{1/2}}
         {(\pi/m)^{3/2}}
    =\frac{m}{\pi}\tau^{(1-\alpha)/2}\rightarrow0.
\end{align*}
%
Consequently, for all sufficiently large $m$, the supremum
of $f_\tau$ on $[\tau,\pi/m]$ equals $(\pi/m)^{3/2}$.
Set
%
\begin{align}
    \label{eq:s4-curvature-scalar-smallness}
    \sigma_m:=\frac{\mathfrak a_m}{\tau}
    \sup_{\tau\leq r\leq\pi/m}f_\tau(r)
    \leq C\delta^{-4-2\beta}\tau^{\alpha/3}m^{-3/2}
    \rightarrow0.
\end{align}
%
In particular, the scalar displacement has size at most
$C\sigma_m\tau$, and its first two derivatives have the
corresponding bounds after multiplication by the appropriate
powers of $\rho_\zeta$. This will make the change of torus
coordinate small relative to $\rho_\zeta$.

\begin{lemma}[Height coordinates away from the neck waists]
    \label{lem:s4-outer-coordinates}
    Fix $R\geq R_0$. For all sufficiently large $m$ and every
    $x\in M_\zeta$ with $\rho_\zeta(x)\geq2R\tau$, the surface
    $G_m(M_\zeta)$ is locally a height graph over $\mathbb T$
    near $G_m(x)$.
    After applying the reflection $Z$ if $x$ lies on the
    lower sheet, there are an open neighborhood
    $U\subset\{q'\in\mathbb T:\rho(q')>R\tau\}$ of
    $q:=\pi(G_m(x))$ and smooth functions $z,b,\eta$ on $U$
    such that the corresponding portions of
    $G_m(M_\zeta)$ and $\Sigma_m$ are parametrized by
    %
    \begin{align}
        \label{eq:s4-outer-graph-form}
        \begin{split}
            X(q')&:=\mathcal F(q',z(q')),\\
            X_b(q')&:=\cos\bigl(b(q')\bigr)X(q')
                        +\sin\bigl(b(q')\bigr)N(q'),\\
            z(q')&=h(q')+\eta(q')
            \qquad\forall\,q'\in U.
        \end{split}
    \end{align}
    %
    For every $q'\in U$ and $j\in\{0,1,2\}$,
    %
    \begin{align}
        \label{eq:s4-outer-coordinate-estimates}
        \begin{split}
            |D^j\eta(q')|
            &\leq C\mathfrak a_m f_\tau(\rho(q'))\rho(q')^{-j},\\
            |D^jb(q')|&\leq C\tau^2\rho(q')^{-j}.
        \end{split}
    \end{align}
    %
    Here $D$ is the flat covariant derivative of
    $g_{\mathbb T}$, and the derivative norms are induced
    by $g_{\mathbb T}$. On the same neighborhood,
    %
    \begin{align}
        \label{eq:s4-outer-height-bounds}
        \begin{split}
            |z(q')|&\leq C\varepsilon,\\
            |Dz(q')|&\leq C\tau/\rho(q'),\\
            |D^2z(q')|&\leq C\tau/\rho(q')^2
            \qquad\forall\,q'\in U.
        \end{split}
    \end{align}
    %
    The constants are independent of $R\geq R_0$, $m$, and $x$.
\end{lemma}

\begin{proof}
    We first record the derivatives of the initial upper
    height. Increasing the fixed threshold $R_0$ if necessary,
    we have
    %
    \begin{align}
        \label{eq:s4-initial-flat-derivatives}
        |D^kh(p)|\leq C_k\tau\rho(p)^{-k}
        \qquad\text{whenever }\rho(p)\geq R_0\tau/2,
    \end{align}
    %
    for every $p\in\mathbb T$ in the stated region and every
    fixed integer $k\geq1$. To verify this, consider
    the three parts of the definition of $h$.
    On a catenoidal part, $h=C_\tau(r)$, where
    $r$ is distance to the corresponding lattice point and
    %
    \begin{align*}
        C_\tau(r)&=\tau\operatorname{arcosh}(r/\tau),\\
        C_\tau'(r)&=\frac{\tau}{\sqrt{r^2-\tau^2}},\qquad
        C_\tau''(r)=-\frac{\tau r}{(r^2-\tau^2)^{3/2}}
        \qquad\forall\,r>\tau.
    \end{align*}
    %
    Differentiating these expressions, or the logarithmic
    formula for $C_\tau$, gives
    $|D^kC_\tau|\leq C_k\tau r^{-k}$ whenever
    $r\geq R_0\tau/2$.
    On the exterior part,
    $h=\widetilde\psi_\zeta
       =\tau\Phi_m-\tau\zeta\widehat v_m$.
    The estimate for $\Phi_m$ is
    \eqref{eq:s4-physical-derivatives}. The cutoff derivatives
    in $\widehat v_m$ have size at most $C_k\delta^{-k}$
    and are supported where $\rho$ is comparable to $\delta$;
    where the cutoff is constant, the derivatives of
    $J_0(\sqrt2r)$ satisfy the same required bounds.
    Finally, on an interpolation annulus
    $2r_c\leq r\leq3r_c$, the estimates in
    Section~\ref{sec:s4-scalar-family} give
    %
    \begin{align*}
        |D^k(C_\tau-\widetilde\psi_\zeta)|
        &\leq C_k\bigl(\tau r_c^2|\log\tau|
                         +\tau^3r_c^{-2}\bigr)r_c^{-k}=C_k\tau r_c^{-k}
              \bigl(r_c^2|\log\tau|+\tau^2r_c^{-2}\bigr)
    \end{align*}
    %
    for every fixed $k\geq0$.
    The last factor tends to zero because $r_c=\tau^\alpha$
    and $0<\alpha<1$. Each derivative of an interpolation
    cutoff contributes at most $C_jr_c^{-j}$. The product
    rule therefore proves
    \eqref{eq:s4-initial-flat-derivatives} on these annuli too.

    Suppose now that $x$ lies on the upper initial sheet.
    Put $p_0:=\pi(x)$ and $r:=\rho(p_0)$, so that
    $r\geq2R\tau$. The map  $X_+:\{p\in\mathbb T:\rho(p)>\tau\}\rightarrow M_\zeta$ given by
    %
    \begin{align*}
        X_+(p)&:=\mathcal F(p,h(p)) \qquad\forall\,p\in\mathbb T\text{ with }\rho(p)>\tau
    \end{align*}
    %
    parametrizes this sheet away from the waists.
    When an ambient point is needed below, we view
    $X_+(p)$ in $\mathcal N$ through the inclusion
    $\jmath_\zeta$.
    Choose a fixed small $c_1>0$ and a flat coordinate disk
    $B(p_0,4c_1r)$ that does not meet $L_m$.
    On this disk, $\rho(p)$ is comparable to $r$ because
    $\rho$ is $1$-Lipschitz. The powers in the definition
    of $f_\tau$ then show that $f_\tau(\rho(p))$ is
    comparable to $f_\tau(r)$ as well.
    Write
    %
    \begin{align*}
        A:=\mathfrak a_m f_\tau(r),\qquad
        \widetilde\phi:=\phi_m\circ X_+,\qquad
        \widetilde b:=b_m\circ X_+.
    \end{align*}
    %
    In the flat coordinates, the induced metric of $X_+$ is
    %
    \begin{align*}
        (X_+^*g_\zeta)_{ij}
        =\cos^2h\,\delta_{ij}+\partial_i h\,\partial_j h
        \qquad\forall\,i,j\in\{1,2\}.
    \end{align*}
    %
    It is uniformly equivalent to the flat metric, since
    $|h|\to0$ and $|Dh|\leq C\tau/r$.
    Its first derivatives have size at most $C/r$, so its
    Christoffel symbols have the same bound.
    Applying \eqref{eq:s4-final-two-jets} gives
    %
    \begin{align*}
        |D^j\widetilde\phi|\leq CAr^{-j},\qquad
        |D^j\widetilde b|\leq C\tau^2r^{-j}
        \qquad\forall\,j\in\{0,1,2\}.
    \end{align*}
    %
    For $j=2$, this uses the coordinate identity
    $\partial_{ij}v=(\nabla^2v)_{ij}
       +\Gamma_{ij}^k\partial_kv$ for a scalar function
    $v$, applied with the metric $X_+^*g_\zeta$.

    We next estimate the change of coordinates caused by
    $\widetilde\phi$. For $p$ in a smaller concentric disk
    and $|t|<c_2r$, where $c_2>0$ is fixed and small, define
    %
    \begin{align*}
        \mathcal E(p,t):=
        \mathcal F^{-1}\!\left(
        \exp_{X_+(p)}^{\mathcal N}
             \bigl(t\nu_\zeta(X_+(p))\bigr)\right).
    \end{align*}
    %
    The values of $\mathcal E$ are expressed in the chosen
    flat coordinates and the coordinate $z$.
    These geodesics remain in the coordinate neighborhood
    for this range of $t$: their speed is one, the metric
    is uniformly equivalent to the Euclidean coordinate
    metric, and their initial points lie in a smaller disk.
    Their heights also remain in $(-z_0,z_0)$ because
    $|h|\to0$ and $r\leq\pi/m$.

    The required dependence on $p$ and $t$ can be read
    directly from the geodesic equations. If
    $\mathcal E(p,t)=(P_p(t),Z_p(t))$, then
    %
    \begin{align*}
        \ddot P_p
            &=2\tan Z_p\,\dot Z_p\,\dot P_p,\\*
        \ddot Z_p
            &=-\sin Z_p\cos Z_p\,|\dot P_p|^2.
    \end{align*}
    %
    Here dots denote derivatives in $t$, and the norm in
    the second equation is the flat norm on the base.
    The initial position is $(p,h(p))$. Writing
    $c(p):=\cos h(p)$, the formula for the upward unit
    normal of a height graph gives the initial velocity
    %
    \begin{align*}
        \bigl(\dot P_p(0),\dot Z_p(0)\bigr)
        =\left(
          -\frac{Dh(p)}{c(p)\sqrt{c(p)^2+|Dh(p)|^2}},
          \frac{c(p)}{\sqrt{c(p)^2+|Dh(p)|^2}}
          \right).
    \end{align*}
    %
    After the changes of variables $p=p_0+ry$ and $t=ru$,
    and division of the coordinate displacement by $r$,
    the equations have uniformly bounded smooth coefficients
    on fixed coordinate neighborhoods. The initial position
    and velocity have bounded derivatives in $y$ through
    order two by \eqref{eq:s4-initial-flat-derivatives}
    with $k\leq3$.
    Differentiating the first-order system for position and
    velocity in $y$ once and twice gives linear ordinary
    differential equations for the corresponding derivatives;
    the inhomogeneous terms in the second equation are
    products of bounded first derivatives. Integration on
    the fixed interval $|u|<c_2$ gives uniform bounds for
    these derivatives. Derivatives in $u$ are then bounded
    by the system itself.

    More explicitly, set
    %
    \begin{align*}
        \mathcal V(p,t):=\int_0^1\partial_t\mathcal E(p,\theta t)
                    \,d\theta.
    \end{align*}
    %
    The preceding bounds, after returning to the original
    coordinates, imply
    %
    \begin{align*}
        |D_p^j\partial_t^k\mathcal V(p,t)|\leq C_{j,k}r^{-j-k}
        \qquad\forall\,j,k\geq0\text{ with }j+k\leq2,
        \quad |t|<c_2r.
    \end{align*}
    %
    The fundamental theorem of calculus gives
    $\mathcal E(p,t)=(p,h(p))+t\mathcal V(p,t)$.
    Consequently, writing
    %
    \begin{align*}
        G_m(X_+(p))&=\mathcal F(Q(p),H(p)),\\
        Q(p)&:=p+e(p),\qquad H(p):=h(p)+a(p),
    \end{align*}
    %
    we have the vector identity
    %
    \begin{align*}
        (e(p),a(p))
        =\widetilde\phi(p)\mathcal V(p,\widetilde\phi(p)).
    \end{align*}
    %
    The product and chain rules now give
    %
    \begin{align*}
        |e|+|a|&\leq CA,\\
        |De|+|Da|&\leq C(A/r+A^2/r^2),\\
        |D^2e|+|D^2a|
            &\leq C(A/r^2+A^2/r^3+A^3/r^4).
    \end{align*}
    %
    For example, the first derivative of
    $\widetilde\phi\mathcal V(p,\widetilde\phi)$ is
    %
    \begin{align*}
        D\widetilde\phi\,\mathcal V
        +\widetilde\phi\bigl(D_p\mathcal V
                   +\partial_t\mathcal V\,D\widetilde\phi\bigr),
    \end{align*}
    %
    with every occurrence of $\mathcal V$ and its derivatives
    evaluated at $(p,\widetilde\phi(p))$.
    Differentiating this expression once more gives the
    stated second derivative bound; no third derivative
    of $\widetilde\phi$ is required.
    Since
    %
    \begin{align*}
        A\leq\sigma_m\tau,\qquad
        \frac{A}{r}\leq\sigma_m\rightarrow0,
    \end{align*}
    %
    the substitution $t=\widetilde\phi(p)$ is allowed for
    large $m$, and the last three estimates imply
    %
    \begin{align}
        \label{eq:s4-outer-projection-estimates}
        |D^je|+|D^ja|
        \leq C\mathfrak a_m f_\tau(r)r^{-j}
        \qquad\forall\,j\in\{0,1,2\}.
    \end{align}

    In particular, $DQ=I+De$ is invertible for large $m$.
    The inverse function theorem provides a local inverse
    $Q^{-1}$ near $q=Q(p_0)$. Shrinking its domain if
    necessary, differentiation of $Q\circ Q^{-1}
    =\operatorname{id}$ gives
    %
    \begin{align*}
        |DQ^{-1}|&\leq2,\\
        D^2Q^{-1}
        &=-(DQ\circ Q^{-1})^{-1}
          (D^2Q\circ Q^{-1})[DQ^{-1},DQ^{-1}],\\
        |D^2Q^{-1}|&\leq CA/r^2.
    \end{align*}
    %
    Moreover,
    %
    \begin{align*}
        |\rho(Q(p))-\rho(p)|
        \leq|Q(p)-p|=|e(p)|\leq CA.
    \end{align*}
    %
    Thus $\rho(Q(p))/\rho(p)=1+o_m(1)$ uniformly on the
    smaller disk. At its center,
    $\rho(q)\geq\tfrac34r\geq\tfrac32R\tau$ for large
    $m$. We may therefore choose the domain $U$ of the
    local inverse inside $\{\rho>R\tau\}$.

    On $U$, define
    %
    \begin{align*}
        z:=H\circ Q^{-1},\qquad
        b:=\widetilde b\circ Q^{-1},\qquad
        \eta:=z-h.
    \end{align*}
    %
    The definition of $\widehat X_m$ gives exactly
    \eqref{eq:s4-outer-graph-form} in these coordinates.
    To estimate $\eta$, put $k:=h-h\circ Q$, so that
    $\eta\circ Q=a+k$. The mean-value formula and
    \eqref{eq:s4-initial-flat-derivatives} give
    %
    \begin{align*}
        |k|&\leq C\tau A/r,\\
        |Dk|&\leq C\tau A/r^2,\\
        |D^2k|&\leq C\tau A/r^3+C\tau A^2/r^4.
    \end{align*}
    %
    Indeed, the first inequality follows by integrating
    $Dh$ along the segment from $p$ to $Q(p)$.
    These segments lie in the larger disk chosen above.
    For the other two inequalities, use
    %
    \begin{align*}
        Dk&=Dh-(Dh\circ Q)(I+De),\\
        D^2k&=D^2h-(D^2h\circ Q)[I+De,I+De]
                   -(Dh\circ Q)D^2e.
    \end{align*}
    %
    The differences $Dh-Dh\circ Q$ and
    $D^2h-D^2h\circ Q$ are bounded by integrating
    $D^2h$ and $D^3h$, respectively, along the same
    segments. Since $\tau/r\leq1$ and $A/r\leq1$, it
    follows that $|D^j(a+k)|\leq CAr^{-j}$ for
    $j\in\{0,1,2\}$.
    The chain rule for $(a+k)\circ Q^{-1}$ now gives
    %
    \begin{align*}
        |\eta|&\leq CA,\qquad |D\eta|\leq CA/r,\\
        |D^2\eta|&\leq CA/r^2+C(A/r)(A/r^2)
                        \leq CA/r^2.
    \end{align*}
    %
    The same calculation for $b=\widetilde b\circ Q^{-1}$
    gives
    %
    \begin{align*}
        |b|\leq C\tau^2,\qquad |Db|\leq C\tau^2/r,\qquad
        |D^2b|\leq C\tau^2/r^2
                  +C(\tau^2/r)(A/r^2)
                  \leq C\tau^2/r^2.
    \end{align*}
    %
    On $U$, $\rho(q')$ is comparable to $r$, and
    $f_\tau(\rho(q'))$ is comparable to $f_\tau(r)$.
    These estimates prove
    \eqref{eq:s4-outer-coordinate-estimates}.
    The bound $|z|\leq C\varepsilon$ follows from
    Lemma~\ref{lem:s4-full-family-geometry}.
    The derivative bounds in
    \eqref{eq:s4-outer-height-bounds} follow by adding
    the estimates for $h$ and $\eta$, using
    $\mathfrak a_m f_\tau(\rho)\leq\sigma_m\tau$.

    Finally, $Z$ is an isometry preserving both surfaces,
    and $G_m$ is equivariant under $Z$.
    It sends the lower initial sheet to the upper one
    and leaves the torus coordinate unchanged.
    The preceding argument therefore applies after this
    reflection to every point on the lower sheet.
\end{proof}

The next lemma estimates the difference between $h$ and
$\psi$ on every part of the sheet, including the catenoidal
part where they are different functions. The term $R^{-2}$
comes from the difference between a catenoidal height and
its logarithmic approximation at distances at least
$R\tau$ from a lattice point. It becomes small by choosing
$R$ large. The other errors become small as $m\to\infty$
with $R$ fixed.

\begin{lemma}[Comparison of the glued and Green heights]
    \label{lem:s4-glued-Green-comparison}
    For every $R\geq R_0$ and all sufficiently large $m$,
    %
    \begin{align}
        \label{eq:s4-glued-Green-comparison}
        \sum_{j=0}^2\rho(q)^j|D^j(h-\psi)(q)|
        \leq C\tau\bigl(|\zeta|+m^2\tau^{2\alpha}+R^{-2}\bigr)
    \end{align}
    %
    for every $q\in\mathbb T$ with $\rho(q)\geq R\tau$.
    Consequently, on every upper graph neighborhood $U$
    from Lemma~\ref{lem:s4-outer-coordinates},
    %
    \begin{align}
        \label{eq:s4-final-Green-comparison}
        \begin{split}
            \sum_{j=0}^2\rho(q)^j|D^j(z-\psi)(q)|
            &\leq C\mu_m(R)\tau
                     \qquad\forall\,q\in U,\\
            \mu_m(R)&:=|\zeta|+m^2\tau^{2\alpha}
                            +\sigma_m+R^{-2}.
        \end{split}
    \end{align}
    %
    The constant $C$ is independent of $R\geq R_0$ and $m$.
    For every fixed $R$, $\mu_m(R)\to R^{-2}$ as
    $m\to\infty$.
\end{lemma}

\begin{proof}
    Since $\widetilde\psi_\zeta
       =\psi-\tau\zeta\widehat v_m$, the definition
    of the glued height and its catenoidal extension gives
    %
    \begin{align}
        \label{eq:s4-height-difference-decomposition}
        h-\psi=-\tau\zeta\widehat v_m
          +\sum_{p\in L_m}\eta_{3r_c,2r_c}(r_p)
                  (C_\tau(r_p)-\widetilde\psi_\zeta)
        \qquad\text{on }\{\rho>\tau\}.
    \end{align}
    %
    Each summand is defined where $\tau<r_p<3r_c$
    and extended by zero where $r_p\geq3r_c$.
    This extension is smooth because the cutoff vanishes
    in a neighborhood of its outer endpoint.
    The disks of radius $3r_c$ about distinct lattice
    points are disjoint for large $m$. Thus at most one
    summand is nonzero at each point, and estimating the
    sum does not introduce a factor $m^2$.

    Fix $p\in L_m$ and a point with
    $R\tau\leq r:=r_p\leq3r_c$.
    On this disk, $\rho=r$ and
    $\widehat v_m=J_0(\sqrt2r)$.
    The local expansion \eqref{eq:s4-physical-expansion}
    and the identity $\log\tau=\log2-A_m+\zeta$ yield
    %
    \begin{align*}
        \widetilde\psi_\zeta
               -\tau\log(2r/\tau)
        &=\tau\bigl(\log r+A_m+\mathcal E_{m,p}
                     -\zeta J_0(\sqrt2r)
                     -\log r-\log2+\log\tau\bigr)\\
        &=\tau\mathcal E_{m,p}
                    +\tau\zeta(1-J_0(\sqrt2r)).
    \end{align*}
    %
    For every $j\in\{0,1,2\}$, the Green remainder
    estimate and the power series for $J_0$ give
    %
    \begin{align*}
        r^j|D^j\mathcal E_{m,p}|
            &\leq Cr^2(m^2+|\log r|),\\
        r^j|D^j(1-J_0(\sqrt2r))|&\leq Cr^2.
    \end{align*}
    %
    The derivatives here are flat derivatives at the
    selected point. The annular estimate
    \eqref{eq:s4-physical-expansion}, used at scale $r$,
    gives the first line. In the second line, the series
    begins with a constant multiple of $r^2$.
    Since $\tau\leq r<1$,
    $|\log r|\leq|\log\tau|\leq Cm^2$.

    To compare the catenoidal height with the same logarithm,
    define
    %
    \begin{align*}
        F(s):=\log\left(\frac{1+\sqrt{1-s}}2\right)
        \qquad\forall\,s\in[0,R_0^{-2}].
    \end{align*}
    %
    Taking $R_0>1$, this is a smooth function with $F(0)=0$
    and bounded first and second derivatives. The exact
    catenoidal formula gives
    %
    \begin{align*}
        C_\tau(r)-\tau\log(2r/\tau)
        =\tau F(\tau^2/r^2).
    \end{align*}
    %
    If $s(r):=\tau^2/r^2$, then
    $s'(r)=-2s/r$ and $s''(r)=6s/r^2$.
    Therefore the last expression, its first radial
    derivative, and its second radial derivative have
    respective bounds
    $C\tau s$, $C\tau s/r$, and $C\tau s/r^2$.
    For a radial function $v=v(r)$ in two flat dimensions,
    %
    \begin{align*}
        Dv&=v'(r)e_r,\\
        D^2v&=v''(r)e_r\otimes e_r
                 +\frac{v'(r)}r(I-e_r\otimes e_r),
        \qquad e_r:=\frac{q-p}{|q-p|}.
    \end{align*}
    %
    Here the calculation is in the flat disk about $p$,
    and $I$ is the identity matrix in its parallel
    orthonormal coordinates. It follows that
    %
    \begin{align*}
        r^j\left|D^j\bigl(C_\tau-\tau\log(2r/\tau)\bigr)\right|
        \leq C\tau^3r^{-2}
        \qquad\forall\,j\in\{0,1,2\}.
    \end{align*}
    %
    Combining the last bounds with $|\zeta|\leq1$ gives
    %
    \begin{align}
        \label{eq:s4-log-profile-error}
        r^j|D^j(C_\tau-\widetilde\psi_\zeta)|
        \leq C\tau\left(m^2r^2+\frac{\tau^2}{r^2}\right)
        \qquad\forall\,j\in\{0,1,2\}.
    \end{align}

    Derivatives of $\eta_{3r_c,2r_c}(r)$ are supported
    where $r$ is comparable to $r_c$ and satisfy
    $|D^j\eta_{3r_c,2r_c}|\leq C_jr_c^{-j}$.
    In a term where $i$ of the $j$ derivatives fall on
    the cutoff, multiplication by $r^j$ leaves a factor
    $r^ir_c^{-i}$ times the scaled derivative of order
    $j-i$ in \eqref{eq:s4-log-profile-error}.
    This factor is uniformly bounded. Thus multiplication
    by the cutoff preserves the bound in that equation.
    On its support, $r\leq3r_c=3\tau^\alpha$ and
    $r\geq R\tau$, so every such term is bounded by
    %
    \begin{align*}
        C\tau\bigl(m^2\tau^{2\alpha}+R^{-2}\bigr).
    \end{align*}

    It remains to estimate the first term of
    \eqref{eq:s4-height-difference-decomposition}.
    The definition of $\widehat v_m$ gives
    %
    \begin{align*}
        \sum_{j=0}^2\rho(q)^j|D^j\widehat v_m(q)|\leq C
        \qquad\forall\,q\in\mathbb T\smallsetminus L_m.
    \end{align*}
    %
    Inside a disk where its cutoff is one, the assertion
    follows from the bounded derivatives of
    $J_0(\sqrt2r_p)$. On each cutoff annulus, $\rho=r_p$
    is comparable to $\delta$, and the factors $\rho^j$
    cancel the factors $\delta^{-j}$ from cutoff
    differentiation. Outside these disks the function is
    zero. The first term in
    \eqref{eq:s4-height-difference-decomposition} is
    consequently bounded by $C\tau|\zeta|$ in the
    stated sum of scaled derivatives.
    Adding the estimates proves
    \eqref{eq:s4-glued-Green-comparison}.

    Finally, on an upper graph neighborhood,
    \eqref{eq:s4-outer-coordinate-estimates} gives
    %
    \begin{align*}
        \sum_{j=0}^2\rho(q)^j|D^j(z-h)(q)|
        \leq C\mathfrak a_m f_\tau(\rho(q))
        \leq C\sigma_m\tau
        \qquad\forall\,q\in U.
    \end{align*}
    %
    The triangle inequality proves
    \eqref{eq:s4-final-Green-comparison}.
    By \eqref{eq:s4-final-balance},
    \eqref{eq:s4-curvature-scalar-smallness}, and the
    exponential decay of $\tau$, the quantities
    $|\zeta|$, $\sigma_m$, and $m^2\tau^{2\alpha}$
    all tend to zero. This proves the last assertion.
\end{proof}

\subsection{Negative curvature away from the neck waists}

The preceding comparison controls the height of the surface
$G_m(M_\zeta)\subset\mathcal N$ and its first two derivatives.
We now relate those quantities to its Gaussian curvature.
We first compute the curvature of an arbitrary height graph
inside $\mathcal N$. We then use the equation satisfied by
$\psi$ to prove that the graphs supplied by
Lemma~\ref{lem:s4-outer-coordinates} have strictly negative
curvature. At this stage we consider the parametrization
$X$ in \eqref{eq:s4-outer-graph-form}; the effect of replacing
$X$ by the final parametrization $X_b$ will be estimated
separately.

For an open set $U\subset\mathbb T$ and a function
$z\in C^2(U,(-z_0,z_0))$, put
%
\begin{align}
    \label{eq:s4-quadratic-curvature}
    \mathcal K_2(z):=\det D^2z+z\Delta_{\mathbb T}z+|Dz|^2.
\end{align}
%
Here the Hessian $D^2z$ and its determinant are computed
with respect to the flat metric $g_{\mathbb T}$, and
$\Delta_{\mathbb T}z=\operatorname{tr}_{g_{\mathbb T}}D^2z$.
The quantity $\mathcal K_2(z)$ consists of the terms of
order two in the height and its first two derivatives
that occur in the curvature formula. The exact curvature
also contains higher-order terms and a positive denominator.
The following lemma bounds their contribution by a small
multiple of the weight $\mathcal W$. This is the comparison
needed to obtain a strict sign from the Green estimates.

\begin{lemma}[Exact curvature formula for a height graph]
    \label{lem:s4-exact-height-curvature}
    Let $U\subset\mathbb T$ be open, let
    $z\in C^2(U,(-z_0,z_0))$, and define the embedding
    $X:U\to\mathcal N$ by
    $X(q):=\mathcal F(q,z(q))$ for every $q\in U$.
    Denote by $K_X$ the Gaussian curvature of the induced
    metric $g_X:=X^*g_{\mathcal N}$, regarded as a function
    on $U$.
    At each point of $U$, write
    %
    \begin{align*}
        c:=\cos z,\qquad s_z:=\sin z,\qquad
        p:=Dz,\qquad H:=D^2z,\qquad Q:=c^2+|p|^2.
    \end{align*}
    %
    Then
    %
    \begin{align}
        \label{eq:s4-exact-height-curvature}
        Q^2K_X=\det H+c s_z\,\operatorname{tr}H+|p|^2+|p|^4
        +2\tan z\,\operatorname{cof}(H)(p,p).
    \end{align}
    %
    Here $\operatorname{cof}(H):=(\operatorname{tr}H)I-H$
    is the cofactor matrix of the symmetric $2\times2$
    matrix $H$, $I$ is the identity matrix, and
    $\operatorname{cof}(H)(p,p):=p^T\operatorname{cof}(H)p$.
    Suppose, in addition, that $R\geq R_0$, that
    $\rho(q)\geq R\tau$ for every $q\in U$, and that
    \eqref{eq:s4-outer-height-bounds} holds on $U$.
    For all sufficiently large $m$, one then has
    %
    \begin{align}
        \label{eq:s4-height-quadratic-error}
        |K_X-\mathcal K_2(z)|
        \leq C(\varepsilon^2+R^{-2})\mathcal W
        \qquad\mbox{on }U.
    \end{align}
    %
    The constant $C$ depends on the constants in
    \eqref{eq:s4-outer-height-bounds} and on $z_0$, but
    is independent of $m$ and $R\geq R_0$.
\end{lemma}

\begin{proof}
    Fix $q\in U$, and choose flat coordinates near $q$.
    Write $\partial_1:=\partial_s$ and
    $\partial_2:=\partial_t$, and put
    $p_i:=\partial_i z$ and $H_{ij}:=\partial_{ij}z$
    for $i,j\in\{1,2\}$. All scalar quantities in the
    following computation are evaluated at $q$, and the
    frame fields $E_1,E_2,E_3$ are evaluated at $X(q)$.
    Differentiating $X$ gives the tangent vectors
    %
    \begin{align*}
        T_i:=dX_q(\partial_i)=cE_i+p_iE_3
        \qquad\forall\,i\in\{1,2\}.
    \end{align*}
    %
    Consequently the induced metric has matrix
    %
    \begin{align*}
        g_{ij}&:=g_{\mathcal N}(T_i,T_j)
        =c^2\delta_{ij}+p_ip_j,\\
        \det(g_{ij})
        &=(c^2+p_1^2)(c^2+p_2^2)-p_1^2p_2^2
        =c^2(c^2+|p|^2)=c^2Q.
    \end{align*}
    %
    In particular, $X$ is an immersion. Its injectivity
    follows from $\pi\circ X=\operatorname{id}_U$, and
    the same identity gives a continuous inverse on its
    image. Thus $X$ is an embedding.

    The unit normal to this graph whose $E_3$ component
    is positive is
    %
    \begin{align*}
        \nu=\frac{-p_1E_1-p_2E_2+cE_3}{\sqrt Q}.
    \end{align*}
    %
    Indeed, its numerator has squared norm $Q$ and is
    orthogonal to both $T_1$ and $T_2$. In the next
    calculation, $h_X^\nu$ is the scalar second fundamental
    form of $X$ as an immersion into $\mathcal N$.
    The connection identities from
    Lemma~\ref{lem:s4-suspension-geometry} give
    %
    \begin{align*}
        \nabla^{\mathcal N}_{T_i}E_j
        &=c\tan z\,\delta_{ij}E_3,\\
        \nabla^{\mathcal N}_{T_i}E_3
        &=-c\tan z\,E_i
        \qquad\forall\,i,j\in\{1,2\}.
    \end{align*}
    %
    Applying the product rule to $T_j=cE_j+p_jE_3$,
    with the connection taken along $X$, yields
    %
    \begin{align*}
        \nabla^{\mathcal N}_{T_i}T_j
        &=-s_z p_iE_j+c\nabla^{\mathcal N}_{T_i}E_j
          +H_{ij}E_3+p_j\nabla^{\mathcal N}_{T_i}E_3\\
        &=-s_z(p_iE_j+p_jE_i)
          +(H_{ij}+c s_z\delta_{ij})E_3.
    \end{align*}
    %
    Taking the inner product with $\nu$, we obtain
    %
    \begin{align*}
        h_X^\nu(\partial_i,\partial_j)
        &=\frac{2s_z p_ip_j+c(H_{ij}+c s_z\delta_{ij})}{\sqrt Q}\\
        &=\frac{c}{\sqrt Q}
        \bigl(H_{ij}+c s_z\delta_{ij}+2\tan z\,p_ip_j\bigr).
    \end{align*}
    %
    Define the symmetric matrix
    $B:=H+c s_z I+2\tan z\,p\otimes p$.
    Since the matrix of the shape operator $S_X^\nu$
    is the product of $(g_{ij})^{-1}$ with the matrix
    of $h_X^\nu$, it follows that
    %
    \begin{align*}
        \det S_X^\nu
        =\frac{\det\bigl(h_X^\nu(\partial_i,\partial_j)\bigr)}
        {\det(g_{ij})}
        =\frac{(c^2/Q)\det B}{c^2Q}
        =\frac{\det B}{Q^2}.
    \end{align*}
    %
    We also need the sectional curvature of the tangent
    plane $\operatorname{span}\{T_1,T_2\}$ in $\mathcal N$.
    The formulas for the shape operator $S$ of
    $\mathcal N\subset\mathbb S^4$ give
    %
    \begin{align*}
        S(T_1)&=-E_1,\\
        S(T_2)&=E_2,
    \end{align*}
    %
    and hence
    %
    \begin{align*}
        g_{\mathcal N}(S(T_1),T_1)&=-c,\\
        g_{\mathcal N}(S(T_2),T_2)&=c,\\
        g_{\mathcal N}(S(T_1),T_2)&=0.
    \end{align*}
    %
    The Gauss equation for $\mathcal N\subset\mathbb S^4$
    therefore gives
    %
    \begin{align*}
        \operatorname{sec}_{g_{\mathcal N}}
        (\operatorname{span}\{T_1,T_2\})
        &=1+\frac{(-c)c}{c^2Q}
        =1-\frac1Q
        =\frac{|p|^2-s_z^2}{Q}.
    \end{align*}
    %
    Applying the Gauss equation once more, now to
    $X:U\to\mathcal N$, yields
    %
    \begin{align*}
        Q^2K_X=Q(|p|^2-s_z^2)+\det B.
    \end{align*}
    %
    To expand $\det B$, recall that for a symmetric
    $2\times2$ matrix $A$, a vector $p\in\mathbb R^2$
    and a scalar $t\in\mathbb R$,
    %
    \begin{align*}
        \det(A+t p\otimes p)
        &=\det A+t(A_{22}p_1^2-2A_{12}p_1p_2+A_{11}p_2^2)\\
        &=\det A+t\operatorname{cof}(A)(p,p).
    \end{align*}
    %
    The terms of order two in $t$ cancel in the
    determinant. Apply this identity with $A=H+c s_z I$
    and $t=2\tan z$. Since
    %
    \begin{align*}
        \det(H+c s_z I)&=\det H+c s_z\operatorname{tr}H+c^2s_z^2,\\
        \operatorname{cof}(H+c s_z I)&=\operatorname{cof}(H)+c s_z I,
    \end{align*}
    %
    we obtain
    %
    \begin{align*}
        \det B
        =\det H+c s_z\operatorname{tr}H+c^2s_z^2
        +2\tan z\operatorname{cof}(H)(p,p)+2s_z^2|p|^2.
    \end{align*}
    %
    Finally,
    %
    \begin{align*}
        Q(|p|^2-s_z^2)
        =c^2|p|^2+|p|^4-c^2s_z^2-s_z^2|p|^2.
    \end{align*}
    %
    Adding these two identities cancels $c^2s_z^2$ and
    leaves $(c^2+s_z^2)|p|^2=|p|^2$ as the terms of
    degree two in $p$. This proves
    \eqref{eq:s4-exact-height-curvature}.

    We now assume the bounds in the second part of the
    lemma. For large $m$, $\rho\leq\pi/m<1$ and
    %
    \begin{align*}
        |p|\leq C\tau\rho^{-1}\leq C/R,\qquad
        |H|\leq C\tau\rho^{-2},\qquad
        |z|\leq C\varepsilon\leq C\tau\rho^{-2}.
    \end{align*}
    %
    The last inequality uses $m^2\rho^2\leq\pi^2$.
    Consequently
    %
    \begin{align*}
        |H|^2&\leq C\tau^2\rho^{-4}\leq C\mathcal W,\\
        |z||H|&\leq C\tau^2\rho^{-4}\leq C\mathcal W,\\
        |p|^2&\leq C\tau^2\rho^{-2}\leq C\mathcal W.
    \end{align*}
    %
    Since $|\det H|\leq|H|^2/2$ and
    $|\operatorname{tr}H|\leq\sqrt2|H|$, these bounds imply
    $|\mathcal K_2(z)|\leq C\mathcal W$.
    Moreover $Q$ is bounded above and below by positive
    constants independent of $m$ and $R\geq R_0$. Indeed
    $Q\geq\cos^2z_0>0$, and $|p|\leq C/R_0$.
    The identities $Q-1=-\sin^2z+|p|^2$ and
    $Q^{-2}-1=(1-Q)(1+Q)/Q^2$ give
    %
    \begin{align*}
        |Q^{-2}-1|\leq C(z^2+|p|^2).
    \end{align*}
    %
    The difference between the numerator in
    \eqref{eq:s4-exact-height-curvature} and
    $\mathcal K_2(z)$ is exactly
    %
    \begin{align*}
        E:=(c s_z-z)\operatorname{tr}H+|p|^4
        +2\tan z\operatorname{cof}(H)(p,p).
    \end{align*}
    %
    For $|z|<z_0$, Taylor's theorem gives
    $|\cos z\sin z-z|\leq C|z|^3$ and
    $|\tan z|\leq C|z|$. Moreover,
    $|\operatorname{cof}(H)(p,p)|\leq|H||p|^2$.
    Hence
    %
    \begin{align*}
        |E|&\leq C\bigl(|z|^3|H|+|p|^4
              +|z||H||p|^2\bigr)\\
        &\leq C(z^2+|p|^2)\mathcal W.
    \end{align*}
    %
    Each term in the second line is controlled by one of
    the three preceding bounds: for example,
    $|z|^3|H|=z^2(|z||H|)$ and
    $|z||H||p|^2=|p|^2(|z||H|)$.
    It follows that
    %
    \begin{align*}
        |K_X-\mathcal K_2(z)|
        &\leq |Q^{-2}-1|\,|\mathcal K_2(z)|+Q^{-2}|E|\\
        &\leq C(z^2+|p|^2)\mathcal W
        \leq C(\varepsilon^2+R^{-2})\mathcal W,
    \end{align*}
    %
    as required.
\end{proof}

We apply this formula first to the Green height $\psi$.
Because $\Delta_{\mathbb T}\psi=-2\psi$ away from the lattice,
the terms involving $\psi^2$ cancel in
$\mathcal K_2(\psi)$. What remains is
$|D\psi|^2-\tfrac12|D^2\psi|^2$, whose negative upper
bound was proved in Section~\ref{sec:s4-scalar-family}.
We then estimate the effect of replacing $\psi$ by the
actual graph height $z$. The comparison must be made
relative to $\mathcal W$, since this negative upper
bound tends to zero on the parts of the sheets furthest
from the necks.

\begin{lemma}[Negative curvature of the graphs inside $\mathcal N$]
    \label{lem:s4-scalar-negative-curvature}
    There exists $c_*>0$ such that, after choosing $R$
    sufficiently large and then $m$ sufficiently large,
    every scalar graph $X:U\to\mathcal N$ from
    Lemma~\ref{lem:s4-outer-coordinates} satisfies
    %
    \begin{align}
        \label{eq:s4-scalar-negative-curvature}
        K_X(q)\leq-c_*\mathcal W(q)
    \end{align}
    %
    at each point in its domain with $\rho(q)\geq R\tau$.
\end{lemma}

\begin{proof}
    On $\mathbb T\smallsetminus L_m$ we have
    $\Delta_{\mathbb T}\psi=-2\psi$. For any symmetric
    $2\times2$ matrix $H$,
    $\det H=((\operatorname{tr}H)^2-|H|^2)/2$.
    Therefore
    %
    \begin{align*}
        \mathcal K_2(\psi)
        &=\tfrac12\bigl((\Delta_{\mathbb T}\psi)^2
            -|D^2\psi|^2\bigr)
            +\psi\Delta_{\mathbb T}\psi+|D\psi|^2\\
        &=2\psi^2-\tfrac12|D^2\psi|^2
            -2\psi^2+|D\psi|^2\\
        &=|D\psi|^2-\tfrac12|D^2\psi|^2.
    \end{align*}
    %
    Recall that $\psi=\varepsilon U_m$. At each point
    of $\mathbb T\smallsetminus L_m$, set
    $r:=\sqrt2m\rho$. Since $\varepsilon$ is constant on
    $\mathbb T$, multiplication of
    \eqref{eq:s4-height-quadratic-margin} by
    $\varepsilon^2$ gives
    %
    \begin{align*}
        \mathcal K_2(\psi)
        &=\varepsilon^2\bigl(|DU_m|^2-\tfrac12|D^2U_m|^2\bigr)\leq-c\varepsilon^2(1+r^{-4})
        =-c\left(\varepsilon^2
             +\frac{\tau^2\rho^{-4}}{16\pi^2}\right).
    \end{align*}
    %
    In the last equality we used
    $\varepsilon^2r^{-4}=\tau^2\rho^{-4}/(16\pi^2)$.
    Since $\mathcal W=\varepsilon^2+\tau^2\rho^{-4}$,
    reducing the positive constant if necessary gives
    %
    \begin{align}
        \label{eq:s4-Green-curvature-margin}
        \mathcal K_2(\psi)\leq-c_0\mathcal W,
    \end{align}
    %
    where $c_0>0$ is independent of $m$ and $R$.
    On $\rho\geq R\tau$, the Green estimates also give
    %
    \begin{align*}
        |\psi|\leq C\varepsilon\leq C\tau\rho^{-2},\qquad
        |D\psi|\leq C\tau\rho^{-1},\qquad
        |D^2\psi|\leq C\tau\rho^{-2}.
    \end{align*}
    %
    The derivative bounds follow from the Green estimates
    in Section~\ref{sec:s4-scalar-family}. For the value
    of $\psi$, use
    $\Phi_m=m^2/(2\pi)+\log(\sqrt2m\rho)+O(1)$.
    Since $R\geq R_0>1$, $\rho\geq R\tau$ and
    $\rho\leq\pi/m<1$, we have
    $|\log\rho|\leq|\log\tau|\leq Cm^2$ for large $m$.
    Thus $|\psi|\leq C\tau m^2=C(2\pi)\varepsilon$.
    The further bound $\varepsilon\leq C\tau\rho^{-2}$
    follows from $m^2\rho^2\leq\pi^2$.

    Consider an upper graph neighborhood. All estimates
    below are at its points with $\rho\geq R\tau$.
    Put $e:=z-\psi$ and $\mu:=\mu_m(R)$. By
    Lemma~\ref{lem:s4-glued-Green-comparison},
    %
    \begin{align*}
        |e|\leq C\mu\tau,\qquad
        |De|\leq C\mu\tau\rho^{-1},\qquad
        |D^2e|\leq C\mu\tau\rho^{-2}.
    \end{align*}
    %
    Using
    $\det(H+E)-\det H=\operatorname{cof}(H):E+\det E$
    for symmetric $2\times2$ matrices, we obtain
    %
    \begin{align*}
        \mathcal K_2(z)-\mathcal K_2(\psi)
        ={}&\operatorname{cof}(D^2\psi):D^2e+\det D^2e
             +\psi\Delta_{\mathbb T}e
             +e\Delta_{\mathbb T}\psi
             +e\Delta_{\mathbb T}e\\
        &+2\langle D\psi,De\rangle+|De|^2.
    \end{align*}
    %
    The colon denotes the sum of the products of corresponding
    matrix entries. For symmetric $2\times2$ matrices,
    $|\operatorname{cof}(H)|=|H|$, and
    $|\det H|\leq|H|^2/2$. The Hessian terms therefore
    satisfy
    %
    \begin{align*}
        |\operatorname{cof}(D^2\psi):D^2e|
        &\leq C\mu\tau^2\rho^{-4},\\
        |\det D^2e|&\leq C\mu^2\tau^2\rho^{-4}.
    \end{align*}
    %
    For the terms involving the heights, the preceding
    bounds give
    %
    \begin{align*}
        |\psi\Delta_{\mathbb T}e|
        &\leq C(\tau\rho^{-2})(\mu\tau\rho^{-2})
        =C\mu\tau^2\rho^{-4},\\
        |e\Delta_{\mathbb T}\psi|
        &\leq C(\mu\tau)(\tau\rho^{-2})
        =C\mu\tau^2\rho^{-2},\\
        |e\Delta_{\mathbb T}e|
        &\leq C(\mu\tau)(\mu\tau\rho^{-2})
        =C\mu^2\tau^2\rho^{-2}.
    \end{align*}
    %
    Finally, the gradient terms satisfy
    %
    \begin{align*}
        2|\langle D\psi,De\rangle|+|De|^2
        \leq C(\mu+\mu^2)\tau^2\rho^{-2}.
    \end{align*}
    %
    Since $\rho<1$, every occurrence of $\rho^{-2}$
    is bounded by $\rho^{-4}$. Summing these estimates
    and using $\tau^2\rho^{-4}\leq\mathcal W$ yields
    %
    \begin{align*}
        |\mathcal K_2(z)-\mathcal K_2(\psi)|
        \leq C(\mu+\mu^2)\mathcal W.
    \end{align*}
    %
    Combining this with
    \eqref{eq:s4-Green-curvature-margin} and
    \eqref{eq:s4-height-quadratic-error} yields
    %
    \begin{align*}
        K_X\leq
        \bigl[-c_0+C(\mu+\mu^2+\varepsilon^2+R^{-2})\bigr]
        \mathcal W.
    \end{align*}
    %
    To make the choice of parameters explicit, write
    $\mu=t_m+R^{-2}$, where
    $t_m:=|\zeta|+m^2\tau^{2\alpha}+\sigma_m\to0$.
    First choose $R\geq R_0$ so large that
    $C(2R^{-2}+R^{-4})<c_0/4$.
    Keeping this $R$ fixed, choose $m$ so large that
    $C(t_m+t_m^2+2t_mR^{-2}+\varepsilon^2)<c_0/4$.
    We then have $K_X\leq-(c_0/2)\mathcal W$, proving
    the claim with $c_*:=c_0/2$.
    Reflection by $Z$ preserves the induced metric, so the
    same curvature bound holds on the lower graphical parts.
\end{proof}

We now estimate the change in Gaussian curvature produced by
$b$, which moves the graph $X$ in the direction normal to
$\mathcal N$ within $\mathbb S^4$. The bound $|b|\leq C\tau^2$
alone does not control this change: curvature also depends on
the first and second derivatives of the parametrization. In
fact, the bound for $|D^2b|$ is only $C/R^2$ when
$\rho=R\tau$, and need not tend to zero for fixed $R$.
The following estimate compares the curvature change directly
with the positive function $\mathcal W$ in the preceding
lemma. It will therefore allow us to preserve the negative
upper bound for $K_X$.

\begin{lemma}[Curvature change under the deformation normal to $\mathcal N$]
    \label{lem:s4-transverse-curvature-comparison}
    Let $X,X_b:U\to\mathbb S^4$ be the parametrizations in
    Lemma~\ref{lem:s4-outer-coordinates}. For every $q\in U$
    with $\rho(q)\geq R\tau$, one has
    %
    \begin{align}
        \label{eq:s4-transverse-curvature-comparison}
        |K_{X_b}(q)-K_X(q)|
        \leq C(m^{-2}+\tau)\mathcal W(q).
    \end{align}
    %
    Here $K_X$ and $K_{X_b}$ are the Gaussian curvatures of
    the metrics induced on $U$ by the indicated
    parametrizations. The constant $C$ is independent of
    $R\geq R_0$ and of sufficiently large $m$.
\end{lemma}

\begin{proof}
    We use local flat coordinates $(s,t)$ on $U$ and regard
    $X$ and $X_b$ as maps into $\mathbb R^5$. Derivatives
    and matrix norms in this proof are Euclidean. All
    estimates are evaluated at a point $q\in U$ as in the
    statement, and we write $\rho$ for $\rho(q)$. We may
    assume $R_0\geq1$, so that
    %
    \begin{align*}
        \tau\leq\rho\leq\frac{\pi}{m}<1.
    \end{align*}
    %
    Put $\Lambda:=|D^2z(q)|$. Since all derivatives of the
    fixed maps $f$ and $N$ are bounded, differentiating
    $X=\cos z\,f+\sin z\,e_5$ and using
    \eqref{eq:s4-outer-height-bounds} gives
    %
    \begin{align*}
        |DX|&\leq C,
        & |D^2X|&\leq C(1+\Lambda),
        & \Lambda&\leq C\tau\rho^{-2}.
    \end{align*}
    %
    The matrix $g_X:=(DX)^T DX$ of the induced metric
    satisfies the exact identity
    %
    \begin{align*}
        g_X=\cos^2z\,I+Dz\otimes Dz.
    \end{align*}
    %
    In particular, $g_X\geq\cos^2z_0\,I$, and its
    eigenvalues are bounded above and below by positive
    constants independent of $m$ and $R$.

    We first estimate the change in the first two
    derivatives of the parametrization. Write
    $\partial_1:=\partial_s$ and $\partial_2:=\partial_t$.
    For every $i,j\in\{1,2\}$, differentiating
    $X_b=\cos b\,X+\sin b\,N$ gives
    %
    \begin{align*}
        \partial_iX_b-\partial_iX
        ={}&(\cos b-1)\partial_iX+\sin b\,\partial_iN
             +b_i(-\sin b\,X+\cos b\,N),\\
        \partial_{ij}X_b-\partial_{ij}X
        ={}&(\cos b-1)\partial_{ij}X
             +\sin b\,\partial_{ij}N\\
        &+b_i(-\sin b\,\partial_jX+\cos b\,\partial_jN)
         +b_j(-\sin b\,\partial_iX+\cos b\,\partial_iN)\\
        &+b_{ij}(-\sin b\,X+\cos b\,N)
         -b_ib_jX_b,
    \end{align*}
    %
    where $b_i:=\partial_i b$ and
    $b_{ij}:=\partial_{ij}b$. Define the nonnegative
    quantities
    %
    \begin{align*}
        E_1:=\tau^2\rho^{-1},\qquad
        E_2:=\tau^2\rho^{-2}.
    \end{align*}
    %
    The coordinate estimates
    \eqref{eq:s4-outer-coordinate-estimates} imply
    %
    \begin{align}
        \label{eq:s4-transverse-E12}
        |DX_b-DX|\leq CE_1,\qquad
        |D^2X_b-D^2X|\leq CE_2.
    \end{align}
    %
    To check the second bound, the first term in the
    second derivative formula is bounded by
    %
    \begin{align*}
        C|b|^2(1+\Lambda)
        &\leq C\tau^4(1+\tau\rho^{-2})
        \leq C\tau^2\rho^{-2}.
    \end{align*}
    %
    Each of the remaining terms is bounded by a constant
    times one of $|b|$, $|Db|$, $|D^2b|$, or $|Db|^2$.
    Their sum is at most $CE_2$, since $\tau<1$ and
    $\rho<1$. The first derivative formula similarly gives
    %
    \begin{align*}
        |DX_b-DX|
        \leq C\bigl(|b|+|Db|\bigr)
        \leq C(\tau^2+\tau^2\rho^{-1})
        \leq CE_1.
    \end{align*}
    %
    Moreover, $E_1\leq\tau/R\leq\tau\to0$.
    Thus the eigenvalues of $g_{X_b}$ also have uniform
    positive lower and upper bounds for large $m$.

    We next express curvature in terms of these
    derivatives. For either immersion $Y=X$ or $Y=X_b$,
    let $P_Y$ be the Euclidean orthogonal projection onto
    the normal space of $Y$ in $\mathbb R^5$. In matrix
    notation,
    %
    \begin{align*}
        P_Y=I-DY(g_Y)^{-1}(DY)^T,
        \qquad g_Y:=(DY)^TDY.
    \end{align*}
    %
    Its Euclidean second fundamental form has coordinate
    components $P_Y\partial_{ij}Y$. The Gauss equation
    therefore gives
    %
    \begin{align}
        \label{eq:s4-euclidean-Gauss-comparison}
        K_Y=
        \frac{\langle P_Y\partial_{ss}Y,
                         P_Y\partial_{tt}Y\rangle
        -|P_Y\partial_{st}Y|^2}{\det g_Y}.
    \end{align}
    %
    This computes the Gaussian curvature of the metric
    induced by $Y$ also when $Y$ takes values in
    $\mathbb S^4$: the Euclidean normal projection includes
    the radial direction of the sphere, so that no
    additional ambient-curvature term is to be added to
    this formula.

    Write $g:=g_X$ and $\widetilde g:=g_{X_b}$.
    Expanding the products of first derivatives in the
    metric matrices gives
    $|\widetilde g-g|\leq CE_1$. Their uniform positive
    lower bounds and the identity
    %
    \begin{align*}
        \widetilde g^{-1}-g^{-1}
        =\widetilde g^{-1}(g-\widetilde g)g^{-1}
    \end{align*}
    %
    then give $|\widetilde g^{-1}-g^{-1}|\leq CE_1$.
    Substitution into the projection formula yields the
    first of the following estimates; the second follows
    by expanding the $2\times2$ determinants and using
    their positive lower bounds:
    %
    \begin{align*}
        |P_{X_b}-P_X|&\leq CE_1,\\
        |(\det g_{X_b})^{-1}-(\det g_X)^{-1}|
        &\leq CE_1.
    \end{align*}
    %
    For $i,j\in\{1,2\}$, put
    %
    \begin{align*}
        T_{ij}:=P_X\partial_{ij}X,\qquad
        \widetilde T_{ij}:=P_{X_b}\partial_{ij}X_b.
    \end{align*}
    %
    Since orthogonal projections have operator norm one,
    %
    \begin{align*}
        |T_{ij}|&\leq C(1+\Lambda),\\
        |\widetilde T_{ij}-T_{ij}|
        &\leq|P_{X_b}(\partial_{ij}X_b-\partial_{ij}X)|
               +|(P_{X_b}-P_X)\partial_{ij}X|\leq C\bigl(E_2+E_1(1+\Lambda)\bigr).
    \end{align*}
    %
    Define
    $\mathcal Q:=\langle T_{11},T_{22}\rangle-|T_{12}|^2$
    and let $\widetilde{\mathcal Q}$ be the corresponding
    expression with $\widetilde T_{ij}$. Subtracting these
    quadratic expressions gives
    %
    \begin{align*}
        |\widetilde{\mathcal Q}-\mathcal Q|
        \leq C\bigl[(1+\Lambda)
                       \bigl(E_2+E_1(1+\Lambda)\bigr)+\bigl(E_2+E_1(1+\Lambda)\bigr)^2\bigr].
    \end{align*}
    %
    For large $m$ we have $E_1\leq1$, and hence
    %
    \begin{align*}
        \bigl(E_2+E_1(1+\Lambda)\bigr)^2
        \leq2E_2^2+2E_1(1+\Lambda)^2.
    \end{align*}
    %
    Using also $|\mathcal Q|\leq C(1+\Lambda)^2$,
    we obtain from \eqref{eq:s4-euclidean-Gauss-comparison}
    %
    \begin{align*}
        |K_{X_b}-K_X|
        &\leq
        \frac{|\widetilde{\mathcal Q}-\mathcal Q|}
             {\det\widetilde g}
        +|\mathcal Q|\,
             |(\det\widetilde g)^{-1}-(\det g)^{-1}|\leq C\bigl((1+\Lambda)E_2
                 +E_1(1+\Lambda)^2+E_2^2\bigr).
    \end{align*}
    %
    Substituting the definitions of $E_1,E_2$ and
    $\Lambda\leq C\tau\rho^{-2}$ gives
    %
    \begin{align*}
        |K_{X_b}-K_X|
        \leq C\left(
        \frac{\tau^2}{\rho^2}+\frac{\tau^3}{\rho^4}
        +\frac{\tau^2}{\rho}+\frac{\tau^3}{\rho^3}
        +\frac{\tau^4}{\rho^5}+\frac{\tau^4}{\rho^4}
        \right).
    \end{align*}
    %
    We compare each term with
    $\mathcal W\geq\tau^2\rho^{-4}$. Division by this
    lower bound yields
    %
    \begin{align*}
        \frac{|K_{X_b}-K_X|}{\mathcal W}
        &\leq C\left(
             \rho^2+\tau+\rho^3+\tau\rho
             +\frac{\tau^2}{\rho}+\tau^2\right)\leq C(m^{-2}+\tau).
    \end{align*}
    %
    In the last inequality we used
    $\rho\leq\pi/m<1$, $\tau\leq\rho$, and $\tau<1$.
    All constants depend only on the fixed geometry and
    the constants in Lemma~\ref{lem:s4-outer-coordinates},
    and are therefore independent of $R\geq R_0$.
\end{proof}

\subsection{Curvature near the neck waists and the area limit}

At a neck waist, the projection of the initial surface onto
$\mathbb T$ is not a local diffeomorphism. Thus the height
coordinates used in the preceding subsections do not cover
that circle. We instead use the catenoidal parametrization
from Section~\ref{sec:s4-scalar-family}. Dividing lengths by
$\tau$ transforms the region $\rho_\zeta\leq2R\tau$ into a
fixed portion of the unit catenoid. We will prove that the
final immersions converge to this catenoid in $C^2$ in the
rescaled coordinates. Its curvature is bounded above by a
strictly negative constant on this fixed portion, which gives
the remaining curvature estimate. We then combine the two
regions and compare the area elements before and after both
deformations.

\begin{lemma}[Negative curvature near the neck waists]
    \label{lem:s4-core-negative-curvature}
    For every fixed $R\geq R_0$, there exists $m_R\in\mathbb N$
    such that, whenever $m\geq m_R$,
    %
    \begin{align}
        \label{eq:s4-core-negative-curvature}
        K_{\Sigma_m}(\widehat X_m(x))
        \leq-\frac{1}{2(2R)^4\tau^2}
        \qquad\forall\,x\in M_\zeta
        \mbox{ such that }\rho_\zeta(x)\leq2R\tau.
    \end{align}
\end{lemma}

\begin{proof}
    Fix $R\geq R_0$. Since $\tau^{1-\alpha}\to0$, we have
    $3R\tau<2r_c$ for all sufficiently large $m$. The part
    of $M_\zeta$ with $\rho_\zeta\leq3R\tau$ therefore lies
    entirely in the catenoidal pieces defined by
    \eqref{eq:s4-catenoid-parametrization}.

    We first specify the ambient coordinates and the metric
    after rescaling. Fix $p\in L_m$, let
    $B_R:=(-5R,5R)^4\subset\mathbb R^4$, and define
    %
    \begin{align*}
        \mathcal P_{p,\tau}(y_1,y_2,v,w)
        :=\Psi\bigl(p+\tau(y_1,y_2),\tau v,\tau w\bigr)
        \qquad\forall\,(y_1,y_2,v,w)\in B_R.
    \end{align*}
    %
    The addition to $p$ is in flat coordinates on $\mathbb T$.
    This map is defined and is a diffeomorphism onto its
    image for sufficiently large $m$. Indeed, the horizontal
    displacement is at most $5\sqrt2R\tau<\delta$ for large
    $m$, and the remaining two coordinates have absolute
    value at most $5R\tau$. Thus its values lie in the
    domain of the tubular embedding $\Psi$.
    The metric that measures lengths in units of $\tau$ is
    %
    \begin{align*}
        \overline g_{p,\tau}
        :=\tau^{-2}\mathcal P_{p,\tau}^*g_{\mathbb S^4}.
    \end{align*}
    %
    Write $\widetilde g_{ab}$ for the coefficients of
    $\Psi^*g_{\mathbb S^4}$ in the coordinates $(s,t,z,b)$.
    The chain rule gives
    %
    \begin{align*}
        (\overline g_{p,\tau})_{ab}(y_1,y_2,v,w)
        =\widetilde g_{ab}
          \bigl(p+\tau(y_1,y_2),\tau v,\tau w\bigr)
        \qquad\forall\,a,b\in\{1,2,3,4\}.
    \end{align*}
    %
    At $(p,0,0)$, the four coordinate vectors of $\Psi$
    are $E_1,E_2,E_3,n$ at $f(p)$, so
    $\widetilde g_{ab}(p,0,0)=\delta_{ab}$. The coefficients
    $\widetilde g_{ab}$ and all their derivatives are bounded
    on a fixed smaller tubular band containing the coordinate
    images under consideration. Taylor's theorem consequently
    bounds the difference of the displayed coefficients from
    $\delta_{ab}$ by $C_R\tau$ on $B_R$. A derivative
    of positive order $j$ in the rescaled variables introduces
    a factor $\tau^j$. Thus $\overline g_{p,\tau}$ converges
    to the Euclidean metric in every fixed $C^k$ norm on
    $B_R$, uniformly for $p\in L_m$.

    For $A\in\{2R,3R\}$, put
    %
    \begin{align*}
        D_A:=[-\operatorname{arcosh}A,\operatorname{arcosh}A]
                  \times\bigl(\mathbb R/(2\pi\mathbb Z)\bigr).
    \end{align*}
    %
    In the coordinates just defined, the initial catenoidal
    parametrization is the fixed map
    %
    \begin{align*}
        c(\xi,\theta)
        :=\bigl(\cosh\xi\cos\theta,\cosh\xi\sin\theta,\xi,0\bigr)
        \qquad\forall\,(\xi,\theta)\in D_{3R}.
    \end{align*}
    %
    More precisely, if $\mathcal C_{\zeta,p}$ denotes the
    parametrization in \eqref{eq:s4-catenoid-parametrization},
    then
    %
    \begin{align*}
        \mathcal C_{\zeta,p}=\mathcal P_{p,\tau}\circ c,
        \qquad
        \rho_\zeta(\mathcal C_{\zeta,p}(\xi,\theta))
        =\tau\cosh\xi
        \qquad\forall\,(\xi,\theta)\in D_{3R}.
    \end{align*}
    %
    The restriction of $\overline g_{p,\tau}$ to $w=0$
    is
    %
    \begin{align*}
        dv^2+\cos^2(\tau v)(dy_1^2+dy_2^2).
    \end{align*}
    %
    Therefore the metric induced on $D_{3R}$ is exactly
    %
    \begin{align*}
        \tau^{-2}\mathcal C_{\zeta,p}^*g_\zeta
        =\bigl(1+\cos^2(\tau\xi)\sinh^2\xi\bigr)d\xi^2+\cos^2(\tau\xi)\cosh^2\xi\,d\theta^2.
    \end{align*}
    %
    Its coefficients converge smoothly to those of
    $\cosh^2\xi(d\xi^2+d\theta^2)$, and their eigenvalues
    are bounded above and below by positive constants
    depending only on $R$.

    We next estimate both displacement functions in these
    fixed parameters. For a smooth function $u$ on the
    initial catenoidal piece, the preceding metric formula
    implies
    %
    \begin{align*}
        \left\|\frac{u\circ\mathcal C_{\zeta,p}}{\tau}
                \right\|_{C^2(D_{3R})}
        \leq \frac{C_R}{\tau}
        \sup_{\mathcal C_{\zeta,p}(D_{3R})}
        \bigl(|u|+\tau|\nabla u|_{g_\zeta}
                    +\tau^2|\nabla^2u|_{g_\zeta}\bigr).
    \end{align*}
    %
    Here the $C^2$ norm is the sum of the supremum norms of
    the ordinary $(\xi,\theta)$ derivatives through order two.
    To see the second-order assertion explicitly, the Hessian
    in these parameters satisfies
    %
    \begin{align*}
        \partial_{ij}(u\circ\mathcal C_{\zeta,p})
        = (\mathcal C_{\zeta,p}^*\nabla^2u)_{ij}
          +\Gamma^k_{ij}
             \partial_k(u\circ\mathcal C_{\zeta,p}),
        \qquad i,j\in\{1,2\},
    \end{align*}
    %
    with summation over $k\in\{1,2\}$. The Christoffel
    symbols $\Gamma^k_{ij}$ are those of the displayed
    rescaled metric and are uniformly bounded on $D_{3R}$.
    Its comparison with the product metric bounds the two
    terms by $C_R\tau^2|\nabla^2u|_{g_\zeta}$ and
    $C_R\tau|\nabla u|_{g_\zeta}$, respectively.

    Since $\rho_\zeta=\tau\cosh\xi$ on this piece,
    %
    \begin{align*}
        f_\tau(\rho_\zeta)
        =\max\{\tau^{3/2}\cosh^{3/2}\xi,
                 \tau^{1-\alpha/2}\cosh^{1/2}\xi\}
        \leq C_R\tau^{1-\alpha/2}.
    \end{align*}
    %
    Define functions on $D_{3R}$ by
    %
    \begin{align*}
        \overline\phi_{p,m}
        :=\frac{\phi_m\circ\mathcal C_{\zeta,p}}\tau,
        \qquad
        \overline b_{p,m}
        :=\frac{b_m\circ\mathcal C_{\zeta,p}}\tau.
    \end{align*}
    %
    Applying \eqref{eq:s4-final-two-jets} in the preceding
    norm comparison gives
    %
    \begin{align*}
        \|\overline\phi_{p,m}\|_{C^2(D_{3R})}
        &\leq C_R\delta^{-4-2\beta}
            \tau^{1+\alpha/3}\tau^{1-\alpha/2}\tau^{-1}
          =C_R\delta^{-4-2\beta}\tau^{1-\alpha/6}
          \rightarrow0,\\
        \|\overline b_{p,m}\|_{C^2(D_{3R})}
        &\leq C_R\tau\rightarrow0.
    \end{align*}
    %
    Both limits are uniform in $p$. The first follows because
    $1-\alpha/6>0$ and every positive power of $\tau$ tends
    to zero faster than every fixed inverse power of $m$.

    These are estimates for scalar functions on the initial
    parametrization. We now relate them to the final
    immersion. Let $\overline g^\mathcal N_{p,\tau}$ be the
    restriction of $\overline g_{p,\tau}$ to $w=0$, and let
    $\overline\nu_{p,\tau}$ be the unit normal to $c$ in that
    metric corresponding to $\nu_\zeta$. Equivalently,
    %
    \begin{align*}
        d\mathcal P_{p,\tau}|_{c(\xi,\theta)}
          \bigl(\overline\nu_{p,\tau}(\xi,\theta)\bigr)
        =\tau\nu_\zeta(\mathcal C_{\zeta,p}(\xi,\theta)).
    \end{align*}
    %
    Define the normal exponential map in these coordinates by
    %
    \begin{align*}
        \mathcal E_{p,\tau}(y,t)
        :=\exp_{c(y)}^{\overline g^\mathcal N_{p,\tau}}
              \bigl(t\overline\nu_{p,\tau}(y)\bigr),
        \qquad y=(\xi,\theta)\in D_{3R}.
    \end{align*}
    %
    This map is defined for $|t|$ less than a fixed positive
    constant depending only on $R$, on a slightly larger
    parameter strip as well. Its derivatives needed below
    are uniformly bounded, because the metrics converge
    smoothly and the initial normal is determined smoothly
    by its unit-length and orthogonality conditions.
    Multiplication of a metric by a positive constant does
    not change its Levi-Civita connection or exponential map.
    The definitions of $G_m$ and $\mathcal P_{p,\tau}$ therefore
    give the exact formula
    %
    \begin{align*}
        \mathcal P_{p,\tau}^{-1}
          \circ G_m\circ\mathcal C_{\zeta,p}(y)
        =\mathcal E_{p,\tau}
           \bigl(y,\overline\phi_{p,m}(y)\bigr).
    \end{align*}
    %
    Both sides lie in $w=0$.

    For completeness, write $a:=\overline\phi_{p,m}$ and
    $\mathcal E:=\mathcal E_{p,\tau}$. The first two
    derivatives of this composition are
    %
    \begin{align*}
        \partial_i[\mathcal E(y,a(y))]
        &=\mathcal E_i+\mathcal E_t\partial_i a,\\
        \partial_{ij}[\mathcal E(y,a(y))]
        &=\mathcal E_{ij}
          +\mathcal E_{it}\partial_j a
          +\mathcal E_{jt}\partial_i a
          +\mathcal E_{tt}\partial_i a\partial_j a
          +\mathcal E_t\partial_{ij}a,
    \end{align*}
    %
    for $i,j\in\{1,2\}$, where all derivatives of
    $\mathcal E$ on the right are evaluated at $(y,a(y))$.
    Since $\mathcal E(y,0)=c(y)$, subtracting the same
    formulas at $a=0$ and applying the mean value theorem
    yields
    %
    \begin{align*}
        \|\mathcal E(\,
          \cdot,a(\cdot))-c\|_{C^2(D_{3R})}
        \leq C_R\|a\|_{C^2(D_{3R})}
                    \bigl(1+\|a\|_{C^2(D_{3R})}\bigr)^2
        \rightarrow0.
    \end{align*}
    %
    In tubular coordinates, the second deformation leaves
    these three coordinates unchanged and sets the fourth
    coordinate equal to $\overline b_{p,m}$. Thus, writing
    $\mathcal E^1,\mathcal E^2,\mathcal E^3$ for the first
    three components of $\mathcal E$, we have
    %
    \begin{align*}
        \mathcal P_{p,\tau}^{-1}
          \circ\widehat X_m\circ\mathcal C_{\zeta,p}(y)
        =\bigl(\mathcal E^1(y,a(y)),
                \mathcal E^2(y,a(y)),
                \mathcal E^3(y,a(y)),
                \overline b_{p,m}(y)\bigr).
    \end{align*}
    %
    It follows that these final rescaled immersions converge
    to $c$ in $C^2$, uniformly for $p\in L_m$. The larger
    strip $D_{3R}$ provides neighborhoods of all the points
    in the smaller strip $D_{2R}$, including its boundary.

    We compute the curvature of the limiting immersion.
    Its first derivatives satisfy
    %
    \begin{align*}
        \langle c_\xi,c_\xi\rangle
        =\langle c_\theta,c_\theta\rangle=\cosh^2\xi,
        \qquad\langle c_\xi,c_\theta\rangle=0.
    \end{align*}
    %
    A Euclidean unit normal in the hyperplane $w=0$ is
    %
    \begin{align*}
        \nu_c(\xi,\theta)
        :=\left(-\frac{\cos\theta}{\cosh\xi},
                -\frac{\sin\theta}{\cosh\xi},
                 \tanh\xi,0\right).
    \end{align*}
    %
    Differentiation gives
    %
    \begin{align*}
        c_{\xi\xi}
        &=(\cosh\xi\cos\theta,\cosh\xi\sin\theta,0,0),\\
        c_{\theta\theta}
        &=(-\cosh\xi\cos\theta,-\cosh\xi\sin\theta,0,0),\\
        c_{\xi\theta}
        &=(-\sinh\xi\sin\theta,\sinh\xi\cos\theta,0,0).
    \end{align*}
    %
    Pairing these vectors with $\nu_c$ gives $-1$, $1$,
    and $0$, respectively. The second normal is the constant
    vector $(0,0,0,1)$, and all three pairings with it vanish.
    Hence the Euclidean Gauss equation gives
    %
    \begin{align*}
        K_c
        =\frac{(-1)\cdot1-0^2}{\cosh^4\xi}
        =-\cosh^{-4}\xi.
    \end{align*}
    %
    All these formulas hold for every $(\xi,\theta)\in D_{3R}$.

    By the Gauss equation, Gaussian curvature depends
    continuously on the immersion through its first two
    derivatives and on the ambient metric through its first
    two derivatives, provided the induced metric remains
    positive definite. The limiting induced metric above
    has this property uniformly on $D_{3R}$. The $C^2$
    convergence of the immersions and of the ambient metrics
    therefore implies uniform convergence of their Gaussian
    curvatures on $D_{2R}$.
    The rescaled induced metric is $\tau^{-2}$ times the
    pullback of the final induced metric. Its Gaussian
    curvature is consequently $\tau^2 K_{\Sigma_m}$.
    We have proved
    %
    \begin{align*}
        \sup_{p\in L_m}\ 
        \sup_{(\xi,\theta)\in D_{2R}}
        \left|\tau^2 K_{\Sigma_m}
          \bigl(\widehat X_m(\mathcal C_{\zeta,p}(\xi,\theta))\bigr)
             +\cosh^{-4}\xi\right|
        \rightarrow0.
    \end{align*}
    %
    On $D_{2R}$, one has $\cosh^{-4}\xi\geq(2R)^{-4}$.
    Choosing $m$ large enough that the displayed supremum
    is at most $(2R)^{-4}/2$ proves
    \eqref{eq:s4-core-negative-curvature}.
\end{proof}

\begin{proof}[Proof of Theorem~\ref{thm:s4-main}]
    Choose $R\geq R_0$ as required by
    Lemma~\ref{lem:s4-scalar-negative-curvature}, and keep
    it fixed. Then choose $m$ sufficiently large that
    all the estimates above hold for this $R$ and
    $C(m^{-2}+\tau)<c_*/2$ in
    \eqref{eq:s4-transverse-curvature-comparison}.
    If $x\in M_\zeta$ satisfies $\rho_\zeta(x)\geq2R\tau$,
    put $q:=\pi(G_m(x))$ and use the parametrizations
    $X,X_b$ of Lemma~\ref{lem:s4-outer-coordinates}.
    Reflection by $Z$, when needed to use the upper sheet,
    preserves the induced metric and Gaussian curvature.
    The two curvature comparisons therefore give
    %
    \begin{align*}
        K_{\Sigma_m}(\widehat X_m(x))
        &=K_{X_b}(q)\\
        &\leq K_X(q)+|K_{X_b}(q)-K_X(q)|\\
        &\leq\bigl[-c_*+C(m^{-2}+\tau)\bigr]\mathcal W(q)\\
        &\leq-\frac{c_*}{2}\mathcal W(q)<0.
    \end{align*}
    %
    If instead $\rho_\zeta(x)\leq2R\tau$,
    Lemma~\ref{lem:s4-core-negative-curvature} gives
    strict negativity. These two conditions use the same
    function $\rho_\zeta$ on the initial surface and cover
    all of $M_\zeta$. Thus no region is omitted because of
    the change from $\pi(x)$ to $\pi(G_m(x))$ during the
    first deformation. Since $\widehat X_m$ is an embedding,
    we have proved $K_{\Sigma_m}<0$ everywhere.
    Proposition~\ref{prop:s4-coupled-existence} already gives
    smoothness, closedness, connectedness, embeddedness,
    minimality, and genus $m^2+1$.

    We next compute the area of the initial surface. Set
    %
    \begin{align*}
        \mathcal U_m
        :=\mathbb T\smallsetminus\bigcup_{p\in L_m}B(p,3r_c),
    \end{align*}
    %
    where the balls are taken in the flat metric
    $g_{\mathbb T}$. Over this set, $M_\zeta$ consists of
    the two graphs $q\mapsto\mathcal F(q,h(q))$ and
    $q\mapsto\mathcal F(q,-h(q))$. The initial height
    estimates give
    %
    \begin{align*}
        |h(q)|\leq C\varepsilon,
        \qquad |Dh(q)|\leq C\tau/r_c
        \qquad\forall\,q\in\mathcal U_m.
    \end{align*}
    %
    The metric of either graph has matrix
    $\cos^2h\,I+Dh\otimes Dh$ in the flat coordinates.
    Its determinant is
    $\cos^4h+\cos^2h\,|Dh|^2$, so the area density
    relative to $d\operatorname{vol}_{g_{\mathbb T}}$ is
    %
    \begin{align*}
        J_h=\cos h\sqrt{\cos^2h+|Dh|^2}.
    \end{align*}
    %
    We used $|h|<z_0<\pi/4$ to take the positive square
    root of $\cos^2h$. For $|h|+|Dh|$ small,
    %
    \begin{align*}
        |J_h-1|
        &\leq|\cos^2h-1|
             +\cos h\left|\sqrt{\cos^2h+|Dh|^2}-\cos h\right|\\
        &=\sin^2h+
          \frac{\cos h\,|Dh|^2}
               {\sqrt{\cos^2h+|Dh|^2}+\cos h}\\
        &\leq C(h^2+|Dh|^2)
         \leq C(\varepsilon^2+\tau^2r_c^{-2}).
    \end{align*}
    %
    The last bound tends to zero, because
    $\varepsilon=m^2\tau/(2\pi)\to0$ and
    $\tau/r_c=\tau^{1-\alpha}\to0$.
    For sufficiently large $m$, the $m^2$ disks removed
    from $\mathbb T$ are disjoint embedded flat disks.
    Their total area is
    %
    \begin{align*}
        \operatorname{Area}(\mathbb T\smallsetminus\mathcal U_m)
        =9\pi m^2r_c^2=9\pi m^2\tau^{2\alpha}
        \rightarrow0.
    \end{align*}
    %
    Integrating the estimate for $J_h$ over $\mathcal U_m$
    shows that the sum of the areas of the two graph pieces
    tends to $2\operatorname{Area}(\mathbb T)=4\pi^2$.

    It remains to bound the initial pieces above the removed
    disks. For each $p\in L_m$, these consist of a catenoid
    with horizontal radius at most $2r_c$ and two graphs
    over the annulus $2r_c\leq r_p\leq3r_c$. The metric
    $dz^2+\cos^2z\,(ds^2+dt^2)$ is bounded above, as a
    quadratic form, by $dz^2+ds^2+dt^2$. Thus the area of
    the catenoidal piece is no larger than its Euclidean
    area. Put $a_c:=\operatorname{arcosh}(2r_c/\tau)$.
    The Euclidean metric induced by this catenoid is
    $\tau^2\cosh^2\xi(d\xi^2+d\theta^2)$, and hence
    %
    \begin{align*}
        \operatorname{Area}_{\mathrm{Eucl}}
        &=\int_{-a_c}^{a_c}\int_0^{2\pi}
             \tau^2\cosh^2\xi\,d\theta\,d\xi\\
        &=2\pi\tau^2(a_c+\sinh a_c\cosh a_c)\\
        &=2\pi\left(
             2r_c\sqrt{4r_c^2-\tau^2}
             +\tau^2\operatorname{arcosh}(2r_c/\tau)
             \right)\\
        &\leq Cr_c^2.
    \end{align*}
    %
    Indeed, $2r_c\sqrt{4r_c^2-\tau^2}\leq4r_c^2$,
    and $\operatorname{arcosh}u\leq u^2$ for $u\geq1$
    gives the same bound for the second term in parentheses.
    On each of the two annuli, $|Dh|\leq C\tau/r_c$,
    so the graph area density is bounded by a fixed
    constant. Their areas are consequently at most $Cr_c^2$
    as well. Summing over the $m^2$ lattice points bounds
    the total omitted area by
    $Cm^2r_c^2=Cm^2\tau^{2\alpha}\to0$. We obtain
    %
    \begin{align}
        \label{eq:s4-initial-area-limit}
        \operatorname{Area}(M_\zeta)\rightarrow4\pi^2.
    \end{align}

    Finally, we prove that the two deformations change area
    by a factor tending uniformly to one. Write
    $\widetilde g_m:=G_m^*g_{\mathcal N}$ for the metric
    on $M_\zeta$ after the deformation within $\mathcal N$.
    We will use the following dimensionless bound for the
    displacement and its first derivatives:
    %
    \begin{align*}
        \vartheta_m&:=\sup_{\tau\leq r\leq\pi/m}
              \frac{\tau^{1+\alpha/4}f_\tau(r)}r\\
        &=\tau^{1+\alpha/4}
          \sup_{\tau\leq r\leq\pi/m}
          \max\{r^{1/2},\tau^{(1-\alpha)/2}r^{-1/2}\}\\
        &\leq C\left(
             \tau^{1+\alpha/4}m^{-1/2}
             +\tau^{1-\alpha/4}\right)
          \rightarrow0.
    \end{align*}
    %
    In the last inequality, the maximum of $r^{1/2}$ is
    attained at $r=\pi/m$, and that of $r^{-1/2}$ at
    $r=\tau$. The use of the bound
    $\tau^{1+\alpha/4}f_\tau$ for $\phi_m$ is justified
    by the uniform estimates on $\mathcal B_m$; the displayed
    bound is \eqref{eq:s4-dimensionless-phi} applied to the
    fixed point.

    To obtain a metric estimate, fix $x\in M_\zeta$, set
    $r:=\rho_\zeta(x)$, and use a rescaled local parametrization
    $\kappa:D\to M_\zeta$ as in the proof of
    Lemma~\ref{lem:s4-full-family-geometry}. Regard the
    initial and deformed maps as maps into $\mathbb R^5$,
    and put
    %
    \begin{align*}
        I(y)&:=r^{-1}\bigl(\jmath_\zeta(\kappa(y))
                               -\jmath_\zeta(x)\bigr),\\
        I_\phi(y)&:=r^{-1}\bigl(G_m(\kappa(y))
                               -\jmath_\zeta(x)\bigr)
        \qquad\forall\,y\in D.
    \end{align*}
    %
    The chain-rule estimates for the exponential map in that
    proof give $\|dI_\phi-dI\|_{C^0(D)}\leq C\vartheta_m$.
    The initial derivatives $dI$ are uniformly bounded, and
    $r^{-2}\kappa^*g_\zeta$ has a uniform positive lower
    bound. The coefficients of the two rescaled induced
    metrics are Euclidean inner products of the derivatives
    of $I$ and $I_\phi$. For $i,j\in\{1,2\}$, their
    difference is
    %
    \begin{align*}
        \langle\partial_iI_\phi,\partial_jI_\phi\rangle
             -\langle\partial_iI,\partial_jI\rangle
        ={}&\langle\partial_i(I_\phi-I),\partial_jI\rangle
          +\langle\partial_iI,\partial_j(I_\phi-I)\rangle\\
          &+\langle\partial_i(I_\phi-I),
                     \partial_j(I_\phi-I)\rangle.
    \end{align*}
    %
    It is bounded by $C\vartheta_m(1+\vartheta_m)$. The
    lower bound for the initial metric converts this
    coefficient estimate into a relative estimate. Since
    $x$ was arbitrary, we obtain the quadratic-form bounds
    %
    \begin{align*}
        (1-C\vartheta_m(1+\vartheta_m))g_\zeta
        \leq\widetilde g_m
        \leq(1+C\vartheta_m(1+\vartheta_m))g_\zeta.
    \end{align*}

    We estimate the second change of metric directly.
    Regard $\iota_{\mathcal N}\circ G_m$ as an immersion
    into $\mathbb S^4$ with induced metric $\widetilde g_m$,
    where $\iota_{\mathcal N}$ is the inclusion of
    $\mathcal N$ into $\mathbb S^4$. Its unit normal in
    the direction orthogonal to $\mathcal N$ is $n_m$.
    Let
    $\mathcal S_m:=S_{\iota_{\mathcal N}\circ G_m}^{n_m}$
    be its shape operator in that direction, and let
    $\omega_m$ be the normal-connection one-form for the
    frame consisting of its chosen unit normal within
    $\mathcal N$ and $n_m$.
    The bounds in \eqref{eq:s4-full-family-geometry},
    \eqref{eq:s4-final-two-jets}, and the first metric
    comparison give
    %
    \begin{align*}
        |\mathcal S_m|_{\widetilde g_m}
             +|\omega_m|_{\widetilde g_m}&\leq C,\\
        |b_m|&\leq C\tau^2,\\
        |db_m|_{\widetilde g_m}
             &\leq C\tau^2/\rho_\zeta\leq C\tau
    \end{align*}
    %
    everywhere on $M_\zeta$.
    The exact metric formula
    \eqref{eq:s4-graph-metric-exact} gives, for every
    $x\in M_\zeta$ and $v\in T_xM_\zeta$,
    %
    \begin{align*}
        (\widehat X_m^*g_{\mathbb S^4})_x(v,v)
        =\bigl|\cos(b_m(x))v
                 -\sin(b_m(x))\mathcal S_m(v)
             \bigr|_{\widetilde g_m}^2+\sin^2(b_m(x))\omega_m(v)^2+db_m(v)^2.
    \end{align*}
    %
    Expanding the first square and subtracting
    $|v|_{\widetilde g_m}^2$ yields the terms
    %
    \begin{align*}
        &-\sin^2(b_m(x))|v|_{\widetilde g_m}^2
         -2\cos(b_m(x))\sin(b_m(x))
             \widetilde g_m(v,\mathcal S_m(v))\\
        &\qquad
         +\sin^2(b_m(x))|\mathcal S_m(v)|_{\widetilde g_m}^2
         +\sin^2(b_m(x))\omega_m(v)^2+db_m(v)^2.
    \end{align*}
    %
    Using the preceding bounds and $|\sin t|\leq|t|$,
    their sum has absolute value at most
    %
    \begin{align*}
        C\bigl(|b_m(x)|+|b_m(x)|^2
                  +|db_m(x)|_{\widetilde g_m}^2\bigr)
                  |v|_{\widetilde g_m}^2
        \leq C\tau^2|v|_{\widetilde g_m}^2.
    \end{align*}
    %
    Consequently,
    %
    \begin{align*}
        (1-C\tau^2)\widetilde g_m
        \leq\widehat X_m^*g_{\mathbb S^4}
        \leq(1+C\tau^2)\widetilde g_m.
    \end{align*}
    %
    Combining the two comparisons gives numbers
    $\epsilon_m\to0$, with $0\leq\epsilon_m<1$ for
    large $m$, such that
    %
    \begin{align*}
        (1-\epsilon_m)g_\zeta
        \leq\widehat X_m^*g_{\mathbb S^4}
        \leq(1+\epsilon_m)g_\zeta.
    \end{align*}
    %
    At each point, the two eigenvalues of the final metric
    relative to $g_\zeta$ lie in
    $[1-\epsilon_m,1+\epsilon_m]$. The ratio of the area
    elements is the square root of their product, so it
    lies in the same interval. Since $\widehat X_m$ is an
    embedding, integration gives
    %
    \begin{align*}
        (1-\epsilon_m)\operatorname{Area}(M_\zeta)
        \leq\operatorname{Area}(\Sigma_m)
        \leq(1+\epsilon_m)\operatorname{Area}(M_\zeta).
    \end{align*}
    %
    The area limit follows from
    \eqref{eq:s4-initial-area-limit}, completing the proof.
\end{proof}

\let\c\textcedilla
\bibliographystyle{amsalpha}
\bibliography{main}
\end{document}
